%!TEX root = /Users/kbj/Dropbox/*classes/601book/601master.tex
\setcounter{chapter}{5}
\chapter{Theories of Movement}
\label{ch:movement}
\setcounter{exx}{0}
% %\setcounter{page}{131}

\noindent
Our present grammar defines a sentence as a series of parses.
\begin{exe}
	\raggedright
	\ex \label{ex6:derivation} 
	$\mathcal{D}$ = (P_1, P_2\ldots,P_n) is a derivation iff:
	\begin{xlist}
		\ex P_1 is formed by the \ibar{X} Skeleton and Lexical Insertion, and it is semantically interpreted,
		\ex Each P_{i>1} is formed from P_{i-1} by a rule,
		\ex P_n is subject to the linearization algorithm.
	\end{xlist}
	\marginnote[0pt]{The cardinality of a derivation is the number of Ps in it.}
	\ex A sentence is the set of well-formed derivations which have the same P_1 (=D-structure). A grammatical sentence is a derivation whose cardinality isn't greater than the cardinality of any other derivation for that sentence.
\end{exe}
One of the rule types that creates parses is movement:
\begin{exe}
	\raggedright
	\ex \label{ex6:movedef} $\alpha$, a node in a phrase marker tree, \underline{moves} to $\beta$, a position in a phrase marker tree, means:
	\begin{xlist}
		\ex $\alpha$ is put in $\beta$, and
		\ex $\alpha$ is removed from its original position.
	\end{xlist}
\end{exe}
We've developed a system for movement rules that characterizes the changes of position that move creates in solely morphological terms. Argument Movement arises, for instance, because the \obar{X} that assigns Case to a DP has a $\bullet$ associated with it. That is, the agreement relation that establishes Case assignment comes with the instruction that the Case marked DP should be pronounced in proximity to the Case assigner. And Head Movement arises under similar circumstances. When a tense morpheme, for instance, is in an agreement relation with a verb---causing that verb to expone the tense morpheme---and one or the other of the terms in that agreement relation has a $\bullet$ associated with it, Head Movement is required to bring the two heads together into one head.
\pagebreak

Let's reconsider how this system works by way of~(\ref{ex6:firstder}).
\begin{exe}
	\raggedright
	\ex \label{ex6:firstder}
	\small
\begin{forest}
[TP [\ibar{T}
	[\obar{T} [nom\\ $\bullet u \phi$\\ \emph{pres}]]
	[vP 
		[DP\\ $u$nom\\ $\phi$ [we, roof]]
		[\ibar{v}
			[\obar{v} [$\bullet$\emph{agent}]]
			[$\mu$P [\ibar{$\mu$}
				[\obar{$\mu$} [$\bullet$acc]]
					[VP [\ibar{V}
						[\obar{V} [$\surd$buy\\ $u \phi$\\ $u$pres\\ $u$agent\\ $u$acc]]
						[$u$acc\\DP [dogs, roof]]
					]]
			]]
		]
	]
]]
\end{forest}
\end{exe}
This structure can be created by the \ibar{X} Skeleton, and it can be semantically evaluated. It isn't a well-formed representation otherwise, however, as it does not meet the requirement that features \textsc{agree}, nor does it honor the requirements imposed by~$\bullet$.

To meet these requirements, a derivation must be created which includes parses that do meet these requirements. This produces the steps in~(\ref{ex6:dersecond}). (I've omitted from this derivation some intermediary parses that don't create useful parses.)
\begin{exe}
	\ex\label{ex6:dersecond}
	\footnotesize
	a.
\begin{forest}
[TP[\ibar{T}
	[\obar{T} [nom\\ $\bullet$\ding{51}$\phi$\\ \emph{pres}] ]
	[vP 
		[DP\\ \ding{51}nom\\ $\phi$ [we, roof]]
		[\ibar{v} 
			[\obar{v} [$\bullet$\emph{agent}] ]
			[$\mu$P [\ibar{$\mu$} 
				[\obar{$\mu$} [$\bullet$acc] ]
				[VP [\ibar{V} 
					[\obar{V} 
						[\obar{V} [$\surd$buy\\ $u \phi$\\ $u$pres\\ $u$agent\\ \ding{51}acc] ]
						[up]
					 ]
					[DP\\ \ding{51}acc [dogs, roof]]
				]]
			]
		]
	]
]]]
\end{forest}
b.
\begin{forest}
[TP [\ibar{T}
	[\obar{T} [nom\\ $\bullet$\ding{51}$\phi$\\ \emph{pres}] ]
	[vP
	[DP\\ \ding{51}nom\\ $\phi$ [we, roof] ]
	[\ibar{v} 
			[\obar{v} [$\bullet$\emph{agent}] ]
			[$\mu$P 
			[DP\\ \ding{51}acc [dogs, roof]]
			[\ibar{$\mu$} 
				[\obar{$\mu$} [$\bullet$acc] ]
				[VP [\ibar{V} 
					[\obar{V} 
						[\obar{V} [$\surd$buy\\ $u \phi$\\ $u$pres\\ $u$agent\\ \ding{51}acc] ]
						[up]
					 ]
				]]
			]
		]
	]
]]]
\end{forest}

c.
\begin{forest}
	[TP [\ibar{T}
		[\obar{T} [nom\\ $\bullet$\ding{51}$\phi$\\ \emph{pres}] ]
		[vP
		[DP\\ \ding{51}nom\\ $\phi$ [we, roof] ]
		[\ibar{v} 
				[\obar{v} [$\bullet$\emph{agent}] ]
				[$\mu$P 
				[DP\\ \ding{51}acc [dogs, roof]]
				[\ibar{$\mu$} 
					[\obar{$\mu$}  [\obar{V} [$\surd$buy\\ $u \phi$\\ $u$pres\\ \ding{51}agent\\ \ding{51}acc] ]
					[\obar{$\mu$} [$\bullet$acc] ]]
					[VP [\ibar{V} 
						[\obar{V} 
							[{\ }]
							[up]
						 ]
					]]
				]
			]
		]
	]]]
\end{forest}
	d.
	\begin{forest}
		[TP [\ibar{T}
			[\obar{T} [nom\\ $\bullet$\ding{51}$\phi$\\ \emph{pres}] ]
			[vP
				[DP\\ \ding{51}nom\\ $\phi$ [we, roof] ]
				[\ibar{v} 
					[\obar{v} 
						[\obar{$\mu$} 
						[\obar{V} [$\surd$buy\\ $u \phi$\\ $u$pres\\ \ding{51}agent\\ \ding{51}acc] ]
							[\obar{$\mu$} [$\bullet$acc] ]
						]
						[\obar{v} [$\bullet$\emph{agent}]]
					]
					[$\mu$P
						[DP\\ \ding{51}acc [dogs, roof]]
						[\ibar{$\mu$}
							[{\ }]
							[VP [\ibar{V}
								[\obar{V}
									[{\ }]
									[up]
								]
							]
							]
						]
						]
					]
				]
			]
		]
	\end{forest}
	
	e.
	\begin{forest}
		[TP [\ibar{T}
			[\obar{T} [nom\\ $\bullet$\ding{51}$\phi$\\ \emph{pres}], name=probe ]
			[vP
				[DP\\ \ding{51}nom\\ $\phi$ [we, roof] ]
				[\ibar{v} 
					[\obar{v} 
						[\obar{$\mu$} 
						[\obar{V} [$\surd$buy\\ \ding{51}$\phi$\\ \ding{51}pres\\ \ding{51}agent\\ \ding{51}acc], name=goal ]
							[\obar{$\mu$} [$\bullet$acc] ]
						]
						[\obar{v} [$\bullet$\emph{agent}]]
					]
					[$\mu$P
						[DP\\ \ding{51}acc [dogs, roof]]
						[\ibar{$\mu$}
							[{\ }]
							[VP [\ibar{V}
								[\obar{V}
									[{\ }]
									[up]
								]
							]
							]
						]
						]
					]
				]
			]
		]
		% \draw[|-|, dotted] (goal) to[out=south west,in=south west] node[pos=0.5,xshift=-5mm]{\emph{agree}} (probe);
	\end{forest}
\end{exe}
These four instances of movement and \textsc{agree} have satisfied the $\bullet$ attribute on \emph{acc}, and have brought $\surd$buy into a position which allows it to value the person/number/tense features and produce a terminal that corresponds to \emph{buy}. The only thing left to do to make this a well-formed s-structure is to satisfy the $\bullet$ attribute on the $\phi$ feature in \obar{T}. This is done by moving the subject so that it merges with a projection of~\obar{T}.

I pointed out that (\ref{ex6:firstder}) is a structure that our semantics can interpret, and which gives us the correct interpretation. But none of the representations produced from (\ref{ex6:firstder}) can be semantically interpreted however. This is a general property of the derivations we've considered so far. A reasonable hypothesis is that this defines derivations. That is, that the first parse in a derivation is the one that is semantically interpreted, and no other parse is. The D-structure parse is the one that represents how the morphemes in a sentence are grouped semantically, and the S-structure parse is the one from which the morphological and prosodic grouping is determined. All the other parses in a derivation are just are model of the function that maps a D-structure onto an S-structure. 

That derivations are defined in this way was a popular view in the 1960's. In 1964, Jerrold Katz and Paul Postal, formulated what's come to be known as the Katz-Postal Hypothesis.\nocite{Katz:1964} That thesis maintains that the denotations of the elements that make up a sentence are semantically composed in a way that the syntax of what we're calling D-structures corresponds to. The rest of the forms in a derivation, including the S-structure, do not reflect the way in which the denotations of the component parts are composed. Syntax furnishes, on the one hand, a way of structuring how the denotations are combined and, on the other hand, a way of structuring the morphology: bringing together the information that make up lexical items. A theory of derivations is a theory of the relation between these two representations. I don't know where the $\bullet$ attribute's requirements fit into this picture, but since their requirements are ones that hold of the morphological connection that is established by \textsc{agree}, it doesn't seem too wild to hope that it too will come from some aspect of morphology. Syntax, then, could be seen as a way of mediating the independent needs that the morphology and the semantics place on the items that are being arranged into sentences. Derivations are models of the solution space to the open equations that the morphological and semantic requirements produce. It's this conception of syntax that we've been using up to now.

In the derivation we've just looked at, the movement operations are all ones that cannot feed the semantic interpretation. In this chapter we're going to confront evidence that this is incorrect. In particular, we'll see that movement can create structures that are semantically interpreted. As it stands, our definition of movement does nothing more than change the syntactic position a term has in order to make sentences comply with morphological requirements. But there are examples which seem to require that movement can itself create meanings. We'll need to develop two things together, then. We'll need a new definition of derivation, one that allows more than just the D-structure to be semantically interpreted. And we'll need a new definition of movement, one that weds a change in meaning to the change in position that it presently expresses.

\section{Trace Theory}
\label{sec:tracetheory}

To see that parses other than D-structures are semantically interpreted, we're going to have to beef up our tools for seeing what the meanings of sentences are. There's a kind of anaphora that provides one such tool, and I'll use it in the following sections to give us a better picture of the meanings of sentences that involve movement. There are also argument DPs whose meanings are not as simple as those we've been considering up till now. These too provide us a better image of how the syntactic parts of a sentence relate to the meaning that sentence has. I'll show how the picture that results from looking at these cases can be modeled by a theory of movement that allows it to insert variables which are bound by the moved term. Those variables are called ``traces,'' and the theory of movement that makes use of them is called ``Trace theory.'' But first, we must discover what variables are, and how they are expressed in syntax. Variables are part of the system of anaphora that natural languages use. We'll begin, then, by introducing some preliminaries about anaphora.

\subsection{Deep and Surface anaphors}
\label{sec:deepsurfaceanapors}

In one kind of anaphora, the denotation some expression has is matched by the denotation of another linguistic expression. \citet{Hankamer:1976} calls this ``surface anaphora,'' and it is to be distinguished from ``deep anaphora.'' In deep anaphora, the term that the anaphor refers to is determined by virtue of how another linguistic expression refers. For example, the pronoun \emph{she} in (\ref{ex6:da1}) is most naturally understood to refer to the same individual that \emph{Rodica} refers to.
\begin{exe}
	\raggedright
	\ex \label{ex6:da1}
	Rodica said that she would explain this.
\end{exe}
Personal pronouns refer to individuals which the context has made salient. In sentences that have been isolated from a context of use---like (\ref{ex6:da1})---this is typically done linguistically. In (\ref{ex6:da1}), the presence of the name \emph{Rodica} has made salient in the relevant sense an individual with that name, and \emph{she} is most naturally understood to refer to that individual. We can describe this situation as one in which \emph{Rodica} and \emph{she} \underline{corefer}, and we'll call the term that establishes the requisite saliency in these cases the pronoun's \underline{antecedent}. Coreference is an instance of deep anaphora. The anaphoric term, the pronoun, has its reference established with the aid of an antecedent. In (\ref{ex6:da1}), the anaphora arises indirectly; the DP \emph{Rodica} makes an individual named Rodica salient, and this makes Rodica a possible referent for \emph{she}.

Another instance of deep anaphora arises when pronouns are understood to be bound variables. In these kinds of examples, illustrated by~(\ref{ex6:da2}), the anaphoric connection between a pronoun and another DP is more direct.
\begin{exe}
	\raggedright
	\ex \label{ex6:da2}
	Every woman remembers her first paper.
\end{exe}
As in (\ref{ex6:da1}), the most natural interpretation of (\ref{ex6:da2}) gives the pronoun \emph{her} an interpretation that depends on the meaning of \emph{every woman}. But in this case, the pronoun doesn't corefer with \emph{every woman}. In fact, the meaning of \emph{every woman} is complex; it doesn't refer. %
%
\marginnote[0pt]{The determiner \emph{every} is called a ``quantifier,'' and the DP it heads is called a ``quantificational DP.'' Other quantificational determiners in English are \emph{some}, \emph{a}, \emph{no}, and \emph{each}. The semantics I will adopt for quantifiers comes from what is known as ``generalized quantifier'' theory. The central idea of this theory is that quantifiers denote relations between sets of things that meet certain restrictions. In (\ref{ex6:da2}), those sets of things are (a) women and (b) the rememberers of her first paper: \emph{every} says that the first set is a subset of the second. This way of thinking about quantificational DPs originates with \citet{Frege:1879}, and is formalized in \citet{Mostowski:1957} and \citet{Lindstrom:1966}. Good introductions to the linguistics of generalized quantifiers are \citet{Barwise:1981} and \citet{Keenan:1986}.}%
%
Instead, it says something more general about the relation between women and rememberers. One can paraphrase the effect that \emph{every woman} has on the meaning of (\ref{ex6:da2}), and include the way in which \emph{her} is connected to that meaning with~(\ref{ex6:da2r}).
\begin{exe}
	\raggedright
	\ex \label{ex6:da2r}
	For all the women, 1 \ldots n, under consideration:
	\begin{xlist}
		\ex woman \#1 remembers woman \#1's first paper, and
		\ex woman \#2 remembers woman \#2's first paper, and
		\ex woman \#3 remembers woman \#3's first paper, and
		\sn \hspace*{1.25in}\vdots
		\ex woman \#n remembers woman \#n's first paper.
	\end{xlist}
\end{exe}
An interesting property of the meaning of sentences like (\ref{ex6:da2}) is that they can be matched to a family of sentence meanings, like that in (\ref{ex6:da2r}). Each of the sentences in the family varies by virtue of a value given to certain terms within. One of these terms in (\ref{ex6:da2}) is \emph{her}, and the other is the external argument, \emph{every woman}. The pronoun's interpretation varies with the values given to \emph{every woman}; it is a variable bound to \emph{every woman}.

Let's unpack what's happening in (\ref{ex6:da2}). I'll begin by considering how to model the meaning (\ref{ex6:da2}) gets when the pronoun \emph{her} is not interpreted as a variable bound by \emph{every woman}. Let's imagine, for instance, that (\ref{ex6:da2}) is uttered as part of a conversation about Thui. In that context, (\ref{ex6:da2}) gets an interpretation that can be paraphrased by (\ref{ex6:da3}), where \emph{her} refers to Thui.
\pagebreak
\begin{exe}
\raggedright
	\ex \label{ex6:da3}
	For all the women, 1 \ldots n, under consideration:
	\begin{xlist}
		\ex woman \#1 remembers woman her first paper, and
		\ex woman \#2 remembers woman her first paper, and
		\ex woman \#3 remembers woman her first paper, and
		\sn \hspace*{1.25in}\vdots
		\ex woman \#n remembers woman her first paper.
	\end{xlist}
\end{exe}
What's required is to find a way of giving a denotation to \emph{every woman} that allows (\ref{ex6:da2}) to get a meaning like that (\ref{ex6:da3}) expresses. 

One way of achieving this meaning is to let the determiner \emph{every} have a denotation that causes it to relate the denotation of \emph{woman} to the denotation of \emph{remembers her first paper} in the way that (\ref{ex6:da3}) describes. Suppose that we take the denotation of \emph{woman} to describe things which are women. We can do this by letting \emph{woman} be a function that maps things to truth just in case they are women. Using the $\lambda$ operator, we can model that with~(\ref{ex6:cnd}).
\begin{exe}
	\raggedright
	\ex \label{ex6:cnd}
	\denotes{woman} =_{\text{def}} $\lambda x$ woman($x$)=T
\end{exe}\marginnote[-15pt]{Remember that \denotes{$\alpha$} is the denotation of~$\alpha$.}
We can now capture the meaning of (\ref{ex6:da2}) with the denotation for \emph{every} in~(\ref{ex6:dno}).%
%
\marginnote[-5pt]{The ``$\rightarrow$'' in (\ref{ex6:dno}) is a logical connective that comes close to the meaning of the word ``implies.'' $\alpha$ $\rightarrow$ $\beta$ is True unless $\alpha$ is True and $\beta$ is False.}%
%
\begin{exe}
	\raggedright
	\ex \label{ex6:dno}
	\denotes{every} =_{\text{def}} $\lambda P\ \lambda Q\ \lambda \epsilon\ \forall x [P(x) \rightarrow Q(x)(\epsilon)]$.
\end{exe}
This gives the vP of (\ref{ex6:da2r}) the denotation in~(\ref{ex6:da2rd}).
\medskip
\begin{fullwidth}
\begin{exe}
	\raggedright
	\ex \label{ex6:da2rd}
	\hspace*{-15pt}
	\small
\begin{forest}
[$\lambda \epsilon\ \forall x$ \textsc{woman}($x$) $\rightarrow$ \textsc{agent}($x$)($\epsilon$) $\wedge$ \textsc{theme}(\denotes{her first paper})($\epsilon$) $\wedge$ \textsc{remember}($\epsilon$) \\vP
	[$\lambda Q \lambda \epsilon \neg \exists x. $ \textsc{woman}($x$) $\wedge$ $Q(x)(\epsilon)$\\DP 
	[\ibar{D}
		[$\lambda P \lambda Q \lambda \epsilon \forall x P(x) \rightarrow Q(x)(\epsilon)$\\\obar{D} [every]]
		[$\lambda x$ \textsc{woman}($x$)\\ NP [woman, roof] ]
	] 
	]
	[$\lambda x \lambda \epsilon$ \textsc{agent}($x$)($\epsilon$) $\wedge$ \textsc{theme}(\denotes{her first paper})($\epsilon$) $\wedge$ \textsc{remembers}($\epsilon$) \\ \ibar{v}
		[remembers her first paper, roof ]
	]
]
\end{forest}
\end{exe}
\end{fullwidth}
\medskip
When \emph{every} semantically combines with \emph{woman}, a denotation is formed that is a function whose domain is relations between entities and events. Our semantics makes the denotation of \ibar{v} a function of that sort, and so the function that the DP denotes can take the denotation of the \ibar{v} as its argument. The result is \denotes{vP}, which is a function from events to the statement that for every x if x is a woman then x is the agent of that event, where that event is a remembering whose theme is the thing that \emph{her first paper} refers to. The vP describes, then, events which are rememberings of her (=Thui's) first paper by x if for every x, x is a woman. This is a very close approximation of the meaning expressed by~(\ref{ex6:da3}). %
%
\marginnote[-10pt]{We'll see later that there is a subtle, but important, mismatch between what this semantics delivers and the meaning that (\ref{ex6:da2}) actually has. Our semantics will require some tweaking.}%
%
An event of the kind that \denotes{vP} describes requires that the sequence of conjunctions in (\ref{ex6:da2r}) hold of that event. That is, for an event to make the function that \denotes{vP} is map to True, it must be an event in which woman \#1 remembers Thui's first paper and woman \#2 remembers Thui's first paper, and woman \#3 remembers Thui's first paper, etc., until every woman (in the context) is listed among the rememberers.

The bound variable interpretation of the pronoun \emph{her} in (\ref{ex6:da2r}) can be captured if it is possible to make the denotation of \emph{her} one of the things that the ``''$\forall$'' binds in the denotation delivered to the vP. That is, we need to find a procedure that allows \emph{her} to get interpreted as the variable $x$ bound by $\forall$ in the way that is indicated in~(\ref{ex6:da2rd2}).
\medskip
\begin{fullwidth}
\begin{exe}
	\raggedright
	\ex \label{ex6:da2rd2}
	\hspace*{-22pt}
	\small
\begin{forest}
[$\lambda \epsilon \forall x$ \textsc{woman}($x$) $\rightarrow$ \textsc{agent}($\epsilon$)($x$) $\wedge$ \textsc{theme}(\denotes{$x$'s first paper})($\epsilon$) $\wedge$ \textsc{remember}($\epsilon$)\\ vP
	[$\lambda Q \lambda \epsilon \forall x. \textsc{woman}(x) \rightarrow\ Q(x)(\epsilon)$\\ DP 
		[\ibar{D}
			[$\lambda P\ \lambda Q \lambda \epsilon \forall x. P(x) \rightarrow\ Q(x)(\epsilon)$\\ \obar{D} [every]]
			[NP [woman, roof]]
		]
	]
	[$\lambda x \lambda \epsilon$ \textsc{agent}($x$)($\epsilon$) $\wedge$ \textsc{theme}(\denotes{$x$'s first paper})($\epsilon$) $\wedge$ \textsc{remember}($\epsilon$)\\ \ibar{v} [remember her first paper, roof]]
]
\end{forest}
\end{exe}
\end{fullwidth}
\medskip
When the domain of the function denoted by \emph{woman}, i.e., $\lambda x$ \textsc{woman}($x$)=T, is restricted to just those individuals about which the sentence is making a claim---that is, it includes only those women under consideration---(\ref{ex6:da2rd2}) is a model of the meaning for (\ref{ex6:da2}) described in~(\ref{ex6:da2r}). One of our aims, then, will be to formulate a syntax that produces a way of interpreting a pronoun as a variable bound in the way that (\ref{ex6:da2rd2}) indicates.

These are instances of deep anaphora. Surface anaphora arises when the anaphor's denotation is equated with the denotation of another term: its antecedent. An example of this is in~(\ref{ex6:one1}).
\begin{exe}
	\raggedright
	\ex \label{ex6:one1}
	Rodica read a bad \underline{paper} after she read a good one.\\ Rodica read a bad \underline{book} after she read a good one.\\ Rodica read a bad \underline{poem} after she read a good one.
\end{exe}
The meaning of \emph{one} in (\ref{ex6:one1}) matches the meaning of the underlined noun. In these cases, \emph{one} is not referring to the same thing that the underlined noun is referring to. Instead, the semantic contribution \emph{one} makes to the DP that it stands within is the same as the semantic contribution that the underlined noun makes to the DP it's within. We can capture this by ensuring that the denotation of \emph{one} is identical to the denotation of the underlined noun. Whereas in Deep Anaphora, what's at stake is how the anaphor and antecedent refer to things, in Surface Anaphora, the relationship between antecedent and anaphor concerns the denotations of the two terms. Surface anaphora is of interest to us because it provides a way of determining the denotation of expressions. We can use it to see if the syntax we have given to phrases correctly corresponds the way the denotations of terms inside that phrase semantically compose.

The surface anaphor we'll use is expressed prosodically. It arises when one phrase is offered as a partial recapitulation of another phrase. (\ref{ex6:deacc1}) illustrates.
\begin{exe}
	\raggedright
	\ex \label{ex6:deacc1}
	\begin{xlist}
		\exi{A:} Thuy eats natto.
		\exi{B:} <Thuy eats> SPINACH (too).
	\end{xlist}
\end{exe}
In this short conversation, B can be understood as offering a response to A's sentence that either corrects it, or adds to it, depending on whether \emph{too} is included. In those circumstances, B's utterance carries a particular prosody. The word \emph{spinach} is accented, and the pitch range of \emph{Thuy eats} is compressed. Pitch accent on \emph{spinach} signals that the DP it makes is focused. Where relevant, I will mark portions of a sentence that have the compressed pitch range indicative of this prosody with angled brackets. These strings are said to be ``de-accented.'' I will also signal the term that bears accent by putting it in all caps. If the focus marked term is different from the accented term, which happens when a multiword phrase is focused, I will mark that phrase with ``[~]_{\text{F}}.'' Phrases that contain this pairing of de-accented string and focus-marked term are a kind of surface anaphor. Their denotation must match---in the same kind of way that a surface anaphor does---the denotation of a phrase uttered elsewhere.

This can be seen by considering the similar conversation in~(\ref{ex6:deacc2}).
\begin{exe}
	\raggedright
	\ex \label{ex6:deacc2}
	\begin{xlist}
		\exi{A:} When Rodica cooks, Thuy eats natto.
		\exi{B:} <She eats> SPINACH (too).
		\exi{B$'$:} Did you know that about her?
	\end{xlist}
\end{exe}
In B$'$'s reply, the referent of \emph{her} can be either Rodica or Thuy, making this question awkwardly ambiguous. This establishes that A's remark makes both Rodica and Thuy salient enough for the pronoun to find them as a referent; either \emph{Rodica} or \emph{Thuy} can be the antecedent for \emph{her}. The response in B, however, does not have this freedom of interpretation. Here \emph{she} can only refer to Thuy. The prosody attendant with B's reply requires that its meaning correspond to the meaning of \emph{Thuy eats natto}, and a consequence of this is that \emph{she} and \emph{Thuy} must corefer.

A similar thing is demonstrated by (\ref{ex6:deacc3}).
\begin{exe}
	\raggedright
	\ex \label{ex6:deacc3}
	\begin{xlist}
		\exi{A:}[]{Thuy will eat natto.}
		\exi{B:}[*]{<Thuy will> COOK <spinach> (too).}
		\exi{B$'$:}[]{<Thuy will> COOK <natto> (too).}
	\end{xlist}
\end{exe}
The prosody in the replies by B and B$'$ in this conversation require that the meaning of the de-accented portion match the meaning in the corresponding positions of A's remark. The ungrammaticality of B's response traces back to the fact that \emph{spinach} doesn't have the same meaning that \emph{natto} has; when that is corrected, the result is grammatical, as B$'$'s response demonstrates.

Let's try to characterize how the surface anaphora that involves this prosody works. Assume that both the de-accenting and the accent that marks focus expone the anaphor. Let's call a phrase that contains both of these prosodic cues a ``focus anaphor.''  %
%
\marginnote[0pt]{The account I am describing here is roughly that of \citet{Rooth:1992}. Much of the foundation for this theory, including the relationship between focus and prosody, is in \citet{Jackendoff:1972}. The idea that these prosodies expone anaphors can be found in \citet{Ladd:1980}. \Citet{Selkirk:1984} is a good introduction to the complex relationship between the prosody of focus marking and the item that is focused.}%
%
The phrase that a focus anaphor must match will be called a ``focus antecedent.'' In (\ref{ex6:deacc1})-(\ref{ex6:deacc3}), it is the entire TP that B utters that is a focus anaphor, and A's remark that is its focus antecedent. Other phrasal types can be focus anaphors and focus antecedents as well, as~(\ref{ex6:deacc4}) indicates.
\begin{exe}
	\raggedright
	\ex \label{ex6:deacc4}
	{}[\DP Thuy's dissertation] will be finished in July\\ and [\DP RODICA'S <dissertation>] arrives soon thereafter.
\end{exe}
Here, the DP \emph{Rodica's dissertation} has the prosody of a focus anaphor, and it finds its focus antecedent in the DP that starts the sentence: \emph{Thuy's dissertation}. I'll assume that any phrasal type can be a focus anaphor or a focus antecedent. A description of the relationship between anaphor and antecedent that the examples we've seen illustrate is~(\ref{ex6:fanaphora}).
\begin{exe}
	\raggedright
	\ex \label{ex6:fanaphora} Let XP be a focus anaphor, and $\alpha$ be the focus-marked term within XP. A \underline{focus value} of XP (=\denotes{XP}^{F}) is created by semantically interpreting the XP that results from replacing $\alpha$ with a \underline{focus alternative} of $\alpha$ (=\textsc{fa}($\alpha$)). $\beta$ is a \textsc{fa}($\alpha$) if it is the same syntactic type as~$\alpha$, and contextually salient.\\[4pt]
	If XP is a focus anaphor, then there must be a YP and a \denotes{XP}^F such that \denotes{YP}=\denotes{XP}^F.
\end{exe}
This gives an account of (\ref{ex6:deacc1})-(\ref{ex6:deacc3}), as indicated below.
\begin{exe}
	\raggedright
	\ex \label{ex6:d1} Because \emph{natto} is a \textsc{fa}(\emph{spinach}), \denotes{Thuy eats SPINACH}^F=\denotes{Thuy eats natto}, and therefore (\ref{ex6:deacc1}) works.
	\ex \label{ex6:d2} Because \emph{natto} is a \textsc{fa}(\emph{spinach}), \denotes{she eats SPINACH}^F=\denotes{Thuy eats natto} when \emph{she} and \emph{Thuy} corefer, and therefore (\ref{ex6:deacc2}) works.
	\ex \label{ex6:d3} Because \emph{eat} is a \textsc{fa}(\emph{cook}), \denotes{Thuy will COOK natto}^F=\denotes{Thuy will eat natto}, and therefore (\ref{ex6:deacc3}) works.
\end{exe}

In the cases we've reviewed, a focus anaphor is indicated by giving a phrase in a sentence a certain prosody, a prosody that combines pitch accent with de-accenting. There is another way that a focus anaphor can be indicated, that is slightly different. To see this, consider first the cases of focus anaphora in~(\ref{ex6:ellip1}), which deploys the combination of de-accenting and pitch accent.
\begin{exe}
	\raggedright
	\ex \label{ex6:ellip1}
	\begin{xlist}
		\ex Rodica eats natto when THUY <eats natto>.
		\ex Rodica eats natto when <she eats> RICE.
	\end{xlist}
\end{exe}
The TP in the \emph{when} clause is a focus anaphor, and meets the conditions of (\ref{ex6:fanaphora}), as (\ref{ex6:whenc}) shows.
\pagebreak
\begin{exe}
	\raggedright
	\ex \label{ex6:whenc}
	\begin{xlist}
		\ex Because Rodica is a \textsc{fa}(Thuy), \denotes{THUY eats natto}^F = \denotes{Rodica eats natto}, and (\ref{ex6:ellip1}a) is grammatical.
		\ex Because natto is a \textsc{fa}(rice), \denotes{she eats RICE}^F = \denotes{Thuy eats natto}, when \emph{she} corefers with \emph{Thuy}, and therefore (\ref{ex6:ellip1}b) is grammatical.
	\end{xlist}
\end{exe}

There is an alternative pronunciation of (\ref{ex6:ellip1}a)---it's in (\ref{ex6:ellip2}a)---but not of~(\ref{ex6:ellip1}b).
\begin{exe}
	\raggedright
	\ex \label{ex6:ellip2}
	\begin{xlist}
		\ex[]{Rodica eats natto when THUY does.}
		\ex[*]{Rodica eats natto when RICE.}
	\end{xlist}
\end{exe}
What's happened in (\ref{ex6:ellip2}a) is that the de-accented words have not been pronounced. That isn't possible for (\ref{ex6:ellip1}), however. A common conclusion about what makes these cases different is that the de-accented words form a phrase in (\ref{ex6:ellip1}a), but not in~(\ref{ex6:ellip1}b). We should engineer the process that allows de-accented words to go unspoken so that it applies only to those words that form a constituent.

A popular way of characterizing this gives to (\ref{ex6:ellip2}a) nearly the same syntax as (\ref{ex6:ellip1}a), but allows for a way of preventing the words in the VP from showing up in the string that is associated with that syntax. The difference in the syntax between (\ref{ex6:ellip1}a) and (\ref{ex6:ellip2}a) is connected to the requirement that the tense feature in \obar{T} be exponed. If we understand the requirement that the \textsc{agree} relation which tense has with the morphology on a verb be satisfied only if the verb with that morphology is pronounced, then the failure to pronounce the verb in (\ref{ex6:ellip2}) will result in ungrammaticality. Instead, then, a sentence that has both tense and the process that unpronounces the verb must also provide a way for the tense feature to be exponed. This is achieved in (\ref{ex6:ellip2}a) by including the auxiliary verb \emph{do} in the syntax. The syntactic representation for (\ref{ex6:ellip2}a) approximates~(\ref{ex6:ellip2t}).
\begin{exe}
	\raggedright
	\ex \label{ex6:ellip2t}
	\small
\begin{forest}
[TP
	[DP [Thuy, roof]]
	[\ibar{T} 
		[\obar{T} 
		[\obar{V} [do, name=binder]]
		[\obar{T} [pres]]
		]
		[VP
		[\ , name=trace]
		[VP [\ibar{V}
			[\obar{V} [\sout{eat}]]
			[DP [\sout{natto}, roof]]
		]]
	]
]]
\draw[->,>=stealth', thin, gray] (trace) to[out=south,in=south east,looseness=1.3] node[pos=0.75,xshift=-10mm]{} (binder);
\end{forest}
\end{exe}
I've used \sout{strikeouts} to indicate the words that haven't been pronounced. These words are said to be ``elided,'' and the process that achieves this in (\ref{ex6:ellip2t}) is called ``VP Ellipsis.'' %
%
\marginnote[-40pt]{There are several ways of modeling ellipsis. We might imagine, for instance, that there is a designated anaphor---like a personal pronoun---for VPs, and that its phonology is silence. \citet{Hardt:1993, Hardt:1999} lays out an account of this type. But we shall see evidence later that supports the view that an ellipsis contains syntactic material that is unspoken.}%
%
This is an unfortunate name for the process, as it can elide the words in phrases that aren't VPs, as (\ref{ex6:apdel}) demonstrates, and being a VP is not sufficient to allow the process, as (\ref{ex6:causedel}) shows.
\begin{exe}
	\raggedright
	\ex[]{\label{ex6:apdel} Thuy must be [\AP happy] and RODICA must be [\AP \sout{happy}] too.}
	\ex[*]{\label{ex6:causedel} She made Thuy [\VP work all night] because he made RODICA [\VP \sout{work} \sout{all} \sout{night}].}
\end{exe}
I'll persist in using this name, as it's what you're most likely to find in the literature, but you should remember that the expression ``VP Ellipsis'' doesn't have a compositional meaning.

If the presence of \emph{does} adds no meaning to these sentences, then the conditions on focus anaphora are still met in (\ref{ex6:ellip2}). We can describe VP Ellipsis with~(\ref{ex6:vpedef}).
\begin{exe}
	\raggedright
	\ex \label{ex6:vpedef} VP Ellipsis\\
	If $\alpha$ is the complement to \obar{T} or an auxiliary verb, then a word in $\alpha$ can go unspoken iff:
	\begin{xlist}
		\ex every other word in $\alpha$ is unspoken,
		\ex the word is not part of something that is focus marked, and
		\ex $\alpha$ is dominated by a focus anaphor.
	\end{xlist}
\end{exe}
I've framed the conditions under which the words in a phrase can elide so that the cases in (\ref{ex6:apdel}) and~(\ref{ex6:causedel}) are captured.

There is an interesting feature of the condition on focus anaphora in (\ref{ex6:fanaphora}) that distinguishes instances where VP ellipsis has applied from those where it hasn't. To see this difference, consider first~(\ref{ex6:lakoff}).%
%
\marginnote[0pt]{This contrast was first discussed in Lakoff?}%
%
\begin{exe}
	\raggedright
	\ex \label{ex6:lakoff}
	Thuy is an accomplished syntactician and RODICA <is smart> too.
\end{exe}
The conditions in (\ref{ex6:fanaphora}) require of the focus anaphor in the second conjunct that its meaning match that of the first conjunct.
\begin{exe}
	\raggedright
	\ex \label{ex6:lakofff}
	Because \emph{Thuy} is an \textsc{fa}(\emph{Rodica}), (\ref{ex6:fanaphora}) requires that\\ \denotes{RODICA is smart}^F = \denotes{Thuy is an accomplished syntactician}
\end{exe}
Because \emph{Thuy is smart} (which is what \emph{RODICA is smart} means when the \textsc{fa}(Rodica) is Thuy) does not have the same denotation that \emph{Thuy is an accomplished syntactician}, it is not clear why the focus anaphor is allowed in (\ref{ex6:lakoff}). To account for this, Chris Tancredi suggested that we think of (\ref{ex6:fanaphora}) as a condition which can be used by a speaker to invite inferences.\cite{Tancredi:1992} By making the second conjunct a focus anaphor whose antecedent is the first conjunct, the speaker leads the hearer to infer that being an accomplished syntactician is equivalent to being smart. The proposal, then, is that (\ref{ex6:fanaphora}) is not a condition that distinguishes the grammatical from the ungrammatical sentences by virtue of their meaning. Instead, it is a condition that shapes the interpretation of the message being conveyed. We won't develop a framework in which conditions shape meanings, rather than factor sentences into the grammatical and ungrammatical bins. %
%
\marginnote[-20pt]{See \citet[]{Fox:2000} for a sketch. A slightly different way of expressing the difference between ellipsis and de-accenting is offered in \citet{Merchant:2001}.}% 
It will be sufficient for us to have an informal understanding of this.

Interestingly, this function of (\ref{ex6:fanaphora}) collapses when the material in the focus anaphor relevant for drawing the inference is elided. An alternative pronunciation of (\ref{ex6:lakoff}) cannot involve VP ellipsis.
\begin{exe}
	\raggedright
	\ex[*]{\label{ex6:lakoffvpe} Thuy is an accomplished syntactician and RODICA is \sout{smart} too.}
	\sn \emph{compare:}
	\sn[]{Thuy is an accomplished syntactician and RODICA is \sout{an} \sout{accomplished} \sout{syntactician} too}
\end{exe}
Using cases in which the focus anaphor employs ellipsis has a certain utility for us then. It allows us to see what the denotation of the antecedent of the focus anaphor is without the haze that arises when inferences of the sort demonstrated in (\ref{ex6:lakoff}) are tolerated.

Focus anaphora, especially in its ellipsis guise, gives us a tool to investigate the denotations of parts of sentences. In this chapter, we will use that tool to investigate the meanings of sentences that involve movement. We start by investigating a particular way in which focus anaphora interacts with the bound variable interpretation that we've seen pronouns can receive.

\subsection{Sloppy Anaphora}
\label{sec:slopana}

Recall that one kind of deep anaphora arises when a personal pronoun gets an interpretation that depends on the value given to some other DP in the same sentence. This is to be contrasted with the interpretation a personal pronoun can get that depends on a context making some individual salient enough to be the referent of that pronoun. (\ref{ex6:corebnd}) illustrates.
\begin{exe}
	\raggedright
	\ex \label{ex6:corebnd}
	As they lunched in the Wellesley cafeteria, Carolyn told Thuy that no student understood her paper.
\end{exe}
To appreciate the range of meanings that (\ref{ex6:corebnd}) gets, it is helpful to know a few things about Carolyn and Thuy. They are both women. They are both published linguists. Carolyn teaches at Wellesley, a college in Cambridge, Massachusetts, and Thuy teaches at Nguyen Tat Thanh University in Ho Chi Minh City. The pronoun \emph{her} in (\ref{ex6:corebnd}) can be understood in one of three ways. It can refer to the same individual that \emph{Carolyn} does: it can corefer with \emph{Carolyn}. It can also refer to the same individual that \emph{Thuy} does: it can corefer with \emph{Thuy}. Finally, it can be understood as a variable bound by \emph{no student}. On the first interpretation, (\ref{ex6:corebnd}) entails that Carolyn's report was that no student understood Carolyn's paper. And on the second, (\ref{ex6:corebnd}) entails that Carolyn's report was that no student understood Thuy's paper. On the final interpretation, (\ref{ex6:corebnd}) says that Carolyn's report was that no student understood her own paper.

Consider now (\ref{ex6:corebnd1}), which includes a focus anaphor. (I give the focus anaphor in both its ellipsis and non-ellipsis versions.)
\begin{exe}
	\raggedright
	\ex \label{ex6:corebnd1}
	\begin{xlist}
	\ex As they lunched at the Wellesley cafeteria, Carolyn told Thuy that no freshman understood her paper and that <no> SOPHOMORE <understood her paper> either.
	\ex As they lunched at the Wellesley cafeteria, Carolyn told Thuy that no freshman understood her paper and that <no> SOPHOMORE did either.
	\end{xlist}
\end{exe}
There are two instances of \emph{her} in this sentence, but because the second is de-accented, it is part of a focus anaphor whose meaning is going to be determined by the clause that has the first instance of \emph{her}. What we should expect to find, then, is that the interpretation of the two instances of \emph{her} should match. This is almost true. It's completely true when the interpretations of the pronouns involve coreference. When the first \emph{her} corefers with \emph{Carolyn} then so must the second \emph{her}; and when the first \emph{her} corefers with \emph{Thuy}, then so must the second \emph{her}. (\ref{ex6:corebnd1}) can have one of the two interpretations paraphrased by (\ref{ex6:corebnd1I}a,b), but neither of the interpretations paraphrased by~(\ref{ex6:corebnd1I}c,d).
\begin{exe}
	\raggedright
	\ex \label{ex6:corebnd1I} As they lunched at the Wellesley cafeteria, Carolyn told Thuy that\ldots
	\begin{xlist}
		\ex[]{no freshman understood her (=Carolyn's) paper and\\ no sophomore understood her (=Carolyn's) paper either.}
		\ex[]{no freshman understood her (=Thuy's) paper and\\ no sophomore understood her (=Thuy's) paper either.}
		\ex[*]{no freshman understood her (=Carolyn's) paper and\\ no sophomore understood her (=Thuy's) paper either.}
		\ex[*]{no freshman understood her (=Thuy's) paper and\\ no sophomore understood her (=Carolyn's) paper either.}
	\end{xlist}
\end{exe}
We can see why this is expected from the conditions on focus anaphora in (\ref{ex6:fanaphora}) by considering how they apply to the last TP---the focus anaphor---in~(\ref{ex6:corebnd1}).
\begin{exe}
	\raggedright
	\ex \label{ex6:corebnd1s}
	Because freshman is a \textsc{fa}(sophomore), \denotes{no SOPHOMORE understands her paper}^F = \denotes{no freshman understands her paper} when \emph{her} refers to the same individual in both.
\end{exe}

What's unexpected in (\ref{ex6:corebnd1}) is how the interpretation of the pronouns arises when they are understood as bound variables. We might expect that when the first \emph{her} is interpreted as a variable bound by \emph{no freshman}, so must the second \emph{her}. In fact, this interpretation is not possible. It isn't possible to interpret (\ref{ex6:corebnd1}) in the way that (\ref{ex6:ccombnd1}) describes.
\begin{exe}
	\raggedright
	\ex \label{ex6:ccombnd1} As they lunched in the Wellesley cafeteria, Carolyn told Thuy that\ldots
	\sn no freshman, $x$, understood $x$'s paper and\\
	no sophomore understood $x$'s paper either.
\end{exe}
It's useful when judging the meanings of complicated sentences to frame those sentences with a context that biases towards the target meaning. A context that favors giving (\ref{ex6:corebnd1}) the interpretation described by (\ref{ex6:ccombnd1}) is provided in~(\ref{ex6:context1}).
\begin{exe}
	\raggedright
	\sn 
	In Carolyn's last class, her students were to report their interpretation of each other's final papers. She noticed an interesting pattern. If none of the sophomores understood a freshman's paper, then even the freshman author of that paper didn't either.
	\ex[\#]{\label{ex6:context1} As they lunched that afternoon, Carolyn told Thuy that no freshman understood her paper and no SOPHOMORE did either.}
\end{exe}
%
\marginnote[-20pt]{The symbol ``\#'' indicates that the sentence has an anomalous semantic interpretation.}%
%
If Carolyn has correctly reported the situation in her class to Thuy, then (\ref{ex6:context1}) is not the right way of describing that. Unexpectedly, (\ref{ex6:corebnd1}) cannot have the meaning in~(\ref{ex6:ccombnd1}).

This turns out to be due to a condition on the bound variable interpretation that limits the syntactic positions that the binder may have relative to the variable. Let's take a moment to learn about this condition.

The method we adopted for providing the semantics of a bound variable interpretation crucially involved the meaning of the determiner responsible for introducing the ``family of sentences'' part of the meaning. We saw how that worked for pronouns that get bound by DPs headed by \emph{every} earlier, but it's true in a parallel way for pronouns that get bound by DPs headed by other quantificational determiners as well. For instance, in (\ref{ex6:no1}), it is crediting \emph{no} with a denotation that involves variable binding that is responsible for giving \emph{her} its bound variable interpretation.
\begin{exe}
	\raggedright
	\ex \label{ex6:no1} No woman reads her first paper
\end{exe}
The denotation of the quantificational determiner \emph{no} (given in (\ref{ex6:nodenote})) takes the NP it combines with, and the \ibar{v} that follows, and creates a meaning from the denotations of these two constituents, as shown in~(\ref{ex6:no2}).
\medskip
\begin{fullwidth}
\begin{exe}
	\raggedright
	\ex \label{ex6:nodenote} \denotes{no} = $\lambda P\ \lambda Q\ \lambda \epsilon\ \neg \exists x\ P(x) \wedge\ Q(x)(\epsilon)$
	\ex \label{ex6:no2}
	\hspace*{-22pt}
	\small
\begin{forest}
[$\lambda \epsilon \neg \exists x$ \textsc{woman}($x$) $\wedge$ \textsc{agent}($x$)($\epsilon$) $\wedge$ \textsc{theme}(\denotes{her first paper})($\epsilon$) $\wedge$ \textsc{read}($\epsilon$)\\ vP
	[$\lambda Q \lambda \epsilon \neg \exists x$ \textsc{woman}($x$) $\wedge\ Q(x)(\epsilon)$ \\ DP 
	[\ibar{D}
		[$\lambda P\ \lambda Q \lambda \epsilon \neg \exists x. P(x) \wedge\ Q(x)(\epsilon)$\\ \obar{D} [no]]
		[$\lambda x$ \textsc{woman}($x$) \\ NP [woman, roof] ]
	]
	]
	[$\lambda x \lambda \epsilon$ \textsc{agent}($x$)($\epsilon$) $\wedge$ \textsc{theme}(\denotes{her first paper})($\epsilon$) $\wedge$ \textsc{read}($\epsilon$) \\ \ibar{v} [read her first paper, roof]]
]
\end{forest}
\end{exe}
\end{fullwidth}
\medskip
Informally, \emph{no} creates a meaning from these two parts which describes reading-eventualities in which the things that are the agents of that eventuality are not things which are a woman. The quantificational determiner \emph{every} has a different meaning, as we've seen, but it binds a variable in a parallel way. The quantificational determiner \emph{some} also has a meaning that can be expressed as a family of sentences, and it too can be given a denotation that achieves this meaning through variable binding. I've given examples of \emph{every} and \emph{some} below, so that all three cases---the situation with \emph{no}, \emph{every}, and \emph{some}---can be compared.
\begin{exe}
	\raggedright
	\ex \label{ex6:dano} No woman reads her paper.
	\begin{xlist}
		\sn For all women under consideration, 1, \ldots,n, it's the case that:
		\begin{xlist}
			\ex woman \#1 doesn't read woman \#1's paper, and
			\ex woman \#2 doesn't read woman \#2's paper, and
			\ex woman \#3 doesn't read woman \#3's paper, and
			\sn \hspace*{1in}\vdots
		\end{xlist}
	\end{xlist}
	\ex \label{ex6:daevery} Every woman reads her paper.
	\begin{xlist}
		\sn For all women under consideration, 1,\ldots,n, it's the case that
		\begin{xlist}
			\ex woman \#1 reads woman \#1's paper, and
			\ex woman \#2 reads woman \#2's paper, and
			\ex woman \#3 reads woman \#3's paper, and
			\sn \hspace*{1in}\vdots
			\ex woman \#n reads woman \#n's paper.
		\end{xlist}
		\end{xlist}
		\ex \label{ex6:dasome} Some woman reads her paper.
		\begin{xlist}
			\sn For some woman under consideration, 1,\ldots,n, it's the case that
		\begin{xlist}
			\ex woman \#1 reads woman \#1's paper, or
			\ex woman \#2 reads woman \#2's paper, or
			\ex woman \#3 reads woman \#3's paper, or
			\sn \hspace*{1in}\vdots
			\ex woman \#n reads woman \#n's paper.
		\end{xlist}
	\end{xlist}
\end{exe}
The meaning for \emph{no} is captured by giving \emph{no} the denotation in (\ref{ex6:nodenote}), as (\ref{ex6:no2}) shows. The meaning for \emph{every} can be captured by giving \emph{every} the denotation repeated in (\ref{ex6:everyd}). And giving \emph{some} the denotation in (\ref{ex6:denose}b) provides a way of modeling its meaning too. The compositions in (\ref{ex6:denoset}) show how.
\medskip
\begin{fullwidth}
\begin{exe}
	\raggedright
	\ex \label{ex6:denose}
	\begin{xlist}
		\ex \label{ex6:everyd} \denotes{every} = $\lambda P \lambda Q \lambda \epsilon \forall x [P(x) \rightarrow Q(x)(\epsilon)]$
		\ex \denotes{some} = $\lambda P \lambda Q \lambda \epsilon \exists x [P(x) \wedge Q(x)(\epsilon)]$
	\end{xlist}
	\ex \label{ex6:denoset}
	\hspace*{-22pt}
	\small
\begin{forest}
[$\lambda \epsilon \forall x$ {[}\textsc{woman}($x$) $\rightarrow$ {[}\textsc{agent}($x$)($\epsilon$) $\wedge$ \textsc{theme}(\denotes{her first paper})($\epsilon$) $\wedge$ \textsc{read}($\epsilon$){]]}\\ vP
	[$\lambda Q \lambda \epsilon \forall x$ {[}\textsc{woman}(x) $\rightarrow Q(x)(\epsilon){]}$\\DP 
	[\ibar{D}
		[$\lambda P\ \lambda Q \lambda \epsilon \forall x {[}P(x) \rightarrow Q(x)(\epsilon){]}$\\\obar{D} [every]]
		[$\lambda x$ \textsc{woman}($x$)\\NP [woman, roof] ]
	]]
	[$\lambda x \lambda \epsilon$ \textsc{agent}($x$)($\epsilon$) $\wedge$ \textsc{theme}(\denotes{her first paper})($\epsilon$) $\wedge$ \textsc{read}($\epsilon$)\\ \ibar{v} [read her first paper, roof]]
]
\end{forest}
\vspace*{20pt}

\hspace*{-22pt}
\begin{forest}
[$\lambda \epsilon \exists x$ {[}\textsc{woman}($x$) $\wedge$ {[}\textsc{agent}($x$)($\epsilon$) $\wedge$ \textsc{theme}(\denotes{her first paper})($\epsilon$) $\wedge$ \textsc{read}($\epsilon$){]]}\\vP
	[$\lambda Q \lambda \epsilon \exists x$ {[}\textsc{woman}($x$) $\wedge\ Q(x)(\epsilon)${]}\\DP 
		[\ibar{D}
			[$\lambda P\ \lambda Q \lambda \epsilon \exists x {[}P(x) \wedge Q(x)(\epsilon){]}$\\ \obar{D} [some]]
			[$\lambda x$ \textsc{woman}($x$)\\ NP [woman, roof]]
		]
	]
	[$\lambda x \lambda \epsilon$ \textsc{agent}($x$)($\epsilon$) $\wedge$ \textsc{theme}(\denotes{her first paper})($\epsilon$) $\wedge$ \textsc{read}($\epsilon$) \\ \ibar{v} [read her first paper, roof]]
]
\end{forest}
\end{exe}
\end{fullwidth}
\medskip
In each case, it is the fact that the second argument of the determiner's denotation is a \ibar{v} whose meaning binds a variable in the \textsc{agent} position that is responsible for causing the determiner to end up binding a variable there.

To get the bound variable interpretation of the pronoun \emph{her} in these cases involves finding a way of letting those pronouns also be variables bound in the denotation of the \ibar{v}. As we saw earlier, the bound variable reading for \emph{her} can be achieved for (\ref{ex6:da2rd}) if we can produce the meaning for the \ibar{v} indicated in (\ref{ex6:da2rd2}). And the same change to the denotations of the \ibar{v}s in (\ref{ex6:denoset}) will bring the bound variable interpretation to the pronoun, as (\ref{ex6:denoseta}), and~(\ref{ex6:denosetb}) demonstrate.
\medskip
\begin{fullwidth}
\begin{exe}
	\raggedright
	\exr{ex6:da2rd2}
	\hspace*{-22pt}
	\small
\begin{forest}
[$\lambda \epsilon \neg \exists x$ {[}\textsc{woman}($x$) $\wedge$ \textsc{agent}($x$)($\epsilon$) $\wedge$ \textsc{theme}(\denotes{$x$'s first paper})($\epsilon$) $\wedge$ \textsc{like}($\epsilon$){]}\\vP
	[$\lambda Q \lambda \epsilon \neg \exists x {[}$\textsc{woman}($x$) $\wedge\ Q(x)(\epsilon){]}$\\DP [\ibar{D} 
		[$\lambda P\ \lambda Q \lambda \epsilon \neg \exists x {[}P(x) \wedge\ Q(x)(\epsilon){]}$\\\obar{D} [no] ]
		[$\lambda x$ \textsc{woman}($x$) \\NP [woman, roof]]
	]]
	[$\lambda x \lambda \epsilon$ \textsc{agent}($x$)($\epsilon$) $\wedge$ \textsc{theme}(\denotes{$x$'s first paper})($\epsilon$) $\wedge$ \textsc{like}($\epsilon$) \\\ibar{v} [like her first paper, roof] ]
]
\end{forest}
	\ex \label{ex6:denoseta}
	\hspace*{-22pt}
	\small
\begin{forest}
[$\lambda \epsilon \forall x$ {[}\textsc{woman}($x$) $\rightarrow$ {[}\textsc{agent}($x$)($\epsilon$) $\wedge$ \textsc{theme}(\denotes{$x$'s first paper})($\epsilon$) $\wedge$ \textsc{read}($\epsilon$){]]}\\vP
	[$\lambda Q \lambda \epsilon \forall x$ {[}\textsc{woman}(x) $\rightarrow Q(x)(\epsilon){]}$\\DP [\ibar{D}
		[$\lambda P\ \lambda Q \lambda \epsilon \forall x {[}P(x) \rightarrow Q(x)(\epsilon){]}$\\ \obar{D} [every] ]
		[$\lambda x$ \textsc{woman}($x$)\\NP [woman, roof]]
	]]
	[$\lambda x \lambda \epsilon$ \textsc{agent}($x$)($\epsilon$) $\wedge$ \textsc{theme}(\denotes{$x$'s first paper})($\epsilon$) $\wedge$ \textsc{read}($\epsilon$)\\ \ibar{v} [read her first paper, roof] ]
]
\end{forest}
\ex \label{ex6:denosetb}
	\hspace*{-22pt}
	\small
\begin{forest}
[$\lambda \epsilon \exists x$ {[}\textsc{woman}($x$) $\wedge$ {[}\textsc{agent}($x$)($\epsilon$) $\wedge$ \textsc{theme}(\denotes{$x$'s first paper})($\epsilon$) $\wedge$ \textsc{read}($\epsilon$){]]}\\vP
	[$\lambda Q \lambda \epsilon \exists x$ {[}\textsc{woman}($x$) $\wedge\ Q(x)(\epsilon){]}$\\DP [\ibar{D}
		[$\lambda P\ \lambda Q \lambda \epsilon \exists x {[}P(x) \wedge Q(x)(\epsilon){]}$\\ \obar{D} [some] ]
		[$\lambda x$ \textsc{woman}($x$) \\NP [woman, roof]]
	]]
	[$\lambda x \lambda \epsilon$ \textsc{agent}($x$)($\epsilon$) $\wedge$ \textsc{theme}(\denotes{$x$'s first paper})($\epsilon$) $\wedge$ \textsc{read}($\epsilon$) \\ \ibar{v} [read her first paper, roof] ]
]
\end{forest}
\end{exe}
\end{fullwidth}
\medskip

What these examples indicate is that for a personal pronoun to be interpreted as a variable bound by a quantificational DP, it must be part of the constituent that the quantifier takes as its second argument. Framing this in syntactic terms, this means that for a quantificational DP, $\alpha$, to bind a variable, $\beta$, $\beta$ must be part of the phrase that is the sister to $\alpha$. This relation is usually expressed in terms of c-command, as in~(\ref{ex6:ccombind}).
\begin{exe}
	\raggedright
	\ex \label{ex6:ccombind}
	For $\beta$ to be a variable bound by $\alpha$, $\alpha$ must c-command~$\beta$.
	\ex $\alpha$ c-commands $\beta$ iff:
	\begin{xlist}
		\ex every phrase that (non-reflexively) dominates $\alpha$, dominates $\beta$ too, and
		\ex $\alpha$ does not (reflexively) dominate $\beta$.
	\end{xlist}
\end{exe}
The law in (\ref{ex6:ccombind}) is a prediction made by the semantics we have adopted for expressing the meanings of the quantifiers \emph{no}, \emph{some}, and \emph{every}. To a certain extent, (\ref{ex6:ccombind}) is confirmed. The bound variable interpretation of \emph{her} in (\ref{ex6:predconf}) is, as (\ref{ex6:ccombind}) predicts, unavailable.
\begin{exe}
	\raggedright
	\ex \label{ex6:predconf}
	\begin{xlist}
		\ex That no woman [left] bothers her friend.
		\ex People that no woman [has talked to] like her friends.
		\ex Her friends [like] no woman.
	\end{xlist}
\end{exe}
The quantificational DP \emph{no woman} doesn't c-command \emph{her} in any of these cases. The sister to \emph{no woman} is marked with square brackets, and \emph{her} is not found within this constituent. There are apparent counterexamples to (\ref{ex6:ccombind}) too, however, and so there are reasons to remain skeptical that we have found the correct method of expressing the meaning of quantificational DPs. We will examine these apparent counterexamples in this chapter. One method of addressing them involves the new theory of movement that will be developed here.

But we can see that to the extent that (\ref{ex6:ccombind}) is true, it provides an account for the fact that an interpretation for both \emph{her}s as a variable bound to \emph{no freshman} in~(\ref{ex6:corebnd1}) isn't possible. The ellipsis version of (\ref{ex6:corebnd1}) is repeated here.
\begin{exe}
	\raggedright
	\exr{ex6:corebnd1}
	As they lunched at the Wellesley cafeteria, Carolyn told Thuy that no freshman understood her paper and that <no> SOPHOMORE did either.
\end{exe}
The DP \emph{no freshman} doesn't c-command the second \emph{her}, as the rudimentary parse in (\ref{ex6:pcorebnd1}) shows.
\medskip
\begin{fullwidth}
\begin{exe}
	\raggedright
	\ex \label{ex6:pcorebnd1}
	\small
\begin{forest}
[\ibar{V}
	[\obar{V} [tell]]
	[CP 
		[CP 
		[\ibar{C}
			[\obar{C} [that] ]
			[TP 
				[DP [no freshman, roof]]
				[\ibar{T} [understood her paper, roof] ]
			]
		]
		]
		[and ]
		[CP 
		[\ibar{C}
			[\obar{C} [that] ]
			[TP [no sophomore understood her paper, roof]]
		]
		]
	]
]
\end{forest}
\end{exe}
\end{fullwidth}
\medskip
If we consider examples which minimally differ from (\ref{ex6:corebnd1}), but put the second pronoun in a position that is c-commanded by \emph{no freshman}, the bound variable interpretation emerges. Consider, for example, the interpretation of~(\ref{ex6:corebnd2}).
\begin{exe}
	\raggedright
	\ex \label{ex6:corebnd2}
	Carolyn told Thuy that no freshman praised her paper if <no> SOPHOMORE did.
\end{exe}
A possible meaning for (\ref{ex6:corebnd2}) gives to both \emph{her}s an interpretation in which they are bound by \emph{no freshman}; that meaning can be paraphrased with~(\ref{ex6:corebnd2para}).
\begin{exe}
	\raggedright
	\ex \label{ex6:corebnd2para}
	Carolyn told Thuy that there is no $x$, $x$ a freshman, who praised $x$'s paper if no sophomore praised $x$'s paper.
\end{exe}
The context in (\ref{ex6:corebnd2con}) helps make that meaning salient.
\begin{exe}
	\raggedright
	\sn
	In Carolyn's last class, the students gave short assessments of each others' final papers. Carolyn noticed a pattern. If no sophomore praised a freshman's paper, then that freshman didn't praise her own paper either. At lunch that afternoon,\ldots
	\ex \label{ex6:corebnd2con} Carolyn told Thuy that no freshman praised her paper if no SOPHOMORE did.
\end{exe}
What makes (\ref{ex6:corebnd2}) relevantly different from (\ref{ex6:corebnd1}), is that there is a parse---it's sketched in (\ref{ex6:corebnd2par})---that causes the second pronoun to be c-commanded by \emph{no freshman}.
\begin{exe}
	\raggedright
	\ex \label{ex6:corebnd2par}
	\small
\begin{forest}
[vP
	[DP [no freshman, roof]]
	[\ibar{v} 
		[\ibar{v} [praised her paper, roof]]
		[PP [if no sophomore praised her paper, roof]]
	]
]
\end{forest}
\end{exe}

So far, then, the condition on focus anaphora correctly captures the meanings of the pronouns in (\ref{ex6:corebnd1}), repeated here.
\begin{exe}
	\raggedright
	\exr{ex6:corebnd1}
	\begin{xlist}
	\ex As they lunched at the Wellesley cafeteria, Carolyn told Thuy that no freshman understood her paper and that <no> SOPHOMORE <understood her paper> either.
	\ex As they lunched at the Wellesley cafeteria, Carolyn told Thuy that no freshman understood her paper and that <no> SOPHOMORE did either.
	\end{xlist}
\end{exe}
It accounts for the fact that the second \textit{her} is coreferent with whatever the first \textit{her} is coreferent with. And it predicts, with the aid of the c-command condition on bound variables, that when the first her is a variable bound to \emph{no freshman}, the second \textit{her} cannot similarly be bound to \emph{no freshman}. Interestingly, though, when the first \textit{her} is bound by \emph{no freshman}, there is a possible interpretation for (\ref{ex6:corebnd1}). When the first pronoun \textit{her} \textit{her} to \emph{no freshman}, the second \textit{her} can be bound to \emph{no sophomore}. (\ref{ex6:corebnd1}) can have the interpretation expressed in~(\ref{ex6:borebnd1slop}).
\begin{exe}
	\raggedright
	\ex \label{ex6:borebnd1slop}
	As they lunched at the Wellesley cafeteria, Carolyn told Thuy that no $x$, $x$ a freshman, understood $x$'s paper and that no $y$, $y$ a sophomore, understood $y$'s~paper.
\end{exe}
This meaning is brought out in the context provided below.
\begin{exe}
	\raggedright
	\sn In Carolyn's last class, the students were responsible with reporting their interpretation of each of the final papers. Carolyn noticed a pattern. No one seemed to understand their own paper, though they understood everyone else's. Later that afternoon, \ldots
	\ex \label{ex6:slopcon} as they lunched at the Wellesley cafeteria, Carolyn told Thuy that no freshman understood her paper, and <no> SOPHOMORE did either.
\end{exe}
John Ross\cite{Ross:1967} called this a ``sloppy interpretation'' of the pronoun. It too follows from the condition on focus anaphora in~(\ref{ex6:fanaphora}).

To see how, consider the denotations of the first and second conjuncts of~(\ref{ex6:corebnd1}).
\begin{exe}
	\raggedright
	\ex \label{ex6:denslop}
	\begin{xlist}
		\ex \denotes{no freshman understood her paper} =\\
		$\lambda \epsilon\ \neg \exists y$ [\textsc{freshman}($y$) $\wedge$ \textsc{agent}($y$)($\epsilon$) $\wedge$ \textsc{theme}(\denotes{$y$'s paper})($\epsilon$) $\wedge$ \textsc{understand}($\epsilon$)]
		\ex \denotes{no sophomore understood her paper} =\\
		$\lambda \epsilon\ \neg \exists y$ [\textsc{sophomore}($y$) $\wedge$ \textsc{agent}($y$)($\epsilon$) $\wedge$ \textsc{theme}(\denotes{$y$'s paper})($\epsilon$) $\wedge$ \textsc{understand}($\epsilon$)]
	\end{xlist}
\end{exe}
Because \emph{freshman} is a focus alternative to \emph{sophomore}, the focus value of (\ref{ex6:denslop}b) can be:
\begin{exe}
	\raggedright
	\ex \label{ex6:denslop1}
	\denotes{no SOPHOMORE understood her paper}^F =\\
			$\lambda \epsilon\ \neg \exists y$ [\textsc{freshman}($y$) $\wedge$ \textsc{agent}($y$)($\epsilon$) $\wedge$ \textsc{theme}(\denotes{$y$'s paper})($\epsilon$) $\wedge$ \textsc{understand}($\epsilon$)]
\end{exe}
Since (\ref{ex6:denslop1}) is the same as (\ref{ex6:denslop}a), the conditions on focus anaphora are met. And since \emph{her} in each conjunct is c-commanded by the DP that binds it, the c-command requirement on the bound variable interpretation is met as well.

Sloppy anaphora plays a role in the new theory of movement developed here.

\subsection{Antecedent Contained Anaphora} 
\label{sec:acd}

We now have enough background to confront the problems that suggest our present theory of movement is wrong. The first that we will examine is in (\ref{ex6:acda}).
\begin{exe}
	\raggedright
	\ex \label{ex6:acda}
	\begin{xlist}
		\ex Thuy will visit every student who [RODICA <will visit>].
		\ex Thuy will visit every student who [RODICA <will>].
	\end{xlist}
\end{exe}
The material in brackets is a focus anaphor; (\ref{ex6:acda}b) uses ellipsis. A problematic property of these examples is that the focus anaphor must find an antecedent from the very VP that contains it. The focus anaphor must therefore have a meaning that matches the meaning of some other constituent in this sentence, where by ``match,'' I mean the condition on focus anaphora that (\ref{ex6:fanaphora}) expresses. Our present syntax makes D-structures the sole parse in a derivation from which meanings are derived. But if we consider The D-structure parse that our present syntax provides for (\ref{ex6:acda}), there is no constituent that has the meaning that the focus anaphor requires. These sentences include a construction-type that we have not yet examined, a kind of modifier of nouns that is called a ``relative clause.'' The relative clauses in (\ref{ex6:acda}) are what contain the focus anaphor: they are the \emph{who Rodica will visit} part. To see what the D-structures are for these sentences requires understanding how relative clauses work. We'll turn to that task immediately below, but in brief, they are CPs headed by the complementizer \emph{that}, and their surface form involves movement of \emph{who} plus deletion of the complementizer \emph{that}. The D-structure representation of (\ref{ex6:acda}) is~(\ref{ex6:acdat}).
\begin{exe}
	\raggedright
	\ex \label{ex6:acdat}
	\footnotesize
\begin{forest}
[TP [\ibar{T}
	[\obar{T} [will] ]
	[vP 
		[DP [Thuy, roof]]
		[\ibar{v} 
			[\obar{v} [agent] ]
			[VP [\ibar{V} 
				[\obar{V} [$\surd$visit] ]
				[DP [\ibar{D}
					[\obar{D} [every] ]
					[NP [\ibar{N}
						[\ibar{N} [student, roof] ]
						[CP [\ibar{C} 
							[\obar{C} [that] ]
							[TP [\ibar{T}
								[\obar{T} [will] ]
								[vP 
									[DP [Rodica, roof]]
									[ \ibar{v}
										[\obar{v} [agent] ]
										[VP [\ibar{V}
											[\obar{V} [$\surd$visit] ]
											[DP [who, roof] ]
										]]
									]
								]
							]]
						]]
					]]
				]]
			]]
		 ]
	]
]]
\end{forest}
\end{exe}
The problem is simply that any constituent that includes \emph{Thuy will visit}, the focus antecedent, also includes the focus anaphor, \emph{Rodica will visit}. \citet{Bouton:1970}, who discovered the ellipsis version of these cases, called them instances of Antecedent Contained Deletion.
% //todo Check what Bouton actually called ACD.

It is very difficult to see how the theory that explains what the antecedent for a focus anaphor is could be manipulated to allow the parse in (\ref{ex6:acdat}) to get the interpretation it does. As long as the focus anaphor is part of the constituent whose meaning provides the antecedent meaning, there is the paradox of making the interpretation of the focus anaphor depend on its own meaning. The problem seems to be with our syntax. It forces (\ref{ex6:acdat}) to be the parse from which the denotation is derived, and this leads to the paradox of antecedent contained deletion. Interestingly, our syntax does allow this sentence to be associated with another parse, but it isn't a parse that the semantic interpretation is read off of. That parse employs Complex NP Shift, and is roughly~(\ref{ex6:cnpcacd1}).
\begin{exe}
	\raggedright
	\ex \label{ex6:cnpcacd1}
	\footnotesize
\begin{forest}
[TP
	[TP 
		[DP [Thuy, roof]]
		[\ibar{T} 
			[\obar{T} [will] ]
			[vP [\ibar{v} 
				[\obar{v} [agent] ]
				[VP [\ibar{V} [\obar{V} [$\surd$visit]]]]
			]]
		]
	]
	[DP [\ibar{D}
		[\obar{D} [every] ]
		[NP [\ibar{N}
			[\ibar{N} [student, roof] ]
			[CP [\ibar{C}
				[\obar{C} [that] ]
				[TP 
					[DP [Rodica, roof]]
					[ \ibar{T}
						[\obar{T} [will] ]
						[vP [\ibar{v}
							[\obar{v} [agent] ]
							[VP [\ibar{V}
								[\obar{V} [$\surd$visit] ]
								[DP [who, roof] ]
							]]
						]]
					]
				]
			]]
		]]
	]]
]
\end{forest}
\end{exe}
This parse also doesn't provide the material necessary to satisfy the conditions on the focus anaphor. Indeed, this parse doesn't provide the right configuration for our semantics to give it an interpretation at all. In particular, the verb \emph{visit} in the root clause cannot semantically compose with the object DP, and so this parse does not allow the object DP to get the $\theta$-role we understand it to have. But at least this syntax does not put the focus anaphor inside the material that could furnish it with an antecedent. It overcomes the largest problem with deriving the requisite meaning from~(\ref{ex6:acdat}). The solution I will lay out here seeks a way of allowing (\ref{ex6:cnpcacd1}) to get a semantic interpretation, and thereby furnish the antecedent for the focus anaphor.

 This will require a change to our syntax, but also a change to how our syntax interacts with the semantics. It will no longer be the case that only the D-structure parse in a derivation gets a semantic interpretation. Instead, we will have to allow parses that are created from the D-structure by a transformational rule be able to give rise to a semantic interpretation. Let's change now our definition of derivation in anticipation of this new architecture.
\begin{exe}
	\raggedright
		\ex \label{ex6:derivation2} Lebeaux Model\\
			$\mathcal{D}$ = (P_1, P_2\ldots,P_n) is a derivation iff:
			\begin{xlist}
				\ex P_1 is formed by the \ibar{X} Skeleton and Lexical Insertion,
				\ex Each P_{i>1} is formed from P_{i-1} by a rule,
				\ex Some P_i is semantically interpreted, and
				\ex P_n is subject to the linearization algorithm.
			\end{xlist}
			\sn A sentence is the set of well-formed derivations which have the same P_1 (=D-structure). A grammatical sentence is a derivation whose cardinality isn't greater than the cardinality of any other derivation for that sentence.
	\end{exe}
%
\marginnote[-20pt]{The cardinality of a derivation is the number of Ps in it.}%
%
This definition of derivation is identical to the one we started this chapter with, except that the parse relevant for providing a semantic interpretation has not been restricted to D-structure. Instead, this definition allows any parse in a derivation to supply the semantic interpretation. We'll investigate other possible ways of giving up D-structures as the semantically interpretable parse; another popular option is to associate another specific element of the derivation---like the S-structure---as the semantically interpretable parse. The view expressed in (\ref{ex6:derivation2}) is often identified with David Lebeaux,\cite{Lebeaux:1988,Lebeaux:2009} who argued for it in his dissertation. I'll call it the Lebeaux model.

The move to the Lebeaux model removes constraints on D-structures that our former model had. If only D-structures are semantically interpreted, then they will have to have a form that delivers a meaning. Therefore the requirements our semantics has will play a role in determining which D-structures are well-formed and which aren't. The Lebeaux model changes this. If we persist with the assumption that a grammatical sentence must provide a parse that can be semantically interpreted, the Lebeaux model lets a sentence satisfy this if any of the elements of its derivation is semantically interpretable. On the Lebeaux model, the semantics requires of D-structures, then, merely that they, or some parse derivable from them, be semantically interpretable. It's presently an open question whether this weakening of the conditions on D-structures brings unwelcome consequences.

To see the kind of issues involved, consider for instance how the Theta Criterion applies to derivations.
\begin{exe}
\raggedright
	\ex \label{ex6:tc} Theta Criterion
	\begin{xlist}
		\ex Every phrase that is interpreted as an argument must be in a $\theta$-marked position, and
		\ex Every $\theta$-marked position must contain a phrase that serves as an argument.
	\end{xlist}
\end{exe}
Our present model derives that the Theta Criterion is a generalization about D-structures from (\ref{ex6:vptype}) and the system of semantic rules and denotations that were developed in chapter~\ref{sec:reconsidered}.
\begin{exe}
\raggedright
	\ex \label{ex6:vptype} A vP must have a denotation that makes it a function whose domain is eventualities.
\end{exe}
Consider, for instance, what would happen if a verb were not supplied with enough arguments by the D-structure (as in (\ref{ex6:tcvio1})) or more arguments are put into a vP than there are $\theta$-roles (as in (\ref{ex6:tcvio2})).
\medskip
\begin{fullwidth}
\begin{exe}
\raggedright
	\ex \label{ex6:tcvio1}
	\begin{forest}
	[$\lambda y\ \lambda \epsilon'$ \textsc{make}($\epsilon'$) $\wedge$ \textsc{patient}(\denotes{it})($\epsilon'$) $\wedge$ \textsc{agent}($y$)($\epsilon'$)\\vP [\ibar{v}
		[$\lambda y\ \lambda \epsilon'$ \textsc{make}($\epsilon'$) $\wedge$ \textsc{patient}(\denotes{it})($\epsilon'$) $\wedge$ \textsc{agent}($y$)($\epsilon'$)\\\obar{v} [agent] ]
		[$\lambda \epsilon$ \textsc{make}($\epsilon$) $\wedge$ \textsc{patient}(\denotes{it})($\epsilon$)\\VP [$\lambda \epsilon$ \textsc{make}($\epsilon$) $\wedge$ \textsc{patient}(\denotes{it})($\epsilon$)\\\ibar{V}
			[$\lambda x\ \lambda \epsilon$ \textsc{make}($\epsilon$) $\wedge$ \textsc{patient}($x$)($\epsilon$)\\\obar{V} [$\surd$make] ]
			[DP [it, roof] ]
		] ]
	]]
	\end{forest}
	\vspace*{-30pt}
	\ex \label{ex6:tcvio2}
	\begin{forest}
	[vP
		[DP [she, roof] ]
		[\ibar{v}
			[\obar{v} [agent] ]
			[VP [??\\\ibar{V}
				[$\lambda \epsilon$ \textsc{arrive}($\epsilon$) \\\obar{V} [$\surd$arrive] ]
				[DP [it, roof] ]
			] ]
		 ]
	]
	\end{forest}
\end{exe}
\end{fullwidth}
\medskip
In (\ref{ex6:tcvio2}) no denotation is delivered to the \ibar{V}, and though our semantics delivers a denotation to vP in (\ref{ex6:tcvio1}), it isn't one that satisfies~(\ref{ex6:vptype}).

If D-structures are no longer the parse in a derivation that must satisfy (\ref{ex6:vptype}), there are new derivations to consider in which (\ref{ex6:vptype}) is satisfied by another parse. For instance, we might consider a derivation whose D-structure is (\ref{ex6:newdstr}) and whose S-structure is~(\ref{ex6:newsstr}).
\begin{exe}
\raggedright
	\ex \label{ex6:newdstr}
	\small
	\begin{forest}
	[TP [\ibar{T}
		[\obar{T} [will] ]
		[vP [\ibar{v}
			[\obar{v} [experiencer] ]
			[VP [\ibar{V} 
				[\obar{V} [$\surd$ believe] ]
				[TP [\ibar{T}
					[\obar{T} [to] ]
					[VP
						[DP [they, roof]]
						[\ibar{V} 
							[\obar{V} [$\surd$seem] ]
							[CP [that pigs fly, roof] ]
						 ]
					 ]
				] ]
			] ]
		] ]
	]]
	\end{forest}
	\ex \label{ex6:newsstr}
	\small
	\begin{forest}
	[TP
		[DP [they, roof] ]
		[\ibar{T}
			[\obar{T} [will] ]
			[vP [\ibar{v}
				[\obar{v} [experiencer] ]
				[VP [\ibar{V}
					[\obar{V} [$\surd$believe] ]
					[TP [\ibar{T}
						[\obar{T} [to] ]
						[VP [\ibar{V}
							[\obar{V} [$\surd$seemed] ]
							[CP [that pigs fly, roof] ]
						] ]
					] ]
				] ]
			] ]
		 ]
	]
	\end{forest}
\end{exe}
This derivation can have the parses indicated in~(\ref{ex6:newder}).
\begin{exe}
\raggedright
	\ex \label{ex6:newder}
	\small
	\begin{forest}
	[TP [\ibar{T}
		[\obar{T} [will] ]
		[vP
			[DP [they, name=binder, roof] ]
			[\ibar{v}
				[\obar{v} [experiencer] ]
				[TP
					[{\ }, name=trace ]
					[\ibar{T}
						[\obar{T} [to] ]
						[VP [\ibar{V}
							[\obar{V} [$\surd$seem] ]
							[CP [that pigs fly, roof] ]
						] ]
					 ]
				 ]
			 ]
		 ]
	]]
	\draw[->,>=stealth', thin, gray] (trace) to[out=south west,in=south,looseness=1.4] node[pos=0.75,xshift=-10mm]{} (binder);
	\end{forest}
\end{exe}
This is a grammatical derivation, and the parse created by movement in (\ref{ex6:newder}) is semantically interpretable. But the sentence it creates---\emph{They will believe to seem that pigs fly}---is ungrammatical. By adopting the Lebeaux model, we will have to ensure that the consequences of letting D-structures be freed from the requirement of semantically interpretability doesn't have untoward results, like admitting the derivation in (\ref{ex6:newdstr})--(\ref{ex6:newder})--(\ref{ex6:newsstr}). 

To exploit the Lebeaux model, we'll change the definition of movement so that the structures created by movement operations have enough information in them to allow the semantics to assemble their meanings correctly. In particular, we'll need a way of allowing the verb \emph{visit} and its object, \emph{every student \ldots}, to semantically combine in (\ref{ex6:cnpcacd1}). We'll also need to have a syntax of relative clauses. And finally, we should make these changes in a way that helps close the problems that arise from letting D-structures be free of the conditions that result from forcing them to be semantically interpreted.

\subsubsection{Traces, Referential Indices, and Quantifiers}
\label{sec:traces}

Let's begin by giving to representations like (\ref{ex6:cnpcacd1}) the ability to semantically combine the verb with the DP that has moved. The method for doing this that we will start with involves adding to the movement operation two additional processes that accompany the change of position function that we've been relying on up to now. The first of these processes inserts into the position that a moved term vacates a term that will be semantically interpreted as a variable. This term is called a ``trace,'' and it is standardly represented by ``\emph{t}.'' %
%
\marginnote{I believe the trace-based system for movement was introduced in \citet{Chomsky:1977}, and developed in \citet{Chomsky:1977b}. The method of expressing how traces are bound that I will sketch here is from \citet{Heim:1998}.}%
%
The second process makes the term that has moved bind that variable. One method of expressing this is to let the movement operation change the meaning of the largest phrase that contains the trace but not the moved phrase, so that it becomes a function whose $\lambda$ operator binds the trace. 

To do this will require a syntax that regulates the connection between variables and their binders. Up to now, the connection between a variable and its binder has been part of the denotation of the morphemes that populate syntactic representations. The fact that pronouns can also be interpreted as variables shows us that we will need a syntax that links expressions that are interpreted as variables to their binders as well. The idea behind Trace theory is that movement also invokes this syntax. We'll add to the morphemes that can be part of our syntactic representations, terms whose job it will be to express the variable-binder connection. These terms are known as ``referential indices,'' and they will be represented as numerical subscripts. To begin with, we will let indices be parts of traces, and nothing else.

With this background, we can now redefine movement with~(\ref{ex6:tracemove}).
\begin{exe}
\raggedright
	\ex \label{ex6:tracemove} $\alpha$ \underline{moves to} $\beta$ =_{\text{def.}} $\alpha$ takes position $\beta$ and if $\alpha$ is an~XP:
	\begin{xlist}
			\ex $\alpha$ is replaced by \emph{t}_n, where $n$ is a unique referential index, and
			\ex The denotation of the sister of $\beta$, $\gamma$, becomes: $\lambda n$\denotes{$\gamma$}.
	\end{xlist}
\end{exe}

This will give to an S-structure that has been created with Complex NP Shift a representation like that illustrated by~(\ref{ex6:tracet1}).
\begin{exe}
\raggedright
	\ex \label{ex6:tracet1}
	\footnotesize
	\begin{forest}
	[TP
		[TP 
			[DP [Thuy, roof]]
			[\ibar{T}	
				[\obar{T} [will]]
				[vP 
					[$t_2$]				
					[\ibar{v}
						[\obar{v} [agent] ]
						[VP 
						[\ibar{V} 
							[\obar{V} [$\surd$visit]]
							[$t_3$]
						]]
				]]
			 ]
		]
		[DP [\ibar{D}
			[\obar{D} [every] ]
			[NP [student, roof] ]
		]]
	]
	\end{forest}
\end{exe}
Note that both Argument Movement of \emph{Thuy} and Complex NP Shift of \emph{every student} have invoked traces, and that their referential indices are distinct. The requirement in (\ref{ex6:tracemove}a) that the referential index on a trace be unique is responsible for ensuring that each trace in a syntactic representation be associated with independent referential indices. To see the effect that (\ref{ex6:tracemove}b) has, we will look at how the semantics assigns denotations to the pieces of (\ref{ex6:tracemove}). We have not developed a semantics for \emph{will}, however, and so we are not equipped with everything necessary to do that. In general terms, we have not developed a semantics for modals or tense---the two kinds of terms that occupy \obar{T} position in finite clauses---and this is why we've not endeavored to give meanings for anything larger than vP in our parses. The meanings of these expressions requires a dive into purely semantic issues that will not, in my estimation, further our understanding of syntax. Let's therefore ignore entirely the semantic contribution that the contents of \obar{T} makes.

If we do this, the semantic composition of the pieces of (\ref{ex6:tracet1}) that we have previously determined is~(\ref{ex6:tracet1den}).
\medskip
\begin{fullwidth}
\begin{exe}
\raggedright
	\ex \label{ex6:tracet1den}	
	\hspace*{-20pt}
	\footnotesize
	\begin{forest}
		[$\lambda \epsilon \forall x$ \textsc{student}($x$) $\rightarrow$ {[}\textsc{visit}($\epsilon$) $\wedge$ \textsc{patient}($x$)($\epsilon$) $\wedge$ \textsc{agent}(\denotes{Thuy})($\epsilon$){]} \\TP
		[$\lambda 3 \lambda \epsilon'$ \textsc{visit}($\epsilon'$) $\wedge$ \textsc{patient}($3$)($\epsilon'$) $\wedge$ \textsc{agent}(\denotes{Thuy})($\epsilon'$)\\TP 
			[DP [Thuy, roof]]
			[$\lambda 2 \lambda \epsilon'$ \textsc{visit}($\epsilon'$) $\wedge$ \textsc{patient}($3$)($\epsilon'$) $\wedge$ \textsc{agent}($2$)($\epsilon'$)\\\ibar{T}	
				[\obar{T} [will]]
				[$\lambda \epsilon'$ \textsc{visit}($\epsilon'$) $\wedge$ \textsc{patient}($3$)($\epsilon'$) $\wedge$ \textsc{agent}($2$)($\epsilon'$)\\ vP 
					[$t_2$]				
					[$\lambda y \lambda \epsilon'$ \textsc{visit}($\epsilon'$) $\wedge$ \textsc{patient}($3$)($\epsilon'$) $\wedge$ \textsc{agent}($y$)($\epsilon'$)\\ \ibar{v}
						[$\lambda P \lambda y \lambda \epsilon' P(\epsilon')$ $\wedge$ \textsc{agent}($y$)($\epsilon'$)\\\obar{v} [agent] ]
						[$\lambda \epsilon$ \textsc{visit}($\epsilon$) $\wedge$ \textsc{patient}($3$)($\epsilon$)\\VP 
						[$\lambda \epsilon$ \textsc{visit}($\epsilon$) $\wedge$ \textsc{patient}($3$)($\epsilon$)\\ \ibar{V} 
							[$\lambda z \lambda \epsilon$ \textsc{visit}($\epsilon$) $\wedge$ \textsc{patient}($z$)($\epsilon$)\\ \obar{V} [$\surd$visit]]
							[$t_3$]
						]]
				]]
			 ]
		]
		[$\lambda Q \lambda \epsilon \forall x$ \textsc{student}($x$) $\rightarrow$ $Q(x)(\epsilon)$\\DP [\ibar{D}
			[\obar{D} [every] ]
			[NP [student, roof] ]
		]]
	]
	\end{forest}
\end{exe}
\end{fullwidth}
\medskip
Note that the denotation of traces is just the variable that their referential index is. When a trace is the argument of a function it contributes its value to the argument position(s) in that function. Thus, when the denotation of \emph{$\surd$visit} combines with the trace in its complement position, the result puts the value of the variable $3$ into the relation that the patient $\theta$-role denotes.%
%
\marginnote[0pt]{This means that the phrases which contain the trace, but not its binder, are technically without a denotation, since we have not given a denotation for a variable except when bound. A more common approach to the semantics of variables lets them have denotations even in contexts where they are not bound.}%
%
The same thing happens with the trace in the Specifier of vP. That trace is the argument of the function that \ibar{v} denotes, and when they compose, the value of that variable is inserted into the agent $\theta$-role relation.

The action of (\ref{ex6:tracemove}c) can be seen in the denotations that \ibar{T} and the lower TP have. If we ignore the semantic contribution made by \emph{will}, the denotation of \ibar{T} should be the same as the denotation of vP. They contain all the same meaningful items. But the denotation of \ibar{T} differs from vP by virtue of~(\ref{ex6:tracemove}c). Because \ibar{T} is the largest phrase that dominates \emph{t}_2 but does not dominate the phrase that \emph{t}_2 replaces, it is the target for (\ref{ex6:tracemove}c). The rule in (\ref{ex6:tracemove}c) adds to the beginning of the denotation that \ibar{T} would otherwise have, a $\lambda$ operator that binds the trace; in this case, that is ``$\lambda 2$.'' A similar step is responsible for the ``$\lambda$'' that starts the denotation for lower the TP as well. That phrase is the largest phrase that contains the trace \emph{t}_3, but not the phrase that the trace replaced, and so it is also the target for~(\ref{ex6:tracemove}c).

In this way, giving movement the ability to change the meanings of the sentences allows the moved DPs in (\ref{ex6:tracet1den}) to be interpreted in the positions they move to. But now that we are looking at the meanings of parts of sentences that are larger than vP, a problem to our semantics is exposed.%
%
\marginnote[0pt]{See \citet{Higginbotham:1985} for an early discussion of this issue.}%
%
 The denotation of vPs in our semantics is a function whose domain is eventualities: the denotation of vP is a description of eventualities. The denotations of clauses constitutes a claim about the existence of an eventuality of the sort that its vP describes. There must consequently be something in a clause that is responsible for changing the denotation of a vP into this claim. If we assume that this change is effected by the denotation of some term in a clause, then we can schematize what is needed with~(\ref{ex6:exclose}).
\begin{exe}
\raggedright
	\ex \label{ex6:exclose}
	\small
	\begin{forest}
	[$\exists \epsilon'$ \textsc{dance}($\epsilon'$) $\wedge$ \textsc{agent}(\denotes{Presley})($\epsilon'$) \\XP
		[$\lambda P\ \exists \epsilon'\ P(\epsilon')$\\ \obar{X} ]
		[$\lambda \epsilon$ \textsc{dance}($\epsilon$) $\wedge$ \textsc{agent}(\denotes{Presley}($\epsilon$)) \\vP [Presley dance, roof] ]
	]
	\end{forest}
\end{exe}
What the \obar{X} in (\ref{ex6:exclose}) does is change the description of eventualities that vP denotes into a sentence that claims that there is such an eventuality. This is the type of thing that we want (non-negated) sentences in general to do. Our semantics so far has been focused on the issue of combining DP arguments and the verbs or adjectives that they are arguments of, and this has allowed us to restrict attention just to how vPs are semantically composed. Now that we must consider the semantics of larger parts of sentences, we will have to acknowledge the existence of something within them that converts the description of eventualities that vPs are into expressions that say that such eventualities exist. Somewhere within a sentence, there must be something that has the function that \obar{X} has in~(\ref{ex6:exclose}).

We can see that this term must be lower than negation in English from~(\ref{ex6:pres}).
\begin{exe}
\raggedright
	\ex \label{ex6:pres}
	Presley did not dance.
\end{exe}
The denotation for this sentence should make it deny the existence of a dancing eventuality that has Presley as its agent. The existential quantification over events should be part of what \emph{not} negates.

A popular speculation is that \obar{X} should be equated with a head that determines the ``aspect'' of a clause. %
%
\marginnote[0pt]{See \citet{Pancheva:2003a} for an accessible introduction.}%
%
Aspectual morphemes have meanings that relate the event described by a vP to tense. By way of illustration, consider the sentences in~(\ref{ex6:aspectintro}), all of which have present tense.
\begin{exe}
\raggedright
	\ex \label{ex6:aspectintro}
	\begin{xlist}
		\ex She sings.
		\ex She is singing.
		\ex She has sung.
	\end{xlist}
\end{exe}
Semantically, these sentences differ in a way that can be described in terms of how the singing eventuality is related to the time at which these sentences are produced. Thus, for instance, (\ref{ex6:aspectintro}b) conveys the information that the time at which the singing eventuality is happening includes the time at which (\ref{ex6:aspectintro}b) is produced. By contrast, (\ref{ex6:aspectintro}c) claims that the time at which the singing eventuality occurs is prior to the time at which (\ref{ex6:aspectintro}c) is produced. Finally, (\ref{ex6:aspectintro}a) indicates that the singing eventuality is a kind of state that holds at the time when (\ref{ex6:aspectintro}a) is produced. In each case, the tense provides a time to which the eventuality is related. This becomes clearer if (\ref{ex6:aspectintro}) is compared to the meanings these sentences have when they contain past tense.
\begin{exe}
\raggedright
	\ex \label{ex6:aspectintro2}
	\begin{xlist}
		\ex She sang.
		\ex She was singing.
		\ex She had sung.
	\end{xlist}
\end{exe}
Just as in (\ref{ex6:aspectintro}), these sentences differ in how the singing eventuality is situated temporally with respect to tense. But rather than that relation holding between the time of the eventuality and the time at which the sentences are produced, the relation holds between the time of the eventuality and a time prior to when the sentences are produced. Thus, for instance, (\ref{ex6:aspectintro2}b) says that the singing eventuality occurs at a time that includes some time prior to when (\ref{ex6:aspectintro2}b) is produced. And (\ref{ex6:aspectintro2}a) conveys that the singing eventuality should be regarded as a kind of state that occurred at a time prior to the time at which (\ref{ex6:aspectintro2}a) is produced. 

We shall not investigate how this semantic information is conveyed by the material in a syntactic representation. As the examples above indicate, these distinctions in English involve the use of auxiliary verbs as well as morphology on the verbs those auxiliaries embed. Because neither of these ingredients are present in (\ref{ex6:aspectintro}a)/(\ref{ex6:aspectintro2}a), it is not obvious that aspectual information is present in these examples. The evidence for an aspect morpheme in these examples comes largely from a semantic theory of aspect. It is also in the semantics of aspect that we find evidence for it being responsible for introducing the existential binder of an event variable. Because I will not develop this semantics, I will not be able to provide you with the evidence for equating \obar{X} with Aspect. Instead, I will rely solely on the evidence provided by (\ref{ex6:pres}) for the conclusion that \obar{X} is situated in a position syntactically lower than \obar{T}---a position I will call~\obar{Asp}. For our purposes, we need not reveal the full semantics of the morpheme that occupies \obar{Asp}. Instead, we can let \obar{Asp} have just the semantics of \obar{X}. That is, we will focus on just the part of its denotation that existentially closes the event variable.

If we adopt the denotation for \emph{asp}, the content of \obar{Asp}, indicated in~(\ref{ex6:existev1}), then \emph{she ran} will have the denotation indicated in~(\ref{ex6:existev2}).%
%
\marginnote[0pt]{I continue to ignore the contribution made to the denotation of a sentence that comes from tense.}%
%
\begin{exe}
\raggedright
	\ex \label{ex6:existev1}
	\denotes{\emph{asp}} = $\lambda P\ \exists \epsilon\ P(\epsilon)$ 
	\ex \label{ex6:existev2}
	\small
	\begin{forest}
	[$\exists \epsilon$ \textsc{run}($\epsilon$) $\wedge$ \textsc{agent}(\denotes{she})($\epsilon$)\\TP
	[\ibar{T}
	[\obar{T} [past]]
	[$\exists \epsilon$ \textsc{run}($\epsilon$) $\wedge$ \textsc{agent}(\denotes{she})($\epsilon$)\\AspP
		[$\exists \epsilon$ \textsc{run}($\epsilon$) $\wedge$ \textsc{agent}(\denotes{she})($\epsilon$)\\\ibar{Asp} 
			[\obar{Asp} [asp] ]
			[$\lambda \epsilon'$ \textsc{run}($\epsilon'$) $\wedge$ \textsc{agent}(\denotes{she})($\epsilon'$) \\vP
				[DP [she, roof]]
				[$\lambda x \lambda \epsilon'$ \textsc{run}($\epsilon'$) $\wedge$ \textsc{agent}($x$)($\epsilon'$) \\\ibar{v}  
					[\obar{v} [agent] ]
					[$\lambda \epsilon'$ \textsc{run}($\epsilon'$) \\VP [\ibar{V} [\obar{V} [$\surd$run]]] ]
				]
			 ]
		]
	]]]
	\end{forest}
\end{exe}
This gives to TP a denotation that claims that there exists an eventuality that has the properties that vP describes. In the case of (\ref{ex6:existev2}), the TP's meaning claims that there is an eventuality of dancing whose agent is whatever \emph{she} refers to. 

Putting the term that existentially closes the eventuality in the position that (\ref{ex6:existev2}) does correctly puts it inside the constituent that \emph{not}, and the other exponents of \emph{neg}, negates. Assume that \emph{neg} is a function which takes the denotation of its complement and returns that denotation negated; this is what (\ref{ex6:negden}) does.%
%
\marginnote[0pt]{Understand the meaning of ``$\neg p$'' to have the paraphrase given by ``not p.'' That is: if $p$ is true, $\neg p$ is false.}%
%
\begin{exe}
\raggedright
	\ex \label{ex6:negden}
	\denotes{neg} = $\lambda p\ \neg p$
\end{exe}
(\ref{ex6:negdent}) illustrates how this will deliver the correct meaning to a sentence with \emph{neg} in it.
\begin{exe}
\raggedright
	\sn she didn't run $\equiv$
	\ex \label{ex6:negdent}
	\small
	\begin{forest}
		[$\neg \exists \epsilon$ \textsc{run}($\epsilon$) $\wedge$ \textsc{agent}(\denotes{she})($\epsilon$)\\TP
		[\ibar{T}
		[\obar{T} [did]]
		[NegP
			[$\neg \exists \epsilon$ \textsc{run}($\epsilon$) $\wedge$ \textsc{agent}(\denotes{she})($\epsilon$)\\\ibar{Neg}
			[$\lambda p\ \neg p$\\\obar{Neg} [nt]]
		[$\exists \epsilon$ \textsc{run}($\epsilon$) $\wedge$ \textsc{agent}(\denotes{she})($\epsilon$)\\AspP
			[$\exists \epsilon$ \textsc{run}($\epsilon$) $\wedge$ \textsc{agent}(\denotes{she})($\epsilon$)\\\ibar{Asp} 
				[\obar{Asp} [asp] ]
				[$\lambda \epsilon'$ \textsc{run}($\epsilon'$) $\wedge$ \textsc{agent}(\denotes{she})($\epsilon'$) \\vP
					[DP [she, roof]]
					[$\lambda x \lambda \epsilon'$ \textsc{run}($\epsilon'$) $\wedge$ \textsc{agent}($x$)($\epsilon'$) \\\ibar{v}  
						[\obar{v} [agent] ]
						[$\lambda \epsilon'$ \textsc{run}($\epsilon'$) \\VP [\ibar{V} [\obar{V} [$\surd$run]]] ]
					]
				 ]
			]
		]]]]]
		\end{forest}
\end{exe}

This change has consequences for the denotations we have adopted for the quantificational determiners \emph{no}, \emph{some}, and \emph{every}. Those denotations, repeated in (\ref{ex6:quantoldden}), are designed so that they can be interpreted in Specifier of vP and complement of \obar{V} positions.
\begin{exe}
\raggedright
	\ex \label{ex6:quantoldden}
	\begin{xlist}
		\ex \denotes{no} = $\lambda P\ \lambda Q\ \lambda \epsilon\ \neg\exists x [P(x) \wedge Q(x)(\epsilon)]$
		\ex \denotes{some} = $\lambda P\ \lambda Q\ \lambda \epsilon\ \exists x [P(x) \wedge Q(x)(\epsilon)]$
		\ex \denotes{every} = $\lambda P\ \lambda Q\ \lambda \epsilon\ \forall x [P(x) \rightarrow Q(x)(\epsilon)]$
	\end{xlist}
\end{exe}
The second argument for these denotations calls for a term that relates things to eventualities. That is the right semantic type for \ibar{v}s, and this is why these denotations allow a DP headed by one of these quantifiers to be semantically interpreted in Specifier of vP; (\ref{ex6:oldqspecvp}) illustrates.
\begin{fullwidth}
\begin{exe}
\raggedright
	\ex \label{ex6:oldqspecvp}
	\hspace*{-20pt}
	\small
	\begin{forest}
	[AspP [\ibar{Asp}
		[\obar{Asp} [asp]]
		[$\lambda \epsilon\ \exists x$ \textsc{woman}($x$) $\wedge$ \textsc{run}($\epsilon$) $\wedge$ \textsc{agent}($x$)($\epsilon$)\\vP 
			[$\lambda Q\ \lambda \epsilon\ \exists x$ \textsc{woman}($x$) $\wedge$ \textsc{run}($\epsilon$) $\wedge$ \textsc{agent}($x$)($\epsilon$)\\DP [\ibar{D}
				[$\lambda P\ \lambda Q\ \lambda \epsilon\ \exists x$ $P(x)$ $\wedge$ $Q(x)(\epsilon)$\\ \obar{D} [some]]
				[$\lambda y$ \textsc{woman}($y$)\\NP [woman, roof] ]
			]]
			[$\lambda x\ \lambda \epsilon$ \textsc{run}($\epsilon$) $\wedge$ \textsc{agent}($x$)($\epsilon$)\\\ibar{v}
				[\obar{v} [agent] ]
				[$\lambda \epsilon$ \textsc{run}($\epsilon$)\\ VP [$\surd$run, roof] ]
			]
		]
	]]
	\end{forest}
\end{exe}
\end{fullwidth}
\medskip
Because transitive verbs are also of this semantic type, DPs headed by one of these quantifiers can also combine with transitive verbs, as (\ref{ex6:oldqcompvp}) illustrates.
\medskip
\begin{fullwidth}
\begin{exe}
\raggedright
	\ex \label{ex6:oldqcompvp}
	\small
	\begin{forest}
	[AspP [\ibar{Asp}
		[\obar{Asp} [asp]]
		[vP
			[DP [she, roof]]
			[$\lambda y\ \lambda \epsilon\ \exists x$ \textsc{beans}($x$) $\wedge$ \textsc{eat}($\epsilon$) $\wedge$ \textsc{patient}($x$)($\epsilon$) $\wedge$ \textsc{agent}($y$)($\epsilon$)\\ \ibar{v} 
				[\obar{v} [agent] ]
				[$\lambda \epsilon\ \exists x$ \textsc{beans}($x$) $\wedge$ \textsc{eat}($\epsilon$) $\wedge$ \textsc{patient}($x$)($\epsilon$)\\VP [$\lambda \epsilon\ \exists x$ \textsc{beans}($x$) $\wedge$ \textsc{eat}($\epsilon$) $\wedge$ \textsc{patient}($x$)($\epsilon$)\\ \ibar{V}
					[$\lambda x\ \lambda \epsilon$ \textsc{eat}($\epsilon$) $\wedge$ \textsc{patient}($x$)($\epsilon$)\\\obar{V} [$\surd$eat] ]
					[$\lambda Q\ \lambda \epsilon\ \exists x$ \textsc{beans}($x$) $\wedge$ $Q(x)(\epsilon)$\\DP [\ibar{D}
						[$\lambda P\ \lambda Q\ \lambda \epsilon\ \exists x$ $P(x)$ $\wedge$ $Q(x)(\epsilon)$\\ \obar{D} [some] ]
						[$\lambda y$ \textsc{beans}($y$)\\NP [beans, roof] ]
					] ]
				] ]
			]
		 ]
	]]
	\end{forest}
\end{exe}
\end{fullwidth}
\medskip
But by introducing AspP, and thereby changing the denotations of \ibar{T} and TP, the position that DPs move to will no longer put them in a position that supplies an argument of the right semantic type for a quantifier. This can be seen by considering how (\ref{ex6:tracet1den}) is changed by this semantics. The TP that the moved object DP is a sister to is no longer the right kind of semantic term, as~(\ref{ex6:tracet1den2}) illustrates.
\medskip
\begin{fullwidth}
\begin{exe}
	\raggedright
		\ex \label{ex6:tracet1den2}	
		\footnotesize
		\begin{forest}
			[??\\TP
			[$\lambda 3 \exists \epsilon'$ \textsc{visit}($\epsilon'$) $\wedge$ \textsc{patient}($3$)($\epsilon'$) $\wedge$ \textsc{agent}(\denotes{Thuy})($\epsilon'$)\\TP 
				[DP [Thuy, roof]]
				[$\lambda 2 \exists \epsilon'$ \textsc{visit}($\epsilon'$) $\wedge$ \textsc{patient}($3$)($\epsilon'$) $\wedge$ \textsc{agent}($2$)($\epsilon'$)\\\ibar{T}	
					[\obar{T} [will]]
					[$\exists \epsilon'$ \textsc{visit}($\epsilon'$) $\wedge$ \textsc{patient}($3$)($\epsilon'$) $\wedge$ \textsc{agent}($2$)($\epsilon'$)\\AspP [\ibar{Asp}
					[\obar{Asp} [asp]]
					[$\lambda \epsilon'$ \textsc{visit}($\epsilon'$) $\wedge$ \textsc{patient}($3$)($\epsilon'$) $\wedge$ \textsc{agent}($2$)($\epsilon'$)\\ vP 
						[$t_2$]				
						[$\lambda y \lambda \epsilon'$ \textsc{visit}($\epsilon'$) $\wedge$ \textsc{patient}($3$)($\epsilon'$) $\wedge$ \textsc{agent}($y$)($\epsilon'$)\\ \ibar{v}
							[$\lambda P \lambda y \lambda \epsilon' P(\epsilon')$ $\wedge$ \textsc{agent}($y$)($\epsilon'$)\\\obar{v} [agent] ]
							[$\lambda \epsilon$ \textsc{visit}($\epsilon$) $\wedge$ \textsc{patient}($3$)($\epsilon$)\\VP 
							[$\lambda \epsilon$ \textsc{visit}($\epsilon$) $\wedge$ \textsc{patient}($3$)($\epsilon$)\\ \ibar{V} 
								[$\lambda z \lambda \epsilon$ \textsc{visit}($\epsilon$) $\wedge$ \textsc{patient}($z$)($\epsilon$)\\ \obar{V} [$\surd$visit]]
								[$t_3$]
							]]
					]]
				 ]
			]]]
			[$\lambda Q \lambda \epsilon \forall x$ \textsc{student}($x$) $\rightarrow$ $Q(x)(\epsilon)$\\DP [\ibar{D}
				[\obar{D} [every] ]
				[NP [student, roof] ]
			]]
		]
		\end{forest}
	\end{exe}
\end{fullwidth}
\medskip
But this is just the syntactic representation that we need to find a denotation for, in order to solve the problem posed by antecedent contained anaphora.

Let's solve this problem by designing the denotations of \emph{no}, \emph{some}, and \emph{every} so that they can be interpreted in their moved positions, but not in their underlying position. The new denotations are~(\ref{ex6:newqden}).
\begin{exe}
\raggedright
	\ex \label{ex6:newqden}
	\begin{xlist}
		\ex \denotes{no} = $\lambda P\ \lambda Q\ \neg\exists x [P(x) \wedge Q(x)]$
		\ex \denotes{some} = $\lambda P\ \lambda Q\ \exists x [P(x) \wedge Q(x)]$
		\ex \denotes{every} = $\lambda P\ \lambda Q\ \forall x [P(x) \rightarrow Q(x)]$
	\end{xlist}	
\end{exe}
The second argument of these denotations now requires something that is a function whose domain is things and whose range is a truth value. That is the right fit for the position that the moved DPs occupy in (\ref{ex6:tracet1den2}), as can be seen from~(\ref{ex6:tracet1den3}).
\medskip
\begin{fullwidth}
\begin{exe}
\raggedright
	\ex \label{ex6:tracet1den3}
	\hspace*{-20pt}
	\footnotesize
		\begin{forest}
			[$\forall x$ \textsc{student}($x$) $\rightarrow \exists \epsilon'$ \textsc{visit}($\epsilon'$) $\wedge$ \textsc{patient}($x$)($\epsilon'$) $\wedge$ \textsc{agent}(\denotes{Thuy})($\epsilon'$)\\TP
			[$\lambda 3 \exists \epsilon'$ \textsc{visit}($\epsilon'$) $\wedge$ \textsc{patient}($3$)($\epsilon'$) $\wedge$ \textsc{agent}(\denotes{Thuy})($\epsilon'$)\\TP 
				[DP [Thuy, roof]]
				[$\lambda 2 \exists \epsilon'$ \textsc{visit}($\epsilon'$) $\wedge$ \textsc{patient}($3$)($\epsilon'$) $\wedge$ \textsc{agent}($2$)($\epsilon'$)\\\ibar{T}	
					[\obar{T} [will]]
					[$\exists \epsilon'$ \textsc{visit}($\epsilon'$) $\wedge$ \textsc{patient}($3$)($\epsilon'$) $\wedge$ \textsc{agent}($2$)($\epsilon'$)\\AspP [\ibar{AsP}
					[\obar{Asp} [asp]]
					[$\lambda \epsilon'$ \textsc{visit}($\epsilon'$) $\wedge$ \textsc{patient}($3$)($\epsilon'$) $\wedge$ \textsc{agent}($2$)($\epsilon'$)\\ vP 
						[$t_2$]				
						[$\lambda y \lambda \epsilon'$ \textsc{visit}($\epsilon'$) $\wedge$ \textsc{patient}($3$)($\epsilon'$) $\wedge$ \textsc{agent}($y$)($\epsilon'$)\\ \ibar{v}
							[$\lambda P \lambda y \lambda \epsilon' P(\epsilon')$ $\wedge$ \textsc{agent}($y$)($\epsilon'$)\\\obar{v} [agent] ]
							[$\lambda \epsilon$ \textsc{visit}($\epsilon$) $\wedge$ \textsc{patient}($3$)($\epsilon$)\\VP 
							[$\lambda \epsilon$ \textsc{visit}($\epsilon$) $\wedge$ \textsc{patient}($3$)($\epsilon$)\\ \ibar{V} 
								[$\lambda z \lambda \epsilon$ \textsc{visit}($\epsilon$) $\wedge$ \textsc{patient}($z$)($\epsilon$)\\ \obar{V} [$\surd$visit]]
								[$t_3$]
							]]
					]]
				 ]
			]]]
			[$\lambda Q \forall x$ \textsc{student}($x$) $\rightarrow$ $Q(x)$\\DP [\ibar{D}
				[\obar{D} [every] ]
				[NP [student, roof] ]
			]]
		]
		\end{forest}
\end{exe}
\end{fullwidth}
\medskip

This change to the semantics of quantifiers is independently required. The semantics I introduced originally for quantifiers harbors a problem, which can be seen by considering~(\ref{ex6:lowno1}).
\begin{exe}
\raggedright
	\ex \label{ex6:lowno1} No reader finishes this book.
\end{exe}
On the denotation for \emph{no} that I gave originally (=(\ref{ex6:quantoldden}a)), (\ref{ex6:lowno1}) should get the meaning indicated in~(\ref{ex6:lowno1t}).%
%
\marginnote[-20pt]{Because we are ignoring the semantic contribution made by tense, the denotation of (\ref{ex6:lowno1}) is the same as the denotation of \ibar{Asp} in~(\ref{ex6:lowno1t}).}%
%

\medskip
\begin{fullwidth}
\begin{exe}
\raggedright
	\ex \label{ex6:lowno1t}
	\footnotesize
	\begin{forest}
	[$\exists \epsilon\ \neg \exists x$ \textsc{reader}($x$) $\wedge$ \textsc{finish}($\epsilon$) $\wedge$ \textsc{patient}(\denotes{this book})($\epsilon$) $\wedge$ \textsc{agent}($x$)($\epsilon$) \\\ibar{Asp}
		[\obar{Asp} ]
		[$\lambda \epsilon\ \neg \exists x$ \textsc{reader}($x$) $\wedge$ \textsc{finish}($\epsilon$) $\wedge$ \textsc{patient}(\denotes{this book})($\epsilon$) $\wedge$ \textsc{agent}($x$)($\epsilon$) \\vP
			[$\lambda Q\ \lambda \epsilon\ \neg \exists x$ \textsc{reader}($x$) $\wedge$ $Q(x)(\epsilon)$ \\DP [no reader, roof]]
			[$\lambda y\ \lambda \epsilon$ \textsc{finish}($\epsilon$) $\wedge$ \textsc{patient}(\denotes{the book})($\epsilon$) $\wedge$ \textsc{agent}($y$)($\epsilon$) \\\ibar{v}
				[\obar{v} [agent]]
				[$\lambda \epsilon$ \textsc{finish}($\epsilon$) $\wedge$ \textsc{patient}(\denotes{this book})($\epsilon$)\\VP [\ibar{V}
					[\obar{V} [$\surd$finish]]
					[DP [this book, roof]]
				] ]
			 ]
		 ]
	]
	\end{forest}
\end{exe}
\end{fullwidth}
\medskip
Informally, what (\ref{ex6:lowno1}) claims is that no one who is a reader is also someone who finishes this book. If we understand \emph{reader} to be restricted to just those individuals who are readers of this book, then (\ref{ex6:lowno1}) amounts to the claim that this book is never read to completion: there are no readers of this book who finish reading it. But that doesn't correspond to the denotation of \ibar{Asp} in (\ref{ex6:lowno1t}). The denotation of \ibar{Asp} in (\ref{ex6:lowno1t}) is that there is an event of finishing this book and that there are no readers who are the agent of that finishing event. It implies that there is a way of finishing this book without being a reader of it, and that, dear reader, is (regrettably) merely wishful thinking.

This problem is remedied by adopting the denotation for \emph{no} in (\ref{ex6:newqden}a). That denotation prevents \emph{no reader} from being interpreted in its D-structure position, and instead requires that it be semantically interpreted in its S-structure position, as in~(\ref{ex6:hino1}).
\medskip
\begin{fullwidth}
\begin{exe}
\raggedright
	\ex \label{ex6:hino1}
	\footnotesize
	\begin{forest}
	[$\neg \exists x$ \textsc{reader}($x$) $\wedge$ $\exists \epsilon$ \textsc{finish}($\epsilon$) $\wedge$ \textsc{patient}(\denotes{this book})($\epsilon$) $\wedge$ \textsc{agent}($x$)($\epsilon$)\\TP
		[$\lambda Q\ \neg \exists x$ \textsc{reader}($x$) $\wedge$ $Q(x)$ \\DP [no reader, roof]]
		[$\lambda 2\ \exists \epsilon$ \textsc{finish}($\epsilon$) $\wedge$ \textsc{patient}(\denotes{this book})($\epsilon$) $\wedge$ \textsc{agent}($2$)($\epsilon$) \\\ibar{T} 
			[\obar{T} [pres]]
			[AspP 
			[$\exists \epsilon$ \textsc{finish}($\epsilon$) $\wedge$ \textsc{patient}(\denotes{this book})($\epsilon$) $\wedge$ \textsc{agent}($2$)($\epsilon$) \\\ibar{Asp}
				[\obar{Asp} ]
				[$\lambda \epsilon$ \textsc{finish}($\epsilon$) $\wedge$ \textsc{patient}(\denotes{this book})($\epsilon$) $\wedge$ \textsc{agent}(2)($\epsilon$) \\vP
					[\emph{t}_2 ]
					[$\lambda x\ \lambda \epsilon$ \textsc{finish}($\epsilon$) $\wedge$ \textsc{patient}(\denotes{this book})($\epsilon$) $\wedge$ \textsc{agent}($x$)($\epsilon$) \\\ibar{v}
						[\obar{v} [agent]]
						[VP [finish this book, roof] ]
					 ]
				 ]
			] ]
		]
	]
	\end{forest}
\end{exe}	
\end{fullwidth}
\medskip
This delivers the right meaning. It produces a sentences that says there is no $x$ which is a reader for which there is an event of finishing this book that has $x$ as agent.

The same problem arises for our earlier denotation for \emph{every} as well (i.e., (\ref{ex6:quantoldden}c)). This can be appreciated by considering (\ref{ex6:everylow1}), whose meaning according to (\ref{ex6:quantoldden}c) would be~(\ref{ex6:everylow1t}).%
%
\marginnote[10pt]{This D-structure assumes that \emph{every class} is the external argument of \emph{end}. It could well be, however, that \emph{end} is an unaccusative verb, and \emph{every class} is its internal argument. This doesn't affect the point being made.}%
%
\medskip
\begin{fullwidth}
\begin{exe}
\raggedright
	\ex \label{ex6:everylow1} Every class ended.
	\ex \label{ex6:everylow1t}
	\begin{forest}
	[$\exists \epsilon\ \forall x$ \textsc{class}($x$) $\rightarrow$ \textsc{end}($\epsilon$) $\wedge$ \textsc{theme}($x$)($\epsilon$)\\\ibar{Asp}
		[\obar{Asp} ]
		[$\lambda \epsilon\ \forall x$ \textsc{class}($x$) $\rightarrow$ \textsc{end}($\epsilon$) $\wedge$ \textsc{theme}($x$)($\epsilon$) \\vP
			[$\lambda Q\ \lambda \epsilon\ \forall x$ \textsc{class}($x$) $\rightarrow$ $Q(x)(\epsilon)$ \\DP [every class, roof] ]
			[$\lambda x\ \lambda \epsilon$ \textsc{end}($\epsilon$) $\wedge$ \textsc{theme}($x$)($\epsilon$)\\\ibar{v}
				[\obar{v} [theme]]
				[$\lambda \epsilon$ \textsc{end}($\epsilon$)\\VP [$\surd$end, roof] ]
			 ]
		 ]
	]
	\end{forest}
\end{exe}
\end{fullwidth}
\medskip
The denotation in (\ref{ex6:everylow1t}) says that there is a single ending event that every class is the theme of. Instead, what we want is each class to be able to have its own ending. If we adopt the new denotation for \emph{every} in (\ref{ex6:newqden}), however, the meaning in (\ref{ex6:everylow1t}) is blocked, and instead the denotation in (\ref{ex6:everylow2}) is derived.
\medskip
\begin{fullwidth}
	\begin{exe}
	\raggedright
		\ex \label{ex6:everylow2}
		\begin{forest}
		[$\forall x$ \textsc{class}($x$) $\rightarrow$ $\exists \epsilon$ \textsc{end}($\epsilon$) $\wedge$ \textsc{theme}($x$)($\epsilon$) \\TP
			[$\lambda Q\ \forall x$ \textsc{class}($x$) $\rightarrow\ Q(x)$ \\DP [every class, roof] ]
			[$\lambda 2\ \exists \epsilon$ \textsc{end}($\epsilon$) $\wedge$ \textsc{theme}($2$)($\epsilon$) \\\ibar{T}
				[\obar{T} [past]]
				[AspP 
				[$\exists \epsilon$ \textsc{end}($\epsilon$) $\wedge$ \textsc{theme}($2$)($\epsilon$) \\\ibar{Asp}
					[\obar{Asp} ]
					[$\lambda \epsilon$ \textsc{end}($\epsilon$) $\wedge$ \textsc{theme}($2$)($\epsilon$) \\vP
						[\emph{t_2} ]
						[$\lambda x\ \lambda \epsilon$ \textsc{end}($\epsilon$) $\wedge$ \textsc{theme}($x$)($\epsilon$) \\\ibar{v}
							[\obar{v} [theme]]
							[$\lambda \epsilon$ \textsc{end}($\epsilon$) \\VP [$\surd$end, roof] ]
						 ]
					 ]
				] 
				]
			 ]
		]
		\end{forest}
	\end{exe}
\end{fullwidth}
\medskip
As desired, (\ref{ex6:everylow2}) says that for all $x$ if $x$ is a class then there is an end-eventuality that it is the theme of.

These facts about the meaning of quantifiers are another indication that we don't want to force D-structures to be the only parse that is semantically interpreted. DPs headed by quantifiers shouldn't be semantically interpreted in their D-structure positions, as the cases in (\ref{ex6:lowno1}) and (\ref{ex6:everylow1}) indicate.

Summarizing, then, we've done three things to allow movement to make new meanings. We've let movement create something---a trace---in the position that the term being moved resides in. That trace employs a new syntactic formative, an index. This has been wedded to a semantic rule that changes the meaning of the phrase that is the sister to the moved term's new position. That change introduces a $\lambda$-expression that binds the trace. And finally, we've adjusted the denotations of quantificational determiners so that they can be interpreted in the positions they move into. This adjustment now prevents the DPs that are headed by quantifiers to be interpreted in their D-structure position. Thus, this last amendment has made D-structures that have quantificational DPs in them incapable of delivering a meaning. Movement is forced to interpret these~DPs.%
%
\marginnote[0pt]{This method of controlling where quantificational DPs are interpreted can be found in \citet{Chung:2004a}.}%
%
 
Note that the revision to movement in (\ref{ex6:tracemove}) puts traces into the syntax only when phrases move. This is because the syntax we've adopted for Head Movement is incompatible with traces. Consider, for instance, the result of moving the verb \emph{aime} to a tensed \obar{T} position in French. If this movement included traces, it would create a representation like that in~(\ref{ex6:notheadm}).
\begin{exe}
\raggedright
	\ex \label{ex6:notheadm}
	\footnotesize
	\begin{forest}
	[TP [\ibar{T}
		[\obar{T} 
			[\obar{V} [aime] ]
			[\obar{T} [present] ]
		 ]
		[$\lambda 2\ \lambda \epsilon$ $2(\epsilon)$ $\wedge$ \textsc{patient}(\denotes{Marie})($\epsilon$) $\wedge$ \textsc{agent}(\denotes{Jean})($\epsilon$) \\vP
			[DP [Jean, roof]] 
			[\ibar{v}
				[\obar{v} [agent] ]
				[VP [\ibar{V} 
					[\emph{t}_2]
					[DP [Marie, roof] ]
				]]
			]
		 ]
	]]
	\end{forest}
\end{exe}
The (\ref{ex6:tracemove}c) part of the new definition of movement causes vP to have a denotation that forces it to take the \emph{aime}+present=\obar{T} term as its argument. I have hidden from our view the meaning of tense, but there is compelling evidence that its denotation does not allow it to combine semantically with the denotation we have given to verbs. And even if there was a means of combining these two pieces, the resulting meaning should not be placed ``into'' the position marked by the variable ``$2$.'' For these reasons, I have prevented the new rule of movement from allowing a head to be semantically interpreted in any except its D-structure position. We will have an opportunity to revisit this decision before this chapter is over.

It is useful to have a way of representing in our syntactic parses which phrases have undergone the change in meaning that (\ref{ex6:tracemove}c) describes. That change has the effect of making the phrase that has been replaced by the trace, be the argument of the function that (\ref{ex6:tracemove}c) creates. Although technically not correct, this phrase is said to be the trace's binder. The standard way of representing which phrase has undergone the change that (\ref{ex6:tracemove}c) describes is to give the ``binder'' a subscript that matches the referential index on the trace. With this notation in place, (\ref{ex6:tracet1}) becomes~(\ref{ex6:tracet1co}).%
%
\marginnote[0pt]{Wherever convenient, I will suppress the presence of AspP in parses. I will reveal their presence only when I judge that doing otherwise would confuse.}%
%

\begin{exe}
\raggedright
	\ex \label{ex6:tracet1co}
	\footnotesize
	\begin{forest}
	[TP
		[TP 
			[DP_2 [Thuy, roof]]
			[\ibar{T}	
				[\obar{T} [will]]
				[vP 
					[$t_2$]				
					[\ibar{v}
						[\obar{v} [agent] ]
						[VP 
						[\ibar{V} 
							[\obar{V} [$\surd$visit]]
							[$t_3$]
						]]
				]]
			 ]
		]
		[DP_3 [\ibar{D}
			[\obar{D} [every] ]
			[NP [student, roof] ]
		]]
	]
	\end{forest}	
\end{exe}
The subscript on \textit{Thuy} indicates, then, that the movement rule has changed the denotation for \ibar{T} into a function that takes as its argument This notation gives us a compact way of indicating which traces are related to which moved phrases. 

Although this notation is standard, I find it confusing and will modify it in these lectures. We will see that DPs other than traces can have referential indices as part of their meaning. When we encounter these cases, it will be convenient to be able to represent them with the same subscript notation that we are using now with traces.  Some of these terms can also be binders; that is, we will want to be able to indicate that one term can have a referential index that is bound and simultaneously be a binder of another referential index. We will therefore want to have a different way of indicating what a term binds and what a term is bound by. The notation illustrated in (\ref{ex6:tracet1co}) to illustrate that \textit{Thuy} is a binder looks just like the way used to indicate that \textit{Thuy} has a referential index and therefore is bound. We should keep these distinct. We need different ways of indicating where a referential index is and where a binder is. For this reason, I will adopt the convention of indicating where the binder of a trace is by appending to it a superscript that matches the referential index on the trace.%
%
\marginnote[0pt]{This convention is introduced in \citet{Heim:1993}. That terms can be both variables and binders of variables is first discussed in \citet{Higginbotham:1980}.}%
%
Instead of (\ref{ex6:tracet1co}), I will use~(\ref{ex6:heimtracecon}).
\begin{exe}
\raggedright
	\ex \label{ex6:heimtracecon}
	\footnotesize
	\begin{forest}
	[TP
		[TP 
			[DP^2 [Thuy, roof]]
			[\ibar{T}	
				[\obar{T} [will]]
				[vP 
					[$t_2$]				
					[\ibar{v}
						[\obar{v} [agent] ]
						[VP 
						[\ibar{V} 
							[\obar{V} [$\surd$visit]]
							[$t_3$]
						]]
				]]
			 ]
		]
		[DP^3 [\ibar{D}
			[\obar{D} [every] ]
			[NP [student, roof] ]
		]]
	]
	\end{forest}
\end{exe}
Keep in mind that the superscripts on binders are not part of the syntax. They are merely a typographical diacritic that serves as a guide to the underlying semantics. In this respect they differ from referential indices, which are morphemes in a syntactic representation.

The machinery we've just reviewed for letting movement create semantically interpretable parses can now be exploited to give a treatment of pronouns when they get a bound variable interpretation. Let pronouns resemble traces in being capable of containing a referential index. We can think of pronouns as consisting of the features that express their person, number and, sometimes, gender information; when they are interpreted as bound variables, they will have a referential index as well. When pronouns are interpreted as referring, we can imagine that their person, number and gender features steer that reference. Informally, then, we can think of pronouns in the way that is portrayed by~(\ref{ex6:pronrefd}).
\begin{exe}
\raggedright
	\ex \label{ex6:pronrefd}
	\footnotesize
	\begin{forest}
		[\denotes{DP} [\ibar{D} [\obar{D} [sing\\ \nth{3}\\ fem]]]]
	\end{forest}
= the $x$ such that $x$ is singular and female and not the speaker or hearer.
\end{exe}
When a pronoun has an index, we can let the value delivered to that index have the role that \emph{the x} does in the formula in~(\ref{ex6:pronrefd}).
\begin{exe}
\raggedright
	\ex \label{ex6:pronvbld}
	\footnotesize
	\begin{forest}
		[\denotes{DP} [\ibar{D} [\obar{D} [sing\\\nth{3}\\ fem\\ 3]]]]
	\end{forest}
= $3$ and $3$ is singular and female and not the speaker or hearer. 
\end{exe}
I will adopt the convention of indicating when a pronoun contains a referential index by appending that referential index as a subscript to the pronoun, parallel to the representation that expresses traces. Thus, (\ref{ex6:pronvbld}) can be represented by:
\begin{exe}
\raggedright
	\ex \label{ex6:pronsub}
	her_3
\end{exe}

Let's take a look at how this will provide a way of deriving the bound variable interpretation of pronouns in~(\ref{ex6:sloprep1}). If we let \emph{her} come with a referential index that matches the one created by Argument movement, the S-structure representation for (\ref{ex6:sloprep1}) will be~(\ref{ex6:sloprept1}).
\medskip
\begin{fullwidth}
\begin{exe}
\raggedright
	\ex \label{ex6:sloprep1} No freshman praised her paper.
	\ex \label{ex6:sloprept1}
	\small
	\begin{forest}
	[$\neg \exists x$ \textsc{freshman}($x$) $\wedge$ $\exists \epsilon$ \textsc{praise}($\epsilon$) $\wedge$ \textsc{patient}($x$'s\_paper)($\epsilon$) $\wedge$ \textsc{agent}($x$)($\epsilon$) \\TP
		[$\lambda Q \neg \exists x$ \textsc{freshman}($x$) $\wedge$ $Q(x)$\\DP^2 [\ibar{D} 
			[\obar{D} [no] ]
			[NP [freshman, roof]]
		]]
		[$\lambda 2\ \exists \epsilon$ \textsc{praise}($\epsilon$) $\wedge$ \textsc{patient}(2's\_paper)($\epsilon$) $\wedge$ \textsc{agent}($2$)($\epsilon$)\\\ibar{T} 
			[\obar{T} [past] ]
			[AspP
			[\ibar{Asp}
			[\obar{Asp}]
			[vP 
				[{\emph{t}_2}]
				[\ibar{v}
					[\obar{v} [agent] ]
					[VP [\ibar{V}
						[\obar{V} [$\surd$praise] ]
						[DP
							[\obar{DP} [\ibar{D} [\obar{D} [her_2]]] ]
							[\ibar{D} [paper, roof] ]
						 ]
					] ]
				 ]
			 ]
		]
	]]]
	\end{forest}
\end{exe}
\end{fullwidth}
\medskip
Because movement of \emph{no freshman} has changed the meaning of \ibar{T} into a function that binds variables with the index ``2,'' the pronoun \emph{her} is construed as one of those variables when it includes that index. The fact that pronouns can be understood as variables corresponds to the proposal that pronouns can include indices as part of the ingredients that make them up. Those indices can be interpreted when movement has created a function that is positioned so that it binds that index.

Trace theory, then, has the useful consequence of providing a syntax for determining how, and where, pronouns can get their bound variable interpretation. Because the rule of movement changes the meaning of the phrase that the moved term becomes a sister to, it creates a meaning that can bind pronouns in that phrase. The movement rule not only creates a meaning for this phrase that allows it to bind the trace, it also creates a meaning that allows it to bind any pronoun that shares the referential index of that trace as well. Independently of providing an ingredient in the still-to-be-given solution to the problem of antecedent contained anaphora, Trace theory gives a model of how pronouns can get their bound variable interpretation.

Trace theory also constrains derivations in a way that removes the problems that arise by adopting the Lebeaux model, and its abandonment of making D-structure the only semantically interpretable element of a derivation. Recall that by removing the requirement that D-structures be the only element of the derivation that is semantically interpreted, the Lebeaux model allows D-structures to violate the Theta Criterion. This wrongly allows derivations like that sketched in~(\ref{ex6:baddertc}).
\medskip
\begin{fullwidth}
\begin{exe}
	\raggedright
		\ex \label{ex6:baddertc}
		\small
		\begin{forest}
		[TP 
		[\ibar{T}
			[\obar{T} [will] ]
			[vP
				[DP [they, name=binder, roof] ]
				[\ibar{v}
					[\obar{v} [experiencer] ]
					[VP 
					[V [$\surd$believe]]
					[TP
						[{\ }, name=trace ]
						[\ibar{T}
							[\obar{T} [to] ]
							[VP [\ibar{V}
								[\obar{V} [$\surd$seem] ]
								[CP [that pigs fly, roof] ]
							] ]
						 ]
					 ]
				 ]
			 ]
		]]]
		\draw[->,>=stealth', thin, gray] (trace) to[out=south west,in=south,looseness=1.4] node[pos=0.75,xshift=-10mm]{} (binder);
		\end{forest}
		\hspace*{-25pt}
		\begin{forest}
			[TP 
			[DP [they, name=binder, roof] ]
			[\ibar{T}
				[\obar{T} [will] ]
				[vP
					[{\ }, name=trace]
					[\ibar{v}
						[\obar{v} [experiencer] ]
						[VP
						[V [$\surd$believe]]
						[TP
							[{\ }]
							[\ibar{T}
								[\obar{T} [to] ]
								[VP [\ibar{V}
									[\obar{V} [$\surd$seem] ]
									[CP [that pigs fly, roof] ]
								] ]
							 ]
						 ]
					 ]
				 ]
			]]]
			\draw[->,>=stealth', thin, gray] (trace) to[out=south west,in=south,looseness=1.4] node[pos=0.75,xshift=-10mm]{} (binder);
			\end{forest}
	\end{exe}
\end{fullwidth}
\medskip
Without Trace theory, the derivation in (\ref{ex6:baddertc}) provides a parse---the one created by the first movement of \emph{they}---that can be semantically interpreted. All the other conditions on derivations, e.g., Upwards, the requirements of bullets and \textsc{agree}, are satisfied as well. Unless this derivation is blocked, our system will characterize \emph{They will believe to seem that pigs fly} as grammatical.

Trace theory closes this problem. Because it requires movement of phrases to leave a trace, and because traces can only be semantically interpreted as variables bound by the $\lambda$-operator that movement invokes, derivations like that in (\ref{ex6:baddertc}) will be blocked. The reason is that the trace left by the movement shown in the first step of (\ref{ex6:baddertc}) cannot get a semantic interpretation. Consider, for instance, what would happen if we let the S-structure, roughly equivalent to the output of the second movement in (\ref{ex6:baddertc}), be subject to our semantic rules.
\medskip
\begin{fullwidth}
\begin{exe}
\raggedright
	\ex \label{ex6:badt}
	\small
	\begin{forest}
		[TP 
		[DP [they, name=binder, roof] ]
		[\ibar{T}
			[\obar{T} [will] ]
			[vP
				[\emph{t_2^3}]
				[\ibar{v}
					[\obar{v} [experiencer] ]
					[VP
					[V [$\surd$believe]]
					[??\\TP
						[\emph{t_3}]
						[$\exists \epsilon$ \textsc{seem}($\epsilon$) $\wedge$ \textsc{cause}(\denotes{that pigs fly})($\epsilon$) \\\ibar{T}
							[\obar{T} [to] ]
							[$\exists \epsilon$ \textsc{seem}($\epsilon$) $\wedge$ \textsc{cause}(\denotes{that pigs fly})($\epsilon$) \\AspP [\ibar{Asp}
							[\obar{Asp} [asp]]
							[$\lambda \epsilon$ \textsc{seem}($\epsilon$) $\wedge$ \textsc{cause}(\denotes{that pigs fly})($\epsilon$) \\VP [\ibar{V}
								[\obar{V} [$\surd$seem] ]
								[CP [that pigs fly, roof] ]
							] ]
						 ]
					 ]
				 ]
			 ]
		]]]]]
		\end{forest}%
		%
		\marginnote[0pt]{We have not developed a semantics that models how verbs take clausal complements. I've somewhat arbitrarily---and probably incorrectly---given a semantics here that uses $\theta$-roles to connect clauses to the verbs they are arguments of.}%
		%
\end{exe}
\end{fullwidth}
\medskip
The complement TP to \emph{$\surd$believe} has no denotation, as there is no semantic rule that can compose the meaning of the trace in Specifier TP with the denotation of its sister:~\ibar{T}. This problem will arise in every parse of the derivation; therefore, no parse in this derivation will be semantically interpretable. The ungrammaticality of this derivation, then, derives from the fact that Trace theory blocks a semantic interpretation of the sentence the derivation models.

In general terms, Trace theory requires every position a moved XP occupies in a derivation to be a position in which it, or its trace, can be semantically interpreted. This requires that the XPs in a D-structure are either semantically interpretable in those positions, or could be replaced by traces that are semantically interpretable in those positions. That is sufficient to ensure that the Theta Criterion holds of D-structures. To my knowledge, it is sufficient to ensure that the constraints on D-structures that trace back to the semantics are preserved in the Lebeaux Model.

Trace theory has three useful ingredients, then. It allows more parses in a derivation to be relevant to the semantic interpretation of a sentence without losing the effects of restricting semantic interpretability to D-structures. It provides a method for allowing pronouns to be understood as bound variables, by allowing movement to change the meaning of the phrase that contains a pronoun. And when this is wedded to the existence of traces, it provides a way of interpreting structures that are formed by movement. 

This last effect brings us a step closer to understanding how to defease the problem posed by cases of antecedent contained anaphora. Our syntax now countenances a syntax for those examples in which the focus anaphor is not inside its focus antecedent, but is semantically interpretable. Those syntaxes arise by way of Complex NP Movement in the way that (\ref{ex6:tracet1den3}) illustrates. The second step in accounting for antecedent contained anaphora requires understanding the syntax of relative clauses. As it happens, relative clauses independently illustrate the need for a syntax that allows movement to create interpretable structures. We turn now to that task.

\subsubsection{Relative Clauses and Wh Movement}
\label{sec:relclause}

Let's start with (\ref{ex6:rel1}) as our model for building an analysis of relative clauses.
\begin{exe}
\raggedright
	\ex \label{ex6:rel1}
	No woman who you talked to will leave.
\end{exe} 
The relative clause is the portion of this sentence that corresponds to the string \emph{who you talked to}. This expression has a form that is parallel to the ``object'' argument of the verb \emph{know} in~(\ref{ex6:relw}).
\begin{exe}
\raggedright
	\ex \label{ex6:relw}
	She knows who you talked to.
\end{exe}
At present, our syntax cannot form either of these types of clauses. Because their form is so close, it makes sense to develop a syntax that takes into consideration both relative clauses and complement clauses of the form in~(\ref{ex6:relw}). 

If we assume that the word \emph{who} is a kind of pronoun, our phrase structure rules can create a clause like that in~(\ref{ex6:rel2}).
\begin{exe}
\raggedright
	\ex \label{ex6:rel2}
	\small
	\begin{forest}
	[CP[\ibar{C}
		[\obar{C} [??]]
		[TP 
			[DP^2 [\ibar{D} [\obar{D} [you]]]]
			[\ibar{T} 
				[\obar{T} [past] ]
				[vP
					[\emph{t}_2 ]
					[\ibar{v} 
						[\obar{v} [agent] ]
						[VP [\ibar{V}
							[\obar{V} [$\surd$talk] ]
							[PP [\ibar{P}
								[\obar{P} [to] ]
								[DP [\ibar{D} [\obar{D} [who]]] ]
							] ]
						] ]
					]
				 ]
			]
		]
	]]
	\end{forest}%
	%
	\marginnote[0pt]{I will suppress the existence of AspP in parses where its existence doesn't matter, as in (\ref{ex6:rel2}).}%
	%
\end{exe}
Let's assume that this is the D-structure representation for the expression \emph{who you talked to} in both (\ref{ex6:rel1}) and~(\ref{ex6:relw}). To achieve the word order that we see in (\ref{ex6:rel1}) and~(\ref{ex6:relw}), we'll formulate a movement rule.

When this CP is the complement to verbs like \emph{know}, it has a semantics that relates it to what a question is. The semantics for (\ref{ex6:relw}) should cause \emph{who you talked to} to characterize a question that (\ref{ex6:relw}) says she knows the answer to. Similarly, in (\ref{ex6:relqex}) the complement should characterize a question that she wonders about, asks, explains, etc.
\begin{exe}
\raggedright
	\ex \label{ex6:relqex}
	\begin{xlist}
		\ex She wonders who you talked to.
		\ex She asked who you talked to.
		\ex She explained who you talked to.
	\end{xlist}
\end{exe}
This isn't the meaning associated with relative clauses. It must be associated, then, with something in these clauses that is not entirely contained in the meaning of the words \emph{who you talked to}. A standard hypothesis is that there is a morpheme in the \obar{C} position that is responsible for giving these CPs their question-like meaning. I'll represent this morpheme with ``Q,'' and I'll assume that it comes with a feature that is an agreement relationship with the phrase that moves. If we call this feature ``wh,'' then the non-relative clause form of (\ref{ex6:rel2}) is~(\ref{ex6:reltq}).
\begin{exe}
\raggedright
	\ex \label{ex6:reltq}
	\footnotesize
	\begin{forest}
	[CP[\ibar{C}
		[\obar{C} [Q\\ $\bullet$wh ]]
		[TP 
			[DP^2 [\ibar{D} [\obar{D} [you]]]]
			[\ibar{T} 
				[\obar{T} [past] ]
				[vP
					[\emph{t}_2 ]
					[\ibar{v} 
						[\obar{v} [agent] ]
						[VP [\ibar{V}
							[\obar{V} [$\surd$talk] ]
							[PP [\ibar{P}
								[\obar{P} [to] ]
								[DP [\ibar{D} [\obar{D} [who\\ \ding{51}\emph{wh}]]] ]
							] ]
						] ]
					]
				 ]
			]
		]
	]]
	\end{forest}
\end{exe}
If there is a $\bullet$ associated with the wh-feature, then our syntax requires a derivation that leads to the correct surface form. \textsc{Agree} will hold between \obar{C} and \emph{who}, and the $\bullet$'s requirements can be satisfied if this is paired with the movement shown in~(\ref{ex6:questmove}).
\begin{exe}
\raggedright
	\ex \label{ex6:questmove}
	\footnotesize
	\begin{forest}
	[CP
		[DP^3 [who\\ \ding{51}\emph{wh}, name=binder, roof]]
		[\ibar{C} 
			[\obar{C} [Q\\ $\bullet$\emph{wh}] ]
			[TP
				[DP^2 [she, roof]]
				[\ibar{T}
					[\obar{T} [past] ]
					[vP
						[\emph{t_2} ]
						[\ibar{v}
							[\obar{v} [agent] ]
							[VP [\ibar{V}
								[\obar{V} [$\surd$talk] ]
								[PP [\ibar{P}
									[\obar{P} [to] ]
									[{t_3}, name=trace ]
								] ]
							] ]
						 ]
					 ]
				 ]
			 ]
		 ]
	]
	\draw[->,>=stealth', ultra thin, gray] (trace) to[out=south,in=-80,looseness=1.2] node[pos=0.75,xshift=-10mm]{} (binder);
	\end{forest}
	\vspace*{-50pt}
\end{exe}
This is a way of using the methods developed here to deliver a syntax for these~CPs.

The class of phrases that can host the wh-feature, and therefore satisfy the bulleted wh-feature in \obar{C} position, all contain a specialized lexical item. In addition to the pronoun \emph{who}, there are other pro-forms, the determiner \emph{which} and the degree word~\emph{how}. Some examples are in~(\ref{ex6:wnphrases}).
\begin{exe}
\raggedright
	\ex \label{ex6:wnphrases} She knows \ldots
	\begin{xlist}
		\ex what/where/when/why you danced.
		\ex which tango you danced.
		\ex how fast you danced.
	\end{xlist}
\end{exe}
Because it is the existence of these words that correlates with the question-like and relative clause constructions, the wh-feature should reside in them. Thus, the D-structures for these constructions should have the form in~(\ref{ex6:pied1}).
\begin{fullwidth}
\begin{exe}
\raggedright
	\ex \label{ex6:pied1}
	\footnotesize
	\begin{forest}
		[CP
			[\obar{C} [Q\\ $\bullet$\emph{wh}]]
			[TP 
				[DP^2 [you, roof] ]
				[\ibar{T} 
					[\obar{T} [past] ]
					[vP
						[\emph{t_2} ]
						[\ibar{v} 
							[\obar{v} [agent] ]
							[VP [\ibar{V}
								[\obar{V} [$\surd$dance] ]
								[DP, circle, draw [\ibar{D} [\obar{D} [what\\ \ding{51}\emph{wh}]]] ]
							] ]
						]
					 ]
				]
			]
		]
		\end{forest}
	\begin{forest}
	[CP
		[\obar{C} [Q\\ $\bullet$\emph{wh}]]
		[TP 
			[DP^2 [you, roof] ]
			[\ibar{T} 
				[\obar{T} [past] ]
				[vP
					[\emph{t_2} ]
					[\ibar{v} 
						[\obar{v} [agent] ]
						[VP [\ibar{V}
							[\obar{V} [$\surd$dance] ]
							[DP, circle, draw [\ibar{D}
								[\obar{D} [which\\ \ding{51}\emph{wh}] ]
								[NP [tango, roof] ]
							] ]
						] ]
					]
				 ]
			]
		]
	]
	\end{forest}
	\begin{forest}
		[CP
			[\obar{C} [Q\\ $\bullet$\emph{wh}]]
			[TP 
				[DP^2 [you, roof] ]
				[\ibar{T} 
					[\obar{T} [past] ]
					[vP
						[\emph{t_2} ]
						[\ibar{v} 
							[\obar{v} [agent] ]
							[VP [\ibar{V}
								[\obar{V} [$\surd$dance] ]
								[PP, circle, draw [\ibar{P} [\obar{P} [where\\ \ding{51}\emph{wh}]]] ]
							] ]
						]
					 ]
				]
			]
		]
		\end{forest}
		\begin{forest}
			[CP
				[\obar{C} [Q\\ $\bullet$\emph{wh}]]
				[TP 
					[DP^2 [you, roof] ]
					[\ibar{T} 
						[\obar{T} [past] ]
						[vP
							[\emph{t_2} ]
							[\ibar{v} 
								[\obar{v} [agent] ]
								[VP [\ibar{V}
									[\obar{V} [$\surd$dance] ]
									[DegP, circle, draw [\ibar{Deg}
										[\obar{Deg} [how\\ \ding{51}\emph{wh}] ]
										[AP [fast, roof] ]
									]]
								] ]
							]
						 ]
					]
				]
			]
			\end{forest}
\end{exe}
\end{fullwidth}
\bigskip
The \textsc{agree} relation holds between \obar{C} and either \obar{D} or \obar{Deg}. The movement relation, however, never involves movement of the heads, but instead always affects a phrase, in (\ref{ex6:pied1}) they are circled, that contains that head. 

In the cases we've seen so far, the phrase that moves is always an \ibar{X} projection of the word that has the \emph{wh}-feature. But this isn't always the case. There are examples---(\ref{ex6:pied2}) illustrates some---in which the phrase that moves is not a projection of the word containing the \emph{wh}-feature.
\pagebreak
\begin{exe}
\raggedright
	\ex \label{ex6:pied2} She asked
	\begin{xlist}
		\ex {}[\PP to whom] you talked.
		\ex {}[\DP how long a race] you ran.
		\ex {}[\DP whose mother] you talked to.
		\ex {}[\DP whose mother's son] you talked to.
	\end{xlist}
\end{exe}
\Citet{Ross:1967} dubbed this phenomenon ``pied-piping.'' The phrases that pied-pipe are known as ``\emph{wh}-phrases'' and the movement operation responsible for relocating \emph{wh}-phrases is called \emph{Wh} Movement.%

\marginnote[0pt]{Ross used the term ``pied-piping'' just for examples in which the phrase that is moved is not an \ibar{X}-projection of the term that has the \emph{wh}-feature. I'll extend this term to all cases in which what moves is a phrase containing the term with the \emph{wh}-feature.}%

Relative clauses employ \emph{Wh} Movement, just like the question-like CPs we've been examining, but there are some differences with respect to the shape of the \emph{wh}-phrases involved. In general, with the exception of \emph{what}, pro-forms can be found in relative clauses (see~(\ref{ex6:relpron})).
\begin{exe}
\raggedright
	\ex \label{ex6:relpron} 
	\begin{xlist}
		\ex[]{the guy who you talked to}
		\ex[]{the place where you danced}
		\ex[]{the time when you left}
		\ex[]{the reason why you left}
		\ex[*]{the thing what you bought}
	\end{xlist}
\end{exe}
Pro-forms can pied-pipe in basically the same way in relative clauses as they do in question~CPs.
\begin{exe}
\raggedright
	\ex \label{ex6:relpron2}
	\begin{xlist}
		\ex the guy to whom you talked
		\ex the lady whose mother you visited
		\ex the child whose mother's hat you stole
	\end{xlist}
\end{exe}

But degree words and the phrases they pied-pipe cannot be involved in relative clauses at all.
\begin{exe}
\raggedright
	\ex \label{ex6:reldeg}
	\begin{xlist}
		\ex[*]{the distance how far you ran}
		\ex[*]{the race how long a trail you ran}
	\end{xlist}
\end{exe}
And when DPs headed by \emph{which} are pied-piped, the NP portion of that DP must be silent and understood to be the same as the NP to which the relative clause is attached. So whereas (\ref{ex6:whichrel1}) is ungrammatical, the form in (\ref{ex6:whichrel2}) is its grammatical variant.
\pagebreak
\begin{exe}
\raggedright
	\ex[*]{\label{ex6:whichrel1} the dog which dog you walked}
	\ex \label{ex6:whichrel2}
	\small
	\begin{forest}
	[DP[\ibar{D}
		[\obar{D} [the]]
		[NP [\ibar{N}
			[\ibar{N} [dog, roof] ]
			[CP
				[DP [\ibar{D}
					[\obar{D} [which] ]
					[NP [\sout{dog}] ]
				] ]
				[\ibar{C}
					[\obar{C} [??] ]
					[TP
						[DP [you, roof] ]
						[\ibar{T} [walked, roof] ]
					 ]
				 ]
			 ]
		]]
	]]
	\end{forest}
\end{exe}

Interestingly, relative clauses have recourse to a wh-word that is not found in the question-like CPs that use \emph{Wh} Movement. This wh-word has the behavior of a pro-form, but is phonologically silent. I'll represent it with ``\emph{wh}.'' An example is in~(\ref{ex6:silrel}).
\begin{exe}
\raggedright
	\ex \label{ex6:silrel}
 the dog you brought 
 \hspace*{12pt}
	\footnotesize
	\begin{forest}
	[DP[\ibar{D}
		[\obar{D} [the]]
		[NP [\ibar{N}
			[\ibar{N} [dog, roof]]
			[CP
				[DP^3 [\ibar{D} [\obar{D} [\emph{wh}, name=binder]]] ]
				[\ibar{C}
					[\obar{C} [??] ]
					[TP 
						[DP^2 [you, roof]]
						[\ibar{T}
							[\obar{T} [past] ]
							[vP
								[\emph{t_2} ]
								[\ibar{v}
									[\obar{v} [agent] ]
									[VP [\ibar{V}
										[\obar{V} [$\surd$bring] ]
										[\emph{t_3}, name=trace ]
									] ]
								 ]
							 ]
						 ]
					 ]
				 ]
			 ]
		]]
	]]
	\draw[->,>=stealth', ultra thin, gray] (trace) to[out=south,in=south,looseness=1.5] node[pos=0.75,xshift=-10mm]{} (binder);
	\end{forest}
\end{exe}
This silent pro-form can be found in forming relative clauses that correspond to PPs as well.
\begin{exe}
\raggedright
	\ex \label{ex6:whpp}
	\begin{xlist}
		\ex the time wh^3 you came \emph{t_3}
		\ex the place wh^3 you left \emph{t_3}
		\ex the reason wh^3 you left \emph{t_3}
	\end{xlist}
\end{exe}
Strangely, in many Englishes it cannot be moved from the most local subject position, as in~(\ref{ex6:localsjb}).\marginnote{Sentences such as these are grammatical in Appalachian, African American, and Hiberno Englishes. See \citet{Sistrunk:2012} and \citet{Doherty:1993}.}
\begin{exe}
\raggedright
	\ex[*]{\label{ex6:localsjb} the dog is barking}
	\sn \emph{compare:}
	\ex[]{the dog which is barking}
\end{exe}

In contexts where the silent \emph{wh} is involved in building a relative clause, the up-till-now mysterious identity of the clause's head is revealed. Alongside (\ref{ex6:silrel}) is~(\ref{ex6:silrelt}).
\bigskip
\begin{fullwidth}
\begin{exe}
	\raggedright
		\ex \label{ex6:silrelt}
	 the dog that you brought 
	 \hspace*{5pt}
		\footnotesize
		\begin{forest}
		[DP[\ibar{D}
			[\obar{D} [the]]
			[NP [\ibar{N}
				[\ibar{N} [dog, roof]]
				[CP
					[DP^3 [\ibar{D} [\obar{D} [\emph{wh}, name=binder]]] ]
					[\ibar{C}
						[\obar{C} [that] ]
						[TP 
							[DP^2 [you, roof]]
							[\ibar{T}
								[\obar{T} [past] ]
								[vP
									[\emph{t_2} ]
									[\ibar{v}
										[\obar{v} [agent] ]
										[VP [\ibar{V}
											[\obar{V} [$\surd$bring] ]
											[\emph{t_3}, name=trace ]
										] ]
									 ]
								 ]
							 ]
						 ]
					 ]
				 ]
			]]
		]]
		\end{forest}
	\end{exe}
\end{fullwidth}
\bigskip
The complementizer \emph{that} has the option of being unpronounced, as can be seen from the alternation in~(\ref{ex6:compdel1}).
\begin{exe}
\raggedright
	\ex \label{ex6:compdel1}
	\begin{xlist}
		\ex Katia knows that Petr loves his family.
		\ex Katia knows Petr loves his family.
	\end{xlist}
\end{exe}
And this option is enforced, it seems, when there is spoken material in the Specifier of the phrase that \emph{that} projects:
\begin{exe}
\raggedright
	\ex \label{ex6:dfcf}
	\begin{xlist}
		\ex[*]{the dog which that you brought}
		\ex[*]{the woman who that left}
		\ex[*]{the time when that you left}
	\end{xlist}
\end{exe}
\Citet{Chomsky:1977a} called this generalization the ``Doubly Filled Comp filter.''
\begin{exe}
\raggedright
	\ex \label{ex6:dfcff} Doubly-filled Comp Filter\\
	The Specifier and head of a CP cannot both contain phonological material.
\end{exe}

Relative clauses can be seen as another case that illustrates the necessity of letting movement create meanings. To see this, consider the pieces of the DP in (\ref{ex6:restrr1}) whose denotations are identified.
\begin{exe}
\raggedright
	\ex \label{ex6:restrr1} no dog that you brought
	\sn
	\small
	% \footnotesize
	\begin{forest}
	[$\lambda Q\ \neg \exists x$ \textsc{dog}($x$) $\wedge$ \textsc{you\_brought}($x$) $\wedge Q(x)$\\ DP [\ibar{D}
		[$\lambda P\ \lambda Q\ \neg \exists x\ P(x) \wedge Q(x)$  \\\obar{D} [no]]
		[NP [\ibar{N}
			[$\lambda y$ \textsc{dog}($y$)\\\ibar{N} [dog, roof]]
			[CP
				[DP^3 [\ibar{D} [\obar{D} [\emph{wh}, name=binder]]] ]
				[\ibar{C}
					[\obar{C} [that] ]
					[TP 
						[DP^2 [you, roof]]
						[\ibar{T}
							[\obar{T} [past] ]
							[vP
								[\emph{t_2} ]
								[\ibar{v}
									[\obar{v} [agent] ]
									[VP [\ibar{V}
										[\obar{V} [$\surd$bring] ]
										[\emph{t_3}, name=trace ]
									] ]
								 ]
							 ]
						 ]
					 ]
				 ]
			 ]
		]]
	]]
	\end{forest}
\end{exe}
We have taken the denotation of common nouns to be functions from things to truth values, and this informs the denotation I have assigned to the \ibar{N} that contains the common noun \emph{dog}. The denotation for \emph{no} is a relation between the functions that the NP it combines with is, and the denotation of the constituent that follows the DP \emph{no} heads. As the aspirational denotation assigned to the DP in (\ref{ex6:restrr1}) indicates, we want to be able to semantically combine the denotation of \emph{dog} with the denotation of the relative clause so that it forms a meaning for the NP that they together comprise, so that the first argument for~\emph{no} is created. Let's take a closer look at how the relative clause's meaning is derived from the pieces that we presently have.
\medskip
\begin{fullwidth}
\begin{exe}
\raggedright
	\ex \label{ex6:reloned}
	\small
	\begin{forest}
	[CP
		[DP [\emph{wh}, roof]]
		[$\lambda 3\ \exists \epsilon'$ \textsc{bring}($\epsilon'$) $\wedge$ \textsc{theme}($3$)($\epsilon'$) $\wedge$ \textsc{agent}(\denotes{you})($\epsilon'$)\\ \ibar{C} 
			[\obar{C} [that]]
			[$\exists \epsilon'$ \textsc{bring}($\epsilon'$) $\wedge$ \textsc{theme}($3$)($\epsilon'$) $\wedge$ \textsc{agent}(\denotes{you})($\epsilon'$)\\TP 
				[DP [you, roof]]
				[$\lambda 2\ \exists \epsilon'$ \textsc{bring}($\epsilon'$) $\wedge$ \textsc{theme}($3$)($\epsilon'$) $\wedge$ \textsc{agent}($2$)($\epsilon'$) \\ \ibar{T}
					[\obar{T} [past] ]
					[$\lambda \epsilon'$ \textsc{bring}($\epsilon'$) $\wedge$ \textsc{theme}($3$)($\epsilon'$) $\wedge$ \textsc{agent}($2$)($\epsilon'$)\\vP
						[\emph{t_2} ]
						[$\lambda x\ \lambda \epsilon'$ \textsc{bring}($\epsilon'$) $\wedge$ \textsc{theme}($3$)($\epsilon'$) $\wedge$ \textsc{agent}($2$)($\epsilon'$) \\ \ibar{v}
							[\obar{v} [agent] ]
							[$\lambda \epsilon$ \textsc{bring}($\epsilon$) $\wedge$ \textsc{theme}($3$)($\epsilon$)\\VP [$\surd$bring \emph{t_3}, roof] ]
						]
					 ]
				 ]
			 ]
		]
	]
	\end{forest}
\end{exe}
\end{fullwidth}
\medskip
I've continued to ignore the contribution that tense makes, and I've made the assumption that the complementizer \emph{that} has no denotation either. 

Note that because of the meaning that movement creates, movement of the silent \emph{wh} phrase causes \ibar{C} to be a function from things to a sentence meaning. It maps things to Truth just in case those things are brought by you. Suppose that this is also the denotation of the relative clause; that is, suppose that the silent \emph{wh} phrase also has no denotation. This would give (\ref{ex6:restrr1}) the pieces in~(\ref{ex6:restrr2}).
\medskip
\begin{fullwidth}
\begin{exe}
\raggedright
	\ex \label{ex6:restrr2}
	\small
	\begin{forest}
	[DP [\ibar{D}
		[$\lambda P\ \lambda Q\ \neg \exists x\ P(x) \wedge Q(x)$  \\\obar{D} [no]]
		[NP [\ibar{N}
			[$\lambda y$ \textsc{dog}($y$)\\\ibar{N} [dog, roof]]
			[$\lambda 3\ \exists \epsilon'$ \textsc{bring}($\epsilon'$) $\wedge$ \textsc{theme}($3$)($\epsilon'$) $\wedge$ \textsc{agent}(\denotes{you})($\epsilon'$) \\CP [that you brought, roof]]
		]]
	]]
	\end{forest}
\end{exe}
\end{fullwidth}
\medskip
To bring the \ibar{N} and CP together in the necessary way, we use a rule of semantic composition rule called Predicate Modification.
\begin{exe}
\raggedright
	\ex \label{ex6:predmod} Predicate Modification\\
	If \denotes{$\alpha$} is $\lambda z P(z)=T$ and \denotes{$\beta$} is $\lambda y Q(y)=T$, then the denotation of their immediate mother is $\lambda x P(x)$ $\wedge\ Q(x)=T$ 
\end{exe}
This produces~(\ref{ex6:restrr3}).
\medskip
\begin{fullwidth}
\begin{exe}
\raggedright
	\ex \label{ex6:restrr3}
	\small
	\begin{forest}
	[$\lambda Q\ \neg\exists x$ \textsc{dog}($x$) $\wedge$ $\exists \epsilon'$ \textsc{bring}($\epsilon'$) $\wedge$ \textsc{theme}($x$)($\epsilon'$) $\wedge$ \textsc{agent}(\denotes{you})($\epsilon'$) $\wedge$ $Q(x)$ \\DP [$\lambda Q\ \neg\exists x$ \textsc{dog}($x$) $\wedge$ $\exists \epsilon'$ \textsc{bring}($\epsilon'$) $\wedge$ \textsc{theme}($x$)($\epsilon'$) $\wedge$ \textsc{agent}(\denotes{you})($\epsilon'$) $\wedge$ $Q(x)$ \\\ibar{D}
		[$\lambda P\ \lambda Q\ \neg \exists x\ P(x) \wedge Q(x)$  \\\obar{D} [no]]
		[$\lambda x$ \textsc{dog}($x$) $\wedge$ $\exists \epsilon'$ \textsc{bring}($\epsilon'$) $\wedge$ \textsc{theme}($x$)($\epsilon'$) $\wedge$ \textsc{agent}(\denotes{you})($\epsilon'$)\\NP [\ibar{N}
			[$\lambda y$ \textsc{dog}($y$)  \\\ibar{N} [dog, roof]]
			[$\lambda 3\ \exists \epsilon'$ \textsc{bring}($\epsilon'$) $\wedge$ \textsc{theme}($3$)($\epsilon'$) $\wedge$ \textsc{agent}(\denotes{you})($\epsilon'$) \\CP [that you brought, roof]]
		]]
	]]
	\end{forest}
\end{exe}
\end{fullwidth}
\medskip
This is the meaning we are after.

If these steps are correct, what we see is that movement has created a meaning that allows CPs to combine with \ibar{N}s in the way required to derive the meaning of the constituent that they form. It's in this sense that relative clauses provide evidence that we want to let movement create new meanings. It is because movement is able to create a function from things to truth values that the CPs which are relative clauses have a denotation that allows them to combine with nominal projections. Trace theory makes movement creates responsible for converting CPs into functions of the right kind to compose by Predicate Modification. Trace theory, then, explains why movement is found in relative clauses.

At least, Trace Theory does this in simple relative clauses. It doesn't generalize to relative clauses in which a more complex \emph{wh}-phrase has moved, however. In cases like (\ref{ex6:piedpden}), for instance, the moved \emph{wh}-phrase makes a semantic contribution, , unlike what we see for the silent wh-phrase in (\ref{ex6:reloned}).
\begin{exe}
\raggedright
	\ex \label{ex6:piedpden}
	the dog whose owner^2 you met \emph{t}_2
\end{exe}
The meaning of the moved wh-phrase should make the relative clause describe properties that the dog has---namely, that it has an owner that you met---while also describing a property that the owner has---namely, being met. The Trace theory of Movement doesn't do that in (\ref{ex6:piedpden}). The conclusion that movement creates new meanings remains. The fact that \emph{whose owner} has moved in (\ref{ex6:piedpden}) that is responsible for allowing the CP to combine semantically with \emph{dog}. And it delivers the fact that the owner is understood to be met. But it doesn't capture how \denotes{dog} is modified by the relative clause. Our present syntax and semantics doesn't yet deliver the right outcome for (\ref{ex6:piedpden}). Let's set aside these cases for the moment, and return to the job at hand: solving the problem of antecedent contained anaphora.%
%
\marginnote[0pt]{A solution to the problem posed by (\ref{ex6:piedpden}) can be found in the Copy theory of movement, section~\ref{sec:piedpipecopies}.}%
%

\subsubsection{Antecedent Contained Anaphora Resolved}
\label{sec:acdresolved}

The puzzle we're in pursuit of solving is how the focus anaphor in (\ref{ex6:acda}), repeated here, can be resolved.
\begin{exe}
	\raggedright
	\exr{ex6:acda}
	\begin{xlist}
		\ex Thuy will visit every student who [RODICA <will visit>].
		\ex Thuy will visit every student who [RODICA <will \sout{visit}>].
	\end{xlist}
\end{exe}
The bracketed portion of these sentences is a focus anaphor, whose antecedent is found in the sentence that seemingly contains the focus antecedent. Focus anaphors, recall, are subject to the requirement in~(\ref{ex6:fanaphora}).
\begin{exe}
	\raggedright
	\exr{ex6:fanaphora} Let XP be a focus anaphor, and $\alpha$ be the focus-marked term within XP. A \underline{focus value} of XP (\denotes{XP}^{F}) is identical to \denotes{XP}, except that $\alpha$ is replaced by a \underline{focus alternative} of $\alpha$ (=\textsc{fa}($\alpha$)). $\beta$ is a \textsc{fa}($\alpha$) if it is the same syntactic type as~$\alpha$.\\[4pt]
	If XP is a focus anaphor, then there must be a YP and a \denotes{XP}^F such that \denotes{YP}=\denotes{XP}^F.
\end{exe}
The puzzle, then, is giving (\ref{ex6:acda}) a syntax that allows (\ref{ex6:fanaphora}) to be satisfied. 

The Trace theory of movement does that. With Trace theory, the S-structure representation for (\ref{ex6:acda}) is~(\ref{ex6:acdres}).
\medskip
\begin{fullwidth}
\begin{exe}
	\raggedright
	\ex \label{ex6:acdres}
	\footnotesize
\begin{forest}
[$\forall x$ {[}\textsc{student}($x$) $\wedge$ $\exists \epsilon$ {[}\textsc{visit}($\epsilon$) $\wedge$ \textsc{patient}($x$)($\epsilon$) $\wedge$ \textsc{agent}(\denotes{Rodica})($\epsilon$){]]} $\rightarrow$ {[}$\exists \epsilon'$ \textsc{visit}($\epsilon'$) $\wedge$ \textsc{patient}($x$)($\epsilon'$) $\wedge$ \textsc{agent}(\denotes{Thuy})($\epsilon'$) {]}\\TP
	[$\lambda 5\ \exists \epsilon'$ \textsc{visit}($\epsilon'$) $\wedge$ \textsc{patient}($5$)($\epsilon'$) $\wedge$ \textsc{agent}(\denotes{Thuy})($\epsilon'$)\\\fbox{TP}
		[DP [Thuy, roof]]
		[\ibar{T}
			[\obar{T} [will] ]
			[vP
				[\emph{t_4} ]
				[\ibar{v}
					[\obar{v} [agent] ]
					[VP [\ibar{V}
						[\obar{V} [$\surd$visit] ]
						[\emph{t_5} ]
					] ]
				 ]
			 ]
		 ]
	 ]
	[DP [\ibar{D}
		[\obar{D} [every] ]
		[NP [\ibar{N} 
			[\ibar{N} [student, roof] ]
			[CP
				[DP [who, roof] ]
				[$\lambda 3\ \exists \epsilon$ \textsc{visit}($\epsilon$) $\wedge$ \textsc{patient}($3$)($\epsilon$) $\wedge$ \textsc{agent}(\denotes{Rodica})($\epsilon$) \\\fbox{\ibar{C}}
					[\obar{C} [\sout{that}] ]
					[TP
						[DP_F [RODICA, roof] ]
						[\ibar{T} 
							[\obar{T} [will] ]
							[vP
								[\emph{t}_2 ]
								[\ibar{v}
									[\obar{v} [agent] ]
									[VP [\ibar{V}
										[\obar{V} [$\surd$visit] ]
										[\emph{t}_3 ]
									] ]
								 ]
							 ]
						]
					 ]
				 ]
			 ]
		] ]
	] ]
]
\end{forest}
\end{exe}
\end{fullwidth}
% \medskip
\vspace*{-18pt}
The object DP of \textit{visit} has moved outside of the TP that provides the antecedent for the focus anaphor, and left a trace in its former position. This provides the right denotation, and also satisfies (\ref{ex6:fanaphora}), the conditions on focus anaphors. The boxed \ibar{C} is the focus anaphor, and the boxed TP is its focus antecedent. There is a focus-alternative to the denotation of \ibar{C} that is identical to the denotation that TP has. In particular, replacing the focus-marked \textit{Rodica} with \textit{Thuy} produces something from the denotation of \fbox{\ibar{C}} that is equivalent to \fbox{TP}. Note that this is partly due to the fact that the traces \textit{t_5} and \textit{t_3} are sloppily bound. That is, they have the same semantics that we've seen sloppily-bound pronouns have in section 6.1.2. 

This solution to the puzzle of antecedent contained focus anaphora predicts that it will arise only in cases where movement can create a trace bound by the term that contains the focus anaphor. This in turn entails that the focus anaphor in such cases must contain a variable in a position that corresponds to that trace, and that this variable and the trace meet the conditions of sloppy binding. I don't know that this prediction has been thoroughly examined, but it appears to be correct. There is a contrast, for example, between (\ref{ex6:acda}) and~(\ref{ex6:slopfana}).
\begin{exe}
\raggedright
	\ex[*]{\label{ex6:slopfana} Thuy will examine every claim that [RODICA <will examine it>].}
\end{exe} 
This is because in (\ref{ex6:slopfana}), the focus anaphor has a pronoun in the position that corresponds to where the trace of moving the \textit{every claim}-DP resides in the antecedent. This pronoun and the trace do not obey the conditions on sloppy binding, and as a consequence the condition in (\ref{ex6:fanaphora}) is not satisfied. To see this, compare (\ref{ex6:slopfanat}), the representation that corresponds to (\ref{ex6:slopfana}), with (\ref{ex6:acdres}), the representation that corresponds to the successful instance of antecedent contained anaphora.
\begin{fullwidth}
\begin{exe}
\raggedright
	\ex \label{ex6:slopfanat}
	\small
	\begin{forest}
	[$\forall x\ {[}\textsc{student}(x) \wedge \exists \epsilon$ {[}\textsc{examine}($\epsilon$) $\wedge$ \textsc{patient}($x$)($\epsilon$) $\wedge$ \textsc{agent}(\denotes{Rodica})($\epsilon$){]]} $\rightarrow$ {[} $\exists \epsilon'$ \textsc{examine}($\epsilon'$) $\wedge$ \textsc{patient}($x$)($\epsilon'$) $\wedge$ \textsc{agent}(\denotes{Thuy})($\epsilon'$){]}\\TP
		[$\lambda 5\ \exists \epsilon'$ \textsc{examine}($\epsilon'$) $\wedge$ \textsc{patient}(5)($\epsilon'$) $\wedge$ \textsc{agent}(\denotes{Thuy})($\epsilon'$)\\TP
			[DP [Thuy, roof]]
			[\ibar{T}
				[\obar{T} [will]
				[vP
					[\textit{t_4} ]
					[\ibar{v}
						[\obar{v} [agent]
						[VP [\ibar{V}
							[\obar{V} [$\surd$examine]
							[\textit{t_5} ]
						] ]
					 ]
				 ]
			 ]
		 ]
		[DP [\ibar{D}
			[\obar{D} [every]
			[NP [\ibar{N}
				[\obar{N} [claim]]
				[CP [$\exists \epsilon$ \textsc{examine}($\epsilon$) $\wedge$ \textsc{patient}(it)($\epsilon$)\\ \textsc{agent}(\denotes{Rodica})($\epsilon$)\\\ibar{C} 
					[\obar{C} [that]
					[TP
						[DP_F [RODICA]
						[\ibar{T}
							[\obar{T} [will]
							[vP
								[\textit{t_2} ]
								[\ibar{v}
									[\obar{v} [agent]
									[VP [\ibar{V} 
										[\obar{V} [$\surd$examine]
										[DP [it, roof]]
									]]
								 ]
							 ]
						 ]
					 ]
				]]
			 ]
		 ]
	]
	\end{forest}
\end{exe}
\end{fullwidth}
\fbox{\ibar{C}} doesn't have a focus alternative that matches the denotation of \fbox{TP} in (\ref{ex6:slopfanat}), and therefore fails (\ref{ex6:fanaphora}). There is no phrase smaller than \fbox{TP} that could be an antecedent for \fbox{\ibar{C}, because \fbox{TP} is the smallest phrase that contains a binder for the trace left by moving the object DP. Our semantics doesn't allow phrases with unbound traces in them to have a denotation. No phrase larger than \fbox{TP} can be an antecedent because it will contain the focus anaphor. Consequently, the only focus anaphors that \fbox{TP} can be the antecedent for are those that are also functions from things to truth values. In general, a focus anaphor whose D-structure position puts it within it antecedent will have a contain a trace that corresponds to the trace left when that focus anaphor moves out of its antecedent.

This solution also makes predictions about the position that the DP containing the focus anaphor will be semantically interpreted in. That DP will have to be interpreted in a position that puts it outside the antecedent for the focus anaphor. When the DP is headed by a quantifier, this can be tested by considering where that quantifier is interpreted. This prediction is examined in \citet{Williams:1977} and \citet{Sag:1976}, and has been confirmed. Let's look at one of the cases that Williams discusses. 

First, consider (\ref{ex6:wantvpe}).
\begin{exe}
\raggedright
	\ex \label{ex6:wantvpe}
	\begin{xlist}
		\ex {}[_a Cherlon wants [_b to beat every contestant that
		\begin{xlist}
			\ex 
			\begin{xlist}
			\ex MARIBEL <wants to beat>.
			\ex MARIBEL <does \sout{want to beat}>.
			\end{xlist}
			\ex
			\begin{xlist}
				\ex MARIBEL <beat>.
				\ex MARIBEL <did \sout{beat}>.
			\end{xlist}
		\end{xlist}
	\end{xlist}
\end{exe}
What (\ref{ex6:wantvpe}) shows is that the sentence \textit{Cherlon wants to beat every contestant} makes available two antecedents for a focus anaphor. One of those antecedents is marked with an opening bracket labeled $a$, and the other is marked with a bracket labeled $b$. The focus anaphors in (\ref{ex6:wantvpe}a) make use of the $a$-antecedent, and the focus anaphors in (\ref{ex6:wantvpe}b) make use of the $b$-antecedent. 

Next, consider the simpler (\ref{ex6:wantvpe1}).
\begin{exe}
\raggedright
	\ex \label{ex6:wantvpe1}
	Cherlon wants to beat every contestant.
\end{exe}
This sentence has two interpretations that depend on where \textit{every contestant} is found in the syntactic representation. If the syntax puts \textit{every contestant} in the position indicated in (\ref{ex6:wantvpe1t}), it will be interpreted within the infinitival complement clause. 
\begin{exe}
\raggedright
	\ex \label{ex6:wantvpe1t}
	\small
	\begin{forest}
	[TP
		[DP^2 [Cherlon, roof]]
		[\ibar{T}
			[\obar{T} [pres]]
			[vP
				[\textit{t_2} ]
				[\ibar{v}
					[\obar{v} [experiencer]
					[VP [\ibar{V} 
						[\obar{V} [$\surd$want]]
						[CP [\ibar{C} 
							[\obar{C} [e]]
							[TP
								[TP 
								[PRO^3]
								[\ibar{T}
									[\obar{T} [to]]
									[vP
										[\textit{t_3} ]
										[\ibar{v}
											[\obar{V} [agent]]
											[VP [\ibar{V}
												[\obar{V} [$\surd$beat]]
												[\textit{t_4} ]
											] ]
										 ]
									 ]
								] ]
								[DP^4 [every contestant, roof]]
							 ]
						] ]
					]]
				 ]
			 ]
		 ]
	]
	\end{forest}
\end{exe}
The meaning that (\ref{ex6:wantvpe1t}) has corresponds to (\ref{ex6:wantvpe1ti}).
\begin{exe}
\raggedright
	\ex \label{ex6:wantvpe1ti}
	It's to beat every contestant that Cherlon wants.
\end{exe}
Note that the desire that (\ref{ex6:wantvpe1ti}) ascribes to Cherlon is consistent with her not wanting to beat each particular contestant. If her beloved sister is among the contestants, she may not want to beat that sister, for instance. (\ref{ex6:wantvpeti}) merely says that what Cherlon desires is to be the overall winner. It could describe a scenario in which Cherlon has conflicting desires. Let's call this the ``narrow scope'' interpretation of \textit{every contestant}. 

If, by contrast, the syntax puts \textit{every contestant} in the position indicated in (\ref{ex6:wantvpe1tt}), then it will be interpreted outside the complement clause.
\begin{exe}
\raggedright
	\ex \label{ex6:wantvpe1tt}
	\small
	\begin{forest}
	[TP
		[TP
			[DP^2 [Cherlon, roof]]
			[\ibar{T}
				[\obar{T} [pres]]
				[vP
					[\textit{t_2} ]
					[\ibar{v}
						[\obar{v} [experiencer]]
						[VP [\ibar{V}
							[\obar{V} [$\surd$want]]
							[CP [\ibar{C}
								[\obar{C} [e]]
								[TP
									[PRO^3 ]
									[\ibar{T}
										[\obar{T} [to]]
										[vP
											[\textit{t_3} ]
											[\ibar{v}
												[\obar{v} [agent]]
												[VP [\ibar{V}
													[\obar{V} ]$\surd$beat]
													[\textit{t_4} ]
												] ]
											 ]
										 ]
									 ]
								 ]
							] ]
						] ]
					 ]
				 ]
			 ]
		 ]
		[DP^4 [every contestant, roof]]
	]
	\end{forest}
\end{exe}
The meaning delivered by (\ref{ex6:wantvpe1tt}) can be paraphrased by (\ref{ex:wantvpe1tti}).
\begin{exe}
\raggedright
	\ex \label{ex6:wantvpe1tti} Every contestant is one that Cherlon wants to beat.
\end{exe}
On this meaning, Cherlon has a desire regarding each contestant, and that desire is to beat them. If the contestants include Cherlon's sister, then the meaning derived from (\ref{ex6:wantvpe1tt}) entails that Cherlon has the desire to beat her as well. Let's call this the ``wide scope'' interpretation of \textit{every contestant}. The account of antecedent contained anaphora provided by Trace theory predicts that (\ref{ex6:wantvpe}) cannot simultaneously have the narrow scope interpretation of \textit{every contestant} and resolve the focus anaphor so that it takes the larger $a$-antecedent. This interpretation would be mapped onto the syntax in (\ref{ex6:acdparse}).
\begin{exe}
\raggedright
	\ex \label{ex6:acdparse}
	\small
	\begin{forest}
	[TP
		[DP^2 [Cherlon, roof]]
		[\ibar{T}
			[\obar{T} [pres]]
			[vP
				[\textit{t_2} ]
				[\ibar{v}
					[\obar{v} [experiencer]]
					[VP [\ibar{V}
						[\obar{V} [$\surd$want]]
						[CP [\ibar{C}
							[\obar{C} [e]]
							[TP
								[TP
									[PRO^3 ]
									[\ibar{T}
										[\obar{T} [to]]
										[vP
											[\textit{t_3} ]
											[\ibar{v}
												[\obar{v} [agent]]
												[VP [\ibar{V}
													[\obar{V} [$\surd$beat]]
													[\textit{t_4} ]
												] ]
											 ]
										 ]
									 ]
								 ]
								[DP^4 [every contestant\\ that MARIBEL <wants to beat>, roof]]
							 ]
						] ]
					] ]
				 ]
			 ]
		 ]
	]
	\end{forest}
\end{exe}
The problem with (\ref{ex6:acdparse}) is that the focus anaphor lies within its antecedent. When \textit{every contestant} takes narrow scope, the focus anaphor can only correspond to the lower clause. (\ref{ex6:acdparse2}) shows how. The denotations aren't provided in (\ref{ex6:acdparse2}), but the relationship between \fbox{\ibar{C}} and \fbox{TP} is parallel to the \fbox{\ibar{C}} and \fbox{TP} in (\ref{ex6:slopfanat}). Because a focus alternative to \textit{Maribel} is PRO, \denotes{TP} is the same as \denotes{\ibar{C}}^F. This makes \denotes{TP} in (\ref{ex6:acdparse2}) an antecedent for \denotes{\ibar{C}}.
\begin{exe}
\raggedright
	\ex \label{ex6:acdparse2}
	\small
	\begin{forest}
	[TP
		[DP^2 [Cherlon, roof]]
		[\ibar{T}
			[\obar{T} [pres]]
			[vP
				[\textit{t_2} ]
				[\ibar{v}
					[\obar{v} [experiencer]]
					[VP [\ibar{V}
						[\obar{V} [$\surd$want]]
						[CP [\ibar{C}
							[\obar{C} [e]]
							[TP
								[TP
									[PRO^3 ]
									[\ibar{T}
										[\obar{T} [to]]
										[vP
											[\textit{t_3} ]
											[\ibar{v}
												[\obar{v} [agent]]
												[VP [\ibar{V}
													[\obar{V} ]$\surd$beat]
													[\textit{t_4} ]
												] ]
											 ]
										 ]
									 ]
								 ]
								[DP^4 [\ibar{D}
									[\obar{D} [every]]
									[NP [\ibar{N}
										[\obar{N} [contestant]]
										[CP
											[wh^5 ]
											[\ibar{C} [that MARIBEL <beat \textit{t_5>}]]
										 ]
									] ]
								]]
							 ]
						] ]
					] ]
				 ]
			 ]
		 ]
	]
	\end{forest}
\end{exe}

These predicted outcomes seem to be correct. In a context that limits the interpretation to a narrow scope one for \textit{every contestant}, like that provided below, only the smaller b-antecedent is available for a focus anaphor.
\begin{exe}
\raggedright
\sn \textit{Maribel and Cherlon are fierce rivals. Neither of them expects to win the tournament, but they each expect that they can outperform the other. So fierce is their rivalry, that Cherlon would even try to beat her sister, who is also a contestant in the tournament, if it would mean besting Maribel. This is surprising, because if Maribel weren't in the tournament, Cherlon would let her sister beat her. Cherlon hates beating her sister. I tell Cherlon's sister:}

	\ex \label{ex6:acdrange} {}[_a Cherlon wants [_b to beat every contestant that\marginnote{``\#'' marks a sentence whose meaning is anomolous.}
	\begin{xlist}
		\ex[\#]{MARIBEL <wants to beat>} 
		\sn[\#]{MARIBEL <wants to>}
		\ex[]{MARIBEL <beats>}
		\sn[]{MARIBEL does}
	\end{xlist}
\end{exe}
It's only in a context which allows the wide scope interpretation for \textit{every contestant}  that the a-antecedent for the focus anaphor becomes available. The context below illustrates.
\begin{exe}
\raggedright
	\sn \textit{Maribel and Cherlon are sisters. When they play a tournament together, they discuss who they hope to beat. It turns out that:}
	\ex \label{ex6:acdrange2} {}[_a Cherlon wants [_b to beat every contestant that
	\begin{xlist}
		\ex MARIBEL <wants to beat>.
		\ex MARIBEL <wants to>.
	\end{xlist}
\end{exe}
The syntax for a wide-scope interpretation of \textit{every contestant} is~(\ref{ex6:widerange}).
\begin{exe}
\raggedright
	\ex \label{ex6:widerange}
	\small
	\begin{forest}
	[TP
		[\fbox{TP}
			[DP^2 [Cherlon, roof]]
			[\ibar{T}
				[\obar{T} [pres]]
				[vP
					[\textit{t_2} ]
					[\ibar{v}
						[\obar{v} [experiencer]]
						[VP [\ibar{V}
							[\obar{V} [$\surd$want]]
							[CP [\ibar{C}
								[\obar{C} [e]]
								[TP
									[PRO^3 ]
									[\ibar{T}
										[\obar{T} [to]]
										[vP
											[\textit{t_3} ]
											[\ibar{v}
												[\obar{v} [agent]]
												[VP [\ibar{V}
													[\obar{V} [$\surd$beat]]
													[\textit{t_4} ]
												] ]
											 ]
										 ]
									 ]
								 ]
							] ]
						] ]
					 ]
				 ]
			 ]
		 ]
		[DP^4 [\ibar{D}
			[\obar{D} [every]]
			[NP [\ibar{N}
				[\obar{N} [contestant]]
				[CP
					[wh^5 ]
					[\ibar{C} [that MARIBEL <wants to beat \textit{t_5}>, roof]]
				 ]
			] ]
		] ]
	]
	\end{forest}
\end{exe}
The a-antecedent is available here because \denotes{\fbox{TP}} is equivalent to a \denotes{\ibar{C}}^F, because \textit{Cherlon} is a focus alternative to \textit{Maribel}.

The solution to antecedent contained anaphora that Trace theory provides gives a good account of the conditions necessary for antecedent contained anaphora to show up. Trace theory's requirements narrowly limit antecedent contained anaphora to environments where movement can put traces in positions where they qualify to be sloppily bound. Because Trace theory also gives a straightforward account of the semantics of relative clauses and the semantics of DPs with quantificational heads, it's widely regarded by syntacticians as a successful theory. It has known problems as well, as I'll show you in \ref{sec:copytheory}, but it has nonetheless become the standard model of movement.

Adopting Trace theory, and the connected idea that parses which movement creates can be ones that provide the denotation for a sentence, complicates the relation between form and meaning that derivations are intended to model. In the simpler models that we've used up till now, we can impose all the requirements that the semantics imposes on a single parse in the derivation: the D-structure. In the new model, a derivation created by movement can meet those requirements and therefore become the parse relevant for a semantic interpretation. This will have consequences for the ordering of the operations which give rise to those parses. Because some sentences will require a non-D-structure parse to provide the denotation, it will also require that derivations be steered so that those parses are created. This means that some derivations will require a specific order of operations to bring about the right parse. We look at this effect in the next section.

\subsection{The Syntax of Quantificational DPs}
\label{sec:QDPs}

One of the consequences of building a denotation for quantificational determiners that allows them to be interpreted in positions outside of vP, is that they cannot be interpreted in positions inside vP. The denotation we've given for vP is as functions whose domains are eventualities. There is some phrase containing that vP, however, where the denotations are no longer functions of this kind. Instead, they are closed formulas---no longer functions. Because quantifiers cannot be interpreted unless the DP they head is a sister to a closed formula, this means that all DPs headed by quantifiers must be semantically interpreted outside vP. This section investigates consequences of this decision.

To put the details back into view, we adopted the denotation for quantificational determiners like that in (\ref{ex6:newqnpden}), which allows the quantificational DPs to be interpreted in their moved positions, as (\ref{ex6:newqnpdent}) illustrates.
\medskip
\begin{fullwidth}
\begin{exe}
\raggedright
	\ex \label{ex6:newqnpden} \denotes{every} = $\lambda P\ \lambda Q\ \forall x$ $P(x)\ \rightarrow Q(x)$ \\[8pt]
	\ex \label{ex6:newqnpdent}
	\footnotesize
	\begin{forest}
		[$\forall x$ \textsc{man}($x$) $\rightarrow$ {[}$\forall y$ \textsc{woman}($y$) $\rightarrow$ {[}$\exists \epsilon$ \textsc{agent}($y$)($\epsilon$) $\wedge$ \textsc{meet}($\epsilon$) $\wedge$ \textsc{patient}($x$)($\epsilon$){]]}\\TP
			[$\lambda 3\ \forall y$ \textsc{woman}($y$) $\rightarrow$ {[}$\exists \epsilon$ \textsc{agent}($y$)($\epsilon$) $\wedge$ \textsc{meet}($\epsilon$) $\wedge$ \textsc{patient}($3$)($\epsilon$){]}\\TP 
				[$\lambda Q\ \forall y$ \textsc{woman}($y$) $\rightarrow$ $Q(y)$\\DP [every woman, roof]]
				[$\lambda 2\ \exists \epsilon$ \textsc{agent}($2$)($\epsilon$) $\wedge$ \textsc{meet}($\epsilon$) $\wedge$ \textsc{patient}($3$)($\epsilon$)\\ \ibar{T}
					[\obar{T} [past] ]
					[$\exists \epsilon$ \textsc{agent}($2$)($\epsilon$) $\wedge$ \textsc{meet}($\epsilon$) $\wedge$ \textsc{patient}($3$)($\epsilon$)\\AspP [\ibar{Asp}
					[\obar{Asp} [asp]]
					[vP
						[\emph{t_2}]
						[\ibar{v} 
							[\obar{v} [agent] ]
							[VP [\ibar{V}
								[\obar{V} [$\surd$meet] ]
								[\emph{t_3}]
							] ]
						]
					 ]
				 ]
			]]]
			[$\lambda Q\ \forall x$ \textsc{man}($x$) $\rightarrow$ $Q(x)$\\DP [every man, roof]]
		]
		\end{forest}	
\end{exe}
\end{fullwidth}
\medskip
This correctly captures the semantic fact that the event variable is existentially bound within the scope of \emph{no} and \emph{every}. The denotation for (\ref{ex6:newqnpden}), for instance, allows there to be a different meeting event for each of the men and women involved. It's also this syntax, and semantics, that provides the account of Antecedent Contained anaphora sketched above.

But this denotation for quantificational determiners does not allow the DPs they project to be semantically interpreted in their D-structure positions, as (\ref{ex6:newqnpdend}) indicates.
\begin{exe}
\raggedright
	\ex \label{ex6:newqnpdend}
	\footnotesize
\begin{forest}
[AspP [\ibar{Asp}
	[\obar{Asp} [asp]]
	[?? \\vP
		[$\lambda Q\ \forall y$ \textsc{woman}($y$) $\rightarrow$ $Q(y)$\\ DP [every woman, roof]]
		[$\lambda y\ \lambda \epsilon$ \textsc{agent}($y$)($\epsilon$) $\wedge$ ?? \\ \ibar{v}
			[\obar{v} [agent] ]
			[??\\VP [??\\\ibar{V}
				[$\lambda y\ \lambda \epsilon$ \textsc{meet}($\epsilon$) $\wedge$ \textsc{patient}($y$)($\epsilon$)\\\obar{V} [$\surd$meet] ]
				[$\lambda Q\ \forall x$ \textsc{man}($x$) $\rightarrow$ $Q(x)$\\DP [every man, roof] ]
			] ]
		 ]
	 ]
]]
\end{forest}	
\end{exe}
As noted, this means that quantificational DPs, unlike DPs that simply refer, must move out of vP to be semantically interpreted.

Our syntax presently independently requires that DPs which bear nominative Case move, and therefore our derivations already include a parse that will allow nominative Case-marked quantificational DPs to be semantically interpreted. The movement of \emph{every woman} in the derivations sketched above exemplifies the general case. 

We also have a syntax that allows accusative Case-marked DPs to move, but this doesn't always provide for a derivation in which accusative Case-marked quantificational DPs will be able to be semantically interpreted. If we adopt the Split VP Hypothesis described in chapter \ref{sec:ordering}, for instance, then Accusative Case-marked DPs move in the way indicated in~(\ref{ex6:objshiftqnp}). But this doesn't move these DPs far enough for them to be semantically interpreted if they are headed by quantificational determiners, as the denotations in (\ref{ex6:objshiftqnp}) indicate.%
%
\marginnote[10pt]{Assume that the contents of \obar{$\mu$} have no denotation, and therefore that \ibar{$\mu$} inherits its denotation unchanged from ~VP. The semantics of the new movement rule change this denotation in~(\ref{ex6:objshiftqnp})}%
%
\medskip
\begin{fullwidth}
\begin{exe}
\raggedright
	\ex \label{ex6:objshiftqnp}
	\begin{xlist}
		\ex 
		\footnotesize
		\begin{forest}
		[\ibar{v}
			[\obar{v} [agent] ]
			[??\\$\mu$P
				[$\lambda Q\ \forall x$ \textsc{child}($x$) $\rightarrow$ $Q(x)$ \\DP [every child, roof]]
				[$\lambda 2\ \lambda \epsilon$ \textsc{meet}($\epsilon$) $\wedge$ \textsc{patient}($2$)($\epsilon$)\\\ibar{$\mu$} 
					[\obar{$\mu$} [$\bullet$\emph{acc}]]
					[$\lambda \epsilon$ \textsc{meet}($\epsilon$) $\wedge$ \textsc{patient}($2$)($\epsilon$)\\VP [\ibar{V}
						[$\lambda x\ \lambda \epsilon$ \textsc{meet}($\epsilon$) $\wedge$ \textsc{patient}($x$)($\epsilon$)\\\obar{V} [$\surd$meet] ]
						[\emph{t_2} ]
					] ]
				 ]
			]
		]
		\end{forest}

		\ex 
		\begin{forest}
		[\ibar{v}
			[\obar{v} [agent] ]
			[??\\ $\mu$P
				[$\lambda Q\ \forall x$ \textsc{child}($x$) $\rightarrow$ $Q(x)$\\DP [every child, roof]]
				[$\lambda 2\ \lambda \epsilon\ \exists \epsilon'$ \textsc{make}($\epsilon$) $\wedge$ $\epsilon$ causes $\epsilon'$ $\wedge$ $\wedge$ \textsc{meet}($\epsilon'$) $\wedge$ \textsc{patient}($2$)($\epsilon'$)\\\ibar{$\mu$}
					[\obar{$\mu$} [$\bullet$\emph{acc}] ]
					[$\lambda \epsilon\ \exists \epsilon'$ \textsc{make}($\epsilon$) $\wedge$ $\epsilon$ causes $\epsilon'$ $\wedge$ \textsc{meet}($\epsilon'$) $\wedge$ \textsc{patient}($2$)($\epsilon'$)\\VP [\ibar{V}
						[\obar{V} [$\surd$make] ]
						[$\lambda \epsilon$ \textsc{meet}($\epsilon$) $\wedge$ \textsc{patient}($2$)($\epsilon$)\\vP
							[\emph{t_2} ]
							[$\lambda x\ \lambda \epsilon$ \textsc{meet}($\epsilon$) $\wedge$ \textsc{agent}($x$)($\epsilon$) \\\ibar{v} [$\surd$run, roof] ]
						 ]
					] ]
				 ]
			]
		]
		\end{forest}
		%
%
\marginnote[-20pt]{I adopt here a denotation for \emph{make} that allows it to combine with a function whose domain is eventualities, and return another function whose domain is eventualities: events that cause the existence of the eventualities their complement describes.}%
%
	\end{xlist}
\end{exe}
\end{fullwidth}
\medskip
The only movement rule we have for accusative Case-marked DPs that will put them in a position where they can be semantically interpreted is Complex NP Shift.

Our semantics, then, predicts that DPs which are headed by quantifiers will necessarily undergo Complex NP Shift when they are accusative Case-marked. And yet, the positions that quantificational DPs can be pronounced in does not fit this prediction. Accusative Case-marked quantificational DPs need not be pronounced in the positions that Complex NP Shift would put them, as (\ref{ex6:nocnps}) shows.
\begin{exe}
\raggedright
	\ex \label{ex6:nocnps}
	\begin{xlist}
		\ex She will read every story to you.
		\ex She will make every child laugh.
		\ex She will look every number up.
	\end{xlist}
\end{exe}
Something in our system is wrong.

But what? There is evidence that the syntax which delivers the word order in (\ref{ex6:nocnps}) should allow the quantificational DPs to be interpreted outside of the vP. That evidence derives from the relative interpretations of the quantifiers in~(\ref{ex6:inverse1}).
\begin{exe}
\raggedright
	\ex \label{ex6:inverse1}
	\begin{xlist}
		\ex A hired actor will read every story to you.
		\ex A funny actor will make every child laugh.
		\ex A different student will look every number up.
	\end{xlist}
\end{exe}
These sentences have an interpretation which allows for the actors to vary with the stories or children involved. These meanings are roughly those paraphrased in~(\ref{ex6:inverse1p}).
\begin{exe}
\raggedright
	\ex \label{ex6:inverse1p}
	\begin{xlist}
		\ex Every story is such that a hired actor will read it to you.
		\ex Every child is such that a funny actor will make them laugh.
		\ex Every number is such that a different student will look it up.
	\end{xlist}
\end{exe}
To get that interpretation requires that the DPs headed by \emph{every} in (\ref{ex6:inverse1}) occupy a syntactic position higher than the subject DPs. That is not where our present syntax puts these DPs. Facts like these, then, indicate that at least one flaw our system has is in how it makes the position that quantifiers are semantically interpreted in correspond to the position they are spoken in. 

We need a semantics which allows a quantificational DP to be semantically interpreted so that it has access to information ``higher'' in the phrase marker relative to it. Our present semantics lets the quantifier only ``see'' the material that is a sister to the DP it heads. Our syntax for (\ref{ex6:inverse1}), for instance, puts \emph{every} inside a DP that is in Specifier of $\mu$P. In this position, the phrase that \emph{every} has as its second argument is the sister to this DP; that is: \ibar{$\mu$}. What (\ref{ex6:inverse1}) shows is that \emph{every} should instead have a constituent that includes the \emph{actor} DP as its second argument. There are ways of doing this that de-couple the syntax and the semantics, allowing a quantificational DP to combine semantically with material that it doesn't syntactically combine with.%
%
\marginnote[0pt]{See \citet{Cooper:1983}, for an example.}%
%
To pursue this direction, we would need to revisit the relationship we have established between quantificational DPs and the variables they bind, among other things.

But there is some evidence that the solution to the problem posed by (\ref{ex6:inverse1p}) should employ the very syntax/semantics we have developed. Perhaps the most direct evidence for preserving the syntax we now have is that antecedent contained anaphora can be resolved for DPs which are not in the positions that Complex NP Shift would give to them. (\ref{ex6:acdacc1}) illustrates.
\begin{exe}
\raggedright
	\ex \label{ex6:acdacc1}
	\begin{xlist}
		\ex She will read every story that I didn't to you.
		\ex She will make every child that I didn't laugh. 
	\end{xlist}
\end{exe}%
%
\marginnote[-70pt]{The non-elliptical versions of these cases are less pretty, but I think similarly grammatical.
\begin{exe}
\raggedright
	\exi{i.} She will read every story that I didn't <read to you> to you.
	\exi{ii.} She will make every child that I didn't <make laugh> laugh.
\end{exe}
My suspicion is that what ails these cases is the availability of the equivalent, but less repetitive, (\ref{ex6:acdacc1}). I am encouraged to have this suspicion by the improved status that examples like (i)-(ii) have when there isn't a elliptical competitor. For example, the following cases sound better to me than (i)-(ii).
\begin{exe}
\raggedright
		\exi{iii.} She will read every story that I didn't <read to HIS child> to YOUR child.
		\exi{iv.} She will make every child that I <made CRY> laugh.
\end{exe}
The focus anaphors in (iii) and (iv) contain focus-marked material which prevents them from having an elliptical version.}%
%
The grammaticality of these sentences rests on their prosody. They require a pronunciation that puts a rising tone on \emph{didn't} followed by a pause and falling tone on the material that follows \emph{didn't}. We'll return to that prosody, but setting it to the side for now, the grammaticality of (\ref{ex6:acdacc1}) encourages the view that we need a syntax that places the accusative Case-marked DPs outside of the vP they seem to reside in. As we've seen, this is the only way to produce an antecedent for the surface anaphor in (\ref{ex6:acdacc1}). I don't know how to avoid the conclusion that the sentences in (\ref{ex6:acdacc1}) should have a syntax like that indicated in~(\ref{ex6:acdacc1t}).
\medskip
\begin{fullwidth}
\begin{exe}
\raggedright
	\ex \label{ex6:acdacc1t}
	\footnotesize
	\begin{forest}
	[TP
		[TP 
		[DP^2 [she, roof]]
		[\ibar{T}
			[\obar{T} [will] ]
			[vP
				[\emph{t_2}]
				[\ibar{v}
					[\obar{v} [agent] ]
					[VP [\ibar{V} 
					[\ibar{V}
						[\obar{V} [$\surd$read] ]
						[\emph{t_3} ]
					]
					[PP [to you, roof]]
					]
					]
				 ]
			 ]
		]]
		[DP^3 [\ibar{D}
			[\obar{D} [every] ]
			[NP[\ibar{N}
				[\ibar{N} [story, roof] ]
				[CP
					[DP^4 [\ibar{D} [\obar{D} [\emph{wh}]]] ]
					[\ibar{C} 
						[\obar{C} [that] ]
						[TP
							[DP^5 [I, roof] ]
							[\ibar{T}
								[\obar{T} [didn't] ]
								[vP
									[\emph{t_5} ]
									[\ibar{v} 
										[\obar{v} [agent] ]
										[VP [\ibar{V}
											[\ibar{V}
												[\obar{V} [\sout{$\surd$read}] ]
												[\sout{\emph{t_4}}]
											 ]
											[PP [\sout{to you}, roof] ]
										] ]
									 ]
								 ]
							 ]
						 ]
					 ]
				 ]
			 ]]
		] ]
	]
	\end{forest}
	\vspace*{-70pt}

	\begin{forest}
	[TP
	[TP
		[DP^2 [she, roof] ]
		[\ibar{T}
			[\obar{T} [will] ]
			[vP
				[\emph{t_2} ]
				[\ibar{v}
					[\obar{v} [agent] ]
					[VP [\ibar{V}
						[\obar{V} [$\surd$make] ]
						[vP
							[\emph{t_3} ]
							[\ibar{v}
								[\obar{v} [agent] ]
								[VP [\ibar{V} [\obar{V} [$\surd$laugh]]] ]
							 ]
						 ]
					] ]
				 ]
			 ]
		 ]
	]
	[DP^3 [\ibar{D}
		[\obar{D} [every] ]
		[NP [\ibar{N} 
			[\ibar{N} [child, roof] ]
			[CP
				[DP^4 [\emph{wh}, roof] ]
				[\ibar{C}
					[\obar{C} [that] ]
					[TP
						[DP^5 [I, roof] ]
						[\ibar{T}
							[\obar{T} [didn't] ]
							[vP
								[\emph{t_5} ]
								[\ibar{v}
									[\obar{v} [agent] ]
									[VP [\ibar{V}
										[\obar{V} [\sout{$\surd$make}] ]
										[vP
											[\sout{\emph{t_4}} ]
											[\ibar{v}
												[\obar{v} [\sout{agent}] ]
												[VP [\sout{$\surd$laugh}, roof] ]
											 ]
										 ]
									] ]
								 ]
							]
						 ]
					 ]
				 ]
			 ]
		] ]
	]]
	]
	\end{forest}
\end{exe}
\pagebreak
These representations correspond to the strings in (\ref{ex6:acdacc2}).
\begin{exe}
\raggedright
	\ex \label{ex6:acdacc2}
	\begin{xlist}
		\ex She will read to you every story that I didn't.
		\ex She will make laugh every child that I didn't.
	\end{xlist}
\end{exe}
\end{fullwidth}
\medskip
These are grammatical; indeed, they don't require the very particular prosody that (\ref{ex6:acdacc1}) does, and they have a more natural cast to them. The puzzle is to understand how these representations can be expressed with the strings in~(\ref{ex6:acdacc1}).

Robert May\cite{May:1977,May:1985} developed a solution to these problems that rethinks the relationship between a derivation and the procedure that interprets a syntactic representation as a string of words. We've just developed a view of derivations which allows for representations created by movement to be semantically interpreted. So, our present syntax characterizes sentences with~(\ref{ex6:derivation2}), repeated here.
\begin{exe}
\raggedright
	\exr{ex6:derivation2}
		$\mathcal{D}$ = (P_1, P_2\ldots,P_n) is a derivation iff:
		\begin{xlist}
			\ex P_1 is formed by the \ibar{X} Skeleton and Lexical Insertion,
			\ex Each P_{i>1} is formed from P_{i-1} by a rule,
			\ex Some P_i is semantically interpreted, and
			\ex P_n is subject to the linearization algorithm.
		\end{xlist}
		\marginnote[20pt]{The cardinality of a derivation is the number of Ps in it.}
		\sn A sentence is the set of well-formed derivations which have the same P_1 (=D-structure). A grammatical sentence is a derivation whose cardinality isn't greater than the cardinality of any other derivation for that sentence.
\end{exe}
May's proposal is to let the linearization algorithm apply to a non-final element of a derivation. A way of fitting May's proposal to the particulars of our syntax is with~(\ref{ex6:derivation3}).
\begin{exe}
\raggedright
	\ex \label{ex6:derivation3}
	$\mathcal{D}$ = (P_1, P_2,\ldots,P_n) is a derivation iff:
	\begin{xlist}
		\ex P_1 is formed by the \ibar{X} Skeleton and Lexical Insertion,
		\ex Each P_{i>1} is formed from P_{i-1} by a rule,
		\ex P_j is subject to the linearization algorithm, 
		\ex P_n is semantically interpreted.
	\end{xlist}
	\sn \label{ex6:early} A sentence is the set of well-formed derivations which have the same P_1 (=D-structure) and P_n (=Logical Form). A grammatical sentence is a derivation which has a (P_1,\ldots,P_j) sub-sequence that is not longer than the (P_1,\ldots,P_j) sub-sequence of any other derivation for that sentence.
\end{exe}
May called the final parse in a derivation, ``Logical Form'' (LF), after the hypothesis that it is the representation that is semantically interpreted. I've not adopted that view in (\ref{ex6:derivation3}), as our syntax of Head Movement won't allow some LFs to be semantically interpreted. %
%
\marginnote[0pt]{See section \ref{sec:ordering}.}%
%
(\ref{ex6:derivation3}) defines S-structure (=P_j) as the element of a derivation that meets the well-formedness conditions necessary to be subject to the linearization algorithm. Chief among those are the conditions that determine where Vocabulary Insertion puts the words that will be linearized. The requirements of \textsc{agree} and $\bullet$ are among those conditions, as they determine where two pieces of a word must be. (\ref{ex6:early}) requires that an S-structure be as close to the D-structure as possible. 

The role that semantic well-formedness plays in determining whether a structure can be linearized is no longer direct. An S-structure need not be the final parse in the series that contains something that is semantically interpretable. If it's not the final parse in that series, though, it must be possible to create an interpretable representation from an S-structure with the rules that define a derivation. Roughly speaking, what this model does is cause those rules which are unnecessary for forming a linearizable string, but necessary for forming a semantically interpretable string, to form the post-S-structure portion of the derivation. 

By separating the parse in a derivation that determines the string from the parse that determines the meaning, May's proposal solves the problem posed by accusative Case-marked quantificational DPs. The parses that determine the semantic interpretation must be related to those that determine the string. But they need not be identical, and that, in May's system, is precisely what arises with accusative Case-marked quantificational~DPs.

A feature of this model is that it allows one S-structure to correspond to more than one LF. That is, it allows the structure that determines the string to be related to more than one semantically interpreted structure. Thus, it predicts that in principle, a single S-structure should be able to be structurally ambiguous. There is evidence that suggests this is a correct prediction, though for a variety of semantic reasons the evidence is less direct than desirable.

To appreciate the difficulty in evaluating this prediction, consider what this syntax predicts for~(\ref{ex6:raiseamba}).
\begin{exe}
\raggedright
	\ex \label{ex6:raiseamba}
	Some student seemed to have read every assignment.
\end{exe}

One LF that is expected for this sentence is~(\ref{ex6:raiseambat}).
\medskip
\begin{fullwidth}
\begin{exe}
\raggedright
	\ex \label{ex6:raiseambat}
\footnotesize
\hspace*{-25pt}
\begin{forest}
[$\forall x$ \textsc{assignment}($x$) $\rightarrow$ {[}$\exists y$ \textsc{student}($y$) $\wedge$ $\exists \epsilon$ \textsc{seem}($\epsilon$) $\wedge$ \textsc{cause}($\exists \epsilon'$ \textsc{read}($e'$) $\wedge$ \textsc{patient}($x$)($\epsilon'$) $\wedge$ \textsc{agent}($y$)($e'$))($\epsilon$) {]}\\TP
[ $\lambda 2\ \exists y$ \textsc{student}($y$) $\wedge\ \exists \epsilon$\textsc{seem}($e$) $\wedge$ \textsc{cause}(\textsc{read}($e'$) $\wedge$ \textsc{patient}($2$)($e'$) $\wedge$ \textsc{agent}($y$)($e'$))($\epsilon$) \\TP
	[$\lambda Q\ \exists y$ \textsc{student}($y$) $\wedge$ $Q(y)$ \\DP [some student, roof]]
	[$\lambda 3\ \exists \epsilon$ \textsc{seem}($\epsilon$) $\wedge$ \textsc{cause}($\exists \epsilon'$ \textsc{read}($\epsilon'$) $\wedge$ \textsc{patient}($2$)($\epsilon'$) $\wedge$ \textsc{agent}($3$)($\epsilon'$))($\epsilon$)\\\ibar{T} 
		[\obar{T} [past] ]
		[VP [\ibar{V}
			[\obar{V} [$\surd$seem] ]
			[TP [\ibar{T}
				[\obar{T} [to] ]
				[vP [ \emph{t_3} have read \emph{t_2}, roof] ]
			] ]
		]]
	]
]
[$\lambda Q\ \forall x$ \textsc{assignment}($x$) $\rightarrow$ $Q(x)$ \\DP [every assignment, roof]]
]
\end{forest}
	%
	%
	\marginnote[100pt]{We do not have the requisite semantic machinery to explicate how \emph{seem} combines with its TP complement. What I've done in rendering the denotations of \emph{seem} and its complement here is a stop-gap; don't take it too seriously. My semantics claims that the existence of the event described by the embedded infinitival causes the seeming-event. I've also ignored the semantic contribution that \emph{have} makes. The point I am aiming to make doesn't depend on these shortcuts.}%
	%
	
\end{exe}
\end{fullwidth}
\medskip
This LF provides an interpretation that allows the students who seem to have read assignments to vary with those assignments. A rough paraphrase of the denotation derived for (\ref{ex6:raiseambat}) is~(\ref{ex6:raiseamb1}).
\begin{exe}
\raggedright
	\ex \label{ex6:raiseamb1}
	For every assignment there is a (possibly different) student who seemed to have read it.
\end{exe}

Our syntax also predicts that the LF in (\ref{ex6:raiseambat2}) should be related to~(\ref{ex6:raiseamba}).
\medskip
\begin{fullwidth}
\begin{exe}
\raggedright
	\ex \label{ex6:raiseambat2}
	\hspace*{-20pt}
	\footnotesize
	\begin{forest}
	[$\exists y$ \textsc{student}($y$) $\wedge$ $\exists \epsilon$ \textsc{seem}($\epsilon$) $\wedge$\\ \textsc{cause}($\forall x$ \textsc{assignment}($x$) $\rightarrow$ $\exists \epsilon'$ \textsc{read}($\epsilon'$) $\wedge$ \textsc{patient}($x$)($\epsilon'$) $\wedge$ \textsc{agent}($y$)($\epsilon'$))($\epsilon$)\\TP
		[$\lambda Q\ \exists y$ \textsc{student}($y$) $\wedge$ $Q(y)$ \\DP [some student, roof]]
		[$\lambda 3\ \exists \epsilon$ \textsc{seem}($\epsilon$) $\wedge$\\ \textsc{cause}($\forall x$ \textsc{assignment}($x$) $\rightarrow$ $\exists \epsilon'$ \textsc{read}($\epsilon'$) $\wedge$ \textsc{patient}($x$)($\epsilon'$) $\wedge$ \textsc{agent}($3$)($\epsilon'$))($\epsilon$)\\\ibar{T} 
			[\obar{T} [past] ]
			[$\exists \epsilon$ \textsc{seem}($\epsilon$) $\wedge$\\ \textsc{cause}($\forall x$ \textsc{assignment}($x$) $\rightarrow$ $\exists \epsilon'$\textsc{read}($\epsilon'$) $\wedge$ \textsc{patient}($x$)($\epsilon'$) $\wedge$ \textsc{agent}($3$)($\epsilon'$))($\epsilon$)\\VP [\ibar{V}
				[\obar{V} [$\surd$seem] ]
				[$\forall x$ \textsc{assignment}($x$) $\rightarrow$ {[}$\exists \epsilon'$ \textsc{read}($\epsilon'$) $\wedge$ \textsc{patient}($x$)($\epsilon'$) $\wedge$\\ \textsc{agent}($3$)($\epsilon'$){]}\\TP
					[$\lambda 2\ \exists \epsilon'$ \textsc{read}($\epsilon'$) $\wedge$ \textsc{patient}($2$)($\epsilon'$)\\ $\wedge$ \textsc{agent}($3$)($\epsilon'$) \\TP [\ibar{T}
						[\obar{T} [to] ]
						[vP [\emph{t_3} have read \emph{t_2}, roof] ]
					] ]
					[$\lambda Q\ \forall x$ \textsc{assignment}($x$) $\rightarrow$ $Q$($x$)\\DP [every assignment, roof] ]
				 ]
			] ]
		]
	]
	\end{forest}
\end{exe}
\end{fullwidth}
\medskip
This LF provides an interpretation that doesn't allow the students to vary with the assignments. A rough paraphrase of the denotation derived in (\ref{ex6:raiseambat2}) is~(\ref{ex6:raiseambat2p}).
\begin{exe}
\raggedright
	\ex \label{ex6:raiseambat2p}
	There is a student who seemed to have read every assignment.
\end{exe}
If we could detect that (\ref{ex6:raiseamba}) had these two denotations, then that would be evidence consistent with the predicted LFs.

Unfortunately, these denotations are not sufficiently distinct to know if they exist. In particular, the denotation delivered by (\ref{ex6:raiseambat}) describes a state of affairs that includes as a special case those described by the denotation delivered by (\ref{ex6:raiseambat2}). This is because (\ref{ex6:entailpara}) holds.
\begin{exe}
\raggedright
	\ex \label{ex6:entailpara}
	If it is true that there is a student who seemed to read every assignment, then it is true that for every assignment there is a student who seemed to read it. 
\end{exe}
The denotation delivered by (\ref{ex6:raiseambat2}) entails the denotation delivered by (\ref{ex6:raiseambat}). Because we are not able to detect whether the denotation delivered by (\ref{ex6:raiseambat2}) exists, we cannot determine whether the LF that would deliver it does either. Insofar as the meaning of (\ref{ex6:raiseamba}) is concerned, the LF in (\ref{ex6:raiseambat}) is sufficient. The LF in (\ref{ex6:raiseambat2}) might exist, but we won't be able to detect it from the meaning of~(\ref{ex6:raiseamba}).

For reasons of this kind, testing the predictions of the LF model requires carefully curated examples. But it can be done, and to a large extent the predictions are confirmed. The string in (\ref{ex6:katia}) provides an illustration.
\begin{exe}
\raggedright
	\sn Katia said that \ldots
	\ex \label{ex6:katia}
	some professor seems to have read every assignment of hers to the class.
\end{exe}
The pronoun in (\ref{ex6:katia}) can get one of two interpretations. It can be understood as coreferent with \emph{Katia}, or it can be understood as a variable bound to \emph{some professor}. %
%
\marginnote[-110pt]{It is also possible for \emph{her} to refer to the entity made salient by \emph{some professor}. It turns out that an existential quantificational DP, like \emph{some professor}, cannot make an entity salient for the purposes of pronominal reference if it is in the scope of a universal quantifier. That can be discovered by considering the short conversation in~(i).\begin{exe}
\raggedright
	\exi{i.}
	\begin{xlist}
		\exi{A:} Everyone read some paper.
		\exi{B:} Yes, it was pretty good.
	\end{xlist}
\end{exe}
The pronoun \emph{it} in B's response is felicitous, and capable of referring to the paper introduced by A's sentence, only if \emph{some paper} is not understood to introduce different papers for every reader. For this reason, \emph{her} can refer to the individual made salient by \emph{some professor} in~(\ref{ex6:katia}) only if there aren't different professors for each assignment. As a consequence, we can safely set aside this case. It is irrelevant to the question we are investigating.}%
%
When the pronoun is coreferent with \emph{Katia}, an LF parallel to (\ref{ex6:raiseambat}) (=(\ref{ex6:katiahigh})) is possible and the sentence has a denotation in which the professors can vary with the assignments.
\begin{exe}
\raggedright
	\ex \label{ex6:katiahigh}
	\footnotesize
	\begin{forest}
	[TP
		[TP
			[DP^3 [some professor, roof]]
			[\ibar{T}
				[\obar{T} [past] ]
				[VP [\ibar{V}
					[\obar{V} [$\surd$seem] ]
					[TP [\ibar{T}
						[\obar{T} [to] ]
						[vP [\emph{t_3} have read \emph{t_2} to the class, roof] ]
					] ]
				] ]
			 ]
		]	
		[DP^2 [every assignment of hers, roof]]
	]
	\end{forest}
	\sn $\approx$ for every assignment of hers (=Katia), there is some professor or other that seemed to have read it to the class.
\end{exe}
But when the pronoun is understood as a variable bound to \emph{some professor}, it will have to be c-commanded by \emph{some professor}. This can be achieved if an LF parallel to (\ref{ex6:raiseambat2}) (=(\ref{ex6:katialow})) is also possible.
\begin{exe}
\raggedright
	\ex \label{ex6:katialow}
	\footnotesize
	\begin{forest}
	[TP
		[DP^3 [some professor, roof]]
		[\ibar{T} 
			[\obar{T} [past] ]
			[VP [\ibar{V}
				[\obar{V} [$\surd$seem] ]
				[TP
					[TP [\ibar{T}
						[\obar{T} [to] ]
						[vP [\emph{t_3} have read \emph{t_2} to the class, roof] ]
					]]
					[DP^2 [every assignment of hers_3 to the class, roof] ]
				 ]
			] ]
		]
	]
	\end{forest}
	\sn $\approx$ There is some professor it seemed has read every assignment of hers (=that professor) to the class.
\end{exe}
As expected, the LF in (\ref{ex6:katialow}) does not allow a meaning in which the professors vary with the assignments. By adding the pronoun, then, we have separated the two meanings that are associated with how the quantificational DPs are interpreted, and revealed that they can be individuated in (\ref{ex6:katia}). This suggests, then, that there are two syntactic ways in which the quantificational DPs can be semantically composed, and a syntax which associates one surface parse with more than one LF does just that. 

Another illustration of how to determine whether one S-structure is associated with more than one semantically interpreted syntactic representation is in comparing the sentences in~(\ref{ex6:subjobjscope}).
\pagebreak
\begin{exe}
\raggedright
	\ex \label{ex6:subjobjscope}
	\begin{xlist}
		\ex Some student read every assignment to the class.
		\ex Every student read some assignment to the class.
	\end{xlist}
\end{exe}
As we've seen, associating (\ref{ex6:subjobjscope}a) with the LF in (\ref{ex6:soscopelf1}) provides the right denotation, when it's coupled with the semantics of Trace Theory.
\begin{exe}
\raggedright
	\ex \label{ex6:soscopelf1}
	\small
	\begin{forest}
	[TP
		[TP 
			[DP^3 [some student, roof]]
			[\ibar{T}
				[\obar{T} [past] ]
				[AspP [\ibar{AsP}
				[\obar{Asp} [asp]]
				[vP
					[\emph{t_3} ]
					[\ibar{v}
						[\obar{v} [agent] ]
						[VP [\ibar{V}
							[\ibar{V}
								[\obar{V} [$\surd$read] ]
								[\emph{t_2} ]
							 ]
							[PP [to the class, roof] ]
						] ]
					 ]
				 ]
			 ]
		]]]
		[DP^2 [every assignment, roof]]
	]
	\end{forest}
	\sn $\approx$ For every assignment, there is some student or other that read it to the class.
\end{exe}
But a parallel LF for (\ref{ex6:subjobjscope}b) does not provide the right denotation.
\begin{exe}
\raggedright
	\ex \label{ex6:soscopelf2}
	\small
		\begin{forest}
			[TP
				[TP 
					[DP^3 [every student, roof]]
					[\ibar{T}
						[\obar{T} [past] ]
						[AspP [\ibar{Asp}
						[\obar{Asp} [asp]]
						[vP
							[\emph{t_3} ]
							[\ibar{v}
								[\obar{v} [agent] ]
								[VP [\ibar{V}
									[\ibar{V}
										[\obar{V} [$\surd$read] ]
										[\emph{t_2} ]
									 ]
									[PP [to the class, roof] ]
								] ]
							 ]
						 ]
					 ]
				]]]
				[DP^2 [some assignment, roof]]
			]
			\end{forest}
			\sn $\approx$ It's the case that there is some assignment that every student read to the class.
\end{exe}
This LF does correspond to an interpretation that (\ref{ex6:subjobjscope}b) has. That is, it is consistent with what (\ref{ex6:subjobjscope}b) means: it describes a scenario in which there is one assignment that was read by all of the students, and (\ref{ex6:subjobjscope}b) can be used to describe that scenario.

But (\ref{ex6:subjobjscope}b) is also a correct description of a situation in which every student read a different assignment to the class. While that meaning is not associated with the LF in (\ref{ex6:soscopelf2}), it can be achieved in our syntax by giving to (\ref{ex6:subjobjscope}b) the LF in~(\ref{ex6:soscopelf3}).
\medskip
\begin{fullwidth}
\begin{exe}
\raggedright
	\ex \label{ex6:soscopelf3}
	\small
	\begin{forest}
	[TP
		[DP^3 [every student, roof] ]
		[\ibar{T}
			[\obar{T} [past] ]
			[AspP
				[AspP [\ibar{Asp}
					[\obar{AsP} [asp] ]
					[vP
						[\emph{t_3} ]
						[\ibar{v}
							[\obar{v} [agent] ]
							[VP [\ibar{V}
								[\ibar{V}
									[\obar{V} [$\surd$read] ]
									[\emph{t_2} ]
								 ]
								[PP [to the class, roof] ]
							] ]
						 ]
					 ]
				] ]
				[DP^2 [some assignment, roof] ]
			 ]
		 ]
	]
	\end{forest}
	\sn $\approx$ For every assignment, there is some student or other that read it to the class.
\end{exe}
\end{fullwidth}
\medskip
In a general way, then, the pair of sentences in (\ref{ex6:subjobjscope}) demonstrates that the S-structures they have must correspond to two different semantic interpretations. That can be achieved on May's proposal with our syntax and semantics.

\section{Copy Theory}
\label{sec:copytheory}

The Trace Theory of movement is perhaps the best known interpretation of movement, and it serves as the standard model. But it houses a variety of well-known shortcomings that we will examine in this section. At the time of writing, there isn't a successful solution to these shortcomings, but there is a rough set of ideas about how to solve them that has gained a certain following. In this section, I will present some of these shortcomings and sketch the ideas that are being explored to solve them. 

\subsection{Reconstruction}
\label{sec:recon}

I'll start by considering a series of special cases that support the view that one S-structure can be related to more than one semantically interpreted syntactic structure. What makes these cases special is that a moved term is semantically interpreted, in part or in whole, in a position from which it moved. These cases are sometimes said to involve ``reconstruction,'' as the moved term's meaning derives from ``reconstructing'' a parse in the derivation that feeds the S-structure. Instances of reconstruction pose problems for the Trace theory of movement, and those problems have been an important source for building a better theory.

\subsubsection{Quantifier Lowering}
\label{sec:QuantLower}

The first case of reconstruction that we'll consider is illustrated by the example in~(\ref{ex6:lowerill}).
\begin{exe}
\raggedright
	\ex \label{ex6:lowerill}
	A student didn't attend the meeting.
\end{exe}
With the introduction of Aspect Phrase, our S-structure parse for (\ref{ex6:lowerill}) is~(\ref{ex6:lowerillss}).
\medskip
\begin{fullwidth}
\begin{exe}
\raggedright
	\ex \label{ex6:lowerillss}
	\footnotesize
	\begin{forest}
	[$\exists x$ \textsc{student}($x$) $\wedge$ $\neg\ \exists \epsilon$ \textsc{attend}($\epsilon$) $\wedge$ \textsc{patient}(\denotes{the meeting})($\epsilon$) $\wedge$ \textsc{agent}($x$)($\epsilon$) \\TP
		[$\lambda Q\ \exists x$ \textsc{student}($x$) $\wedge$ $Q(x)$ \\DP^2 [a student, roof]]
		[$\lambda 2\ \neg \exists \epsilon$ \textsc{attend}($\epsilon$) $\wedge$ \textsc{patient}(\denotes{the meeting})($\epsilon$) $\wedge$ \textsc{agent}($2$)($\epsilon$)\\\ibar{T} 
			[\obar{T} [did] ]
			[$\neg \exists \epsilon$ \textsc{attend}($\epsilon$) $\wedge$ \textsc{patient}(\denotes{the meeting})($\epsilon$) $\wedge$ \textsc{agent}($2$)($\epsilon$)\\NegP [\ibar{Neg}
				[\obar{Neg} [neg] ]
				[$\exists \epsilon$ \textsc{attend}($\epsilon$) $\wedge$ \textsc{patient}(\denotes{the meeting})($\epsilon$) $\wedge$ \textsc{agent}($2$)($\epsilon$)\\AspP [\ibar{AsP}
					[\obar{Asp} [asp] ]
					[$\lambda \epsilon$ \textsc{attend}($\epsilon$) $\wedge$ \textsc{patient}(\denotes{the meeting})($\epsilon$) $\wedge$ \textsc{agent}($2$)($\epsilon$)\\vP
						[{\emph{t_2}} ]
						[\ibar{v}
							[\obar{v} [agent] ]
							[VP [\ibar{V}
								[\obar{V} [$\surd$attend] ]
								[DP [the meeting, roof] ]
							] ]
						 ]
					 ]
				] ]
			] ]
		]
	]
	\end{forest}
\end{exe}
\end{fullwidth}
\medskip
This parse puts negation within the interpretation of the determiner \emph{a}, and consequently the interpretation it delivers can be paraphrased by~(\ref{ex6:widep2}).
\begin{exe}
\raggedright
	\ex \label{ex6:widep2}
	There is a student who didn't attend the meeting.
\end{exe}
But (\ref{ex6:lowerill}) has another interpretation in which \emph{a student} is interpreted as part of what is negated. 

Our present syntax provides this parse in the part of the derivation from which the S-structure is produced. A Movement can form (\ref{ex6:lowerdL1}) by moving \emph{a student} into Specifier of AspP. %
%
\marginnote[0pt]{Alternatively, Complex NP Shift could create a parse parallel to (\ref{ex6:soscopelf3}), right-adjoining \emph{a student} to AspP.}%
%
From this parse, A movement can derive the S-structure position of \emph{a student} by again moving it into Specifier of~TP.
\medskip
\begin{fullwidth}
\begin{exe}
\raggedright
	\ex \label{ex6:lowerdL1}
	\footnotesize
	\begin{forest}
	[TP[\ibar{T}
		[\obar{T} [did]]
		[$\neg \exists y$ \textsc{student}($y$) $\wedge$ $\exists \epsilon$ \textsc{attend}($\epsilon$) $\wedge$ \textsc{patient}(the\_meeting)($\epsilon$) $\wedge$ \textsc{agent}($y$)($\epsilon$) \\NegP [\ibar{Neg}
			[\obar{Neg} [neg] ]
				[$\exists x$ \textsc{student}($x$) $\wedge$ $\exists \epsilon$ \textsc{attend}($\epsilon$) $\wedge$ \textsc{patient}(the\_meeting)($\epsilon$) $\wedge$ \textsc{agent}($x$)($\epsilon$)\\AspP 
				[$\lambda Q\ \exists x$ \textsc{student}($x$) $\wedge$ $Q(x)$\\DP^2 [a student, roof] ]
				[$\lambda 2\ \exists \epsilon$ \textsc{attend}($\epsilon$) $\wedge$ \textsc{patient}(the\_meeting)($\epsilon$) $\wedge$ \textsc{agent}($2$)($\epsilon$)\\\ibar{Asp}
					[\obar{Asp} [asp] ]
					[$\lambda \epsilon$ \textsc{attend}($\epsilon$) $\wedge$ \textsc{patient}(the\_meeting)($\epsilon$) $\wedge$ \textsc{agent}($2$)($\epsilon$)\\vP
						[{\emph{t_2}}]
						[\ibar{v}
							[\obar{v} [agent] ]
							[VP [\ibar{V}
								[\obar{V} [$\surd$attend] ]
								[DP [the meeting, roof] ]
							] ]
						 ]
					 ]
				] ]
			 ]
		]]
	]
	\end{forest}
	\normalsize
	\sn $\approx$ It's not the case that there is a student that attended the meeting.
\end{exe}	
\end{fullwidth}
\medskip
On the Lebeaux model, the presence of this parse in the denotation ensures that the meaning associated with it is available to the interpretation of~(\ref{ex6:lowerill}). For this reason, (\ref{ex6:lowerill}) can also have the denotation indicated: what's negated includes \emph{a student}.

Robert May suggested in \citet{May:1985} a different way of deriving this parse that is popular in the literature. His system abandons the Lebeaux model, as we've seen, and requires the final parse in a derivation---his ``LF''---to be the structure that is semantically interpreted. Because May's system allows the S-structure to be a parse in the derivation from which an LF is derived, he is able to capture the fact that sometimes a DP is semantically interpreted in a position higher than where it is spoken. His proposals are not entirely successful, as we'll see in the next section, as they don't include derivations with Head Movement in them. But May does try to accommodate Reconstruction phenomena to his model. We'll look at his attempt to do so because it has features that might prove useful.

May's proposal has two components. One is to let movement only optionally leave traces when XPs are moved. %
%
\marginnote[0pt]{By letting traces be optional, May's proposal also loses the use Trace theory has in ensuring that the Theta Criterion holds of D-structures. See the discussion on page~\pageref{ex6:baddertc}.}%
%
The second is to abandon the Upwards constraint on derivations, and let the presence of traces be responsible for its effects. Putting these together with his goal of letting only LFs be semantically interpretable, we have what I'll call the May model.
\pagebreak
\begin{exe}
\raggedright
	\ex \label{ex6:maymodel} May Model
	\begin{xlist}
		\ex Only an LF is semantically interpreted, and
		\ex Leaving traces is not obligatory for moved XPs, and
		\ex Derivations aren't constrained by Upwards
	\end{xlist}
\end{exe}
The effects of Upwards, which is repeated here, can be derived from the syntax and semantics of traces when they are present.%
% 
\marginnote[0pt]{See \pageref{ex4:upwards} for discussion.}%
%
\begin{exe}
	\raggedright
	\ex \label{ex6:upwards} Upwards\\
	If $\alpha$ moves from position T to position P, then the sister of P must have a projection that dominates~T.
\end{exe}
The syntax and semantics of traces derives the generalization in (\ref{ex6:traceup}) for movement that involves them.
\begin{exe}
	\raggedright
		\ex \label{ex6:traceup}
		Let $\alpha$ be a term that is moved. $\alpha$ must c-command any trace that arises from its movement and is otherwise required.
\end{exe}
It will therefore prevent derivations that would also be blocked by Upwards, like that shown in~(\ref{ex6:uptraceill}).
\begin{exe}
\raggedright
	\ex \label{ex6:uptraceill}
	\small
	\begin{forest}
	[TP
		[DP [it, roof]]
		[\ibar{T} 
			[\obar{T} [might] ]
			[vP
				[\emph{t_2}, name=trace ]
				[\ibar{v}
					[\obar{v} [agent] ]
					[VP [\ibar{V}
						[\obar{V} [$\surd$say] ]
						[CP [\ibar{C}
							[\obar{C} [that] ]
							[TP
								[DP^2, name=binder [she, roof] ]
								[\ibar{T}
									[\obar{T} [pres] ]
									[VP [\ibar{V}
										[\obar{V} [$\surd$seem] ]
										[CP [that pigs fly, roof] ]
									] ]
								 ]
							 ]
						] ]
					] ]
				 ]
			 ]
		]
	]
	\draw[->,>=stealth', thin, gray] (trace) to[out=south west,in=south west,looseness=1] node[pos=0.75,xshift=10mm]{} (binder);
	\end{forest}
\end{exe}
The trace cannot be interpreted as a variable in this representation, because it is not part of a phrase that combines with the DP that should bind it. It doesn't conform to the generalization in~(\ref{ex6:traceup}). 

In (\ref{ex6:uptraceill}), the option of not leaving a trace when \emph{she} moves is not available, since the agent $\theta$-role would not be assigned. Because the May Model doesn't allow D-structures, or any other non-LF parse, to be semantically interpreted, it requires that a parse created by moving \emph{she} must be interpretable. These two ingredients, then, combine to require moved argument DPs to move to a position where they c-command the position they have at D-structure. For movement of argument DPs, then, the results of Upwards can be derived from the syntax and semantics invoked by Trace theory.

But because movement of DPs from non-argument positions doesn't require a trace to be left in that position, the effects of Upwards will not be present in those cases. This will give (\ref{ex6:lowerill}) a second LF, and this, May suggests, is why it has a second interpretation. Imagine that there is a rule that can produce (\ref{ex6:ql1}) from (\ref{ex6:lowerillss}), the S-structure of~(\ref{ex6:lowerill}).
\medskip
\begin{fullwidth}
\begin{exe}
\raggedright
	\ex \label{ex6:ql1}
	\footnotesize
	\begin{forest}
		[TP
		[{\ }, name=trace]
		[\ibar{T}
			[\obar{T} [did]]
			[$\neg \exists y$ \textsc{student}($y$) $\wedge$ $\exists \epsilon$ \textsc{attend}($\epsilon$) $\wedge$ \textsc{patient}(the\_meeting)($\epsilon$) $\wedge$ \textsc{agent}($y$)($\epsilon$) \\NegP [\ibar{Neg}
				[\obar{Neg} [neg] ]
					[$\exists x$ \textsc{student}($x$) $\wedge$ $\exists \epsilon$ \textsc{attend}($\epsilon$) $\wedge$ \textsc{patient}(the\_meeting)($\epsilon$) $\wedge$ \textsc{agent}($x$)($\epsilon$)\\AspP 
					[$\lambda Q\ \exists x$ \textsc{student}($x$) $\wedge$ $Q(x)$\\DP^2, name=binder [a student, roof] ]
					[$\lambda 2\ \exists \epsilon$ \textsc{attend}($\epsilon$) $\wedge$ \textsc{patient}(the\_meeting)($\epsilon$) $\wedge$ \textsc{agent}($2$)($\epsilon$)\\\ibar{Asp}
						[\obar{Asp} [asp] ]
						[$\lambda \epsilon$ \textsc{attend}($\epsilon$) $\wedge$ \textsc{patient}(the\_meeting)($\epsilon$) $\wedge$ \textsc{agent}($2$)($\epsilon$)\\vP
							[{\emph{t_2}}]
							[\ibar{v}
								[\obar{v} [agent] ]
								[VP [\ibar{V}
									[\obar{V} [$\surd$attend] ]
									[DP [the meeting, roof] ]
								] ]
							 ]
						 ]
					] ]
				 ]
			]]
		]
		\draw[->,>=stealth', thin, gray] (trace) to[out=south,in=south west,looseness=2] node[pos=0.75,xshift=10mm]{} (binder);
		\end{forest}
\end{exe}
\end{fullwidth}
\medskip
This could, for instance, be what Argument Movement looks like if Upwards were not present. (\ref{ex6:ql1}) is the same as (\ref{ex6:lowerdL1}), and so provides the same semantic interpretation. On May's system, that parse can be achieved at LF, however, and this would be consistent with the May model's hypothesis that only LFs provide the parse that is semantically interpreted. May calls the operation illustrated by (\ref{ex6:ql1}) ``quantifier lowering.'' This is how these kinds of examples have come to be known.

May's analysis of quantifier lowering depends on the success of his method of deriving the effects earlier ascribed to Upwards. May's method derives a condition on derivations from the syntax and semantics of interpreting variables. It replaces Upwards with (\ref{ex6:traceup}). (\ref{ex6:traceup}) is narrower than Upwards, however. Upwards is tailored to apply to the syntax we've adopted for Head Movement which, as the contrast in (\ref{ex6:hmc}) indicates, should fall under its purview.
\begin{exe}
\raggedright
	\ex \label{ex6:hmc}
	\begin{xlist}
		\ex[]{Hsin-Lun is proud that he dominates in billiards.}
		\ex[*]{Hsin-Lun must proud that he is dominant in billiards.}
	\end{xlist}
\end{exe}
Upwards is obeyed in the syntax our system gives to (\ref{ex6:hmc}a), sketched in~(\ref{ex6:hmca}).
\begin{exe}
\raggedright
	\ex \label{ex6:hmca}
	\small
	\begin{forest}
	[TP
		[DP [Hsin-Lun, roof]]
		[\ibar{T} 
			[\obar{T}
			[\obar{V} [$\surd$be, name=binder] ]
			[\obar{T} [pres] ]
			]
			[VP [\ibar{V}
				[{\ }, name=trace ]
				[AP
					[\obar{A} [proud] ]
					[CP [that he dominates in billiards, roof] ]
				 ]
			]]
		]
	]
	\draw[->,>=stealth', ultra thin, gray] (trace) to[out=south west,in=south,looseness=1] node[pos=0.75,xshift=10mm]{} (binder);
	\end{forest}
\end{exe}
The moved term is the verb \emph{be}, and its sister---\obar{T}---projects \ibar{T} and TP both of which dominate the position from which \emph{be} moves. This isn't true of the syntax our system gives to (\ref{ex6:hmc}b), however, as (\ref{ex6:hmcb}) shows.
\begin{exe}
\raggedright
	\ex \label{ex6:hmcb}
	\small
	\begin{forest}
	[TP
		[DP [Hsin-Lun, roof]]
		[\ibar{T} 
			[\obar{T} [must] ]
			[VP [\ibar{V} 
				[{\ }, name=trace]
				[AP [\ibar{A}
					[\obar{A} [proud] ]
					[CP [\ibar{C}
						[\obar{C} [that] ]
						[TP
							[DP [he, roof] ]
							[\ibar{T}
								[\obar{T}
									[\obar{V}, name=binder [$\surd$be] ]
									[\obar{T} [pres] ]
								 ]
								[VP [dominate in billiards, roof] ]
							 ]
						 ]
					] ]
				] ]
			] ]
		]
	]
	\draw[->,>=stealth', ultra thin, gray] (trace) to[out=-110,in=west,looseness=1] node[pos=0.75,xshift=10mm]{} (binder);
	\end{forest}
\end{exe}

For these cases to be fit into (\ref{ex6:traceup}), the syntax and semantics we've given for binding and traces requires revision. Because \emph{be} doesn't c-command the position it moves from in either of the cases in (\ref{ex6:hmc}), a trace in that position will not be bound by it. Either the syntax of Head Movement would require revision, or the syntax we've developed for movement would. We'll examine this issue in~\ref{sec:semanticrecon}.

Quantifier lowering provides interesting confirmation of the analysis of infinitives described in chapter~\ref{sec:infinitives}. On that analysis there are two related differences between the sentences in~(\ref{ex6:lowerinf}).
\begin{exe}
\raggedright
	\ex \label{ex6:lowerinf}
	\begin{xlist}
		\ex A student is likely to attend the meeting.
		\ex A student wants to attend the meeting.
	\end{xlist}
\end{exe}
One difference is that \emph{a student} is an argument of \emph{want}, but not of \emph{likely}. Connected to this is the status of the term that serves as an argument of \emph{attend} in these examples. In (\ref{ex6:lowerinf}a), the argument of \emph{attend} is \emph{a student}. But in (\ref{ex6:lowerinf}b), the argument of \emph{attend} is a silent pronominal that is bound by \emph{a student.} If we use \textsc{pro} to represent the silent pronoun, then the D-structure parses for (\ref{ex6:lowerinf}) are shown in~(\ref{ex6:lowerinfp}).
\medskip
\begin{fullwidth}
\begin{exe}
\raggedright
	\ex \label{ex6:lowerinfp}
	\footnotesize
	\begin{forest}
	[TP[\ibar{T}
		[\obar{T} [pres] ]
		[VP [\ibar{V}
			[\obar{V} [$\surd$be] ]
			[AP [\ibar{A}
				[\obar{A} [likely] ]
				[TP [\ibar{T}
					[\obar{T} [to] ]
					[vP
						[DP [a student, roof] ]
						[\ibar{v}
							[\obar{v} [agent] ]
							[VP [attend the meeting, roof] ]
						 ]
					 ]
				] ]
			] ]
		] ]
	]]
	\end{forest}
	\hspace*{-50pt}
	\begin{forest}
	[TP
		[\obar{T} [pres] ]
		[vP 
		[DP^2 [a student, roof] ]
		[\ibar{v}
			[\obar{v} [experiencer] ]
			[VP
				[\ibar{V}
					[\obar{V} [$\surd$want] ]
					[CP [\ibar{C}
						[\obar{C} [$\epsilon$] ]
						[TP [\ibar{T}
							[\obar{T} [to] ]
							[vP
								[DP [\textsc{pro}_2, roof] ]
								[\ibar{v} 
									[\obar{v} [agent] ]
									[VP [attend the meeting, roof] ]
								 ]
							 ]
						] ]
					] ]
				 ]
			 ]
		] ]
	]
	\end{forest}
\end{exe}	
\end{fullwidth}
\medskip
If we adopt the analysis of quantifier lowering that preserves Upwards and adopts the Lebeaux model, then from the D-structure for (\ref{ex6:lowerinf}) can be derived~(\ref{ex6:lowerinfl}).
\begin{exe}
\raggedright
	\ex \label{ex6:lowerinfl}
	\footnotesize
	\begin{forest}
		[TP[\ibar{T}
			[\obar{T} [pres] ]
			[VP [\ibar{V}
				[\obar{V} [$\surd$be] ]
				[AP [\ibar{A}
					[\obar{A} [likely] ]
					[TP 
					[DP^2 [a student, roof] ]			
					[\ibar{T}
						[\obar{T} [to] ]
						[vP
							[{\emph{t_2}}]
							[\ibar{v}
								[\obar{v} [agent] ]
								[VP [attend the meeting, roof] ]
							 ]
						 ]
					] ]
				] ]
			] ]
		]]
		\end{forest}
\end{exe}
If this parse is semantically interpreted, it will deliver an interpretation in which \emph{a student} is part of what the sentence claims is likely. That interpretation can be paraphrased with~(\ref{ex6:lowpara}).
\begin{exe}
\raggedright
	\ex \label{ex6:lowpara}
	It's likely that a student will attend a meeting.
\end{exe}
This, indeed, is an interpretation that (\ref{ex6:lowerinf}a) has.

From (\ref{ex6:lowerinfl}), the S-structure representation in (\ref{ex6:higherinfl}) can be derived.
\medskip
\begin{fullwidth}
\begin{exe}
\raggedright
	\ex \label{ex6:higherinfl}
	\footnotesize
	\begin{forest}
	[TP
		[DP^3 [a student, roof] ]
		[\ibar{T}
			[\obar{T} [pres] ]
			[VP [\ibar{V}
				[\obar{V} [$\surd$be] ]
				[AP [\ibar{A}
					[\obar{A} [likely] ]
					[$\exists \epsilon$ \textsc{attend}($\epsilon$) $\wedge$ \textsc{patient}(the\_meeting)($\epsilon$) $\wedge$ \textsc{agent}($3$)($\epsilon$)\\TP
						[\emph{t_3}^2 ]
						[$\lambda 2\ \exists \epsilon$ \textsc{attend}($\epsilon$) $\wedge$ \textsc{patient}(the\_meeting)($\epsilon$) $\wedge$ \textsc{agent}($2$)($\epsilon$)\\\ibar{T}
							[\obar{T} [to] ]
							[$\lambda \epsilon$ \textsc{attend}($\epsilon$) $\wedge$ \textsc{patient}(the\_meeting)($\epsilon$) $\wedge$ \textsc{agent}($2$)($\epsilon$)\\vP
								[\emph{t}_2 ]
								[\ibar{v}
									[\obar{v} [agent] ]
									[VP [attend the meeting, roof] ]
								 ]
							 ]
						 ]
					 ]
				] ]
			] ]
		 ]
	]
	\end{forest}
\end{exe}	
\end{fullwidth}
\medskip
This parse can also be semantically interpreted, but it puts the denotation of \emph{a student} outside the embedded clause and therefore outside what this sentence claims is likely. A paraphrase of this meaning is~(\ref{ex6:hihgint}).
\begin{exe}
\raggedright
	\ex \label{ex6:hihgint}
	There is a student who is likely to attend the meeting.
\end{exe}
This is also an interpretation that (\ref{ex6:lowerinf}a) has. 

Note that the trace left by moving \emph{a student} to the highest Specifier of TP is both a binder and a variable. It is a variable bound by \emph{a student}, and it binds the trace in Specifier of vP. I've given a partial explication of the semantics for (\ref{ex6:higherinfl}) to indicate how this works. (A full explication of the semantics of these sentences requires understanding how predicates combine with embedded clauses, and this is not in our toolkit.)

This situation for (\ref{ex6:lowerinf}a) is not reproduced in (\ref{ex6:lowerinf}b), because the D-structure for (\ref{ex6:lowerinf}b) puts \emph{a student} outside the embedded infinitive. On a theory that accepts Upwards, this prevents Argument Movement from building a parse in which \emph{a student} is part of what (\ref{ex6:lowerinf}b) claims is likely. This predicts, correctly, that (\ref{ex6:lowerinf}a) is ambiguous in a way that (\ref{ex6:lowerinf}b) is not.
\begin{exe}
\raggedright
	\exr{ex6:lowerinf}
	\begin{xlist}
		\ex A student is likely to attend the meeting.
		\begin{xlist}
			\ex[$\approx$]{It is likely that a student will attend the meeting.}
			\ex[$\approx$]{There is a student that is likely to attend the meeting.}
		\end{xlist}
		\ex A student wants to attend the meeting.
		\begin{xlist} 
			\ex[$\neq$]{It is wanted that a student will attend the meeting.}
			\ex[$\approx$]{There is a student that wants to attend the meeting.}
		\end{xlist}
	\end{xlist}
\end{exe}

This same contrast emerges from a theory that abandons Upwards, and derives its effects from the semantics of Trace theory, as May suggests. On that model, the semantically interpreted parse is derived from S-structure; the S-structures of (\ref{ex6:lowerinf}) are shown in~(\ref{ex6:lowerinfs}).
\medskip
\begin{fullwidth}
\begin{exe}
\raggedright
	\ex \label{ex6:lowerinfs}
	\footnotesize
	\begin{forest}
	[TP
		[DP^2 [a student, roof] ]
		[\ibar{T}
			[\obar{T} [pres] ]
			[VP [\ibar{V}
				[\obar{V} [$\surd$be] ]
				[AP [\ibar{A}
					[\obar{A} [likely] ]
					[TP [\ibar{T}
						[\obar{T} [to] ]
						[vP
							[\emph{t_2} ]
							[\ibar{v}
								[\obar{v} [agent] ]
								[VP [attend the meeting, roof] ]
							 ]
						 ]
					] ]
				] ]
			] ]
		 ]
	]
	\end{forest}
	\hspace*{-50pt}
	\begin{forest}
	[TP
		[DP^2  [a student, roof]]
		[\ibar{T} 
			[\obar{T} [pres] ]
			[vP
				[\emph{t_2} ]
				[\ibar{v} 
					[\obar{v} [experiencer] ]
					[CP [\ibar{C}
						[\obar{C} [$\epsilon$] ]
						[TP [\ibar{T}
							[\obar{T} [to]]
							[vP [attend the meeting, roof] ]
						] ]
					] ]
				 ]
			 ]
		]
	]
	\end{forest}
\end{exe}	
\end{fullwidth}
\medskip
Because the trace \emph{t_2} is required in the semantically interpreted parse, movement cannot bring \emph{a student} to a position that doesn't c-command that trace. This will prevent movement from moving \emph{a student} into the embedded infinitive from the S-structure for (\ref{ex6:lowerinf}b) (the second graph in (\ref{ex6:lowerinfs})). But because the trace is inside the embedded infinitive in (\ref{ex6:lowerinf}a), this is a structure that can be derived from the S-structure for (\ref{ex6:lowerinf}a) (the first graph in (\ref{ex6:lowerinfs})).
% %
% %//todo Add discussion of how quantifier lowering is restricted by the type of quantifier involved. Include here the discussion of Buering's case of lowering a universal beneath negation, and Lasnik/SabineandIvy's discussion of lowering into infinitival clauses.
% %

What we learn from quantifier lowering, then, is that the parse that is semantically interpreted can differ from the parse that is pronounced in having a DP in a lower position. We've entertained two ways of modeling this. One allows movement to put the moved term in a position that doesn't c-command the position it moves from. It relies on the semantics of traces to limit where this can happen. This approach, the May model, raises the possibility that only LFs provide the semantically interpretable parse. The other model, the Lebeaux model, preserves Upwards, and doesn't specify which parse in a derivation provides the semantically interpretable parse.

We'll now look at other evidence that bears on choosing between these two views.

\subsubsection{Head Movement and Derivations: a problem}
\label{sec:ordering}

Our present definition of derivation orders the parses that make up our representation of a sentence. If a parse contains a trace, for instance, the definition of a derivation requires that this parse not be a D-structure, but must be related to a D-structure by a series of parses that includes one in which the trace is created by movement. If we assume that the S-structure---which is visible---and the D-structure---which can be determined given the S-structure---are known, then the space of possible parses derivable from movement that fit those two parameters can be calculated. Any of those parses from which the semantic rules can produce a denotation corresponds to a meaning the sentence should be able to have. A component in the account we have for the meanings that a sentence can have, then, is in what our syntax allows a derivation to be. The number of parses in a sentence's derivation that have a semantic interpretation should correspond to the number of denotations that sentence can have. If that number is zero, then we should expect the corresponding sentence to have no denotation. We could interpret this failure of the syntax to provide a denotation as one way a sentence could be judged ungrammatical. In this section, we're going to see that our syntax and semantics don't successfully furnish derivations for certain kinds of grammatical sentences that deliver a denotation. Our model fails to capture the grammaticality of a certain type of sentence. What causes the problem are derivations in which both a head and a phrase move.

Our present semantics for movement doesn't allow a head, a verb say, that has been moved to be semantically interpreted in the position it moves to. Thus, for instance, the S-structure parse for a German main clause can put the main verb into \obar{C} position in the way that (\ref{ex6:v2ill}) illustrates.
\begin{exe}
\raggedright
	\ex \label{ex6:v2ill} \gll Jedes Buch lesen ein M\"adchen.\\ every book read.pres a girl\\ ``A girl reads every book.''
	\vspace*{-20pt}
	\sn
	\small
	\begin{forest}
	[CP
		[DP^2 [jedes Buch, roof, name=binder] ]
		[\ibar{C}
			[\obar{C} 
				[\obar{T}, name=binderv
					[\obar{V} [$\surd$lesen] ]
					[\obar{T} [pres] ]
				 ]
				[\obar{C} [dec] ]
			 ]
			[TP
				[DP [ein M\"adchen, roof] ]
				[\ibar{T}
					[VP [\ibar{V}
						[{\emph{t_2}}, name=trace ]
						[\obar{V} ]
					]]
					[{\ }, name=tracev ]
				 ]
			 ]
		 ]
	]
	\draw[->,>=stealth', ultra thin, gray] (trace) to[out=south west,in=south west,looseness=2] node[pos=0.75,xshift=-10mm]{} (binder);
	\draw[->,>=stealth', ultra thin, gray] (tracev) to[out=south west,in=south,looseness=1] node[pos=0.75,xshift=-10mm]{} (binderv);
	\end{forest}
\end{exe}
But the syntax of Head movement doesn't permit a verb to bind a trace in the position from which it's moved, and so our definition of movement doesn't allow for traces of verbs. As a consequence, the verb in (\ref{ex6:v2ill}) must be semantically interpreted in the position it resides in. But this isn't possible; the sister of \emph{$\surd$lesen} is not the object its $\theta$-role needs, but instead present tense. (\ref{ex6:v2ill}) is semantically anomalous.

A parallel situation arises in English, if verbs move to \obar{v}. The S-structure in (\ref{ex6:Vmill}), for instance, cannot be semantically interpreted given our present syntax/semantics, because the denotation of \emph{$\surd$give} cannot compose with the denotation of \emph{agent}.
\begin{exe}
\raggedright
	\ex \label{ex6:Vmill}
	She will give every toy to Coral.
	\sn
	\small
	\begin{forest}
	[TP
		[DP^2 [she, roof]]
		[\ibar{T} 
			[\obar{T} [will] ]
			[AspP [\ibar{Asp}
			[\obar{Asp} [asp]]
			[vP
				[\emph{t_2} ]
				[\ibar{v}
					[\obar{v}
						[\obar{V} [$\surd$give, name=binder] ]
						[\obar{v} [agent] ]
					 ]
					[VP [\ibar{V}
						[\ibar{V}
							[{\ }, name=trace ]
							[DP [every toy, roof] ]
						 ]
						[PP [to Coral, roof] ]
					] ]
				 ]
			 ]
		]
	]]]
	\draw[->,>=stealth', ultra thin, gray] (trace) to[out=south west,in=south,looseness=1] node[pos=0.75,xshift=-10mm]{} (binder);
	\end{forest}
\end{exe}
In general terms, sentences that involve Head Movement must have derivations in which a non-S-structure parse provides the denotation. If Head Movement is governed by Upwards---as the evidence we have suggests---then this means that the parse which provides the denotation will have to be D-structure, or some parse between D-structure and S-structure that preserves the D-structure position of the head.

But consider now what's required to give a parse for (\ref{ex6:Vmill}) that allows \emph{every toy} to be semantically interpreted. If our semantics for \emph{every} is correct, then \emph{every toy} cannot be interpreted in the position it occupies in~(\ref{ex6:Vmill}). Instead, it will have to move to a position outside the term that binds the event variable; that term is \emph{asp} in our syntax. A parse that meets these twin goals---putting \emph{$\surd$give} and \emph{every toy} in their semantically interpretable positions---is~(\ref{ex6:Vmillp}).
\begin{exe}
\raggedright
	\ex \label{ex6:Vmillp}
	\small
	\begin{forest}
	[TP
		[DP^2 [she, roof]]
		[\ibar{T} 
			[\obar{T} [will] ]
			[AspP
			[AspP
			[\ibar{Asp}
			[\obar{Asp} [asp]]
			[vP
				[\emph{t_2} ]
				[\ibar{v}
					[\obar{v} [agent] ]
					[VP 
					[\ibar{V} [\ibar{V}
						[\obar{V} [$\surd$give] ]
						[\emph{t_3} ]
					]
					[PP [to Coral, roof]]
					] 
					]
				 ]
			 ]
		]
	]
	[DP^3 [every toy, roof]]
	]]]
	\end{forest}
\end{exe}
The derivation for this sentence must be able to contain this parse if it is to deliver the right denotation.

We get independent evidence for a parse of this kind being part of the derivation from the methods we've developed elsewhere to reveal the denotations of sub-constituents of a sentence. For instance, we can learn what the meaning of the VP is from \emph{again}, which, in (\ref{ex6:againVm}), can invoke a reading in which its presupposition doesn't include the subject.%
%
\marginnote[0pt]{\citet{Bale:2007} is a good introduction to the interaction of \emph{again} and quantifiers.}%
%
\begin{exe}
\raggedright
	\ex \label{ex6:againVm}
	At the beginning of the day, the teacher gave every toy to Coral. But now the other kids have taken them from Coral, and she's unhappy, so someone should give every toy to Coral again. 
\end{exe}
The meaning of \emph{again} in this sentence requires that it combine with a constituent that contains the meaning of \emph{$\surd$give} and its arguments, but not the meaning of \emph{agent} and its arguments. (See chapter \ref{sub:again} for the details.) That constituent is provided by (\ref{ex6:Vmillp})---it's the VP---but there is no constituent of that type in~(\ref{ex6:Vmill}).

Similarly, we learn from (\ref{ex6:Vmwide}), that the \emph{to}-phrase can fall within the scope of \emph{every toy}.
\begin{exe}
\raggedright
	\ex \label{ex6:Vmwide}
	She gave every toy to a student.
\end{exe}
This requires that \emph{every toy} be the sister to a constituent that contains \emph{a student}. A parse like that in (\ref{ex6:Vmillp}) would achieve that goal, but a parse like (\ref{ex6:Vmill}) wouldn't. A sentence that combines these two assays, like (\ref{ex6:Vmillc}), shows that (\ref{ex6:Vmillcp}) must be available.
\begin{exe}
\raggedright
	\sn At the beginning of the day, the teacher gave every toy to a student. But now Coral has taken them all for herself, so 
	\ex \label{ex6:Vmillc} Someone should give every toy to a student again.
	\ex \label{ex6:Vmillcp}%
	%
	\marginnote[0pt]{This parse leaves \emph{a student} inside VP, but our semantics for \emph{a} would require it to also be interpreted above \emph{asp}. A parse that meets that requirement would put \emph{a student} in a position adjoined to AspP that is below \emph{every toy}. I simplify the picture here. The semantics given for \emph{again} in chapter ?? does not include how it interacts with quantifiers and traces, so for a complete picture to emerge here would require examining how these interact.}%
	%
	\small
	\begin{forest}
	[TP
		[DP^2 [someone, roof]]
		[\ibar{T} 
			[\obar{T} [should] ]
			[AspP 
			[AspP [\ibar{Asp}
				[\obar{Asp} [asp] ]
				[vP
					[\emph{t_2} ]
					[\ibar{v}
						[\obar{v} [agent] ]
						[VP
						[VP 
							[\ibar{V}
							[\ibar{V}
								[\obar{V} [$\surd$give] ]
								[\emph{t_3} ]
							 ]
							[PP [to a student, roof] ]
						] ]
						[again]
						]
					 ]
				 ]
			]] 
			[DP^3 [every toy, roof]]
			]
		]
	]
	\end{forest}
\end{exe}

The same problem arises for the German sentence in (\ref{ex6:v2ill}). Just as in the English (\ref{ex6:Vmill}), this sentence has argument DPs that are headed by quantificational determiners. If the semantics of these expressions is the same as in English---and they do seem to have the same denotations---then the derivation for (\ref{ex6:v2ill}) must provide a parse in which both subject and object DP have moved out of the vP. That parse must also, however, put the verb root $\surd$\emph{lesen} in its D-structure position. To give the derivation that parse, and also achieve the S-structure in (\ref{ex6:v2ill}) in a way that preserves Upwards isn't possible. 

The challenge, then, is to see how to design a derivation for (\ref{ex6:Vmill}) that includes the parse in (\ref{ex6:Vmillp}) but provides the S-structure in~(\ref{ex6:Vmill}), and then generalize that solution to all the cases, like (\ref{ex6:v2ill}), that share the problem. If derivations obey Upwards, our present system doesn't allow a derivation of the requisite kind. But if the effects of Upwards are derived entirely from the semantics of traces in the way that May suggests, then there is a derivation that does what's required. In that derivation, QR or Complex NP Shift applies before Movement relocates the verb forming the parse in~(\ref{ex6:Vmillcp}). The S-structure is derived from this by using Movement to move the verb into \obar{v}, and moving  \emph{every toy} to the lower position indicated in~(\ref{ex6:lowerobj1}).
\begin{exe}
\raggedright
	\ex \label{ex6:lowerobj1}
	\small
	\begin{forest}
	[TP
		[DP^2 [she, roof]]
		[\ibar{T} 
			[\obar{T} [will] ]
			[AspP [AspP [\ibar{Asp}
				[\obar{Asp} [asp] ]
				[vP
					[\emph{t_2} ]
					[\ibar{v} 
						[\obar{v} 
							[\obar{V} [$\surd$give, name=binderv] ]
							[\obar{v} [agent] ]
						 ]
						[VP [\ibar{V}
						[\ibar{V}
							[\ibar{V} 
								[{\ }, name=tracev]
								[\emph{t_3} ]
							]
							[DP^3, name=binder [every toy, roof] ]
						] 
						[PP [to Coral, roof]]
						]]
					 ]
				 ]
			]]
			[{\ }, name=trace]
			]
		]
	]
	\draw[->,>=stealth', ultra thin, gray] (trace) to[out=south east,in=north east,looseness=1.5] node[pos=0.75,xshift=-10mm]{} (binder);
	\draw[->,>=stealth', ultra thin, gray] (tracev) to[out=south west,in=south,looseness=1] node[pos=0.75,xshift=-10mm]{} (binderv);
	\end{forest}
\end{exe}

Abandoning Upwards, however, brings its own unwanted consequences. We lose the consequences of Upwards for any movement that doesn't have to leave a trace. Because only DPs can leave traces in our present syntax, this means we will have to find another way of ensuring Upward's consequences for moved heads, as well as moved VPs, PP, and~APs.

This section, like the previous, then, shows that we need to allow for derivations in which a moved term is semantically interpreted in a position lower than where it is pronounced. The previous section discussed cases that involve subjects, and for these we have two options. One is to let a parse in the derivation that feeds the S-structure provide the semantic interpretation. The other is to allow for derivations that move terms downwards. The second option maintains the possibility that the parse which the semantic rules interpret is limited to either the S-structure, or a parse derived from an S-structure. 
	
The cases examined in this section, however, involve interpreting a head in a lower position than it occupies at S-structure, and this complicates matters. If we let a parse that feeds the S-structure be semantically interpreted, we can avoid reworking the semantics of traces so that they fit the syntax of Head Movement. Heads can be semantically interpreted in a parse in which they have not moved. But the situation involving quantificational objects in English requires that movement downwards be allowed for moved phrases in that scenario. Alternatively, we can maintain the view that the semantics is only read off of parses derivable from an S-structure, but equip our syntax and semantics with a way of letting moved heads leave a trace. In this scenario, not just phrases, but heads too, will have to be equipped with the ability to move downwards. 

Together, then, these two sections suggest that we should find a way of giving a derivation the ability to create a semantically interpreted parse that derives from an S-structure but allows a term in that S-structure to be interpreted lower than where it sits at S-structure. I will use the term ``reconstruction'' for any scenario of this kind.
	\begin{exe}
	\raggedright
		\ex \label{ex6:reconstrdef} Let an \underline{Upwards-obeying position} for $x$ (=\textsc{up}($x$)) be a position to which Upwards would allow a term in position $x$ to move to, and let an S-structure position (=``s-position'') be the position that a term occupies in the S-structure parse of a derivation. $\alpha$ \underline{reconstructs} to position $\beta$ if $\alpha$ is semantically interpreted in $\beta$ and:
		\begin{xlist}
			\ex $\beta$ isn't the s-position of $\alpha$ and
			\ex $\beta$ isn't \textsc{up}(s-position of $\alpha$)
		\end{xlist}
	\end{exe}
Reconstruction describes cases in which a term is semantically interpreted in neither its S-structure position nor in a position that could be derived by an Upwards-obeying instance of movement from its S-structure position. Cases involving Quantifier Lowering can be alternatively described as situations in which the quantificational DP reconstructs. Similarly, the cases described in this section, where a verb is semantically interpreted in a position lower than its surface position, can be described as situations in which the verb reconstructs.

A law of Reconstruction seems to be~(\ref{ex6:reconlaw}).
\begin{exe}
\raggedright
	\ex \label{ex6:reconlaw} Law of Reconstruction\\
	Let $\alpha$ occupy position $\beta$ at D-structure. If $\alpha$ reconstructs to $\beta '$, then either:
	\begin{xlist}
		\ex $\beta' = \beta$, or
		\ex $\beta '$ is \textsc{up}($\beta$).
	\end{xlist}
\end{exe}
This law says that Reconstruction is only possible to the D-structure position of the term reconstructing, or to a position that an Upwards-obeying instance of movement could take that term to from its D-structure position. To date, at any rate, I do not know cases of Reconstruction which violate (\ref{ex6:reconlaw}). For instance, it is not possible to interpret (\ref{ex6:lawill1}) in such a way that puts \emph{a student} within the scope of \emph{shouldn't}.
\begin{exe}
\raggedright
	\ex \label{ex6:lawill1} A student is likely to say that this shouldn't be her problem.
\end{exe}
That would be a possibility if the S-structure for (\ref{ex6:lawill1}) could be related to the parse in~(\ref{ex6:badlower}).
\begin{exe}
\raggedright
	\ex \label{ex6:badlower}
	\footnotesize
\begin{forest}
[TP
	[\obar{T} [pres] ]
	[AP [\ibar{A}
		[\obar{A} [likely] ]
		[TP [\ibar{T}
			[\obar{T} [to] ]
			[vP
				[\emph{t_2} ]
				[\ibar{v}
					[\obar{v} [agent] ]
					[VP [\ibar{V}
						[\obar{V} [$\surd$say] ]
						[CP [\ibar{C}
							[\obar{C} [that] ]
							[TP
								[DP [this, roof] ]
								[\ibar{T}
									[\obar{T} [should] ]
									[NegP [\ibar{Neg}
										[\obar{Neg} [nt] ]
										[VP
											[DP^2 [a student, roof] ]
											[VP [\ibar{V}
												[\obar{V} [be] ]
												[DP [her problem, roof] ]
											] ]
										 ]
									] ]
								 ]
							 ]
						] ]
					]]
				 ]
			 ]
		] ]
	] ]
]
\end{forest}
\end{exe}
But this violates the Law of Reconstruction, as the position that \emph{a student} occupies in (\ref{ex6:badlower}) is not \textsc{up}(\emph{t_2}): \emph{t_2} being the position \emph{a student} occupies at D-structure.

The Law of Reconstruction is straightforwardly derived on the Lebeaux model of the syntax/semantics interface, if it is wedded to Upwards. By contrast, the May Model does not derive the Law of Reconstruction. The May Model also fails, as we've seen, to derive Upwards for terms that don't leave traces---like Verbs---and it wrongly allows D-structures to put arguments in non-$\theta$-marked positions. The May Model is presently a live contender in the literature at large. But let us abandon it in these lectures. 

Instead, let's provisionally adopt Upwards and the Lebeaux Model. 
\begin{exe}
\raggedright
	\ex \label{ex6:lebv} Lebeaux+Upwards Model
	\begin{xlist}
		\ex Derivations are constrained by Upwards,
		\ex If P is a semantically interpretable parse in a derivation for S, then \denotes{P} is a possible denotation for~S.
	\end{xlist}
\end{exe}
The Lebeaux+Upwards model also suffers from the ordering problems this section describes. It requires that a derivation include a parse in which quantificational DPs be interpreted in their moved position, but verbs (and other heads) be interpreted in their un-moved positions. If derivations are governed by Upwards, our present syntax does not let derivations have those parses. Unlike the May Model, however, there is an idea about how to fix these problems for the Lebeaux+Upwards model. The solution we'll investigate below involves changing the syntax and semantics of movement. Before looking at the solution, let's first look at another kind of environment where reconstruction arises.

\subsubsection{Reconstructed Variables}
\label{sec:reconvariable}

In the previous two sections we have seen cases in which a term ``reconstructs.'' We've seen that the denotation of the determiner \emph{a} can be interpreted in a position lower than it's pronounced, and we've interpreted that as a situation in which the DP it heads reconstructs. We've also seen that the denotation of a verb can be interpreted in a position lower than where it's pronounced: verbs too can reconstruct. In this section we'll examine another case of reconstruction. Expressions that are interpreted as variables can reconstruct.

One case of this is illustrated by~(\ref{ex6:reconpron}).%
%
\marginnote[0pt]{The first to examine these cases is, perhaps, \citet{Engdahl:1982}. A pretty thorough examination of cases like this is in \citet{Barss:1986b}.}%
%
\begin{exe}
\raggedright
	\ex \label{ex6:reconpron}
	Which of her articles does no woman remember accurately?
\end{exe}
The derivation for (\ref{ex6:reconpron}) includes a parse in which \emph{which of her articles} resides in the complement position of \emph{remember}. Suppose that parse has been created by movement of the subject argument; an abbreviated picture of that parse could be~(\ref{ex6:reconpronp}).
\begin{exe}
\raggedright
	\ex \label{ex6:reconpronp}
	\footnotesize
	\begin{forest}
	[CP [\ibar{C}
		[\obar{C} [Q] ]
		[TP
			[DP^2 [no woman, roof] ]
			[\ibar{T}
				[\obar{T} [does] ]
				[vP
					[\emph{t_2} ]
					[\ibar{v}
						[\obar{v} 
							[\obar{V} [$\surd$remember, name=binder] ]
							[\obar{v} [experiencer] ]
						 ]
						[VP [\ibar{V}
							[\ibar{V}
							[{\ }, name=trace ]
							[DP [which of her_2 articles, roof] ]
						] 
						[AdvP [accurately, roof]]
						]
						]
					 ]
				 ]
			 ]
		 ]
	]]
	\draw[->,>=stealth', thin, gray] (trace) to[out=south west,in=south,looseness=1] node[pos=0.75,xshift=-10mm]{} (binder);
	\end{forest}
\end{exe}
If the pronoun is equipped with the same referential index that the trace of moving the subject has, then this parse allows the pronoun an interpretation in which it is a variable bound by \emph{no woman}. This is indeed an interpretation (\ref{ex6:reconpron}) can get. That interpretation can be paraphrased with~(\ref{ex6:reconpronpara}).
\begin{exe}
\raggedright
	\ex \label{ex6:reconpronpara}
	What is the article such that for no $x$, $x$ is a woman and $x$ remembers accurately that article of $x$'s.
\end{exe}
This question seeks the identity of articles that no woman has written and remembers accurately. The identity of the articles, then, varies with the value given to the possessive pronoun. That can be achieved by making the pronoun a variable bound to \emph{no woman}. 

\Citet{Engdahl:1982} shows that the ability of the pronoun in sentences like these to get the bound variable interpretation depends on there being a parse in the derivation in which that pronoun is c-commanded by the binder. That a pronoun can only be a variable bound by a c-commanding DP derives, recall, from the semantics we have given to movement. When a phrase moves, like \emph{no woman} in (\ref{ex6:reconpronp}), then the phrase it becomes a sister to turns into a $\lambda$-expression. That $\lambda$-expression will bind every referential index within that the $\lambda$-operator is related to. For a pronoun to be a variable, it must bear one of those referential indices. It must, in other words, be within the phrase that is the sister to its binder. This condition holds in (\ref{ex6:reconpronp}): \emph{her} is within the phrase---\ibar{T}---that is the sister to its binder. By contrast, there is no parse in the derivation of (\ref{ex6:reconpron2}) in which that relation holds between \emph{no woman} and \emph{her}.
\begin{exe}
\raggedright
	\ex \label{ex6:reconpron2}
	Which of her articles revealed that no woman had failed.
\end{exe}
The D-structure for (\ref{ex6:reconpron2}) is sketched in~(\ref{ex6:reconpron2p}).
\begin{exe}
\raggedright
\small
	\ex \label{ex6:reconpron2p}
	\footnotesize
	\begin{forest}
	[CP [\ibar{C}
		[\obar{C} [Q] ]
		[TP [\ibar{T}
			[\obar{T} [past] ]
			[vP
				[DP [which of her articles, roof] ]
				[\ibar{v}
					[\obar{v} [cause] ]
					[VP [\ibar{V}
						[\obar{V} [$\surd$reveal] ]
						[CP [\ibar{C}
							[\obar{C} [that] ]
							[TP [\ibar{T}
								[\obar{T} [past] ]
								[VP [\ibar{V}
									[\obar{V} [$\surd$have] ]
									[vP
										[DP [no woman, roof] ]
										[\ibar{v}
											[\obar{v} [agent] ]
											[VP [\ibar{V} [\obar{V} [$\surd$fail]]] ]
										 ]
									 ]
								] ]
							] ]
						] ]
					] ]
				 ]
			 ]
		] ]
	]]
	\end{forest}
\end{exe}
The phrase \emph{which of her articles} is not c-commanded by \emph{no woman} in this parse, and because movement will only move \emph{which of her articles} into higher positions, there is no parse derivable from this one in which \emph{no woman} c-commands \emph{which of her pictures}. The prediction then is that \emph{her} cannot be interpreted as a variable bound by \emph{no woman} in (\ref{ex6:reconpron2}). That prediction is confirmed. If \emph{her} could be a variable bound to \emph{no woman} in (\ref{ex6:reconpron2}), then it should have an interpretation like that paraphrased in~(\ref{ex6:reconpron2i}).
\begin{exe}
\raggedright
	\ex \label{ex6:reconpron2i}
	What is the article such that for no $x$, $x$ is a woman and that article of $x$'s reveals that $x$ had failed..
\end{exe}
That is not what (\ref{ex6:reconpron2}) means. Instead, it gets an interpretation in which \emph{her} refers to some female individual whose identity the context must supply. The resulting meaning is roughly:
\begin{exe}
\raggedright
	\ex \label{ex6:reconpron2ib}
	What is the article such that that article of hers reveals that no woman had failed.
\end{exe}

Another example that demonstrates the same thing is (\ref{ex6:reconprondisamb}).
\begin{exe}
\raggedright
	\ex \label{ex6:reconprondisamb}
	Which of her articles did a student that no teacher recommended remember? 
\end{exe}
The D-structure for (\ref{ex6:reconprondisamb}) is sketched in~(\ref{ex6:reconprondisambp}).
\begin{exe}
\raggedright
	\ex \label{ex6:reconprondisambp}
	\small
	\begin{forest}
	[CP [\ibar{C}
		[\obar{C} [Q] ]
		[TP [\ibar{T}
			[\obar{T} [did] ]
			[vP
				[DP [a student that no teacher recommended, roof]
				 ]
				[\ibar{v}
					[\obar{v} [agent] ]
					[VP [\ibar{V}
						[\obar{V} [$\surd$remember] ]
						[DP [which of her articles, roof] ]
					] ]
				 ]
			 ]
		] ]
	]]
	\end{forest}
\end{exe}
The pronoun \emph{her} is not c-commanded by \emph{no teacher} in (\ref{ex6:reconprondisambp}), and movement is incapable of creating from (\ref{ex6:reconprondisambp}) a parse in which that happens. As expected, then, (\ref{ex6:reconprondisamb}) cannot have an interpretation like that described by~(\ref{ex6:reconprondisambi}).
\begin{exe}
\raggedright
	\ex \label{ex6:reconprondisambi}
	What is the article such that for no $x$, $x$ is a teacher there is a student that $x$ recommended who remembers that article of $x$'s.
\end{exe}
This is the interpretation that would arise if \emph{her} could be a variable bound to \emph{no teacher}. 

Instead, \emph{her} can be interpreted as a variable for the expression \emph{a student no teacher recommended}. %
%
\marginnote[0pt]{If (\ref{ex6:reconprondisamb}) is uttered in a context which made salient a female individual, then (\ref{ex6:reconprondisamb}) could also get an interpretation in which \emph{her} refers to that individual.}%
%
 That's because there is a derivation for (\ref{ex6:reconprondisamb}) that includes the parse in~(\ref{ex6:reconprondisambpp}).
\begin{exe}
\raggedright
	\ex \label{ex6:reconprondisambpp}%
	\small
	\begin{forest}
	[CP [\ibar{C}
		[\obar{C} [Q] ]
		[TP
			[DP^2 
				[a student no teacher recommended, roof]
			 ]
			[\ibar{T}
				[\obar{T} [did] ]
				[vP
					[\emph{t_2} ]
					[\ibar{v} 
						[\obar{v} [agent] ]
						[VP [\ibar{V}
							[\obar{V} [$\surd$remember] ]
							[DP [which of her_2 articles, roof] ]
						] ]
					]
				 ]
			 ]
		 ]
	]]
	\end{forest}
\end{exe}
This parse puts \emph{her} in a position that allows it to be bound by the moved phrase \emph{a student no teacher recommended}.
%
\marginnote[12pt]{Note that \emph{no teacher} still doesn't c-command \emph{her} in this parse.}%
%
Consequently, it gives (\ref{ex6:reconprondisamb}) the interpretation described by~(\ref{ex6:reconprondisambii}).
\begin{exe}
\raggedright
	\ex \label{ex6:reconprondisambii}
	What is the article such that a student no teacher recommended, $y$, remembers that article of~$y$'s.
\end{exe}
The pronoun in (\ref{ex6:reconprondisamb}) can be a variable bound to \emph{a student no teacher recommended}, but not to \emph{no woman}. This is just what is expected if the pronoun's interpretation can be computed from a position it occupies in a non-s-structure parse in the derivation.

These cases of reconstruction can be captured by the Lebeaux+Upwards model, so they do not pose a problem for our present theory of movement. But they can be used to create the same kind of ordering problem that is described in the previous section. Consider, for instance,~(\ref{ex6:reconvro1}).
\begin{exe}
\raggedright
	\ex \label{ex6:reconvro1}
	Which of her friends should no woman read some paper to?
\end{exe}
This can get an interpretation in which \emph{her} is a variable bound by \emph{no woman} and \emph{some paper} is interpreted outside the scope of \emph{no woman}. The parse that will deliver that interpretation puts \emph{which of her friends} in its unmoved position, and \emph{some paper} is a position higher than \emph{no woman}. For instance, (\ref{ex6:reconvro1t}) has the right configuration.
\begin{exe}
\raggedright
	\ex \label{ex6:reconvro1t}
	\small
	\begin{forest}
	[CP
		[\ibar{C}
		[\obar{C} [Q]]
		[TP
			[TP
				[DP^2 [no woman, roof] ]
				[\ibar{T} 
					[\obar{T} [should]]
					[vP
						[\emph{t_2} ]
						[\ibar{v} 
							[\obar{v} [agent] ]
							[VP [\ibar{V}
								[\ibar{V}
									[\obar{V} [$\surd$read] ]
									[\emph{t_3} ]
								 ]
								[PP [to which of her_2 friends, roof] ]
							] ]
						 ]
					 ]
				]
			 ]
			[DP^3 [some paper, roof] ]
		 ]
	]]
	\end{forest}
\end{exe}
There is no derivation that obeys Upwards which has this parse as well as an S-structure that corresponds to~(\ref{ex6:reconvro1}).

\subsubsection{Achieving reconstruction with traces}
\label{sec:semanticrecon}

These cases of reconstruction show that we need a syntax which allows a moved term to be semantically interpreted in a position it occupies before moving. In the previous sections, I've sketched how this is accounted for on the Lebeaux+Upwards Model, which allows the semantic interpretation of a sentence to be derived from any element in its derivation. Because the Lebeaux+Upwards Model allows one to determine the elements of a derivation from the meanings a sentence has, it provides a method for determining what a derivation consists of. If our model gives a sentence a derivation that consists of elements which are semantically interpretable, but those interpretations are not available to the sentence, then we can read this as evidence that we have the wrong model of a derivation. We should seek a model of derivations that prevents parses whose meanings are not associated with the sentence in question. Conversely, if a sentence has an interpretation that doesn't correspond to a parse in its derivation, then we should modify our model of derivations so that it includes that parse. On the Lebeaux+Upwards model, then, we can use the meanings a sentence has as a guide to the parses in its derivation.

A derivation is a syntactic object. It is a set of parses that represents a sentence. Reconstruction provides us a view of that set of syntactic parses through the meanings associated with those parses. The view reconstruction gives us of a derivation therefore rests on how parses are semantically interpreted. As our model of the semantic interpretation of parses matures, so too will our view of the parses that make up a derivation. If, for instance, our model of semantic interpretation associates a denotation to the parse~P but not the parse~Q, then we will only be able to determine whether the derivations we design are correct with respect to~P. Our model of derivations will predict that sentences that have~P in their derivation will also have the meaning associated with~P, and that sentences which don't contain~P in their derivation won't have the meaning associated with~P (unless that meaning can be derived from some other parse in the derivation). A model of semantic interpretation of this kind will not allow us to determine anything about the distribution of~Q in the parses of a sentence, however. As a consequence, we will not be able to determine whether one theory of derivations is better than another by virtue of what those theories predict about~Q. This is one place, then, where the model we have of semantic interpretation and the model we have of syntactic representation are intertwined. They cannot be developed independently.

This section illustrates this relationship between the semantic model used and the kinds of derivations it fits. I will do this by considering a popular extension of Trace Theory that allows traces to be variables of different semantic types. %
%
\marginnote[0pt]{Perhaps the most extensive development of the technique that I'll describe here can be found in \citet{Cresti:1995}.}%
%
A consequence of this alternative semantics for traces changes the way in which reconstruction arises, and therefore gives a different image of how derivations fit to semantic interpretations. As we'll see, this extension to Trace Theory also affects how we can model the syntax/semantics interface. It will provide a way of modeling ambiguity that allows one parse in a derivation to provide two denotations that in our present theory can only be delivered by considering two parses in a derivation. The extension, then, not only gives us a different view of the parses that make up derivations, but also a different view of how those parses need to be sampled in order to get the correct meanings for a sentence. 

In our present theory, traces are interpreted as variables over expressions that are entities. That is, the functions that movement creates has entities as its domain. Let's reconsider how this works in~(\ref{ex6:lowerill}).
\begin{exe}
\raggedright
	\exr{ex6:lowerill}
	A student didn't attend the meeting.
\end{exe}
Argument movement of \emph{a student} allows for a trace to be left in the Specifier of vP, from which it moves, and causes the \ibar{T} that \emph{a student} ends being a sister to to be interpreted as a function with that trace as its variable. This is illustrated in~(\ref{ex6:lowerillhighparse}), which we can think of as approximating the S-structure representation of~(\ref{ex6:lowerill}).
\medskip
\begin{fullwidth}
\begin{exe}
\raggedright
	\ex \label{ex6:lowerillhighparse}
	\small
	\begin{forest}
	[$\exists x$ \textsc{student}($x$) $\wedge\ \neg\ \exists \epsilon$ \textsc{attend}($\epsilon$) $\wedge$ \textsc{patient}(the\_meeting)($\epsilon$) $\wedge$ \textsc{agent}($x$)($\epsilon$)  \\TP
		[$\lambda P\ \exists x$ \textsc{student}($x$) $\wedge$ $P(x)$\\DP [a student, roof] ]
		[$\lambda 2\ \neg \exists \epsilon $ \textsc{attend}($\epsilon$) $\wedge$ \textsc{patient}(the\_meeting)($\epsilon$) $\wedge$ \textsc{agent}($2$)($\epsilon$) \\ \ibar{T}
			[\obar{T} [did] ]
			[NegP [\ibar{Neg}
				[\obar{Neg} [neg] ]
				[AspP [\ibar{Asp}
					[\obar{Asp} [asp] ]
					[vP 
						[\emph{t_2} ]
						[\ibar{v}
							[\obar{v} [agent] ]
							[VP [\ibar{V}
								[\obar{V} [$\surd$attend] ]
								[DP [the meeting, roof] ]
							] ]
						 ]
					 ]
				] ]
			] ]
		 ]
	]
	\end{forest}
\end{exe}
\end{fullwidth}
\vspace*{12pt}
This parse provides a denotation that puts \emph{a student} outside the scope of negation, and its existence in the derivation for (\ref{ex6:lowerill}) is responsible for allowing (\ref{ex6:lowerill}) to have this meaning. Our syntax also allows the derivation for (\ref{ex6:lowerill}) to include~(\ref{ex6:lowerilllowparse}). This parse, which corresponds to neither the S-structure nor the D-structure for (\ref{ex6:lowerill}), has a denotation that puts \emph{a student} within the scope of negation. Its existence, then, is responsible for (\ref{ex6:lowerill}) having this denotation as well. The S-structure representation for (\ref{ex6:lowerill}) can be derived from (\ref{ex6:lowerilllowparse}) by using movement to bring \emph{a student} into the Specifier of TP. And it is, in turn, derived by moving \emph{a student} into Specifier of AspP, from its D-structure position in Specifier of vP. As in (\ref{ex6:lowerillhighparse}), the interpretation of (\ref{ex6:lowerilllowparse}) uses a semantics of movement that leaves a trace that is interpreted as a variable over entities. And for that reason, the function that \ibar{Asp} is converted into by movement is one whose domain is entities.
\medskip
\begin{fullwidth}
\begin{exe}
\raggedright
	\ex \label{ex6:lowerilllowparse}
	\footnotesize
	\begin{forest}
	[$\neg \exists x$ \textsc{student}($x$) $\wedge\ \exists \epsilon$ \textsc{attend}(the\_meeting)($\epsilon$) $\wedge$ \textsc{agent}($x$)($\epsilon$) \\TP [\ibar{T}
		[\obar{T} [did] ]
		[NegP [$\neg \exists x$ \textsc{student}($x$) $\wedge\ \exists \epsilon$ \textsc{attend}(the\_meeting)($\epsilon$) $\wedge$ \textsc{agent}($x$)($\epsilon$) \\\ibar{Neg}
			[\obar{Neg} [neg] ]
			[$\exists x$ \textsc{student}($x$) $\wedge\ \exists \epsilon$ \textsc{attend}(the\_meeting)($\epsilon$) $\wedge$ \textsc{agent}($x$)($\epsilon$) \\AspP
				[$\lambda P\ \exists x$ \textsc{student}($x$) $\wedge$ $P(x)$\\DP [a student, roof] ]
				[$\lambda 2\ \exists \epsilon $ \textsc{attend}($\epsilon$) $\wedge$ \textsc{patient}(the\_meeting)($\epsilon$) $\wedge$ \textsc{agent}($2$)($\epsilon$) \\\ibar{Asp}
					[\obar{Asp} [asp] ]
					[vP
						[\emph{t_2} ]
						[\ibar{v}
							[\obar{v} [agent] ]
							[VP [\ibar{V}
								[\obar{V} [$\surd$attend] ]
								[DP [the meeting, roof] ]
							] ]
						 ]
					 ]
				 ]
			 ]
		] ]
	]
	]
	\end{forest}
\end{exe}
\end{fullwidth}
\medskip

Consider now an alternative semantics that allows a trace to be a variable over things of various semantic types. In addition to being a variable over things, or ``entities,'' for instance, imagine that a trace could be interpreted as a variable over objects of the same semantic kind of thing that the moved phrase is. For instance, if the moved term has a denotation which makes it a function, let the trace left by that term be a variable over the same functions. I'll indicate the semantic type of a trace with a subscript on the index it bears. The symbol in (\ref{ex6:suptype}a), for instance, will indicate a trace which will be interpreted as a variable over things, or ``entities.'' The symbol in (\ref{ex6:suptype}b), by contrast, will indicate a trace whose denotation makes it a variable over functions of the same kind as the term whose movement creates the trace.
\begin{exe}
\raggedright
	\ex \label{ex6:suptype}
	\begin{xlist}
		\ex \emph{t}_{2}_{e}
		\ex \emph{t}_{2}_{f}
	\end{xlist}
\end{exe}

Allowing this flexibility to the semantics of traces doesn't change how things work in (\ref{ex6:lowerillhighparse}) and (\ref{ex6:lowerilllowparse}), because in these parses the trace cannot compose semantically with its sister unless it is of entity type. To see this, consider (\ref{ex6:lowerillhighparse}), but with the type of the trace adjusted so that it matches that of the moved DP, \emph{a student}.
\medskip
\begin{fullwidth}
\begin{exe}
\raggedright
	\ex \label{ex6:lowerillhighflex}
	\small
	\begin{forest}
	[TP
		[$\lambda P\ \exists x$ \textsc{student}($x$) $\wedge$ $P(x)$\\DP [a student, roof] ]
		[\ibar{T}
			[\obar{T} [did] ]
			[NegP [\ibar{Neg}
				[\obar{Neg} [neg] ]
				[AspP [\ibar{Asp}
					[\obar{Asp} [asp] ]
					[vP 
						[\emph{t_{2}_{f}} ]
						[$\lambda y\ \lambda \epsilon$ \textsc{attend}($\epsilon$)  $\wedge$ \textsc{patient}(the\_meeting)($\epsilon$) $\wedge$ \textsc{agent}($y$)($\epsilon$) \\\ibar{v}
							[\obar{v} [agent] ]
							[VP [\ibar{V}
								[\obar{V} [$\surd$attend] ]
								[DP [the meeting, roof] ]
							] ]
						 ]
					 ]
				] ]
			] ]
		 ]
	]
	\end{forest}
\end{exe}
\end{fullwidth}
\medskip
Consider how the trace in Specifier of vP combines semantically with its \ibar{v} sister in (\ref{ex6:lowerillhighflex}). Because the domain of the function \textsc{agent} is entities, the function that \ibar{v} denotes also has that domain. If the semantic composition were to treat the trace as the argument of the function denoted by \ibar{v} the result would be (\ref{ex6:flexbad1}), which gives \emph{agent} an argument of the wrong semantic type.
\begin{exe}
\raggedright
	\ex \label{ex6:flexbad1}
	$\neg \exists \epsilon$ \textsc{attend}($\epsilon$) $\wedge$ \textsc{patient}(the\_meeting)($\epsilon$) $\wedge$ \textsc{agent}(\emph{t_2}_{f})($\epsilon$)
\end{exe}
However, the trace is a variable over functions, and this is therefore also the semantic type of the domain of the function that the denotation of \ibar{v} is. Thus, the trace can take \ibar{v} as its argument. We can represent this with~(\ref{ex6:hightcomb1}).
\begin{exe}
\raggedright
	\ex \label{ex6:hightcomb1}
	\denotes{vP} = [$2$_{f}]($\lambda y\ \lambda \epsilon$ \textsc{attend}($\epsilon$) $\wedge$ \textsc{patient}(the\_meeting)($\epsilon$) $\wedge$ \textsc{agent}($y$)($\epsilon$))
\end{exe}
A snapshot of the parse in (\ref{ex6:lowerillhighflex}) with the relevant denotations revealed is~(\ref{ex6:hightcombt}).
\medskip
\begin{fullwidth}
\begin{exe}
\raggedright
	\ex \label{ex6:hightcombt}
	\small
	\begin{forest}
	[TP
		[$\lambda P\ \exists x$ \textsc{student}($x$) $\wedge$ $P(x)$\\DP [a student, roof] ]
		[\ibar{T}
			[\obar{T} [did] ]
			[NegP [\ibar{Neg}
				[\obar{Neg} [neg] ]
				[AspP [\ibar{Asp}
					[$\lambda P\ \exists \epsilon\ P(\epsilon)$\\\obar{Asp} [asp] ]
					[{[}$2$_{f}{]}($\lambda y\ \lambda \epsilon$ \textsc{attend}($\epsilon$) $\wedge$ \textsc{patient}(the\_meeting)($\epsilon$) $\wedge$ \textsc{agent}($y$)($\epsilon$))\\vP 
						[\emph{t_{2}_{f}} ]
						[\ibar{v}
							[\obar{v} [agent] ]
							[VP [\ibar{V}
								[\obar{V} [$\surd$attend] ]
								[DP [the meeting, roof] ]
							] ]
						 ]
					 ]
				] ]
			] ]
		 ]
	]
	\end{forest}	
\end{exe}
\end{fullwidth}
\medskip
The domain of \emph{asp} is functions from eventualities to truth values. It takes a description of eventualities and creates a sentence that says there is such an eventuality. The semantic type of vP in (\ref{ex6:hightcombt}) isn't the right semantic type for \emph{asp}. It isn't a description of eventualities, but is instead a complicated sentence: whatever value the index $2$ gets applied to the argument enclosed in parentheses. For this reason, taking the \ibar{v} in (\ref{ex6:lowerillhighflex}) to provide an argument to the trace in Specifier of TP, also fails to provide a denotation of this sentence.

Our semantics, then, blocks an interpretation of (\ref{ex6:lowerillhighparse}), where the trace created by moving \emph{a student} is interpreted as a variable over the same functions that \emph{a student} is. For entirely parallel reasons, letting the trace in (\ref{ex6:lowerilllowparse}) have that semantic type will fail to yield a denotation. For these two cases, then, allowing traces to be anything but variables over entities has no consequences. Our semantics doesn't allow traces to exploit that flexibility in these parses. The denotations of the other material in the parse fits together only if the traces in (\ref{ex6:lowerillhighparse}) and (\ref{ex6:lowerilllowparse}) are interpreted as variables over entities.

Where allowing traces to be variables over functions has consequences is when traces aren't arguments of $\theta$-roles. To see this, consider the parse that arises from movement applying to \emph{a student} in (\ref{ex6:lowerilllowparse}). This movement operation can leave a trace in Specifier of AspP, and this trace can successfully be a variable over things of the same semantic type that \emph{a student} is. (\ref{ex6:lowerhitype}) is a sketch of that parse.
\medskip
\begin{fullwidth}
	\begin{exe}
		\raggedright
			\ex \label{ex6:lowerhitype}
			\small
			\begin{forest}
			[TP
				[DP [a student, roof] ]
				[\ibar{T}
					[\obar{T} [did] ]
					[NegP [\ibar{Neg}
						[\obar{Neg} [neg] ]
						[${[}2_{f}{]}$($\lambda 3\ \exists \epsilon$ \textsc{attend}($\epsilon$) $\wedge$ \textsc{patient}(the\_meeting)($\epsilon$) $\wedge$ \textsc{agent}($3$)($\epsilon$)) \\AspP
							[\emph{t_{2}_{f}} ]
							[$\lambda 3_e\ \exists \epsilon$ \textsc{attend}($\epsilon$) $\wedge$ \textsc{patient}(the\_meeting)($\epsilon$) $\wedge$ \textsc{agent}($3_e$)($\epsilon$) \\\ibar{Asp}
								[\obar{Asp} [asp] ]
								[vP
									[\emph{t_{3}_e} ]
									[\ibar{v}
										[\obar{v} [agent] ]
										[VP [\ibar{V}
											[\obar{V} [$\surd$attend] ]
											[DP [the meeting, roof] ]
										] ]
									 ]
								 ]
							 ]
						 ]
					] ]
				 ]
			]
			\end{forest}
		\end{exe}
\end{fullwidth}
Note that movement of \emph{a student} into Specifier of AspP from Specifier of vP has left a trace of entity type. This is going to give the function \textsc{agent} something of the type that is in its domain. When subsequent movement of \emph{a student} into Specifier of TP occurs, it's this movement that leaves a trace of function type. If we assume that the denotation of \emph{neg} is (\ref{ex6:denneg}), then it's possible for the denotation of AspP and \obar{Neg} to semantically combine.
\medskip
\begin{exe}
\raggedright
	\ex \label{ex6:denneg}
	\denotes{neg} = $\lambda p\ \neg\ p$
\end{exe}
This delivers the denotations indicated in~(\ref{ex6:lowerhightyp2}).
\medskip
\begin{fullwidth}
\begin{exe}
\raggedright
	\ex \label{ex6:lowerhightyp2}
	\footnotesize
	\begin{forest}
	[TP
		[$\lambda Q\ \exists x$ \textsc{student}($x$) $\wedge$ $Q(x)$\\DP [a student, roof] ]
		[$\lambda 2_f\ \neg\ {[}2_{f}{]}$($\lambda 3\ \exists \epsilon$ \textsc{attend}($\epsilon$) $\wedge$ \textsc{patient}(the\_meeting)($\epsilon$) $\wedge$ \textsc{agent}($3$)($\epsilon$)) \\ \ibar{T}
			[\obar{T} [did] ]
			[$\neg\ {[}2_{f}{]}$($\lambda 3\ \exists \epsilon$ \textsc{attend}($\epsilon$) $\wedge$ \textsc{patient}(the\_meeting)($\epsilon$) $\wedge$ \textsc{agent}($3$)($\epsilon$)) \\NegP [\ibar{Neg}
				[$\lambda p\ \neg\ p$\\\obar{Neg} [neg] ]
				[${[}2_{f}{]}$($\lambda 3\ \exists \epsilon$ \textsc{attend}($\epsilon$) $\wedge$ \textsc{patient}(the\_meeting)($\epsilon$) $\wedge$ \textsc{agent}($3$)($\epsilon$)) \\AspP
					[\emph{t_{2}_{f}} ]
					[\ibar{Asp}
						[\obar{Asp} [asp] ]
						[vP
							[\emph{t_{3}_e} ]
							[\ibar{v}
								[\obar{v} [agent] ]
								[VP [\ibar{V}
									[\obar{V} [$\surd$attend] ]
									[DP [the meeting, roof] ]
								] ]
							 ]
						 ]
					 ]
				 ]
			] ]
		 ]
	]
	\end{forest}
\end{exe}	
\end{fullwidth}
\medskip
The denotation of \ibar{T} in (\ref{ex6:lowerhightyp2}) is a function that can take the denotation of \emph{a student} as its argument. Applying \denotes{\ibar{T}} to the denotation of \emph{a student} gives to TP the denotation in~(\ref{ex6:lowerhightyp3}).
\medskip
\begin{fullwidth}
\begin{exe}
\raggedright
	\ex \label{ex6:lowerhightyp3}
	\denotes{TP} = $\neg\ [\lambda Q\ \exists x$ \textsc{student}($x$) $\wedge$ $Q(x)$]($\lambda 3\ \exists \epsilon$ \textsc{attend}($\epsilon$) $\wedge$ \textsc{patient}(the\_meeting)($\epsilon$) $\wedge$ \textsc{agent}($3$)($\epsilon$))
\end{exe}	
\end{fullwidth}
\medskip
This formula can be simplified by applying the function in ``[{}]'' to the argument in ``({}).'' The result is~(\ref{ex6:lowerhightyp4}).
\medskip
\begin{fullwidth}
\begin{exe}
\raggedright
	\ex \label{ex6:lowerhightyp4}
	\denotes{TP} = $\neg\ \exists x$ \textsc{student}($x$) $\wedge$ $\exists \epsilon$ \textsc{attend}($\epsilon$) $\wedge$ \textsc{patient}(the\_meeting)($\epsilon$) $\wedge$ \textsc{agent}($x$)($\epsilon$)
\end{exe}
\end{fullwidth}
\medskip
This is a denotation in which \emph{a student} is part of what is negated; it is one of the meanings that we've seen this sentence can get. 

Allowing traces to be interpreted as variables over functions of the same semantic type as the moved phrase, then, provides a way of semantically interpreting (\ref{ex6:lowerhightyp2}), which corresponds (roughly) to the S-structure of (\ref{ex6:lowerill}), and deliver the negation wide scope interpretation. This alternative semantics, then, changes which parses in the derivation we need to say are semantically interpreted. To capture the ambiguity of (\ref{ex6:lowerill}), we needn't require that more than the S-structure parse be semantically interpreted. The ambiguity arises as a function of the flexibility of the semantic type that traces may have. When the derivation puts a trace in Specifier of AspP, and that trace is interpreted as a variable over functions of the kind that \emph{a student} is, then the denotation of \emph{a student} falls within the scope of negation. When, by contrast, the trace in Specifier of AspP is interpreted as a variable over entities, or alternatively the derivation doesn't put a trace in Specifier of AspP, the result will put the denotation of \emph{a student} outside the scope of negation.
%
%
\marginnote[-30pt]{The stepwise explication of how the semantics works if the trace in Specifier of AspP is interpreted as a varaible over entitites isn't provided here. I recommend it as an exercise for the student.}%
%

In our present syntax, only movement of DPs is capable of leaving a trace, and this corresponds to the decision to let traces be interpreted as variables over things. If the semantics of traces is broadened to allow them to be variables over other kinds of things, then the syntax of movement can be broadened to allow for traces in other contexts. Consider, for instance, the case of moving PPs that modify VPs. For concreteness, suppose that some PPs have denotations that allow them to combine with VPs in the way shown in~(\ref{ex6:ppmodbase}).
\medskip
\begin{fullwidth}
\begin{exe}
\raggedright
	\ex \label{ex6:ppmodbase}
	\small
	\begin{forest}
	[$\exists \epsilon$ \textsc{agent}(she)($\epsilon$) $\wedge$ \textsc{dance}($\epsilon$) $\wedge$ \textsc{on\_this}($\epsilon$) \\TP
		[DP [she, roof]]
		[$\lambda 2_e\ \exists \epsilon$ \textsc{agent}($2_e$)($\epsilon$) $\wedge$ \textsc{dance}($\epsilon$) $\wedge$ \textsc{on\_this}($\epsilon$) \\\ibar{T} 
			[\obar{T} [past] ]
			[$\exists \epsilon$ \textsc{agent}($2_e$)($\epsilon$) $\wedge$ \textsc{dance}($\epsilon$) $\wedge$ \textsc{on\_this}($\epsilon$) \\AspP [\ibar{Asp}
				[\obar{AsP} [asp] ]
				[$\lambda \epsilon$ \textsc{agent}($2_e$)($\epsilon$) $\wedge$ \textsc{dance}($\epsilon$) $\wedge$ \textsc{on\_this}($\epsilon$) \\vP
					[\emph{t_2_e} ]
					[$\lambda x\ \lambda \epsilon$ \textsc{agent}($x$)($\epsilon$) $\wedge$ \textsc{dance}($\epsilon$) $\wedge$ \textsc{on\_this}($\epsilon$) \\\ibar{v}
						[$\lambda P\ \lambda x\ \lambda \epsilon$ \textsc{agent}($x$)($\epsilon$) $\wedge$ $P$($\epsilon$)\\\obar{v} [agent] ]
						[$\lambda \epsilon$ \textsc{dance}($\epsilon$) $\wedge$ \textsc{on\_this}($\epsilon$)\\VP [$\lambda \epsilon$ \textsc{dance}($\epsilon$) $\wedge$ \textsc{on\_this}($\epsilon$)\\\ibar{V}
							[$\lambda \epsilon$ \textsc{dance}($\epsilon$)\\\ibar{V} [\obar{V} [$\surd$dance, roof]] ]
							[$\lambda P\ \lambda \epsilon\ P(\epsilon)$ $\wedge$ \textsc{on\_this}($\epsilon$)\\ PP [on this, roof] ]
						] ]
					 ]
				 ]
			] ]
		]
	]
	\end{forest}
\end{exe}	
\end{fullwidth}
\medskip
Consider now the S-structure that could be created from movement of these PPs in a syntax that allows traces to be created and interpreted more broadly.
\pagebreak
\medskip
\begin{fullwidth}
\begin{exe}
\raggedright
	\ex \label{ex6:ppmodmove}
	On this, she danced.
	\sn
	\small
	\begin{forest}
	[$\exists \epsilon$ \textsc{agent}(she)($\epsilon$) $\wedge$ {[}$\lambda P\ \lambda \epsilon\ P(\epsilon)$ $\wedge$ \textsc{on\_this}($\epsilon$){]}($\lambda \epsilon$ \textsc{dance}($\epsilon$)){]}($\epsilon$) \\TP
		[$\lambda P\ \lambda \epsilon'\ P(\epsilon')$ $\wedge$ \textsc{on\_this}($\epsilon'$) \\PP [on this, roof] ]
		[$\lambda 3_f\ \exists \epsilon$ \textsc{agent}(she)($\epsilon$) $\wedge$ {[}{[}$3_f${]}($\lambda \epsilon$ \textsc{dance}($\epsilon$)){]}($\epsilon$) \\TP
			[DP [she, roof] ]
			[$\lambda 2_f\ \exists \epsilon$ \textsc{agent}($2_e$)($\epsilon$) $\wedge$ {[}{[}$3_f${]}($\lambda \epsilon$ \textsc{dance}($\epsilon$)){]}($\epsilon$) \\\ibar{T}
				[\obar{T} [past] ]
				[$\exists \epsilon$ \textsc{agent}($2_e$)($\epsilon$) $\wedge$ {[}{[}$3_f${]}($\lambda \epsilon$ \textsc{dance}($\epsilon$)){]}($\epsilon$) \\AspP [\ibar{Asp}
					[\obar{Asp} [asp] ]
					[$\lambda \epsilon$ \textsc{agent}($2_e$)($\epsilon$) $\wedge$ {[}{[}$3_f${]}($\lambda \epsilon$ \textsc{dance}($\epsilon$)){]}($\epsilon$) \\vP
						[\emph{t_2_e} ]
						[$\lambda x\ \lambda \epsilon$ \textsc{agent}($x$)($\epsilon$) $\wedge$ {[}{[}$3_f${]}($\lambda \epsilon$ \textsc{dance}($\epsilon$)){]}($\epsilon$) \\\ibar{v}
							[\obar{v} [agent] ]
							[{[}$3_f${]}($\lambda \epsilon$ \textsc{dance}($\epsilon$)) \\VP [\ibar{V}
								[\ibar{V} [\obar{V} [$\surd$dance, roof]] ]
								[\emph{t_3_f} ]
							] ]
						 ]
					 ]
				] ]
			 ]
		 ]
	]
	\end{forest}
\end{exe}	
\end{fullwidth}
\vspace*{-60pt}
Let's unpack the denotation this gives to the root TP:
\begin{exe}
\raggedright
	\ex \label{ex6:compreduc1}
	$\exists \epsilon$ \textsc{agent}(she)($\epsilon$) $\wedge$ {[}$\lambda P\ \lambda \epsilon'\ P(\epsilon')$ $\wedge$ \textsc{on\_this}($\epsilon'$){]}($\lambda \epsilon$ \textsc{dance}($\epsilon$)){]}($\epsilon$)
\end{exe}
This contains the function ``$\lambda P\ \lambda \epsilon'\ P(\epsilon')$ $\wedge$ \textsc{on\_this}($\epsilon'$)'' which takes as its argument ``$\lambda \epsilon$ \textsc{dance}($\epsilon$).'' When these two combine, we get~(\ref{ex6:compreduc2}).
\begin{exe}
\raggedright
	\ex \label{ex6:compreduc2}
	$\exists \epsilon$ \textsc{agent}(she)($\epsilon$) $\wedge$ {[}$\lambda \epsilon'$ {[}$\lambda \epsilon$ \textsc{dance}($\epsilon$){]}($\epsilon'$) $\wedge$ \textsc{on\_this}($\epsilon'$){]}($\epsilon$)
\end{exe}
Applying ``$\lambda \epsilon$ \textsc{dance}($\epsilon$)'' to ``$\epsilon'$'' produces:
\begin{exe}
\raggedright
	\ex \label{ex6:compreduc3}
	$\exists \epsilon$ \textsc{agent}(she)($\epsilon$) $\wedge$ {[}$\lambda \epsilon'$ \textsc{dance}($\epsilon'$) $\wedge$ \textsc{on\_this}($\epsilon'$){]}($\epsilon$)
\end{exe}
This expression contains the function ``$\lambda \epsilon'$ \textsc{dance}($\epsilon'$) $\wedge$ \textsc{on\_this}($\epsilon'$)'' which is given the argument ``$\epsilon$.'' These combine to form:
\begin{exe}
\raggedright
	\ex \label{ex6:compreduc4}
	$\exists \epsilon$ \textsc{agent}(she)($\epsilon$) $\wedge$ \textsc{dance}($\epsilon$) $\wedge$ \textsc{on\_this}($\epsilon$)
\end{exe}
This is, arguably, the correct denotation. At least it puts the pieces together in such a way that (\ref{ex6:ppmodmove}) matches the meaning of its counterpart in~(\ref{ex6:ppmodbase}).

Note that this alternative semantics doesn't change the model we need for derivations; at least not for the cases illustrated by (\ref{ex6:lowerill}). Both semantics require derivations that include a parse that has \emph{a student} occupy Specifier of AspP. The old semantics requires a model of the syntax/semantics interface that lets that parse provide a denotation for the sentence. The new semantics allows for a model of the syntax/semantics interface that doesn't require that parse to be semantically interpreted. Instead, the S-structure parse of the derivation has sufficient information in it to provide both interpretations. Letting the trace in Specifier of AspP be interpreted as a variable over functions of the right kind, when combined with the definition of $\lambda$-expressions, provides a semantic computation that evaluates the denotation of the phrase in Specifier of TP in Specifier of AspP. It provides for a method of expressing reconstruction that relies on this feature of the rules of semantic composition. This method is sometimes called ``semantic reconstruction,'' and it is contrasted with ``syntactic reconstruction,'' which hypothesizes instead that the ambiguities involved are provided by different syntactic parses. Note that because the syntax and semantics employed by these two treatments are very similar, it is difficult to find evidence that distinguishes them. They are not so much alternatives as different executions of the same project. 

In section \ref{sec:ordering}, we saw that our present syntax struggles to deliver a derivation for certain sentences that have both movement of a verb and a quantificational DP in them. Embracing this alternative semantics can be combined with a slightly different syntax of head movement to allow for derivations that make these cases work smoothly.

The problem is illustrated by (\ref{ex6:lowerprob}), whose S-structure representation puts the direct object and verb in the positions indicated in~(\ref{ex6:lowerprobt}).
\begin{exe}
\raggedright
	\ex \label{ex6:lowerprob} She will give every toy to Coral.
	\ex \label{ex6:lowerprobt}
	\small
	\begin{forest}
	[TP
		[DP^2 [she, roof]]
		[\ibar{T} 
			[\obar{T} [will] ]
			[AspP [\ibar{Asp}
				[\obar{Asp} [asp] ]
				[vP
					[\emph{t_2} ]
					[\ibar{v} 
						[\obar{v} 
							[\obar{V} [$\surd$give, name=binderv] ]
							[\obar{v} [agent] ]
						 ]
						[VP [\ibar{V}
							[\ibar{V} 
								[{\ }, name=tracev]
							[DP, [every toy, roof] ]
						] 
						[PP [to Coral, roof]]
						]]
					 ]
				 ]
			]]
			]
		]
	\draw[->,>=stealth', ultra thin, gray] (tracev) to[out=south west,in=south,looseness=1] node[pos=0.75,xshift=-10mm]{} (binderv);
	\end{forest}
\end{exe}
This S-structure isn't semantically interpretable for two reasons: the quantificational DP must be moved outside AspP for it to be interpreted and neither $\surd$\emph{give} nor [_{\obar{v}} agent] are in positions where they can be semantically combined with its surrounding material. Moreover, if derivations obey Upwards, then there is no derivation that contains (\ref{ex6:lowerprobt}) which is semantically interpretable. Derivations that obey Upwards require the parse in~(\ref{ex6:lowerprobg}) to be part of the derivation, for that is the one our semantics provides the right denotation from.
\begin{exe}
\raggedright
	\ex \label{ex6:lowerprobg}
	\small
	\begin{forest}
	[TP
		[DP^2 [she, roof]]
		[\ibar{T} 
			[\obar{T} [will] ]
			[AspP [AspP [\ibar{Asp}
				[\obar{Asp} [asp] ]
				[vP
					[\emph{t_2} ]
					[\ibar{v} 
							[\obar{v} [agent] ]
						[VP [\ibar{V}
							[\ibar{V}
							[\obar{V} [$\surd$give, name=binderv] ]
								[\emph{t_3} ]
							]
						[PP [to Coral, roof]]
						]
						]]
					 ]
				 ]]
			[DP^3, name=binder [every toy, roof] ]
			]
			]
			]
	\end{forest}
\end{exe}
Our present rules do not make this a parse that can be achieved from the S-structure in (\ref{ex6:lowerprobt}), nor can it be part of the parses that lead to~(\ref{ex6:lowerprobt}).

Consider now the parses that make up a derivation for this sentence in a theory that lets traces be able to be interpreted as variables of the same semantic type as the moved term that they are related to, and a syntax that includes the following changes.
\begin{exe}
\raggedright
	\ex \label{ex6:newsyntax}
	\begin{xlist}
		\ex Moved heads can leave a trace.
		\ex Moved heads can adjoin to a phrase.
	\end{xlist}
\end{exe}
This syntax allows the S-structure of (\ref{ex6:lowerprob}) to have the representation in~(\ref{ex6:lowerprobt2}).
\begin{exe}
\raggedright
	\ex \label{ex6:lowerprobt2}
	\small
	\begin{forest}
	[TP
		[DP^2 [she, roof]]
		[\ibar{T}
			[\obar{T} [will] ]
				[AspP [\ibar{Asp} 
				[\obar{Asp} [asp]]
					[vP
					[\obar{V}^3 [$\surd$give]]
					[vP 
					[\emph{t_2} ]
					[\ibar{v}
						[\obar{v} [agent] ]
						[VP [\ibar{V}
							[\ibar{V}
								[\emph{t_3} ]
								[DP [every toy, roof] ]
							 ]
							 [PP [to Coral, roof]]
						] ]
					]
				 ]
			 ]
		 ]
	]
	]
	]
	\end{forest}
\end{exe}
The derivation that leads to this parse exploits the option of moving a head, here $\surd$\emph{give}, to a phrase and leaving a trace. In this case, \emph{$\surd$give} has moved to a position left-adjoined to vP, and left a trace in its D-structure position. Like (\ref{ex6:lowerprobt}), (\ref{ex6:lowerprobt2}) isn't semantically interpretable because \emph{every toy} isn't in a position where the rules of semantics allow it to be combined with its surrounding material. But unlike (\ref{ex6:lowerprobt}), it's possible to create a parse from (\ref{ex6:lowerprobt2}) that preserves Upwards and puts \emph{every toy} in a position that it can be semantically interpreted. That parse is~(\ref{ex6:lowerproblf}).
\medskip
\begin{fullwidth}
\begin{exe}
\raggedright
	\ex \label{ex6:lowerproblf}
	\small
	\begin{forest}
	[TP
		[DP [she, roof]]
		[\ibar{T} 
			[\obar{T} [will] ]
			[AspP 
				[AspP [\ibar{AsP}
					[\obar{Asp} [asp] ]
					[vP
						[$\lambda x\ \lambda y\ \lambda \epsilon'$ \textsc{theme}($x$)($\epsilon'$) $\wedge$ \textsc{goal}($x$)($\epsilon'$) \\ \obar{V} [$\surd$give] ]
						[$\lambda 3_f\ \lambda \epsilon$ \textsc{agent}($2_e$)($\epsilon$) $\wedge$ {[[[} 3_{f}{]}(4_e){]}(to\_Coral){]}($\epsilon$)\\vP
							[\emph{t_2_e} ]
							[$\lambda x\ \lambda \epsilon$ \textsc{agent}($x$)($\epsilon$) $\wedge$ {[[[} 3_{f}{]}(4_e){]}(to\_Coral){]}($\epsilon$)\\\ibar{v}
								[\obar{v} [agent] ]
								[VP [{[}{[} 3_{f}{]}(4_e){]}(to\_Coral)\\\ibar{V}
									[{[} 3_{f}{]}(4_e)\\\ibar{V}
										[\emph{t_3_f} ]
										[\emph{t_4_e} ]
									 ]
									[PP [to Coral, roof] ]
								] ]
							 ]
						 ]
					 ]
				] ]
				[DP [every toy, roof] ]
			 ]
		]
	]
	\end{forest}
\end{exe}
\end{fullwidth}
\medskip
Note that the denotation shown for the lowest vP in (\ref{ex6:lowerproblf}) includes the change that Movement has wrought; Movement causes the sister to a moved term to turn into a $\lambda $-expression that binds the variable created by the trace. In (\ref{ex6:lowerproblf}), that trace is left by movement of $\surd$\emph{give}. The function that is created by this is one whose domain includes functions of the kind that \denotes{$\surd$\emph{give}} is. Therefore, the way that this vP and $\surd$\emph{give} combine semantically is by taking the denotation of $\surd$\emph{give} to be the argument of the function that vP is. The result is~(\ref{ex6:lowerproblfi}).
\medskip
\begin{fullwidth}
\begin{exe}
\raggedright
	\ex \label{ex6:lowerproblfi}
	\denotes{vP} = $\lambda \epsilon$ \textsc{agent}($2_e$)($\epsilon$) $\wedge$ {[[[}$\lambda x\ \lambda y\ \lambda \epsilon'$ \textsc{theme}($x$)($\epsilon'$) $\wedge$ \textsc{goal}($y$)($\epsilon'$){]}($4_e$)](to\_Coral){]}($\epsilon$)
\end{exe}
\end{fullwidth}
\medskip
By applying the function ``$\lambda x\ \lambda y\ \lambda \epsilon$ \textsc{theme}($x$)($\epsilon$) $\wedge$ \textsc{goal}($y$)($\epsilon$)'' recursively to its arguments, we can reduce this formula to:
\medskip
\begin{fullwidth}
\begin{exe}
\raggedright
	\ex \label{ex6:lowerproblfii}
	\denotes{vP} = $\lambda \epsilon$ \textsc{agent}($2_e$)($\epsilon$) $\wedge$ \textsc{theme}($4_e$)($\epsilon$) $\wedge$ \textsc{goal}(to\_Coral)($\epsilon$)
\end{exe}	
\end{fullwidth}
\medskip
Plugging this into the syntactic representation gives~(\ref{ex6:lowerproblfiii}).
\medskip
\begin{fullwidth}
\begin{exe}
\raggedright
	\ex \label{ex6:lowerproblfiii}
	\small
	\begin{forest}
	[$\forall y$ \textsc{toy}($y$) $\rightarrow$ $\exists \epsilon$ \textsc{agent}(she) $\wedge$ \textsc{theme}($y$)($\epsilon$) $\wedge$ \textsc{goal}(to\_Coral)($\epsilon$)\\TP
		[DP [she, roof]]
		[$\lambda 2_e$ $\forall y$ \textsc{toy}($y$) $\rightarrow$ $\exists \epsilon$ \textsc{agent}($2_e$)($\epsilon$) $\wedge$ \textsc{theme}($y$)($\epsilon$) $\wedge$ \textsc{goal}(to\_Coral)($\epsilon$) \\\ibar{T} 
			[\obar{T} [will] ]
			[$\forall y$ \textsc{toy}($y$) $\rightarrow$ $\exists \epsilon$ \textsc{agent}($2_e$)($\epsilon$) $\wedge$ \textsc{theme}($y$)($\epsilon$) $\wedge$ \textsc{goal}(to\_Coral)($\epsilon$)\\AspP 
				[$\lambda 4_e\ \exists \epsilon$ \textsc{agent}($2_e$)($\epsilon$) $\wedge$ \textsc{theme}($4_e$)($\epsilon$) $\wedge$ \textsc{goal}(to\_Coral)($\epsilon$)\\AspP [$\exists \epsilon$ \textsc{agent}($2_e$)($\epsilon$) $\wedge$ \textsc{theme}($4_e$)($\epsilon$) $\wedge$ \textsc{goal}(to\_Coral)($\epsilon$)\\\ibar{AsP}
					[$\lambda P\ \exists \epsilon\ P(\epsilon)$\\\obar{Asp} [asp] ]
					[$\lambda \epsilon$ \textsc{agent}($2_e$)($\epsilon$) $\wedge$ \textsc{theme}($4_e$)($\epsilon$) $\wedge$ \textsc{goal}(to\_Coral)($\epsilon$)\\vP
						[\obar{V} [$\surd$give] ]
						[vP
							[\emph{t_2_e} ]
							[\ibar{v}
								[\obar{v} [agent] ]
								[VP [\ibar{V}
										[\emph{t_3_f} ]
										[\emph{t_4_e} ]
									 ]
									[PP [to Coral, roof] ]
								] ]
							 ]
						 ]
					 ]
				]
				[$\lambda Q\ \forall y$ \textsc{toy}($y$) $\rightarrow$ $Q(y)$ \\DP [every toy, roof] ]
			 ]]
		]
	]
	\end{forest}
\end{exe}	
\end{fullwidth}
\medskip
In calculating the denotation of \ibar{T}, I've ignored the contribution that \emph{will} makes, and this absence is therefore present in the denotation given for the resulting TP. But with that caveat, the syntax and semantics of this parse delivers the right result.

Adopting this alternative semantics, then, can be wedded to a different syntax of Head Movement that allows for derivations that obey Upwards and overcomes the problem posed by examples like (\ref{ex6:lowerprob}). But the syntax it requires of Head Movement is problematic. For a moved verb to bind its trace, it must move to a position that c-commands that trace. In the cases we've just worked through, this involves moving \obar{V} so that it left-adjoins to a projection of \obar{v}. But the same kind of syntax would be required of moved verbs whenever there is a quantificational DP in the sentence. So for instance, it would require that an auxiliary verb in English move to some projection of \obar{T} in a finite clause, perhaps in the way sketched in~(\ref{ex6:auxvm}).
\begin{exe}
\raggedright
	\ex \label{ex6:auxvm}
	\small
	\begin{forest}
	[TP
		[DP^2 [you, roof]]
		[\ibar{T} 
			[\obar{V}^3 [$\surd$have, name=binder] ]
			[\ibar{T}
				[\obar{T} [pres] ]
				[VP [\ibar{V}
					[\emph{t}_3, name=trace ]
					[vP
						[\emph{t_2} ]
						[\ibar{v}
							[\obar{v} [agent] ]
							[VP [read everything, roof] ]
						 ]
					 ]
				] ]
			 ]
		]
	]
	\draw[->,>=stealth', thin, gray] (trace) to[out=south west,in=south,looseness=1] node[pos=0.75,xshift=-10mm]{} (binder);
	\end{forest}
\end{exe}
We've also learned that in question contexts, the contents of \obar{T} move to \obar{C}. In a syntax that allows Heads to move to c-commanding positions, this might produce representations like~(\ref{ex6:tmp}).
\begin{exe}
\raggedright
	\ex \label{ex6:tmp}
	\small
	\begin{forest}
	[CP
		[\ibar{C}
		[\obar{T}^3 [will, name=binder]]
		[\ibar{C}
			[\obar{C} [Q] ]
			[TP
				[DP^2 [you, roof]]
				[\ibar{T}
					[\emph{t_3}, name=trace ]
					[VP [have read everything, roof] ]
				 ]
			 ]
		]
		]
	]
	\draw[->,>=stealth', thin, gray] (trace) to[out=south west,in=south,looseness=1] node[pos=0.75,xshift=-10mm]{} (binder);
	\end{forest}
\end{exe}
We've seen that when these two operations are both present, the result should put the auxiliary verb in the position to which \obar{T} moves.%
%
\marginnote[0pt]{See chapter \ref{ch:features} for details.}%
%
\begin{exe}
\raggedright
	\ex \label{ex6:combmp}
	Have you read everything?
\end{exe}
The requires a syntax that causes the moved auxiliary verb to combine with the tense features in \obar{T} so that they are one \obar{X}. That is, they must have the same status that the \obar{T} \emph{will} has in (\ref{ex6:tmp}). That is not what the syntax in (\ref{ex6:auxvm}) delivers. A syntax that requires a moved verb to c-command the position it moves from is in conflict with the need to form a new head from the moved verb and the \obar{X} whose needs are being satisfied by the movement. If the syntax/semantics sketched in this section is to solve the problem posed by head movement that is discussed in section \ref{sec:ordering}, then we must find a way of resolving this incompatibility.%
%
\marginnote[0pt]{See \citet{Matushansky:2002} for a suggestion.}%
%

As it stands, the Trace theory sketched in this section also doesn't provide a straightforward solution to the cases of reconstructed variables, discussed in section~\ref{sec:reconvariable}. These cases arise when a moved phrase contains a term that is interpreted as a variable bound by a lower phrase. (\ref{ex6:reconvar1}) illustrates.
\begin{exe}
\raggedright
	\ex \label{ex6:reconvar1}
	Which discussion about her_2 will no woman^2 remember?
\end{exe}
A sketch of the relevant parse for (\ref{ex6:reconvar1}) is in~(\ref{ex6:reconvar1t}).
\begin{exe}
\raggedright
	\ex \label{ex6:reconvar1t}
	\small
	\begin{forest}
	[CP
		[DP^3 [which discussion of her_2, roof]]
		[\ibar{C} 
			[\obar{C} [Q]]
			[TP
				[DP^2 [no woman, roof]]
				[\ibar{T} 
					[\obar{T} [will] ]
					[vP
						[\emph{t_{2e}} ]
						[\ibar{v}
							[\obar{v} [agent] ]
							[VP [\ibar{V}
								[\obar{V} [$\surd$remember] ]
									[\emph{t_{3f}} ]
								] 
							] ]
						 ]
					 ]
				]
			 ]
		]
	\end{forest}
\end{exe}
We have not determined a semantics for wh-phrases like \emph{which discussion of her}. Let us suppose, for the purposes of illustration, that it refers to an entity---it is of semantic type~e. If so, even if its trace matches it in semantic type (as indicated in (\ref{ex6:reconvar1t})), it will be of type e, and this gives us the semantics in~(\ref{ex6:reconvar1s}).
\medskip
\begin{fullwidth}
\begin{exe}
\raggedright
	\ex \label{ex6:reconvar1s}
	\footnotesize
	\begin{forest}
		[$\neg \exists y$ \textsc{woman}($y$) $\wedge$ $\exists \epsilon$ \textsc{partake}($\epsilon$) $\wedge$\\ \textsc{theme}($\epsilon$)(which\_discussion\_about\_her_2) $\wedge$ \textsc{agent}($\epsilon$)($y$)\\CP
			[DP^3 [which discussion of her_2, roof]]
			[$\lambda 3 \neg \exists y$ \textsc{woman}($y$) $\wedge$ $\exists \epsilon$ \textsc{remember}($\epsilon$) $\wedge$\\ \textsc{theme}($\epsilon$)($3$) $\wedge$ \textsc{agent}($\epsilon$)($y$)\\\ibar{C} 
				[\obar{C} [Q]]
				[$\neg \exists y$ \textsc{woman}($y$) $\wedge$ $\exists \epsilon$ \textsc{remember}($\epsilon$) $\wedge$ \textsc{theme}($\epsilon$)($3$) $\wedge$ \textsc{agent}($\epsilon$)($y$)\\TP
					[DP^2 [no woman, roof]]
					[$\lambda 2\ \exists \epsilon$ \textsc{remember}($\epsilon$) $\wedge$ \textsc{theme}($\epsilon$)($3$) $\wedge$ \textsc{agent}($\epsilon$)($2$)\\\ibar{T} 
						[\obar{T} [will] ]
						[$\lambda \epsilon$ \textsc{remember}($\epsilon$) $\wedge$ \textsc{theme}($\epsilon$)($3$) $\wedge$ \textsc{agent}($\epsilon$)($2$)\\vP
							[\emph{t_{2e}} ]
							[\ibar{v}
								[\obar{v} [agent] ]
								[VP [\ibar{V}
									[\obar{V} [$\surd$remember] ]
										[\emph{t_{3f}} ]
								] ]
							 ]
						 ]
					]
				 ]
			]
		]
		\end{forest}
\end{exe}
\end{fullwidth}
\medskip
Note that the index on \emph{her} is not part of the constituent that is a sister of ``$\lambda 2$.'' For this reason, it is not bound by $\lambda 2$, and so we do not get a well-formed result.%
%
\marginnote[0pt]{Moreover, the index on \emph{her} is prevented from getting bound by ``$\lambda 2$'' because of the law that prevents variables inside the domain of a $\lambda$-operator from being bound by that $\lambda$-operator. See page~\pageref{lambdalaw} for discussion.}%
%

The way of bringing these cases in line with the Trace theory of this section is to give to wh-phrases a more complex semantics. Suppose, for instance, that wh-phrases could be functions from entities to entities. The entities of their domain are the values that a variable they contain has, and the entities of their range are the entities that they refer to. For instance, imagine that \emph{which discussion of her_2} could have the denotation in~(\ref{ex6:funcwh1}).
\begin{exe}
\raggedright
	\ex \label{ex6:funcwh1}
	\denotes{which discussion of her_1} = $\lambda 1$ $a$ if \textsc{discussion\_of\_$1$}($a$)
\end{exe}
The function defined in (\ref{ex6:funcwh1}) applies to entities, $x$, and returns an entity, $a$, just in case it is a discussion of~$x$. Further, suppose that traces can be more than simple variables of various semantic types. Imagine that they can themselves take arguments when they are the type of a function. Suppose, that is, that \emph{t_f(n)}, where \emph{n} is a referential index, can be put into the position from which a phrase moves. This allows (\ref{ex6:reconvar1}) to have the representation in~(\ref{ex6:reconvar1tt}).
\medskip
\begin{fullwidth}
\begin{exe}
\raggedright
	\ex \label{ex6:reconvar1tt}
	\footnotesize
	\begin{forest}
		[$\neg \exists y$ \textsc{woman}($y$) $\wedge$ $\exists \epsilon$ \textsc{partake}($\epsilon$) $\wedge$\\ \textsc{theme}($\epsilon$)($a$ if {[discussion\_about\_$y$]}($z$)) $\wedge$ \textsc{agent}($\epsilon$)($y$)\\CP
			[DP^3 [which discussion of her_1, roof]]
			[$\lambda 3 \neg \exists y$ \textsc{woman}($y$) $\wedge$ $\exists \epsilon$ \textsc{remember}($\epsilon$) $\wedge$\\ \textsc{theme}($\epsilon$)($3$($y$)) $\wedge$ \textsc{agent}($\epsilon$)($y$)\\\ibar{C} 
				[\obar{C} [Q]]
				[$\neg \exists y$ \textsc{woman}($y$) $\wedge$ $\exists \epsilon$ \textsc{remember}($\epsilon$) $\wedge$ \textsc{theme}($\epsilon$)($3$($y$)) $\wedge$ \textsc{agent}($\epsilon$)($y$)\\TP
					[DP^2 [no woman, roof]]
					[$\lambda 2\ \exists \epsilon$ \textsc{remember}($\epsilon$) $\wedge$ \textsc{theme}($\epsilon$)($3$($2$)) $\wedge$ \textsc{agent}($\epsilon$)($2$)\\\ibar{T} 
						[\obar{T} [will] ]
						[$\lambda \epsilon$ \textsc{remember}($\epsilon$) $\wedge$ \textsc{theme}($\epsilon$)($3$($2$)) $\wedge$ \textsc{agent}($\epsilon$)($2$)\\vP
							[\emph{t_{2e}} ]
							[\ibar{v}
								[\obar{v} [agent] ]
								[VP [\ibar{V}
									[\obar{V} [$\surd$remember] ]
										[\emph{t_{3f}}($2$) ]
								] ]
							 ]
						 ]
					]
				 ]
			]
		]
		\end{forest}
\end{exe}	
\end{fullwidth}

\vspace*{-30pt}
This delivers the desired meaning; let's see how. The trace of the wh-phrase in (\ref{ex6:reconvar1tt}) (=\emph{t_{3f}}) comes with a variable that is bound by \emph{no woman} (=$2$). The wh-phrase which binds the trace that combines with this variable will semantically combine with the variable, if the trace has the same semantic type as the wh-phrase. Schematically, then, (\ref{ex6:analf}a) yields~(\ref{ex6:analf}b).
\begin{exe}
\raggedright
	\ex \label{ex6:analf}
	\begin{xlist}
		\ex which discussion of her_2 $\lambda 3$\ldots no woman $\lambda 2$\ldots remember $3$($2$)
		\ex \ldots no woman $\lambda 2$\ldots remember [$\lambda 1$ $z$ if \textsc{discussion\_of}\_$1$($z$)]($2$)
	\end{xlist}
\end{exe}
If we reduce the the ``[$\lambda 1$ $z$ if \textsc{discussion\_of}\_$1$]($2$)'' expression by applying the function to the argument, we get from (\ref{ex6:analf}b) what's schematized in~(\ref{ex6:analyff}).
\begin{exe}
\raggedright
	\ex \label{ex6:analyff}
	\ldots no woman $\lambda 2$\ldots remember $z$, if \textsc{discussion
	_of}\_$2$
\end{exe}
In this way, then, we can model cases of reconstructed variables in Trace theory by letting traces contain variables that can be bound by other terms in the sentence.

While this method technically brings reconstructed variables within range of the Trace theory of movement, it doesn't do so in a way that can capture an important fact about these cases. That fact is that the morphological form of the reconstructed variable determines the conditions under which it can be bound under reconstruction. First, let me illustrate that the morphological form of a variable can determine the conditions under which it is bound. This is illustrated by the contrast of interpretations that are available to the sentences in~(\ref{ex6:localpron}).
\begin{exe}
\raggedright
	\ex \label{ex6:localpron}
	\begin{xlist}
		\ex No woman will remember someone's discussion of her.
		\ex No woman will remember someone's discussion of herself.
	\end{xlist}
\end{exe}
Both \emph{her} and the reflexive pronoun \emph{herself} can be interpreted as bound variables. In (\ref{ex6:localpron}), for instance, \emph{herself} can be interpreted as a variable bound by \emph{someone}. But only in (\ref{ex6:localpron}a) can \emph{no woman} bind the pronoun; that is, \emph{her} but not \emph{herself} can be interpreted as a variable bound by \emph{no woman}. This is because reflexive pronouns, unlike non-reflexive pronouns, can only be bound by DPs that are sufficiently close. One way of describing that condition is with the requirement that a reflexive pronoun be bound by the closest DP that can bind it. Because \emph{someone} can bind \emph{herself} in (\ref{ex6:localpron}b), this requirement would force \emph{herself} to be understood as a variable bound by \emph{someone} rather than \emph{no woman}. It would also require \emph{herself} to be interpreted as a variable bound by \emph{someone} rather than finding a contextually salient person to refer to, as \emph{her} can. That is also a difference between (\ref{ex6:localpron}a) and~(\ref{ex6:localpron}b). And, finally, it would explain why both \emph{her} and \emph{herself} can be bound by \emph{no woman} in~(\ref{ex6:ssc2}).
\begin{exe}
\raggedright
	\ex \label{ex6:ssc2}
	\begin{xlist}
		\ex No woman will remember this discussion of her.
		\ex No woman will remember this discussion of herself.
	\end{xlist}
\end{exe}
Because there is no genitive DP \emph{someone's} in (\ref{ex6:ssc2}b), the closest possible binder for \emph{herself} is \emph{no woman} and this is what \emph{herself} can, and must, be bound by.

Consider now the contrast in~(\ref{ex6:reconanap}).
\begin{exe}
\raggedright
	\ex \label{ex6:reconanap}
	\begin{xlist}
		\ex[]{Whose discussion of her_2 will no woman^2 remember?}
		\ex[*]{Whose discussion of herself_2 will no woman^2 remember?}
	\end{xlist}
\end{exe}
Parallel to (\ref{ex6:localpron}), only (\ref{ex6:reconanap}a) allows the pronoun to be bound by \emph{no woman}. In (\ref{ex6:reconanap}b), \emph{herself} must be bound to \emph{whose}, as this is the closest possible binder. Just as in (\ref{ex6:ssc2}), when \emph{whose} is absent, \emph{herself} can be bound by \emph{no woman}:
\begin{exe}
\raggedright
	\ex[]{\label{ex6:ssc3} Which discussion of herself_2 will no woman^2 remember?}
\end{exe}%
%
\marginnote[0pt]{Paradigms like these are examined, perhaps for the first time, by \citet{Barss:1986b}.}%
%
These contrast aren't explicable under Trace theory, even when embellished in the way just described. What's needed is a theory that allows binding of reconstructed variables by way of a syntactic representation that puts the phrase containing the variable in the position it needs to be bound. In a Trace theory that allows a trace to contain a variable which, indirectly, does the job of binding the pronoun in the moved expression---as, for example, sketched above---there is no way of capturing the fact that the constraints on binding that come with the pronoun matter. We need a way of putting the pronoun itself into the relevant position for binding in order to capture this fact.

What we've seen in these sections, then, is that Reconstruction phenomena pose a kind of ordering problem for Trace theory. There is no way of capturing these effects with Trace theory that is consistent with Upwards. If derivations obey Upwards, they cannot provide the parses needed to capture certain cases of Reconstruction and still deliver the S-structures that represent those derivations. As we've seen in this section, it is possible to preserve Upwards using a version of Trace theory that allows the trace to have the same denotation as the binder. This allows more of the parses in a derivation to be semantically interpretable and one of the consequences of this is a treatment of reconstruction that fits Upwards constrained derivations. But broadening the interpretation of traces in this way has its own problems. It doesn't fit well the syntax we've given Head Movement, which requires a head to move to a position from which it cannot bind a trace. And it doesn't capture the cases of Reconstruction of pronominal variables in a way that enforces the constraints imposed by the pronouns that are reconstructed. 

What's needed to shore up these problems is a way of thinking of traces as having the same meaning as the term that moved, whether or not they are variables bound by that moved term. That will allow for Upwards obeying derivations in the way that the Trace theory of this section does, but fit the syntax of Head Movement better. And we need to understand trace to be able to have the same syntactic content as the moved phrase. That will capture the cases of Reconstructed pronominal variables in a way that preserves the conditions on binding that come with the choice of pronoun. It's providing these two attributes that the Copy Theory of movement is designed for.

\subsection{Copies}
\label{sec:copies}

Instead of ``traces,'' this theory of movement creates ``copies,'' and it is named after this feature of the theory.\marginnote{That movement involves putting copies of the moved term in different positions is adopted in \citet{Ross:1967}, and given perhaps its first defense in \citet{Ross:1969b}. But it is revived to handle these semantic problems in \citet{Chomskey:1993c}, and developed in \citet{Chomsky:1995}. It's Chomsky's version of this theory that will be discussed here.}

There are various ways of formulating the Copy theory of movement, but I'll use what strikes me as the most common of those ways. In this formulation, movement is factored into three parts. The first is an operation that creates the ``copy.''
\pagebreak
\begin{exe}
\raggedright
	\ex \label{ex6:copydef} \textsc{Copy}\\
	For $\alpha$ a constituent in a phrase marker, \textsc{copy}($\alpha$)=$\alpha'$, where $\alpha$ and $\alpha'$ are syntactically, semantically, and morphologically exactly the same.
\end{exe}
The \textsc{copy} operation is used by the move operation.
\begin{exe}
\raggedright
	\ex \label{ex6:movedeff} Move\\
	$\alpha$ moves to $\beta$ means that \textsc{copy}($\alpha$) is put into the position~$\beta$.
\end{exe}
On this definition, Head Movement creates representations like~(\ref{ex6:headcmoveill}).
\begin{exe}
\raggedright
	\ex \label{ex6:headcmoveill}
	\small
	\begin{forest}
	[CP [\ibar{C}
		[\obar{C} 
			[\obar{T$'$} [will, name=binder] ]
			[\obar{C} [Q] ]
		]
		[TP 
			[DP [this, roof] ]
			[\ibar{T}
				[\obar{T} [will, name=trace] ]
				[vP [make sense, roof] ]
			 ]
		]
	]]
	\draw[->,>=stealth', thin, gray] (trace) to[out=south west,in=south,looseness=1] node[pos=0.75,xshift=-10mm]{} (binder);
	\end{forest}
\end{exe}
Argument Movement creates representations like~(\ref{ex6:amovecmoveill}).
\begin{exe}
\raggedright
	\ex \label{ex6:amovecmoveill}
	\small
	\begin{forest}
	[TP
		[DP$'$ [this, name=binder, roof]]
		[\ibar{T} 
			[\obar{T} [will] ]
			[vP
				[DP [this, name=trace, roof] ]
				[\ibar{v} [make sense, roof] ]
			 ]
		]
	]
	\draw[->,>=stealth', thin, gray] (trace) to[out=south west,in=south,looseness=1] node[pos=0.75,xshift=-10mm]{} (binder);
	\end{forest}
\end{exe}
Wh Movement creates representations like~(\ref{ex6:whmovecmoveill}).
\begin{exe}
\raggedright
	\ex \label{ex6:whmovecmoveill}
	\small
	\begin{forest}
	[CP
		[DP$'$ [which problem, name=binder, roof]]
		[\ibar{C} 
			[\obar{C} [Q] ]
			[TP
				[DP [this idea, roof]]
				[\ibar{T}
					[\obar{T} [will] ]
					[vP
						[DP [this idea, roof] ]
						[\ibar{v}
							[\obar{v} [cause] ]
							[VP [\ibar{V}
								[\obar{V} [$\surd$solve] ]
								[DP [which problem, name=trace, roof] ]
							] ]
						 ]
					 ]
				 ]
			 ]
		]
	]
	\draw[->,>=stealth', thin, gray] (trace) to[out=south west,in=south,looseness=1] node[pos=0.75,xshift=-10mm]{} (binder);
	\end{forest}
\end{exe}
And, finally, Complex NP Shift creates representations like~(\ref{ex6:complexcmoveill}).
\begin{exe}
\raggedright
	\ex \label{ex6:complexcmoveill}
	\small
	\begin{forest}
	[TP
		[TP
			[DP$'$ [this, roof] ]
			[\ibar{T}
				[\obar{T} [might] ]
				[vP
					[DP [this, roof]  ]
					[\ibar{v}
						[\obar{v} [cause] ]
						[VP [\ibar{V}
							[\obar{V} [$\surd$solve] ]
							[DP [the problems we love, name=trace, roof] ]
						] ]
					 ]
				 ]
			 ]
		 ]
		[DP$'$ [the problems we love, name=binder, roof] ]
	]
	\draw[->,>=stealth', thin, gray] (trace) to[out=south east,in=south east,looseness=1.5] node[pos=0.75,xshift=10mm]{} (binder);
	\end{forest}
\end{exe}

When Movement happens to a term more than once, it will invoke additional instances of \textsc{copy}. Thus, for instance, Head Movement can produce derivations like~(\ref{ex6:scycheadc}).
\medskip
\begin{fullwidth}
\begin{exe}
\raggedright
	\ex \label{ex6:scycheadc}
	\small
	\begin{forest}
	[CP [\ibar{C}
		[\obar{C} ]
		[TP
			[DP [this, roof]]
			[\ibar{T}
				[\obar{T} 
					[\obar{V$'$} [$\surd$be, name=binder] ]
					[\obar{T} [past] ]
				 ]
				[VP [\ibar{V}
					[\obar{V} [$\surd$be, name=trace] ]
					[AP [unclear, roof]]
				] ]
			 ]
		 ]
	]]
	\draw[->,>=stealth', thin, gray] (trace) to[out=south west,in=south,looseness=1] node[pos=0.75,xshift=-10mm]{} (binder);
	\end{forest}
	\begin{forest}
		[CP [\ibar{C}
		[\obar{C} 
		[\obar{T$'$}, name=binder 
			[\obar{V$'$} [$\surd$be] ]
			[\obar{T} [past] ]
		]
		[\obar{C} [Q]]
		]
		[TP
			[DP [this, roof]]
			[\ibar{T}
				[\obar{T}, name=trace 
					[\obar{V} [$\surd$be] ]
					[\obar{T} [past] ]
				 ]
				[VP [\ibar{V}
					[\obar{V} [$\surd$be] ]
					[AP [unclear, roof]]
				] ]
			 ]
		 ]
	]]
	\draw[->,>=stealth', thin, gray] (trace) to[out=south west,in=south,looseness=1] node[pos=0.75,xshift=-10mm]{} (binder);
	\end{forest}
\end{exe}
\end{fullwidth}
\medskip
And Argument Movement can produce derivations like~(\ref{ex6:amovescc}).
\medskip
\begin{fullwidth}
\begin{exe}
\raggedright
	\ex \label{ex6:amovescc}
	\small
	\begin{forest}
	[TP
		[\ibar{T}
		[\obar{T} ]
		[VP [\ibar{V}
			[\obar{V} [$\surd$seem] ]
			[TP 
				[DP$'$ [this, name=binder, roof]]
				[\ibar{T}
				[\obar{T} [to] ]
				[vP
					[DP [this, name=trace, roof] ]
					[\ibar{v} [be complicated, roof] ]
				 ]
			] ]
		] ]
	]]
	\draw[->,>=stealth', thin, gray] (trace) to[out=south west,in=south,looseness=1] node[pos=0.75,xshift=-10mm]{} (binder);
	\end{forest}
	\begin{forest}
		[TP
			[DP$''$ [this, name=binder, roof]]
			[\ibar{T}
			[\obar{T} ]
			[VP [\ibar{V}
				[\obar{V} [$\surd$seem] ]
				[TP 
					[DP$'$ [this, name=trace, roof]]
					[\ibar{T}
					[\obar{T} [to] ]
					[vP
						[DP [this, roof] ]
						[\ibar{v} [be complicated, roof] ]
					 ]
				] ]
			] ]
		]]
		\draw[->,>=stealth', thin, gray] (trace) to[out=south west,in=south,looseness=1.2] node[pos=0.75,xshift=-10mm]{} (binder);
		\end{forest}
\end{exe}
\end{fullwidth}
\medskip

In what follows, I will not distinguish the copy from the thing it is copied from in parses of movement. (That is, I will represent ``XP$'$'' with ``XP''.) And I will adopt a correspondingly sloppy habit of speech: I will call all the phrases that are related by \textsc{copy} ``copies,'' even though technically one is a copy of the other. 

Obviously, this rule incorrectly results in two pronunciations of the moved term. So, it must be wedded to a rule that prevents pronunciation of more than one copy. Let's know this rule as Suppress Copy.
\pagebreak
\begin{exe}
\raggedright
	\ex \label{ex6:supcopy} Suppress Copy\\
	Exactly one copy can be in the linearization.
\end{exe}
In all of the examples above, Suppress Copy is satisfied by suppressing all but the highest copy. It would be reasonable to strengthen Suppress Copy so that it captured this generalization; it could be expressed as allowing only the highest copy to be linearized. When we discuss QR, we will see reasons for adopting this weaker version. But in the end, a means for determining which copy is linearized will be necessary, as in the normal case, it doesn't appear to be free. In none of the examples above, for instance, can anything other than the highest copy be linearized. I will adopt the convention of putting the copies that are not pronounced into a \shade{shaded} font. (\ref{ex6:headcmoveill}), for instance, shows the pronunciation in~(\ref{ex6:headcmoveillb}).
\begin{exe}
\raggedright
	\ex \label{ex6:headcmoveillb}
	\small
	\begin{forest}
	[CP [\ibar{C}
		[\obar{C} 
			[\obar{T$'$} [will, name=binder] ]
			[\obar{C} [Q] ]
		]
		[TP 
			[DP [this, roof] ]
			[\ibar{T}
				[\obar{T} [\shade{will}, name=trace] ]
				[vP [make sense, roof] ]
			 ]
		]
	]]
	\draw[->,>=stealth', thin, gray] (trace) to[out=south west,in=south,looseness=1] node[pos=0.75,xshift=-10mm]{} (binder);
	\end{forest}
\end{exe}

\subsubsection{Copies and Reconstruction}
\label{sec:copyrecon}

The Copy Theory of movement has features like that of Trace theory when it is extended in the way described in the previous section. That version of Trace theory allows traces to be variables over the kind of things that the moved term is. If a verb moves, for instance, the trace left by that verb can be a variable over verb meanings. And if a quantificational DP moves, the trace it leaves can be a variable over quantificational DP meanings. The Copy Theory of movement creates representations which have similar information in them. If a verb moves, for instance, the lower position of the moved verb remains occupied by that verb. And if a DP moves, again the lower position remains occupied by that DP. It's this feature of the Copy Theory that allows it to model reconstruction.

We'll need a new semantic interface procedure for the Copy Theory of movement to do so, however. To begin with, let us adopt the view that no more than one copy is required to be semantically interpreted. Like Suppress Copy, this lets copies be in principle invisible.
\begin{exe}
\raggedright
	\ex \label{ex6:semcopyint} No more than one copy needs to be semantically interpreted. 
\end{exe}
In a situation where verbs move---as in French finite clauses or English vPs---Head Movement creates representations that can be semantically interpreted. In those representations, the lower copy of the verb will be semantically interpreted, but the higher copy won't be. This is illustrated by the English case of verb movement in~(\ref{ex6:reconcopyv1}).
\medskip
\begin{fullwidth}
\begin{exe}
\raggedright
	\ex \label{ex6:reconcopyv1}
	\begin{forest}
	[$\lambda \epsilon'$ \textsc{eat}($\epsilon'$) $\wedge$ \textsc{patient}(them)($\epsilon'$) $\wedge$ \textsc{agent}(she)($\epsilon'$) \\vP
		[DP [she, roof]]
		[$\lambda y \lambda \epsilon'$ \textsc{eat}($\epsilon'$) $\wedge$ \textsc{patient}(them)($\epsilon'$) $\wedge$ \textsc{agent}($y$)($\epsilon'$) \\\ibar{v} 
			[$\lambda P \lambda y \lambda \epsilon'$ $P$($\epsilon'$) $\wedge$ \textsc{agent}($y$)($\epsilon'$) \\\obar{v}
				[\obar{V} [$\surd$eat, name=binder] ]
				[$\lambda P \lambda y \lambda \epsilon'$ $P$($\epsilon'$) $\wedge$ \textsc{agent}($y$)($\epsilon'$) \\\obar{v} [agent] ]
			 ]
			[$\lambda \epsilon$ \textsc{eat}($\epsilon$) $\wedge$ \textsc{patient}(them)($\epsilon$)  \\VP [$\lambda \epsilon$ \textsc{eat}($\epsilon$) $\wedge$ \textsc{patient}(them)($\epsilon$)  \\\ibar{V}
				[$\lambda x \lambda \epsilon$ \textsc{eat}($\epsilon$) $\wedge$ \textsc{patient}($x$)($\epsilon$) \\\obar{V} [\shade{$\surd$eat}, name=trace] ]
				[DP [them, roof] ]
			] ]
		]
	]
	\draw[->,>=stealth', thin, gray] (trace) to[out=south west,in=south,looseness=1] node[pos=0.75,xshift=-10mm]{} (binder);
	\end{forest}
\end{exe}	
\end{fullwidth}
\medskip
The denotation of the \obar{v} that contains \emph{eat} is the same as the denotation of the \obar{v} that doesn't. The meaning of \emph{eat} has been ignored in computing the denotation of the larger~\obar{v}.

Note that (\ref{ex6:reconcopyv1}) provides a way of capturing reconstruction that, unlike the use of traces in the previous section, does not require that the moved term c-command the position in moves from. In general terms, what the copy theory of movement does is allow a term to be put into a variety of positions in a phrase marker. That term can be pronounced in one of those positions and semantically interpreted in another. The positions that a term can be distributed in is given entirely by movement; they are the positions that \textsc{Upwards} and the other conditions on derivations permit. The parses in a derivation that are semantically interpretable grow as a consequence. For a parse derived by movement to be semantically interpretable does not require, as it does on Trace theory, that the moved term be able to bind the variable it leaves in the lower position. This allows the Copy theory to provide a less problematic model of reconstruction than does Trace theory. In particular, the problems associated with Head Movement that are chronicled in the previous sections are unproblematic on the Copy theory of movement. 

To see this we'll need to add the semantic component that allows the Copy theory of movement to give the variable binding interpretation that the Trace theory of movement is designed to achieve directly. One place where that interpretation is needed is in~(\ref{ex6:copyhigh1}), where Argument movement brings a quantificational DP into a position higher than negation.
\begin{exe}
\raggedright
	\ex \label{ex6:copyhigh1}
	Some linguist won't finish this chapter.
\end{exe}
The Copy theory of movement provides this sentence with the representation in~(\ref{ex6:copyhigh1t}).
\begin{exe}
\raggedright
	\ex \label{ex6:copyhigh1t}
	\small
	\begin{forest}
	[TP
		[DP [some linguist, name=binder, roof]]
		[\ibar{T} 
			[\obar{T} [will] ]
			[NegP[\ibar{Neg}
				[\obar{Neg} [$\neg$] ]
				[AspP
				[\ibar{Asp}
				[\obar{Asp} [asp]]
				[vP 
					[DP [\shade{some linguist}, name=trace, roof] ]
					[\ibar{v} [finish this chapter, roof]]
				 ]
			] ]
		]
	]]]
	\draw[->,>=stealth', thin, gray] (trace) to[out=south west,in=south,looseness=1] node[pos=0.75,xshift=-10mm]{} (binder);
	\end{forest}
\end{exe}
The Trace theory of movement coupled movement with a set of interpretive rules that, in addition to inserting a variable in the lower position, changed the denotation of the phrase that the moved DP becomes a sister to. For that effect to be mimicked in the Copy theory of movement, we'll need a similar set of rules. I'll present here an approach to these rules that has been made popular by Danny Fox.\cite{Fox:1999}

To understand Fox's proposal, it will first be useful to know a few things about DPs that are headed by the determiner \emph{the}---so-called, ``definite descriptions.'' One popular conception of the meaning of \emph{the} is that it creates a DP which, like a pronoun, refers to something salient in the context. For instance, \emph{the man} and \emph{he} are roughly interchangeable in the scenario described below.
\begin{exe}
\raggedright
	\ex \label{ex6:dem1} \emph{My sister and I are watching people exit an event at school. After most people have left the building, a lone couple walks out: a man in a suit and a woman in a simple dress. My sister says:}
	\begin{xlist}
		\ex He used to be my neighbor.
		\ex The man used to be my neighbor.
	\end{xlist}
\end{exe}
Neither of these expressions would be appropriate in the scenario described below.
\begin{exe}
\raggedright
	\ex \label{ex6:dem2} \emph{My sister and I are watching people exit an event at school. After most people have left the building, a lone couple walks out: a man in a suit and tie and a man in jeans and a t-shirt. My sister says:}
	\begin{xlist}
		\ex He used to be my neighbor.
		\ex The man used to be my neighbor.
	\end{xlist}
\end{exe}
The first context makes salient a single male, but the second doesn't. If a requirement of both \emph{the man} and \emph{he} is that their referent be contextually salient, then we have an explanation for the contrast in (\ref{ex6:dem1}) and~(\ref{ex6:dem2}). 

The conditions under which a definite description can be used are not identical to those that govern where pronouns can be used, but there is a kind of family resemblance to their distribution. Another way in which they are similar is that definite descriptions can be interpreted as variables, like pronouns can. For instance, (\ref{ex6:bounddef1}) can have an interpretation that is roughly like that of~(\ref{ex6:bounddef2}).%
%
\marginnote[0pt]{There is a clear preference for (\ref{ex6:bounddef2}) over (\ref{ex6:bounddef1}). A complete account would include an explanation for this contrast; we'll focus merely on how (\ref{ex6:bounddef1}) can get the relevant interpretation.}%
%
\begin{exe}
\raggedright
	\ex \label{ex6:bounddef}
	\begin{xlist}
		\ex \label{ex6:bounddef1} No man's partner should belittle the man.
		\ex \label{ex6:bounddef2} No man's partner should belittle him.
	\end{xlist}
\end{exe}
In our syntax, bound variable interpretations come by way of the presence of a referential index. Because pronouns can contain a referential index, for instance, pronouns are capable of being interpreted as bound variables. What (\ref{ex6:bounddef1}) suggests, then, is that definite descriptions can also be equipped with referential indices.

Let's put these observations together. The denotation of \emph{the} must allow it to combine with an NP and produce something that refers to a salient individual. Let's take the denotations of NPs, like nouns, to be functions whose domains are things and whose ranges are the values True or False. They are functions, then, that characterize the sets of things that yield truth of the property named by the~NP. For instance, the noun, and NP, \emph{man} can be given the denotation in~(\ref{ex6:mand1}), which is a function that yields truth for any thing that it applies to which makes \textsc{man} true.
\begin{exe}
\raggedright
	\ex \label{ex6:mand1}
	\denotes{man} =_{\text{def}} $\lambda x$ \textsc{man}($x$)=T
\end{exe}
We want \emph{the} to combine with this function and refer to a salient individual that, when fed to this function, yields True. A preliminary way of describing this is illustrated in~(\ref{ex6:thed1}).
\begin{exe}
\raggedright
	\ex \label{ex6:thed1}
	\denotes{the} =_{\text{def.}} $\lambda P\ a$, where $P(a)=T$ and $a$ is a salient individual.
	\sn
	\begin{forest}
	[$a$ where \textsc{man}($a$){=}T and $a$ is salient\\DP
		[$\lambda P\ a$ where $P(a)${=}T and $a$ is salient\\\obar{D} [the] ]
		[$\lambda x$ \textsc{man}($x$){=}T\\ NP [\ibar{N} [\obar{N} [man]]] ]
	]
	\end{forest}
\end{exe}
I'm going to commandeer a symbol from the semantics literature to aid in abbreviating the denotation of \emph{the}. That symbol, \rotiota, I will define informally with~(\ref{ex6:iotadef}).%
%
\marginnote[-170pt]{Understand $a$ to be of semantic type entity: it is an arbitrarily chosen individual that satisfies $P$.

The symbol ``\rotiota'' was made popular by \citet{Whitehead:1910}, where it is used to model the meaning of \emph{the}. The semantics in \citet{Whitehead:1910} is different than (\ref{ex6:iotadef}) however. They treated ``\rotiota'' as a quantifier, which, like other quantifiers, relates the meaning of the NP it combines with to the meaning of the function that the DP headed by that quantifier is a sister to. On their treatment, \rotiota\ is an existential quantifier that imposes a requirement that there be only one thing of the kind that is quantified. This requirement is responsible for the contrast between (\ref{ex6:dem1}) and (\ref{ex6:dem2}). (i) exemplifies. (Note that Whitehead and Russell do not use events or $\theta$-roles to describe the denotations of sentences.)
\begin{exe}
\raggedright
	\exi{i.} \label{ex6:russ}
	\denotes{The man sat} = $\exists x$ \textsc{man}($x$) $\wedge$ \textsc{sit}($x$) $\wedge$ $\forall y$ (\textsc{man}($y$) $\rightarrow$ $y=x$)
\end{exe}
This approach to the meaning of \emph{the} is broached in \citet{Russell:1905} and improved in \citet{Neale:1990}. In (\ref{ex6:iotadef}), the uniqueness requirement is captured by the more general (and vague) requirement that $x$ be salient.}%
%
\begin{exe}
\raggedright
	\ex \label{ex6:iotadef}
	\rotiota $a\ P(a)$ = $a$, if $P(a)$ and $a$ is made salient by the context of use.
\end{exe}
Using this abbreviation, the denotation of \emph{the} can be expressed by~(\ref{ex6:dentheweak}).
\begin{exe}
\raggedright
	\ex \label{ex6:dentheweak}
	\denotes{$\surd$the} =_{\text{def.}} $\lambda P$ \rotiota $a\ P(a)$%
%
\marginnote[10pt]{Note that (\ref{ex6:dentheweak}) gives the denotation for the root of the word that will be exponed by the vocabulary item \emph{the}. This is explained immediately below.}%
%
\end{exe}
With this definition, (\ref{ex6:thed1}) becomes (\ref{ex6:thed2}).
\begin{exe}
\raggedright
	\ex \label{ex6:thed2}
	\begin{forest}
		[\rotiota $a$ \textsc{man}($a$)\\DP
			[$\lambda P\ \rotiota a$ $P(a)$\\\obar{D} [the] ]
			[$\lambda x$ \textsc{man}($x)$\\ NP [man, roof]]]
		\end{forest}
\end{exe}

Now we need to see how to build a definite description so that it includes a referential index. Let's adopt view suggested by Florian Schwarz\cite{Schwarz:2009}: a referential index can be part of the meaning of \emph{the}. The ambiguity of definite descriptions, then, comes from the ambiguity of \emph{the}. In addition to (\ref{ex6:dentheweak}), we'll give \emph{$\surd$the} the denotation in(\ref{ex6:denthestrong}). Assume that any \obar{D} that contains \emph{$\surd$the} is exponed by the word~\emph{the}.
\begin{exe}
\raggedright
	\ex \label{ex6:denthestrong}
	\denotes{$\surd$the} =_{\text{def.}} $\lambda n\ \lambda P$ \rotiota $a$ $a=n$ $\wedge$ $P(a)$
\end{exe}
%
%
\marginnote[-35pt]{There are other ways of expressing how a referential index allows a definite description to get the interpretation of a variable. This is essentially the method advocated by \citet{Schwarz:2009}, who uses it to give an account of the cross-linguistic variation that definite determiners can have. \citet{Elbourne:2005a} also has a good discussion of the issues involved in how the referential index is put into definite descriptions. \citet{Hanink:2018} provides evidence from Washo that helps situate indices in definite descriptions.}%
%
Take ``n'' to be a variable that ranges over referential indices. The denotation in (\ref{ex6:denthestrong}) allows \emph{the} to combine first with a referential index and then an NP to form a definite description. How this works is illustrated in~(\ref{ex6:indexden2}).
\begin{exe}
\raggedright
	\ex \label{ex6:indexden2}
	\begin{forest}
	[DP [\rotiota$a$ $a${=}$2$ $\wedge$ \textsc{man}($a$)\\\ibar{D}
		[$\lambda P$ \rotiota $a$ $a${=}$2$ $\wedge$ $P(a)$\\\obar{D}
		[$\lambda n\ \lambda P$ \rotiota$a$ $a${=}$n$ $\wedge$ $P(a)$\\\obar{D} [$\surd$the] ]
		[$2$]
		]
			[$\lambda y$ \textsc{man}($y$)\\NP [man, roof] ]
			]
			]
	\end{forest}
\end{exe}
The expression in (\ref{ex6:indexden2}) refers to the salient $a$ which is both a man and has the value that $2$ gets. I'll illustrate how this achieves the bound variable reading of \emph{the man} in (\ref{ex6:bounddef1}), falling back on the trace theory of movement. To produce a syntax that gives the right semantic interpretation to (\ref{ex6:bounddef1}) requires a movement rule that we haven't previously encountered. That rule will move \emph{no man} to the position indicated in~(\ref{ex6:qrgen1}). 
If this rule exists in English, it will have to be a rule that is trapped in that portion of the derivation that does not feed the S-structure. That such a rule exists, and why it is constrained to never produce S-structures, is discussed in \citet{Kennedy:2000a}. This rule won't be examined in detail in these lectures, so I will do nothing beyond using it here in an \emph{ad hoc} way.
\medskip
\begin{fullwidth}
\begin{exe}
\raggedright
	\ex \label{ex6:qrgen1}
	\footnotesize
\begin{forest}
[$\neg \exists y$ \textsc{man}($y$) $\wedge$ $\exists \epsilon$  \textsc{belittle}($\epsilon$) $\wedge$ \textsc{patient}($\epsilon$)(\rotiota$a$ $a${=}$y$ $\wedge$ \textsc{man}($a$))\\ $\wedge$ \textsc{agent}($\epsilon$)(\rotiota $a$ \textsc{partner\_of}\_y($a$)\\TP
	[$\lambda Q\ \neg \exists y$ \textsc{man}($y$) $\wedge$ $Q(y)$\\DP^2 [no man, roof]]
	[$\lambda 2\ \exists \epsilon$ \textsc{belittle}($\epsilon$) $\wedge$ \textsc{patient}($\epsilon$)(\rotiota$a$ $a${=}$2$ $\wedge$ \textsc{man}($a$))\\ $\wedge$ \textsc{agent}($\epsilon$)(\rotiota $a$ \textsc{partner\_of\_2}($a$))\\TP
	[\rotiota $a$ \textsc{partner\_of\_2}($a$)\\DP^3
		[\emph{t_2} ]
		[\ibar{D}
			[\obar{D} [gen]]
			[NP [partner, roof] ]
		 ]]
		[$\lambda 3\ \exists \epsilon$ \textsc{belittle}($\epsilon$) $\wedge$ \textsc{patient}($\epsilon$)(\rotiota$a$ $a${=}$2$ $\wedge$ \textsc{man}($a$))\\ $\wedge$ \textsc{agent}($\epsilon$)($3$)\\\ibar{T}
			[\obar{T} [should]]
			[$\lambda \epsilon$ \textsc{belittle}($\epsilon$) $\wedge$ \textsc{patient}($\epsilon$)(\rotiota$a$ $a${=}$2$ $\wedge$ \textsc{man}($a$))\\ $\wedge$ \textsc{agent}($\epsilon$)($3$)\\vP
				[\emph{t_3} ]
				[\ibar{v}
					[\obar{v} [agent]]
					[$\lambda \epsilon$ \textsc{belittle}($\epsilon$) $\wedge$ \textsc{patient}($\epsilon$)(\rotiota$a$ $a${=}$2$ $\wedge$ \textsc{man}($a$))\\VP [\ibar{V}
						[\obar{V} [belittle] ]
						[\rotiota$a$ $a${=}$2$ $\wedge$ \textsc{man}($a$)\\DP [\ibar{D} 
							[$\lambda P$ \rotiota $a$ $a${=}$2$ $\wedge$ $P(a)$\\\obar{D}
							[$\lambda i\ \lambda P$ \rotiota $a$ $a${=}$i$ $\wedge$ $P(a)$ \\\obar{D} [$\surd$the]]
							[2]
							]
								[$\lambda y$ \textsc{man}($y$)\\NP [man, roof] ]
							] ]
						] ]
					] ]
				 ]
			 ]
		]
	]
\end{forest}%
%
\marginnote[0pt]{As usual, I ignore the contribution \obar{T} makes to the denotation. Note that the subject is treated as a definite description. I have obscured how Aspect existentially closes the event variable, and I've also taken shortcuts in expressing how the genitive gets its meaning.}%
%
\end{exe}	
\end{fullwidth}
\vspace*{-50pt}
As desired, the referent of \emph{the man} varies with the value given to $y$ because of the referential index (=$2$) that it contains. A paraphrase of the denotation delivered to this sentence in (\ref{ex6:qrgen1}) is~(\ref{ex6:qrgenpara}).
\begin{exe}
\raggedright
	\ex \label{ex6:qrgenpara}
	There is no $y$ such that $y$ is a man and there is a belittling by the partner of $y$ whose patient is the man which is~$y$.
\end{exe}

Now that we've taken this detour into how definite descriptions work, let's return to the central task: giving the Copy theory of movement a way of expressing the bound variable meanings that movement sometimes creates. If movement creates representations like (\ref{ex6:copyhigh1t}), repeated below, then how do we extract from this representation the meaning in which \emph{some linguist} binds a variable in the Specifier of vP?
\begin{exe}
	\raggedright
		\exr{ex6:copyhigh1t}
		\small
		\begin{forest}
		[TP
			[DP [some linguist, name=binder, roof]]
			[\ibar{T} 
				[\obar{T} [will] ]
				[NegP[\ibar{Neg}
					[\obar{Neg} [$\neg$] ]
					[AspP
					[\ibar{Asp}
					[\obar{Asp} [asp]]
					[vP 
						[DP [\shade{some linguist}, name=trace, roof] ]
						[\ibar{v} [finish this chapter, roof]]
					 ]
				] ]
			]
		]]]
		\draw[->,>=stealth', thin, gray] (trace) to[out=south west,in=south,looseness=1] node[pos=0.75,xshift=-10mm]{} (binder);
		\end{forest}
\end{exe}
Danny Fox's proposal is to let the phrase in the lower position be interpreted in the same way that definite descriptions with a referential index are. %
%
\marginnote[0pt]{This proposal is in its essentials like a proposal in \citet{Engdahl:1980}, whose theory of movement is an early variant of Copy theory. (See also \citet{Engdahl:1986}, for a different implementation of the idea.)}%
%
%
\marginnote[20pt]{See \citet{Fox:1999,Fox:2002, Fox:2003a}. (\ref{ex6:movecopydef1}) deviates from Fox in requiring $\alpha$ to contain a DP for it to undergo Trace Conversion. Fox required $\alpha$ to \underline{be} the DP that Trace Converts. My version extends to cases of Pied-Piping; see \ref{sec:piedpipecopies}. \citet{Sauerland:1998c, Sauerland:2004} also provides a method for interpreting copies as variables that shares features with Fox.}% 
%
Movement will be defined with~(\ref{ex6:movecopydef1}).
\begin{exe}
\raggedright
	\ex \label{ex6:movecopydef1}
	A term $\alpha$ moves to position $\beta$ =_{\text{def.}}
	\begin{xlist}
		\ex $\beta$ becomes occupied by \textsc{copy}($\alpha$), 
		\ex \label{ex6:tracecon}If $\alpha$ contains a DP, then:
		\begin{xlist}
		\ex that DP's head can be interpreted as the \emph{the} which contains referential index~$n$, and
		\ex The sister of $\beta$ goes from \denotes{$\rho$} to $\lambda n$\denotes{$\rho$}.
		\end{xlist}
	\end{xlist}
\end{exe}
Fox calls (\ref{ex6:tracecon}) ``Trace Conversion.'' It mimics closely the definition of traces in the Trace theory of movement, by allowing DPs involved in a movement relation to be interpreted as bound variables. It delivers to (\ref{ex6:copyhigh1}) the denotation shown in~(\ref{ex6:copyhigh1td}).
\medskip
\begin{fullwidth}
\begin{exe}
	\raggedright
		\ex \label{ex6:copyhigh1td}
		\small
		\begin{forest}
		[$\exists y$ \textsc{linguist}($y$) $\wedge$ $\neg \exists \epsilon$ \textsc{finish}($\epsilon$) $\wedge$ \textsc{patient}($\epsilon$)(\denotes{this chapter})\\ $\wedge$ \textsc{agent}($\epsilon$)(\rotiota $a$ $a${=}$y$ $\wedge$ \textsc{linguist}($a$))\\TP
			[$\lambda P\ \exists y$ \textsc{linguist}($y$) $\wedge$ $P(y)$\\DP [some linguist, name=binder, roof]]
			[$\lambda 2\ \neg \exists \epsilon$ \textsc{finish}($\epsilon$) $\wedge$ \textsc{patient}($\epsilon$)(\denotes{this chapter})\\ $\wedge$ \textsc{agent}($\epsilon$)(\rotiota $a$ $a${=}$2$ $\wedge$ \textsc{linguist}($a$))\\\ibar{T} 
				[\obar{T} [will] ]
				[$\neg \exists \epsilon$ \textsc{finish}($\epsilon$) $\wedge$ \textsc{patient}($\epsilon$)(\denotes{this chapter})\\ $\wedge$ \textsc{agent}($\epsilon$)(\rotiota $a$ $a${=}$2$ $\wedge$ \textsc{linguist}($a$))\\NegP[\ibar{Neg}
					[\obar{Neg} [$\neg$] ]
					[AspP
					[\ibar{Asp}
					[\obar{Asp} [asp]]
					[$\lambda \epsilon$ \textsc{finish}($\epsilon$) $\wedge$ \textsc{patient}($\epsilon$)(\denotes{this chapter})\\ $\wedge$ \textsc{agent}($\epsilon$)(\rotiota $a$ $a${=}$2$ $\wedge$ \textsc{linguist}($a$))\\vP 
						[\rotiota $a$ $a${=}$2$ $\wedge$ \textsc{linguist}($a$)\\DP [\shade{some linguist}, name=trace, roof] ]
						[\ibar{v} [finish this chapter, roof]]
					 ]
				] ]
			]
		]]]
		\draw[->,>=stealth', thin, gray] (trace) to[out=south west,in=south,looseness=1] node[pos=0.75,xshift=-10mm]{} (binder);
		\end{forest}
	\end{exe}
\end{fullwidth}
\medskip
This is the desired meaning: it interprets \emph{some linguist} outside negation, and causes it to bind a variable inside the scope of negation. That variable is now a definite description---something semantically equivalent to \emph{the linguist}---rather than a trace.

By letting DPs that are in a movement relation become definite descriptions which get interpreted as variables, the Copy theory of movement provides an account for Reconstruction of pronominal variables that explains why the constraints imposed by the pronoun are preserved. To see how, consider~(\ref{ex6:vrrecon1}).
\begin{exe}
\raggedright
	\ex \label{ex6:vrrecon1}
	Whose articles about her did no woman remember?
\end{exe}
The pronoun \emph{her} in this example can be interpreted as a variable bound by \emph{no woman}. As we saw in section~\ref{sec:reconvariable}, this is possible because \emph{whose articles about her} has moved from a position that is c-commanded by \emph{no woman}, and when the pronoun is in that position it is capable of being bound by \emph{no woman}. The Copy theory of movement produces the representation in~(\ref{ex6:vrrecon1t}).
\medskip
\begin{fullwidth}
\begin{exe}
\raggedright
	\ex \label{ex6:vrrecon1t}
	\small
	\begin{forest}
	[CP
		[DP [whose articles about her, roof]]
		[\ibar{C} 
			[\obar{C} [did]]
			[TP
				[DP^2 [no woman, roof] ]
				[\ibar{T} 
					[\obar{T} [\shade{did}] ]
					[vP
						[DP [\shade{no woman}, roof] ]
						[\ibar{v}
							[\obar{v} [agent] ]
							[VP [\ibar{V}
								[\obar{V} [$\surd$remember] ]
								[DP [\shade{whose articles about her_2}, roof] ]
							] ]
						 ]
					 ]
				]
			 ]
		]
	]
	\end{forest}
\end{exe}	
\end{fullwidth}
\medskip
Note in particular that \emph{her} is capable of being bound in the parallel~(\ref{ex6:vrrecon1b}).
\begin{exe}
\raggedright
	\ex[]{\label{ex6:vrrecon1b} No woman^2 remembered my articles about her_2.}
\end{exe}
By contrast, the reflexive pronoun in (\ref{ex6:vrrecon2}) cannot be bound by \emph{no woman}, because the representation the Copy theory of movement gives to it (in~(\ref{ex6:vrrecon2t})) does not put the reflexive pronoun in a position that allows it to be bound by \emph{no woman}, as the comparison with~(\ref{ex6:vrrecon2b}) indicates.
\medskip
\begin{fullwidth}
\begin{exe}
\raggedright
	\ex[*]{\label{ex6:vrrecon2} Whose article about herself_2 did no woman^2 remember?}
	\ex \label{ex6:vrrecon2t}
	\small
	\begin{forest}
		[CP
			[DP [whose articles about her, roof]]
			[\ibar{C} 
				[\obar{C} [did]]
				[TP
					[DP^2 [no woman, roof] ]
					[\ibar{T} 
						[\obar{T} [\shade{did}] ]
						[vP
							[DP [\shade{no woman}, roof] ]
							[\ibar{v}
								[\obar{v} [agent] ]
								[VP [\ibar{V}
									[\obar{V} [$\surd$remember] ]
									[DP [\shade{whose articles about herself_2}, roof] ]
								] ]
							 ]
						 ]
					]
				 ]
			]
		]
		\end{forest}	
		\ex[*]{\label{ex6:vrrecon2b} No woman^2 remembered someone's article about herself_2.}
\end{exe}	
\end{fullwidth}
\medskip

Note that as these representations make clear, we should adopt an account of the semantics of wh-questions that allows the wh-phrase to be semantically interpreted in just its lower position. We do not want the higher copy to be semantically interpreted, for that would wrongly put \emph{her}, or \emph{herself}, in a position in which it cannot be bound by \emph{no woman}.

The Copy theory of movement is successful at addressing the problems chronicled in this section for the Trace theory of movement. It does so by allowing one parse in a derivation to have a lot of the information that previous parses in that derivation contain. This makes it possible to let one parse give the semantic interpretation mechanisms what is needed to get interpretations that were available only in different parses in the Trace theory of movement. In this way, the Copy theory of movement allows derivations to obey Upwards, but still deliver parses whose semantic interpretations were inconsistent with an S-structure on Trace theory. Let me illustrate that by considering how Trace theory and Copy theory apply to (\ref{ex6:demoder}), which is akin to a case discussed in section~\ref{sec:ordering}.
\begin{exe}
\raggedright
	\ex \label{ex6:demoder}
	Some student solved each problem for me. 
\end{exe}
The S-structure/interpreted parse pair that the Copy theory allows for (\ref{ex6:demoder}) is~(\ref{ex6:demodert}). 
\medskip
\begin{fullwidth}
\begin{exe}
\raggedright
\ex \label{ex6:demodert}
a.
\footnotesize
\begin{forest}
[TP
		[DP^2 [some student, roof] ]
		[\ibar{T}
			[\obar{T} [past] ]
			[vP
				[DP_2 [\shade{some student}, roof] ]
				[\ibar{v}
					[\obar{v} 
					[\obar{V} [$\surd$solve]]
					[\obar{v} [agent]]
					]
					[VP [\ibar{V}
					[\ibar{V}
						[\obar{V} [\shade{$\surd$solve}] ]
						[DP [each problem, roof]]
					] 
						[PP [for me, roof]]]
				 ]
			 ]
		 ]
	 ]]
\end{forest}

\normalsize b. \footnotesize
	\begin{forest}
	[TP
		[TP
			[DP^2 [some student, roof] ]
			[\ibar{T}
				[\obar{T} [past] ]
				[vP
					[DP_2 [\shade{some student}, roof] ]
					[\ibar{v}
						[\obar{v} 
						[\obar{V} [$\surd$solve]]
						[\obar{v} [agent]]
						]
						[VP [\ibar{V}
						[\ibar{V}
							[\obar{V} [\shade{$\surd$solve}] ]
							[DP_3 [\shade{each problem}, roof]]
						] 
							[PP [for me, roof]]]
					 ]
				 ]
			 ]
		 ]]
		[DP^3 [each problem, roof]]
	]
	\end{forest}
\end{exe}
\end{fullwidth}
\medskip
The higher copy of \emph{solve} can go uninterpreted, as it must. And the lower copies of both DPs can be interpreted as variables, giving \emph{each problem} higher scope than \emph{some student}.

Trace theory provides the S-structure/interpreted parse in~(\ref{ex6:demodert2}).
\medskip
\begin{fullwidth}
\begin{exe}
\raggedright
\ex \label{ex6:demodert2}
a.
\footnotesize
\begin{forest}
	[TP
			[DP^2 [some student, roof] ]
			[\ibar{T}
				[\obar{T} [past] ]
				[vP
					[\emph{t_2} ]
					[\ibar{v}
						[\obar{v} 
						[\obar{V} [$\surd$solve]]
						[\obar{v} [agent]]
						]
						[VP [\ibar{V}
						[\ibar{V}
							[\obar{V} ]
							[DP [each problem, roof]]
						] 
							[PP [for me, roof]]]
					 ]
				 ]
			 ]
		 ]]
	\end{forest}
\normalsize b. \footnotesize
	\begin{forest}
	[TP
		[TP
			[DP^2 [some student, roof] ]
			[\ibar{T}
				[\obar{T} [past] ]
				[vP
					[\emph{t_2} ]
					[\ibar{v}
						[\obar{v} 
						[\obar{v} [agent]]
						]
						[VP [\ibar{V}
						[\ibar{V}
							[\obar{V}  [$\surd$solve]]
							[\emph{t_3} ]
						] 
							[PP [for me, roof]]]
					 ]
				 ]
			 ]
		 ]]
		[DP^3 [each problem, roof]]
	]
	\end{forest}
\end{exe}
\end{fullwidth}
\medskip
The interpreted parse can be derived from an Upwards obeying movement of the object \emph{each problem} from the S-structure in (\ref{ex6:demodert}). But there is no way of putting both parses in (\ref{ex6:demodert2}) into an Upwards obeying derivation. The Copy theory of movement provides a solution for all of the problems this section chronicles for Trace theory. It is compatible with the Lebeaux+Upwards model, and inherits from that model its ability to derive the Law of Reconstruction. 

\subsubsection{Pied-Piping}
\label{sec:piedpipecopies}

The Copy theory of movement also has a means for addressing a problem we encountered in the discussion of relative clauses. %
%
\marginnote[0pt]{See section \ref{sec:relclause}.}%
%
That problem is illustrated by the relative clause in~(\ref{ex6:piedpden}), repeated here.
\begin{exe}
	\raggedright
		\exr{ex6:piedpden}
		the dog whose owner^2 you met \emph{t}_2
	\end{exe}
In (\ref{ex6:piedpden}), Wh Movement has moved a constituent larger than the wh-phrase. These are cases of what is known as ``pied-piping'': Wh Movement moves more than just the \obar{X} that contains the wh-feature. It instead moves some phrase dominating that \obar{X}. In (\ref{ex6:piedpden}), it has moved the DP \emph{whose owner}, which contains the wh-feature in \emph{who}. 

Cases of pied-piping like this present a problem for our analysis of relative clauses because it doesn't deliver the right meaning from Trace theory. Relative clauses have a meaning that restricts the range of things that the \ibar{N} they combine with can refer to. In (\ref{ex6:piedpden}), we want the relative clause to restrict the range of things that \emph{dog} can refer to. But Trace theory causes relative clauses to describe the things that could fit semantically into the trace position. In (\ref{ex6:piedpden}), this means that the relative clause describes things that you met. Trace theory therefore causes (\ref{ex6:piedpden}) to have the same meaning that \emph{the dog who you met} has. Because the Copy theory of movement doesn't require the position from which a phrase moves to be a variable, it doesn't have the same consequences in this context. Let's see how.

	Recall that relative clauses are formed with Wh Movement, and that the semantics of movement causes them to become functions from individuals to truth values. In the simpler relative clause of (\ref{ex6:piedpden2}), for instance, a semantics like that in (\ref{ex6:piedpden2t}) is created on the Trace theory of movement. (I take a variety of simplifying shortcuts in~(\ref{ex6:piedpden2t}).)
	\medskip
	\begin{fullwidth}
	\begin{exe}
	\raggedright
		\ex \label{ex6:piedpden2} the dog who you met.
		\ex \label{ex6:piedpden2t}
		\small
		\begin{forest}
		[DP [\ibar{D}
			[\obar{D} [the]]
			[NP [$\lambda x$ \textsc{dog}($x$) $\wedge$ $\exists \epsilon$ \textsc{met}($\epsilon$) $\wedge$ \textsc{agent}($\epsilon$)(\denotes{you}) $\wedge$ \textsc{patient}($\epsilon$)($x$)\\\ibar{N}
				[$\lambda y$ \textsc{dog}($y$)\\\ibar{N} [dog, roof] ]
				[$\lambda 2\ \exists \epsilon$ \textsc{met}($\epsilon$) $\wedge$ \textsc{agent}($\epsilon$)(\denotes{you}) $\wedge$ \textsc{patient}($\epsilon$)($2$)\\CP
					[DP [who, roof] ]
					[$\lambda 2\ \exists \epsilon$ \textsc{met}($\epsilon$) $\wedge$ \textsc{agent}($\epsilon$)(\denotes{you}) $\wedge$ \textsc{patient}($\epsilon$)($2$)\\\ibar{C}
						[\obar{C} ]
						[TP
							[DP [you, roof] ]
							[\ibar{T}
								[\obar{T} [past] ]
								[VP [\ibar{V}
									[\obar{V} [met]]
									[\emph{t_2} ]
								] ]
							 ]
						 ]
					 ]
				 ]
			] ]
		]]
		\end{forest}
	\end{exe}
\end{fullwidth}
\medskip
Movement of \emph{who} causes its sister \ibar{C} to become a $\lambda$ expression in which the trace left by movement is interpreted as a variable. If relative pronouns like \emph{who} are semantically vacuous, which is the analysis we adopted in section \ref{sec:relclause}, then the denotation that \ibar{C} has will be the same as the denotation that CP has. This allows the CP and the \ibar{N} that is its sister to combine semantically in the way indicated in~(\ref{ex6:piedpden2t}), forming the desired meaning. The relative clause in (\ref{ex6:piedpden}) does not get the right meaning on this theory because the variable in (\ref{ex6:piedpden}) is bound by \emph{whose owner}, but what we want is, parallel to (\ref{ex6:piedpden2}), a relative clause whose variable corresponds to dogs.
\pagebreak

This can be remedied by the Copy Theory of movement, which is repeated below.
\begin{exe}
	\raggedright
		\exr{ex6:movecopydef1}
		A term $\alpha$ moves to position $\beta$ =_{\text{def.}}
		\begin{xlist}
			\ex $\beta$ becomes occupied by \textsc{copy}($\alpha$), 
			\ex \label{ex6:tracecon}If $\alpha$ contains a DP, then:
			\begin{xlist}
			\ex that DP's head can be interpreted as the \emph{the} which contains referential index~$n$, and
			\ex The sister of $\beta$ goes from \denotes{$\rho$} to $\lambda n$\denotes{$\rho$}.
			\end{xlist}
		\end{xlist}
	\end{exe}
Let's first consider how the Copy theory applies to (\ref{ex6:piedpden2}), which involves the movement of a wh-pronoun. 

Recall that pronouns are determiners that include information about the thing that their DP refers to. The pronoun \emph{she}, for instance, indicates that the DP it heads will refer to a singular female. The pronoun \emph{who} has information that indicates that its referent will be human. Let's give this a syntactic expression by letting the DPs that \emph{who} heads have the form in~(\ref{ex6:whodp}).
\begin{exe}
\raggedright
	\ex \label{ex6:whodp}
	\footnotesize
	\begin{forest}
	[DP [\ibar{D}
		[\obar{D} 
		[\obar{D} [$\surd$who]]
		[$\phi$ [human]]
		]
	]]
	\end{forest}
\end{exe}
With this analysis of \emph{who}, the Copy theory can assign (\ref{ex6:piedpden2}) the representation in~(\ref{ex6:piedpden2ct}).
\medskip
\begin{fullwidth}
\begin{exe}
\raggedright
	\ex \label{ex6:piedpden2ct}
	\footnotesize
	\begin{forest}
		[DP [\ibar{D}
			[\obar{D} [the]]
			[NP [$\lambda x$ \textsc{dog}($x$) $\wedge$ $\exists \epsilon$ \textsc{met}($\epsilon$) $\wedge$ \textsc{agent}($\epsilon$)(\denotes{you}) $\wedge$ \textsc{patient}($\epsilon$)(\rotiota $a$ $a${=}$x$ $\wedge$ \textsc{human}($a$))\\\ibar{N}
				[$\lambda y$ \textsc{dog}($y$)\\\ibar{N} [dog, roof] ]
				[$\lambda 2\ \exists \epsilon$ \textsc{met}($\epsilon$) $\wedge$ \textsc{agent}($\epsilon$)(\denotes{you}) $\wedge$ \textsc{patient}($\epsilon$)(\rotiota $a$ $a${=}$2$ $\wedge$ \textsc{human}($a$))\\CP
					[DP [who, roof] ]
					[$\lambda 2\ \exists \epsilon$ \textsc{met}($\epsilon$) $\wedge$ \textsc{agent}($\epsilon$)(\denotes{you}) $\wedge$ \textsc{patient}($\epsilon$)(\rotiota $a$ $a${=}$2$ $\wedge$ \textsc{human}($a$))\\\ibar{C}
						[\obar{C} ]
						[TP
							[DP [you, roof] ]
							[\ibar{T}
								[\obar{T} [past] ]
								[VP [\ibar{V}
									[\obar{V} [met]]
									[DP [\ibar{D}
										[\rotiota $a$ $a${=}$2$ $\wedge$ \textsc{human}$(a)$\\\obar{D} 
										[$\lambda P$ \rotiota $a$ $a${=}$2$ $\wedge$ $P(a)$\\\obar{D} [$\lambda n\ \lambda P$ \rotiota $a$ $a${=}$n$ $\wedge$ $P(a)$\\\obar{D}
										[\shade{$\surd$who}]]
										[$2$]]
									[$\phi$ [human]]
									]]
									]
								] ]
							 ]
						 ]
					 ]
				 ]
			] ]
		]]
		\end{forest}
\end{exe}
\end{fullwidth}
\medskip
Note that the lower copy has undergone the ``trace conversion'' described by (\ref{ex6:tracecon}). Trace conversion gives the \obar{D} occupied by \emph{$\surd$who} the meaning that a definite determiner with an index has. This turns the lower copy of \emph{$\surd$who} into a variable that is restricted in the same way that \emph{who} is---that is, it ranges over humans. The result is equivalent, meaning-wise, to~(\ref{ex6:piedpden2t}).

Consider now the representation that the Copy theory will give to the relative clause in~(\ref{ex6:piedpden}).
\medskip
\begin{fullwidth}
	\begin{exe}
\ex \label{ex6:piedpdenct}
\footnotesize
\begin{forest}
[$\lambda 2\ \exists \epsilon$ \textsc{agent}($\epsilon$)(\denotes{you}) $\wedge$ \textsc{meet}($\epsilon$)\\ $\wedge$ \textsc{patient}($\epsilon$)(\rotiota $b$ \textsc{owner}($b$) $\wedge$ \textsc{has}($b$)(\rotiota $a$ $a${=}$2$ $\wedge$ \textsc{human}$(a)$)\\CP
	[DP [whose owner, roof] ]
	[$\lambda 2\ \exists \epsilon$ \textsc{agent}($\epsilon$)(\denotes{you}) $\wedge$ \textsc{meet}($\epsilon$)\\ $\wedge$ \textsc{patient}($\epsilon$)(\rotiota $b$ \textsc{owner}($b$) $\wedge$ \textsc{has}($b$)(\rotiota $a$ $a${=}$2$ $\wedge$ \textsc{human}$(a)$)\\\ibar{C}
		[\obar{C} ]
		[$\exists \epsilon$ \textsc{agent}($\epsilon$)(\denotes{you}) $\wedge$ \textsc{meet}($\epsilon$)\\ $\wedge$ \textsc{patient}($\epsilon$)(\rotiota $b$ \textsc{owner}($b$) $\wedge$ \textsc{has}($b$)(\rotiota $a$ $a${=}$2$ $\wedge$ \textsc{human}$(a)$)\\TP
			[DP [you, roof] ]
			[\ibar{T}
				[\obar{T} [past] ]
				[VP [\ibar{V}
					[\obar{V} [met] ]
					[\rotiota $b$ \textsc{owner}($b$) $\wedge$ \textsc{has}($b$)(\rotiota $a$ $a${=}$2$ $\wedge$ \textsc{human}$(a)$)\\DP
						[DP [\ibar{D}
							[\rotiota $a$ $a${=}$2$ $\wedge$ \textsc{human}$(a)$\\\obar{D}
								[$\lambda P$ \rotiota $a$ $a${=}$2$ $\wedge$ \textsc{human}$(a)$\\\obar{D} [$\lambda P$ \rotiota $a$ $a${=}$2$ $\wedge$ $P(a)$\\\obar{D}
									[$\lambda n\ \lambda P$ \rotiota $a$ $a${=}$n$ $\wedge$ $P(a)$\\\obar{D} [\shade{$\surd$who}]]
									[$2$ ]
								 ]
								[$\phi$ [human] ]
							 ]]]]
							 [\ibar{D}
							[\obar{D} [gen] ]
						[NP [\shade{owner}, roof] ]
							 ]]
					 ]
				] ]
			 ]
		 ]
	 ]
\end{forest}
\end{exe}
\end{fullwidth}
\medskip
The object of \emph{meet} in (\ref{ex6:piedpdenct}) refers to owners, as desired. But the variable bound by the $\lambda$ expression that is the denotation of CP is buried within the object of \emph{meet}: it corresponds to \emph{whose}, which has been converted into a variable by the Trace Conversion part of the Copy Theory of movement.%
%
\marginnote{Pied-piping, on this account, requires a theory which determines which phrases dominating a feature can, or must, move. There is a large literature attempting to provide that theory, but at the time of writing, there is no consensus view that has emerged. A good introduction to the phenomena, and a development of one of the popular theories, is \citet{Heck:2008, Heck:2009}.}%
%
\ Simplifying out the irrelevant details of (\ref{ex6:piedpdenct}), the meaning delivered to the relative clause is~(\ref{ex6:simpledenct}).
\begin{exe}
\raggedright
	\ex \label{ex6:simpledenct} $\lambda x$ you met $x$'s owner
\end{exe}
And that is just the meaning needed for (\ref{ex6:piedpden}):
\begin{exe}
\raggedright
	\ex \label{ex6:piedpdenfinal}
	\small
	\begin{forest}
	[DP
		[\obar{D} [the]]
		[NP 
		[$\lambda z$ \textsc{dog}($z$) $\wedge$ you met $z's$ owner\\\ibar{N}
			[$\lambda y$ \textsc{dog}($y$)\\\ibar{N} [dog] ]
			[$\lambda x$ you met $x$'s owner\\CP [whose owner you met, roof]]
		] ]
	]
	\end{forest}
\end{exe}

The Copy theory of movement can deliver the right semantics for relative clauses when they involve Pied-piping.

\subsubsection{Problem: Antecedent Contained Anaphora}
\label{sec:acdproblem}

As successful as the Copy theory of movement is in solving the problems that reconstruction poses for Trace theory, it harbors a serious problem. It isn't well-equipped to handle the central problem that Trace theory solves: antecedent contained anaphora. An illustration of that problem, and its Trace theoretic solution, is in~(\ref{ex6:acdcopy1}).
\medskip
\begin{fullwidth}
\begin{exe} 
	\raggedright
			\ex \label{ex6:acdcopy1s} She will introduce everyone who she should to me.
			\ex \label{ex6:acdcopy1}
			\footnotesize
			\begin{forest}
				[TP
				[TP
					[DP^2 [she, roof]]
					[\ibar{T} 
						[\obar{T} [will] ]
						[AspP
						[\ibar{Asp}
						[\obar{Asp} [asp]]
						[vP
							[\emph{t_2} ]
							[\ibar{v}
								[\obar{v} [agent]]
								[VP [\ibar{V}
									[\ibar{V}
									[\obar{V} [$\surd$introduce] ]
									[\emph{t_3} ]
									]
									[PP [to me, roof]]
								] ]
							 ]
						 ]]]]]
				[DP^3 [\ibar{D}
					[\obar{D} [every]]
					[NP [\ibar{N}
						[\ibar{N} [one, roof]]
						[CP 
						[DP^5 [who, roof]]
						[\ibar{C}
							[\obar{C} ]
							[TP 
								[DP^4 [she, roof] ]
								[\ibar{T}
									[\obar{T} [SHOULD]]
									[vP
										[\emph{t}_4 ]
										[\ibar{v}
										[\obar{v} [agent]]
										[\fbox{VP} [\ibar{V}
											[\ibar{V} 
												[\obar{V} [$\surd$introduce]]
												[\emph{t_5} ]
											]
											[PP [to me, roof] ]
										] ]]]
									 ]
								 ]
							 ]
						] ]
					] ]
			]]
			\end{forest}%
			%
			\marginnote[0pt]{Movement of the verbal roots into \obar{v}, as well as other details, are ignored in this graph.}%
			%
	\end{exe}
\end{fullwidth}
\medskip
The S-structure that corresponds to (\ref{ex6:acdcopy1s}) can be in a derivation that includes the parse shown in (\ref{ex6:acdcopy1})---QR can form this parse from the S-structure---and this allows the boxed VP to elide. This is possible in part because the boxed VP matches the VP from which the object has QR'd. %
%
\marginnote[0pt]{The condition on ellipsis requires that the elided phrase be part of a constituent that matches the meaning of an antecedent phrase, except for those places where there is contrastive focus. See section \ref{sec:slopana} for details.}%
%
On the Trace theory of movement, movement is capable of creating an antecedent for the elided VP in~(\ref{ex6:acdcopy1s}). By contrast, the Copy theory of movement associates with (\ref{ex6:acdcopy1s}) the representation in (\ref{ex6:acdcopy1a}), and the VPs in this representation do not match in the necessary way.
\medskip
\begin{fullwidth}
	\begin{exe}
	\ex \label{ex6:acdcopy1a}
	\footnotesize
	\begin{forest}
	[TP
		[TP
			[DP^2 [she, roof]]
			[\ibar{T} 
				[\obar{T} [will] ]
				[vP
					[DP_2 [\shade{she}, roof] ]
					[\ibar{v}
						[\obar{v} [agent]]
						[VP [\ibar{V}
							[\ibar{V}
							[\obar{V} [$\surd$introduce] ]
							[DP_3 [everyone who she SHOULD\\ \sout{introduce to me}, roof] ]
							]
							[PP [to me, roof]]
						] ]
					 ]
				 ]
			]
		 ]
		[DP^3 [\ibar{D}
			[\obar{D} [every]]
			[NP [\ibar{N}
				[\ibar{N} [one, roof]]
				[CP 
				[DP^5 [who, roof]]
				[\ibar{C}
					[\obar{C} ]
					[TP 
						[DP^4 [she, roof] ]
						[\ibar{T}
							[\obar{T} [SHOULD]]
							[vP
								[DP_4 [\shade{she}, roof]]
								[\ibar{v}
								[\obar{v} [agent]]
								[\fbox{VP} [\ibar{V}
									[\ibar{V} 
										[\obar{V} [$\surd$introduce]]
										[DP_5 [\shade{who}, roof] ]
									]
									[PP [to me, roof] ]
								] ]
							 ]
						 ]
					 ]
				] ]
			] ]
		] ]
	]]
	\end{forest}%
	%
	\marginnote[0pt]{Assume that this representation is derived from the S-structure by Moving the object DP. Note that the VP \emph{introduce to me} is silent in the S-structure because of VP ellipsis---this is indicated with~\sout{strikeouts}.}%
	%
\end{exe}
\end{fullwidth}
\medskip
The Copy theory of movement doesn't do the job of creating the antecedent for VP ellipsis. 

There is a conflict in what reconstruction phenomena require, and what antecedent contained anaphora require. Reconstruction phenomena seem to need the moved phrase to be in a position different from where it is pronounced. If movement is the process that relates a term to its two (or more) syntactic positions, then reconstruction phenomena recommend a theory of movement that keeps the moved term in its various positions syntactically. The Copy theory is designed for this effect. By contrast, antecedent contained anaphora seems to require that movement remove a term from one of its syntactic positions, and occupy that position with a variable instead. This is what Trace theory is designed for. At the time of writing, the resolution of these conflicting needs has not been fully achieved.

One direction presently under investigation extends that part of the Copy theory of movement that allows one of the copies to be ignored by the semantic component. In cases where a verb moves, for instance, it's necessary that only the lower copy of that verb be semantically interpreted: the higher copy of the verb's semantic contribution can be ignored. This is one of results of the instructions that determine how the syntactic representations created by movement are mapped onto a semantics and a string. Those instructions are repeated in~(\ref{ex6:map1}).
\begin{exe}
\raggedright
	\ex \label{ex6:map1}
	\begin{xlist}
		\ex Only one copy may be linearized.
		\ex Only one copy must be semantically interpreted.
	\end{xlist}
\end{exe}
(\ref{ex6:map1}a) prevents more than one of the copies of some term from being part of the string that corresponds to the syntactic representation of movement. And (\ref{ex6:map1}b) allows copies created by movement to be semantically ignored, as long as just one of them is interpreted. Together, (\ref{ex6:map1}a) and (\ref{ex6:map1}b) mimic some of the information that derivations have. In conjunction with the notion of copies, they permit a term to be pronounced in a position different than where that term is semantically interpreted. In (\ref{ex6:map1demo}), for instance, the verbal root \emph{$\surd$eat} is pronounced in the \obar{v} position, but semantically interpreted in its lower position.
\begin{exe}
\raggedright
	\ex \label{ex6:map1demo}
	\small
	\begin{forest}
	[vP
		[DP [she, roof]]
		[\ibar{v} 
			[\obar{v}
				[\obar{V} [$\surd$eat, name=binder] ]
				[\obar{v} [agent] ]
			 ]
			[VP [\ibar{V}
				[\obar{V} [\shade{$\surd$eat}, name=trace] ]
				[DP [nattoo, roof] ]
			]]
		]
	]
	\draw[->,>=stealth', thin, gray] (trace) to[out=south west,in=south,looseness=1] node[pos=0.75,xshift=-10mm]{} (binder);
	\end{forest}
\end{exe}
Determining which copy is pronounced and which is semantically interpreted is equivalent to determining which parse in a derivation is pronounced and which is semantically interpreted. The solution to antecedent contained anaphora I will describe involves loosening (\ref{ex6:map1}) so that parts of copies are allowed to be semantically and phonologically interpreted.

In place of (\ref{ex6:map1}), then, consider~(\ref{ex6:map2}).%
%
\marginnote[0pt]{This condition is modeled on a proposal in \citet{Sportiche:2019}.}%
%
\begin{exe}
\raggedright
	\ex \label{ex6:map2}
	Let \{$\gamma$, $\gamma '$, $\gamma ''$, \ldots\} be matching phrases or heads dominated by copies \{$\alpha$,~$\alpha '$,~$\alpha ''$, \ldots\}.
	\begin{xlist}
		\ex Only one $\gamma$^n may be linearized.
		\ex Only one $\gamma$^n must be semantically interpreted.
	\end{xlist}
\end{exe}
This allows parts of a moved phrase to be semantically, or phonologically, interpreted. It takes away from movement the job of choosing which things will be pronounced and which semantically interpreted. Instead movement only determines which phrases are eligible to be semantically/phonologically interpreted in different places, and where those places will be. 

I'll adopt the convention of placing within denotation brackets (i.e., \denotes{\ }) the terms that will be semantically interpreted, and I'll continue the practice of indicating copied material that isn't pronounced by putting it in \shade{shaded} font. Reconsider now the representation that movement allows for the instance of antecedent contained anaphora in~(\ref{ex6:acdcopy1s}). That representation (in (\ref{ex6:acdcopy1a})) can now be interpreted along the lines allowed by (\ref{ex6:map2}) in the way indicated in~(\ref{ex6:copy1latea}).
\medskip
\begin{fullwidth}
	\begin{exe}
	\ex \label{ex6:copy1latea}
	\footnotesize
	\begin{forest}
	[TP
		[TP
			[DP^2 [she, roof]]
			[\ibar{T} 
				[\obar{T} [will] ]
				[vP
					[DP_2 [\shade{she}, roof] ]
					[\ibar{v}
						[\obar{v} [agent]]
						[VP [\ibar{V}
							[\ibar{V}
							[\obar{V} [$\surd$introduce] ]
							[\denotes{DP_3} [\denotes{\ibar{D}}
								[\denotes{\obar{D}} [every] ]
								[\denotes{NP} [\denotes{\ibar{N}} 
									[\denotes{\ibar{N}} [one, roof] ]
									[CP [who she SHOULD \sout{introduce to me}, roof] ]
								] ]
							]
							]
							]
							[PP [to me, roof]]
						] ]
					 ]
				 ]
			]
		 ]
		[\denotes{DP^3} [\denotes{\ibar{D}}
			[\denotes{\obar{D}} [every]]
			[\denotes{NP} [\denotes{\ibar{N}}
				[\denotes{\ibar{N}} [one, roof]]
				[\denotes{CP} 
				[DP^5 [who, roof]]
				[\denotes{\ibar{C}}
					[\obar{C} ]
					[\denotes{TP} 
						[\denotes{DP^4} [she, roof] ]
						[\denotes{\ibar{T}}
							[\denotes{\obar{T}} [SHOULD]]
							[\denotes{vP}
								[\denotes{DP_4} [\shade{she}, roof]]
								[\denotes{\ibar{v}}
								[\denotes{\obar{v}} [agent]]
								[\denotes{VP} [\denotes{\ibar{V}}
									[\denotes{\ibar{V}} 
										[\denotes{\obar{V}} [$\surd$introduce]]
										[\denotes{DP_5} [\shade{who}, roof] ]
									]
									[\denotes{PP} [to me, roof] ]
								] ]
							 ]
						 ]
					 ]
				] ]
			] ]
		] ]
	]]
	\end{forest}%
	%
	\marginnote[0pt]{I've only indicated the semantically interpreted and uninterpreted parts of the moved object~DP. Assume that all the phrases that aren't parts of copies are semantically interpreted in~(\ref{ex6:copy1latea})}%
	%
\end{exe}
\end{fullwidth}
\medskip
Note that the relative clause is semantically interpreted in the higher copy but not the lower copy. This will give rise to the desired outcome; let's see how.

In the complement position of the VP headed by \emph{$\surd$introduce} is spoken the DP \emph{everyone who she should}. (The VP \emph{introduce to me} is silent because of ellipsis.) But only the determiner \emph{every} and the \ibar{N} \emph{one} are semantically interpreted in this position. Thus, the phrases that dominate these terms have meanings that are determined only by \emph{every} and \emph{one}: the DP in object position has the meaning that is assigned to \emph{everyone}. It doesn't include the semantic contribution of the relative clause. Because \emph{everyone} in this position is the lower of two copies, it can get the semantic interpretation that involves Trace Conversion; it can be interpreted as equivalent to \emph{the_3 one}. Thus, the smaller TP of (\ref{ex6:copy1latea}) will have an interpretation like that sketched in~(\ref{ex6:antlate1}).
\begin{exe}
\raggedright
	\ex \label{ex6:antlate1}
	$\lambda x$ she will introduce the_x one to me 
\end{exe}

The higher copy of the object DP has a meaning that derives from interpreting everything inside it, including the relative clause. The VP of that relative clause contains the lower copy of the moved relative pronoun \emph{who}. Because it is the lower copy of a movement relation, it too can be given an interpretation by Trace conversion. Its Trace converted interpretation is \emph{the_5 one}. This makes the TP of the relative clause have an interpretation like that sketched in~(\ref{ex6:elllate1}).
\begin{exe}
\raggedright
	\ex \label{ex6:elllate1}
	$\lambda y$ she should_F introduce the_y one to me
\end{exe}
Because (\ref{ex6:antlate1}) matches (\ref{ex6:elllate1}), except where there is a focus-marked term, the conditions on focus anaphora and ellipsis are met. In this way, by ignoring the relative clause's semantic contribution in the lower copy, Copy theory is capable of essentially mimicking Trace theory.

This method of accounting for antecedent contained anaphora showcases that part of (\ref{ex6:map2}) which allows just a part of a copy to be semantically interpreted. But (\ref{ex6:map2}) also allows just a portion of a copy to be phonologically interpreted as well. Jon Nissenbaum and Danny Fox\cite{Fox:1999a} have argued that examples like (\ref{ex6:extracopy1}) arise in this way.
\begin{exe}
\raggedright
	\ex \label{ex6:extracopy1}
	She will introduce everyone to me who she should.
\end{exe}
The relative clause in (\ref{ex6:extracopy1}) is not pronounced together with the DP it is semantically part of. One way of accounting for this case is to posit a Transformational rule that moves relative clauses to the right. Using the Trace theory, this rule---sometimes called ``extraposition from NP''---would permit derivations to contain parses like those in~(\ref{ex6:extrarule}).
\medskip
\begin{fullwidth}
\begin{exe}
\raggedright
	\ex \label{ex6:extrarule}
	\footnotesize
	\begin{forest}
	[TP
		[TP
			[DP^2 [she, roof] ]
			[\ibar{T}
				[\obar{T} [will]]
				[vP
					[\emph{t_2} ]
					[\ibar{v}
						[\obar{v} [agent]]
						[VP [\ibar{V}
							[\ibar{V}
							[\obar{V} [$\surd$introduce]]
							[DP [\ibar{D}
								[\obar{D} [every]]
								[NP [\ibar{N}
									[\ibar{N} [one, roof]]
									[\emph{t_3}, name=trace ]
								] ]
							] ]
						]
						[PP [to me, roof]] ]
						]
					 ]
				 ]
			 ]
		 ]
		[CP^3 [who she should, roof, name=binder] ]
	]
	\draw[->,>=stealth', thin, gray] (trace) to[out=north east,in=south east,looseness=1.2] node[pos=0.75,xshift=10mm]{} (binder);
	\end{forest}
\end{exe}
\end{fullwidth}
\medskip
This is the analysis offered in John Ross's dissertation.\cite{Ross:1967} But there is compelling evidence that a better analysis is one that uses Copy theory, and makes use of that part of (\ref{ex6:map2}) that allows only portions of a copy to be pronounced. That view would give to (\ref{ex6:extracopy1}) a representation like~(\ref{ex6:extracopy1p}).
\medskip
\begin{fullwidth}
\begin{exe}
\raggedright
	\ex \label{ex6:extracopy1p}
	\footnotesize
	\begin{forest}
	[TP
		[TP
			[DP^2 [she, roof] ]
			[\ibar{T}
				[\obar{T} [will]]
				[vP
					[DP_2 [\shade{she}, roof]]
					[\ibar{v} 
						[\obar{v} [agent] ]
						[VP [\ibar{V}
							[\ibar{V}
								[\obar{V} [$\surd$introduce]]
								[DP, name=trace [\ibar{D}
									[\obar{D}_3 [every]]
									[NP 
										[\ibar{N} [one, roof]]
										[CP [\shade{who she should}, roof] ]
									 ]
								] ]
							 ]
							[PP [to me, roof] ]
						] ]
					 ]
				 ]
			 ]
		 ]
		[DP^3, name=binder [\ibar{D} 
			[\obar{D} [\shade{every}]]
			[NP 
				[\ibar{N} [\shade{one}]]
				[CP [who she should, roof] ]
			 ]
		] ]
	]
	\draw[->,>=stealth', thin, gray] (trace) to[out=north east,in=south east,looseness=1.3] node[pos=0.75,xshift=10mm]{} (binder);
	\end{forest}
\end{exe}
\end{fullwidth}
\medskip

A salient difference between the two accounts concerns what they say about the interpretation of the quantificational determiner. In the account that (\ref{ex6:extrarule}) exemplifies, the position that the quantificational determiner is interpreted in is independent of the position of the relative clause. But in the account that (\ref{ex6:extracopy1p}) sketches, the relative clause is in the position that the DP headed by the quantificational determiner occupies. As a consequence, the quantifier must be interpreted in a position that matches where the relative clause is spoken in the second account, but not the first. It's this difference that \citet{Fox:1999a} base their account on.%
%
\marginnote[0pt]{See \citet{Overfelt:2015} for another kind of argument for (\ref{ex6:extracopy1p}). \citet{Baltin:1987} is where the idea originates for taking the syntax of extraposed relative clauses to also be the syntax that resolves antecedent contained anaphora.}%
%
They discover examples which suggest that the position a quantifier is interpreted in is dependent on the position that the relative clause is pronounced in the way that (\ref{ex6:extracopy1p}) predicts.

Let's look at one of their cases. Consider first~(\ref{ex6:extrafn1}).
\begin{exe}
\raggedright
	\ex \label{ex6:extrafn1}
	I must read every paper that I assigned before you do.
\end{exe}
This example has an adjunct clause containing an elided VP whose meaning is determined by the meaning of the other VP in the sentence. We can use this fact to determine what the syntactic form of the other VP is. What we discover is that the VP which forms the antecedent to the ellipsis has two meanings, and therefore two syntactic forms. One of these meanings is, simply, the one associated with \emph{read every paper that I assigned}. On that interpretation, (\ref{ex6:extrafn1}) is synonymous with~(\ref{ex6:extrafn1i}).
\begin{exe}
\raggedright
	\ex \label{ex6:extrafn1i}
	I must read every paper that I assigned before you read every paper that I assigned.
\end{exe}
What (\ref{ex6:extrafn1i}) expresses is the requirement that I finish reading the assigned papers before you finish reading the assigned papers. This requirement can be satisfied even if you read some of the papers before I do. As long as I have read every assigned paper before you have finished reading those papers, the requirement that (\ref{ex6:extrafn1i}) describes will be satisfied.

A parse that corresponds to this interpretation is~(\ref{ex6:extrafn1p}).
\medskip
\begin{fullwidth}
\begin{exe}
\raggedright
	\ex \label{ex6:extrafn1p}
	\footnotesize
	\begin{forest}
	[TP
	[\ibar{T}
		[\ibar{T}
			[\obar{T} [must]]
			[\fbox{AsP} 
				[AsP [\ibar{Asp} 
					[\obar{Asp} ]
					[vP
					[DP^x [I, roof]]
						[\ibar{v}
							[\obar{v} [agent]]
							[VP [\ibar{V}
									[\obar{V} [$\surd$read]]
									[$y$ ]
							] ]
						 ]
					 ]
				]]
				[DP^y [every paper\\ that I assigned, roof]]
			 ]
		]
		[PP [\ibar{P}
			[\obar{P} [before]]
			[TP [\ibar{T}
				[\obar{T} [do]]
				[\fbox{AsP}
					[AspP [\ibar{Asp} 
						[\obar{AsP} ]
						[vP
							[DP [you_F, roof] ]
							[\ibar{V}
								[\obar{v} [agent]]
								[VP [\ibar{V}
									[\obar{V} [\sout{$\surd$read}]]
									[$x$ ]
								] ]
							 ]
						 ]
					] ]
					[DP^x [\sout{every paper}\\ \sout{that I assigned}, roof]]
				 ]
			] ]
		]]
	]]
	\end{forest}
\end{exe}
\end{fullwidth}
\medskip
In (\ref{ex6:extrafn1p}), I've dodged the issue of what is in the position from which the DPs have moved by putting in those positions the variables $x$ and~$y$. What's key is that the object DP \emph{every paper that I assigned} is inside a constituent---here AspP---that serves as the antecedent for the focus anaphor inside the PP headed by \emph{before}. The boxed AspPs are identical, except for constituents that are focus-marked, and this is what is required for the VP inside AspP to elide. The parse in (\ref{ex6:extrafn1p}) can be derived from the S-structure parse for (\ref{ex6:extrafn1}) by moving the two instances of \emph{every paper that I assigned} outside the VP they are part of. Because they are part of those VPs in the S-structure, however, VP ellipsis will cause the second \emph{every paper that I assigned} to be elided. The elided material is shown in strikeouts. The S-structure will have the subject DPs (i.e., \emph{I} and \emph{you}) in Specifier of TP, but they will be semantically interpreted as if they haven't moved from the positions they occupy in~(\ref{ex6:extrafn1p}). I've taken the shortcut here of showing these DPs in the positions that they are interpreted. A more accurate parse would put them in their S-structure position, and use the semantics of movement to give the vPs they've moved from the meanings indicated in~(\ref{ex6:extrafn1p}).

In addition to the meaning expressed by (\ref{ex6:extrafn1i}), (\ref{ex6:extrafn1}) can also have the different meaning paraphrased by~(\ref{ex6:extrafn1ii}).
\begin{exe}
\raggedright
	\ex \label{ex6:extrafn1ii}
	Every paper that I assigned is such that I must read it before you read~it.
\end{exe}
On this interpretation, (\ref{ex6:extrafn1}) expresses the requirement that every assigned paper must be read by me before it is read by you. Unlike the meaning paraphrased by (\ref{ex6:extrafn1i}), the requirement expressed by (\ref{ex6:extrafn1ii}) is not satisfied if there is some assigned paper you read before I read it. (\ref{ex6:extrafn1ii}) requires not just that I finish reading all of the assigned papers before you finish reading them all, but that I finish reading each of the assigned papers before you've read that paper. To achieve this meaning with our semantics, we need a syntax that puts the \emph{before} PP inside a constituent that \emph{every paper} quantifies over. A consequence of that syntax is that \emph{every paper} won't be within VP that is elided; instead only a variable bound by \emph{every paper} will be. A parse that represents this interpretation is~(\ref{ex6:extrafn1pp}).
\medskip
\begin{fullwidth}
\begin{exe}
\raggedright
	\ex \label{ex6:extrafn1pp}
	\footnotesize
\begin{forest}
[TP [\ibar{T}
	[\obar{T} [must]]
	[AspP
		[AspP [\ibar{Asp}
			[\ibar{Asp}
				[\obar{Asp} ]
				[\fbox{vP}
					[DP [I, roof] ]
					[\ibar{v}
						[\obar{v} [agent]]
						[VP [\ibar{V}
							[\obar{V} [$\surd$read]]
							[$x$ ]
						] ]
					 ]
				 ]
			 ]
			[PP [\ibar{P}
				[\obar{P} [before] ]
				[TP [\ibar{T}
					[\obar{T} [do]]
					[AspP [\ibar{Asp}
						[\obar{Asp} ]
						[\fbox{vP}
							[DP [you_F, roof] ]
							[\ibar{v}
								[\obar{v} [agent]]
								[VP [\ibar{V}
									[\obar{V} [\sout{$\surd$read}]]
									[DP [\sout{it_x}, roof] ]
								] ]
							 ]
						 ]
					] ]
				] ]
			] ]
		] ]
		[DP^x [every paper\\ that I assigned, roof] ]
	 ]
]]
\end{forest}
\end{exe}
\end{fullwidth}
\medskip
As in (\ref{ex6:extrafn1p}), I've dodged specifying exactly what's in the position left by movement of the \emph{every paper} DP by putting an $x$ in that position, and I've put the subject DPs in their semantically interpreted positions, rather than their S-structure positions. Unlike (\ref{ex6:extrafn1p}), there is only one instance of \emph{every paper that I assigned} in (\ref{ex6:extrafn1pp}). In the object position of the elided VP is a pronoun that is bound by the the \emph{every paper} DP, rather than another \emph{every paper} DP. The boxed vPs have the same meaning, except for what is focused-marked, and should therefore meet the conditions on ellipsis.\marginnote{In fact, the antecedence conditions on focus anaphors, and therefore ellipsis, are not met in this case. That is partly because the semantics I have given of bound variables does not give to them a denotation until the $\lambda$-operator that binds them is included. As a consequence, there isn't a denotation delivered to any phrase that contains a bound variable, but not the $\lambda$-operator that binds it. What is needed is a semantics that allows variables to have a denotation in phrases in which they are not bound, and then a modification of the antecedence conditions on focus anaphors that lets two variables bound by the same $\lambda$-operator have the same denotation. It would take us too far into the technicalia of variable binding to make these corrections, so I won't. See \citet{Heim:1993} for a direction.}

What \citet{Fox:1999a} observe is that (\ref{ex6:extrafn1p}) can't be the source for an example that is identical to (\ref{ex6:extrafn1}), but with the relative clause in extraposed position. The example in (\ref{ex6:extrafn2}) has only the interpretation that goes with (\ref{ex6:extrafn1pp}), not~(\ref{ex6:extrafn1p}).
\begin{exe}
\raggedright
	\ex \label{ex6:extrafn2}
	I must read every paper before you do that I assigned.
\end{exe}
This is what is expected if extraposition of a relative clause comes by way of selectively pronouncing parts of copies, but not if it comes by way of moving the relative clause. 

To see this, consider what an account which allows relative clauses to move would say about (\ref{ex6:extrafn2})'s derivation. That derivation would include the parse in~(\ref{ex6:extrafn2bp}).
\medskip
\begin{fullwidth}
	\begin{exe}
		\raggedright
			\ex \label{ex6:extrafn2bp}
			\footnotesize
			\begin{forest}
				[TP
				[TP
				[\ibar{T}
					[\ibar{T}
						[\obar{T} [must]]
						[\fbox{AsP} 
							[AsP [\ibar{Asp} 
								[\obar{Asp} ]
								[vP
								[DP [I, roof]]
									[\ibar{v}
										[\obar{v} [agent]]
										[VP [\ibar{V}
												[\obar{V} [$\surd$read]]
												[$y$ ]
										] ]
									 ]
								 ]
							]]
							[DP^y [\ibar{D}
								[\obar{D} [every]]
								[NP [\ibar{N}
									[\ibar{N} [paper, roof]]
									[$z$, name=trace ]
								] ]
							]]
						 ]
					]
					[PP [\ibar{P}
						[\obar{P} [before]]
						[TP [\ibar{T}
							[\obar{T} [do]]
							[\fbox{AsP}
								[AspP [\ibar{Asp} 
									[\obar{AsP} ]
									[vP
										[DP [you_F, roof] ]
										[\ibar{V}
											[\obar{v} [agent]]
											[VP [\ibar{V}
												[\obar{V} [\sout{$\surd$read}]]
												[$x$ ]
											] ]
										 ]
									 ]
								] ]
								[DP^x [\sout{every paper}\\ \sout{that I assigned}, roof]]
							 ]
						] ]
					]]
				]]
				[CP^z [that I assigned, name=binder, roof]]
				]
				\draw[->,>=stealth', thin, gray] (trace) to[out=south east,in=south,looseness=1] node[pos=0.75,xshift=10mm]{} (binder);	
				\end{forest}		
		\end{exe}
\end{fullwidth}
\medskip
Note that (\ref{ex6:extrafn2bp}) presents a problem of reconstruction: the moved relative clause must be semantically interpreted in the position it occupies before movement. However this problem is solved, the fact that the relative clause is interpreted as if it hadn't moved is why the antecedent in (\ref{ex6:extrafn2bp}) behaves as if it contains the relative clause for the purposes of satisfying the conditions on the focus anaphor. That is, the antecedent AspP is equivalent to \emph{I read every paper that I assigned}, and therefore eligible to be the antecedent to the anaphoric AspP. In other words, movement of the relative clause does not affect the ability of VP Ellipsis in~(\ref{ex6:extrafn2bp}). This can be confirmed by observing that the elided VP in~(\ref{ex6:extravpd}) has the interpretation that would be given by (\ref{ex6:extrafn2bp}) to~(\ref{ex6:extrafn2}).
\begin{exe}
\raggedright
	\ex \label{ex6:extravpd}
	I must read every paper before this afternoon that was assigned in class and you must $\triangle$ before tomorrow.\\
	$\triangle$ = read every paper that was assigned in class
\end{exe}
The antecedent in (\ref{ex6:extravpd}) is \emph{read every paper that was assigned in class}, even though the relative clause is pronounced outside that constituent.

An account of extraposed relative clauses that claims they move, therefore, wrongly predicts that (\ref{ex6:extrafn2}), and its accompanying interpretation, should be available to (\ref{ex6:extrafn2}). But an account like that of Fox and Nissenbaum's doesn't. That's because on that account, the relative clause is spoken inside the higher copy of the DP it modifies. When the relative clause is in extraposed position, it is in a higher position than the material that precedes it. In (\ref{ex6:extrafn2}), that means that the \emph{every paper} DP is outside the constituent that serves as antecedent to the focus anaphor; (\ref{ex6:extrafn2gp}) illustrates.
\medskip
\begin{fullwidth}
\begin{exe}
\raggedright
	\ex \label{ex6:extrafn2gp}
	\footnotesize
	\begin{forest}
	[TP [\ibar{T}
		[\obar{T} [must]]
		[AspP
			[AspP [\ibar{Asp}
				[\ibar{Asp}
				[\obar{Asp} ]
				[\fbox{vP}
					[DP [I, roof] ]
					[\ibar{v}
						[\obar{v} [agent]]
						[VP [\ibar{V}
							[\obar{V} [$\surd$read]]
							[DP [\ibar{D}
								[\obar{D}_2 [every]]
								[NP [\ibar{N}
									[\ibar{N} [paper, roof]]
									[CP [\shade{that I assigned}, roof] ]
								] ]
							] ]
						] ]
					 ]
				 ]
			] 
			[PP [\ibar{P}
				[\obar{P} [before]]
				[TP [\ibar{T}
					[\obar{T} [do]]
					[AsP [\ibar{Asp}
						[\obar{Asp} ]
						[\fbox{vP}
							[DP [you_F, roof] ]
							[\ibar{v}
								[\obar{v} [agent]]
								[VP [\ibar{V}
									[\obar{V} [\sout{$\surd$read}]]
									[DP [\sout{it_2}, roof] ]
								] ]
							 ]
						 ]
					] ]
				] ]
			]]
			]]
			[DP^2 [\ibar{D}
				[\obar{D} [\shade{every}]]
				[NP [\ibar{N}
					[\ibar{N} [\shade{paper}, roof]]
					[CP [that I assigned, roof] ]
				] ]
			]]
		 ]
	]]
	\end{forest}
\end{exe}
\end{fullwidth}
\medskip
The higher copy of the \emph{every paper} DP is the one in which \emph{every} is interpreted as the universal quantifier. This higher copy cannot be interpreted inside a constituent to the left of the PP, and cannot therefore be part of the antecedent for the ellipsis in that PP. This is how the Copy theory of movement prevents (\ref{ex6:extrafn2}) from having an interpretation in which the ellipsis contains an \emph{every paper} DP, i.e.,~(\ref{ex6:extrafn1p}).

Instead, (\ref{ex6:extrafn2gp}) produces the meaning that is associated with (\ref{ex6:extrafn1pp}). Let's see how. The antecedent (boxed) vP in (\ref{ex6:extrafn2gp}) is semantically parallel to (\ref{ex6:extrafn2para}a), and the elided (also boxed) vP is parallel semantically to~(\ref{ex6:extrafn2para}b).
\begin{exe}
\raggedright
	\ex \label{ex6:extrafn2para}
	\begin{xlist}
		\ex I read that paper that I assigned
		\ex You_F read it
	\end{xlist}
\end{exe}
The lower copy of the movement relation in (\ref{ex6:extrafn2gp}) gets an interpretation as a definite description with a bound variable meaning, and I've tried to evoke that interpretation in (\ref{ex6:extrafn2para}a).%
%
\marginnote{Again: this isn't technically true in the semantics provided for bound variables in this book. As noted in the previous sidebar, our semantics needs to be modified so that bound variables have an interpretation in phrases that don't include their $\lambda$-operator.}%
%
\ Because the lower copy and the \emph{it} in the elided vP are both bound by the higher copy of \emph{every paper that I assigned}, they have the same meaning. As a consequence, the boxed vPs meet the condition on focus anaphora---their meanings are identical except for what is focus-marked---and the ellipsis indicated with strike-outs is permitted.

There is evidence, then, that extraposition of a relative clause involves the selective pronunciation of parts of two copies. And Antecedent Contained anaphora can be captured by the Copy theory of movement if it is possible to semantically interpret parts of two copies. Together, then, these cases support (\ref{ex6:map2}), repeated here.
\begin{exe}
	\raggedright
		\exr{ex6:map2}
		Let \{$\gamma$, $\gamma '$, $\gamma ''$, \ldots\} be matching phrases or heads dominated by copies \{$\alpha$,~$\alpha '$,~$\alpha ''$, \ldots\}.
		\begin{xlist}
			\ex Only one $\gamma$^n may be linearized.
			\ex Only one $\gamma$^n must be semantically interpreted.
		\end{xlist}
\end{exe}
If Movement works in tandem with (\ref{ex6:map2}), then the derivations it defines have a  different relationship to the sound-meaning mapping. Without (\ref{ex6:map2}), movement provides the range of parses from which a spoken form and a semantically interpreted form can be derived. With (\ref{ex6:map2}), each parse in a derivation has a variety of ways of furnishing a spoken and semantically interpreted form. Each parse in a derivation corresponds to a range of different sound-meaning relationships.

This makes it possible to consider rethinking the relationship between a derivation and the parse that represents its spoken string. It may be possible, for instance, to let S-structures always correspond to the final parse in a derivation. The present definition of derivations allows the S-structure---the parse that corresponds to the spoken string---to feed parses that are semantically interpreted. In many situations involving quantificational DPs, for instance, we have let a semantically interpreted parse be derived from the S-structure by Quantifier Raising. But if the approach to extraposed relative clauses described above is correct, this isn't necessary. The representation that gives rise to extraposed relatives could be interpreted so that all of the lower copy is pronounced and none of the higher copy is. For instance, an alternative to (\ref{ex6:extrafn2gp}) might be (\ref{ex6:singleo1}), where all of the lower copy is pronounced and none of the higher one is.
\medskip
\begin{fullwidth}
\begin{exe}
\raggedright
	\ex \label{ex6:singleo1}
	\footnotesize
	\begin{forest}
	[TP [\ibar{T}
		[\obar{T} [must]]
		[AspP
			[AspP [\ibar{Asp}
				[\ibar{Asp}
				[\obar{Asp} ]
				[\fbox{vP}
					[DP [I, roof] ]
					[\ibar{v}
						[\obar{v} [agent]]
						[VP [\ibar{V}
							[\obar{V} [$\surd$read]]
							[DP [\ibar{D}
								[\obar{D}_2 [every]]
								[NP [\ibar{N}
									[\ibar{N} [paper, roof]]
									[CP [that I assigned, roof] ]
								] ]
							] ]
						] ]
					 ]
				 ]
			] 
			[PP [\ibar{P}
				[\obar{P} [before]]
				[TP [\ibar{T}
					[\obar{T} [do]]
					[AsP [\ibar{Asp}
						[\obar{Asp} ]
						[\fbox{vP}
							[DP [you_F, roof] ]
							[\ibar{v}
								[\obar{v} [agent]]
								[VP [\ibar{V}
									[\obar{V} [\sout{$\surd$read}]]
									[DP [\sout{it_2}, roof] ]
								] ]
							 ]
						 ]
					] ]
				] ]
			]]
			]]
			[DP^2 [\ibar{D}
				[\obar{D} [\shade{every}]]
				[NP [\ibar{N}
					[\ibar{N} [\shade{paper}, roof]]
					[CP [\shade{that I assigned}, roof] ]
				] ]
			]]
		 ]
	]]
	\end{forest}
\end{exe}
\end{fullwidth}
\medskip
The S-structure in (\ref{ex6:singleo1}) has QR built into it. 

The Copy theory of movement allows S-structures to have all of the information that the parses in the derivation leading up to that S-structure have. And amending the Copy theory of movement with (\ref{ex6:map2}) allows an S-structure to also contain the information that the parses which could be derived from that S-structure would have. S-structures, then, can have all the information needed for both the semantic interpretation and the pronunciation of a sentence. Jonathan Bobaljik\cite{Bobaljik:1995a} argues for this theory of derivations, and calls it a ``single output'' model. In a single output model there is one parse that contains all the information necessary to get the correct semantics and phonology. On these models the place for derivations is unclear. Derivations were originally conceived as a way of expressing how the semantic phrases in a sentence don't always correspond to a contiguous string---a phonological phrase. Derivations do this by letting a sentence be mapped onto a series of parses, one of which represents the string and another that represents the semantic composition. On a single output model, there is one parse that has both the words in their spoken positions and the words in their semantically interpreted positions. Derivations are no longer what represents the differing spoken and interpreted positions of phrases and their words. Instead, derivations are just a way of expressing the function that takes a set of morphemes into one syntactic representation. If a single output model should prove to be correct, we should reassess what is required of this function. Thinking of it as a series of operations---Transformations---may not be needed. We may find calling the model of syntax these lectures are charged with explicating ``Transformational Grammar,'' to be misleading.

\section{Summary}
\label{sec:summary}

This chapter has sketched two theories of movement. Both let movement change the meaning of the parse they apply to. In both, the moved term can become a binder of a variable left in the position the term moved from. In Trace theory, this is achieved by letting the position a term moves from be occupied by a silent variable bound by the term from its new position. In the Copy theory of movement, this is achieved by letting a term be turned into a definite description with the meaning of a variable when a copy of it is put into a position from which it binds that definite description. The Copy theory of movement needs to be wedded to a theory that determines which parts of two copies can, and can't, be pronounced and semantically interpreted. The Trace theory of movement needs to be wedded to a theory that explains why a trace is silent and why it arises only in movement contexts.

The Trace theory of movement is demonstrably inadequate. It has not been possible to define derivations correctly under Trace theory. In particular, it's not possible to generate parses in a derivation that correspond to the meanings needed. At least this is not possible unless derivations are allowed to violate Upwards, and without Upwards it's not presently known how to prevent derivations from allowing too many possibilities. 

The Copy theory of movement, on the other hand, does not straightforwardly account for antecedent contained anaphora. At the time of writing, the remedy for this has been to interpret movement as an operation that gives a term two positions from which parts of that term can be semantically and phonologically interpreted. On this view, the parses in a derivation created by movement don't uniquely determine one denotation and one string each, but instead proffer a range of possible denotations and strings. While this successfully captures the cases that are problematic for Trace theory, it allows too many possibilities. Without a model of which parts of the two copies can be semantically and phonologically interpreted, there are too many ways for even one instance of movement to provide a string-meaning correspondence. The Copy theory of movement relies on the success of a program that has yet to provide such a model.

One piece of striking evidence that the Copy theory of movement is on the right track comes from examples in which a moved term can be pronounced in both of its positions. This arises in many languages in cases where a verb, or a verb phrase, are topicalized or clefted. Many languages have a rule that moves a constituent into a clause initial position when that constituent is intended to be contrasted, or otherwise highlighted. In some languages it appears that a verb can undergo this process, leaving its projection behind. In many languages that have this process the moved verb must be pronounced in both of its positions. Some examples are in~(\ref{ex:verbcopy1}).\marginnote{\citet{Nunes:1995,Nunes:2004} pioneers an approach to this phenomenon which involves combining the Copy Theory of Movement with a theory that expresses how syntactic representations are mapped onto strings. See \citet{Landau:2006}, \citet{Trinh:2011}, and \citet{Hein:2018} for more. }
\begin{exe}
\raggedright
	\ex \label{ex:verbcopy1}
	\begin{xlist}
		\ex \gll k{\'a} \textipa{d\=E}-\textipa{k\=a} \textipa{\`al\=\i} \`A\textipa{t\`{\i}m} \textipa{d\`E} \textipa{m\'ang\`o-ku} \textipa{d\=\i\=em}\\ FOC eat-nmlz C Atim ate mango-def yesterday\\ ``It's eating that Atim ate the mango yesterday.''\\ \hfill \citep[\emph{Buli,}][p. 262]{Hiraiwa:2005}\hfill\hspace*{.15in}
		\ex \gll masi-ki-nun Chelsu-ka mayckwu-lul masi-ess-ta\\ drink-nmlz-top Chelsu-nom beer-acc drink-pst-dec\\ ``As for drinking, Chelsu drank beer.''\\ \hfill \citep[\emph{Korean,}][p. 90]{Jo:2000}\hfill\hspace*{.15in}
		\ex \gll doc thi no nen doc sach.\\ read top he should read book\\ ``As for reading this book, I read.''\\
		\hfill \citep[\emph{Vietnamese,}][p. 60]{Trinh:2011}\hfill\hspace*{.15in}
	\end{xlist}
\end{exe}
The verbs in these constructions can have different inflections, but they cannot otherwise be different. They look, in other words, like they are copies of the same verb, but with the morphology that is appropriate for the position they are spoken in. The conditions we normally see of movement---like Upwards---also govern the relationship between the higher and lower verb. This is precisely what would be expected on the Copy theory of movement, if the prohibition against speaking both copies was lifted. By contrast, it's not at all what would be expected on Trace theory. 
